
%%FILE blueprint/src/content.tex
% In this file you should put the actual content of the blueprint.
% It will be used both by the web and the print version.
% It should *not* include the \begin{document}
%
% If you want to split the blueprint content into several files then
% the current file can be a simple sequence of \input. Otherwise It
% can start with a \section or \chapter for instance.

Warning: this document is a draft. Nothing in it is complete. It lacks references and content, its scope is not final, it will likely be reorganized several times and does not reflect the intended final version in any way.

\input{sections/introduction}

\chapter{Testing, Estimation, Risk}

\input{sections/risk}
\input{sections/sample_complexity}

\chapter{Information divergences}

\input{sections/dpi_div}
\input{sections/stat_div}
\input{sections/f_divergence}

\input{sections/kl_divergence}

\input{sections/hellinger_alpha}
\input{sections/renyi_divergence}
\input{sections/hellinger}

\input{sections/chernoff}

\input{sections/jensen_shannon}

\chapter{Mutual information}

\input{sections/mutual_info}

\chapter{The risk of simple binary hypothesis testing}

\input{sections/testing}
\input{sections/sequential}

\bibliographystyle{amsalpha}
\bibliography{biblio}

\appendix

\input{sections/kernels}
\input{sections/rn_deriv}
\input{sections/convex}
%\input{sections/deficiency}

%%FILE blueprint/src/sections/chernoff.tex
\section{Chernoff divergence}

\begin{definition}[Chernoff divergence]
  \label{def:Chernoff}
  \lean{ProbabilityTheory.chernoffDiv}
  \leanok
  \uses{def:Renyi, lem:renyi_properties}
    The Chernoff divergence of order $\alpha > 0$ between two measures $\mu$ and $\nu$ on $\mathcal X$ is
  \begin{align*}
    C_\alpha(\mu, \nu) = \inf_{\xi \in \mathcal P(\mathcal X)}\max\{R_\alpha(\xi, \mu), R_\alpha(\xi, \nu)\} \: .
  \end{align*}
\end{definition}

Note: the name Chernoff divergence is usually reserved for $C_1$, and any reference to $C$ without subscript in this document refers to $C_1$.
The extension to other orders $\alpha$ is motivated by the appearance of $C_\alpha$ in a result further down.
This is not a well established notion and it remains to see if this extension is meaningful.

\begin{lemma}
  \label{lem:chernoff_eq_kl}
  \lean{ProbabilityTheory.chernoffDiv_one}
  \leanok
  \uses{def:Chernoff, def:KL}
  $C_1(\mu, \nu) = \inf_{\xi \in \mathcal P(\mathcal X)}\max\{\KL(\xi, \mu), \KL(\xi, \nu)\}$ .
\end{lemma}

\begin{proof}\leanok
\uses{def:Renyi}
This is $R_1 = \KL$ (by definition).
\end{proof}


\begin{lemma}[Symmetry]
  \label{lem:chernoff_symm}
  %\lean{}
  %\leanok
  \uses{def:Chernoff}
  $C_\alpha(\mu, \nu) = C_\alpha(\nu, \mu)$.
\end{lemma}

\begin{proof}%\leanok
Immediate from the definition.
\end{proof}


\begin{lemma}[Monotonicity]
  \label{lem:chernoff_mono}
  %\lean{}
  %\leanok
  \uses{def:Chernoff}
  The function $\alpha \mapsto C_\alpha(\mu, \nu)$ is monotone.
\end{lemma}

\begin{proof}
\uses{lem:renyi_monotone}
Consequence of Lemma~\ref{lem:renyi_monotone}.
\end{proof}


\begin{lemma}
  \label{lem:chernoff_eq_max_renyi}
  %\lean{}
  %\leanok
  \uses{def:Chernoff, def:Renyi}
  \begin{align*}
  C_1(\mu, \nu) = \max_{\alpha\in [0,1]} (1 - \alpha)R_\alpha(\mu, \nu) \: .
  \end{align*}
\end{lemma}

\begin{proof}%\leanok
\uses{}
\end{proof}



%%FILE blueprint/src/sections/convex.tex
\chapter{Convex functions of Radon-Nikodym derivatives}

Let $f : \mathbb{R} \to \mathbb{R}$ be convex on $[0, +\infty)$. Remark that $f(0)$ is finite.

\begin{definition}
  \label{def:derivAtTop}
  \lean{derivAtTop}
  \leanok
  %\uses{}
  We define $f'(\infty) := \limsup_{x \to + \infty} f(x)/x$. This can be equal to $+\infty$ (but not $-\infty$).
\end{definition}

\begin{lemma}
  \label{lem:integrable_f_rnDeriv_of_derivAtTop_ne_top}
  \lean{MeasureTheory.integrable_f_rnDeriv_of_derivAtTop_ne_top}
  \leanok
  \uses{def:derivAtTop}
  If $\mu$ and $\nu$ are two finite measures and $f'(\infty) < \infty$, then $x \mapsto f\left(\frac{d\mu}{d\nu}(x)\right)$ is $\nu$-integrable. 
\end{lemma}

\begin{proof} \leanok
By convexity and $f(0) < \infty$, $f'(\infty)<\infty$, we can sandwich $f$ between two affine functions of $\frac{d\mu}{d\nu}$.
Those functions are integrable since the measures are finite.
\end{proof}

\begin{theorem}
  \label{thm:condexp_jensen}
  \lean{ConvexOn.map_condExp_le}
  \mathlibok
  %\uses{}
  For $f$ convex, $f\left(\mu [g \mid m]\right) \le \mu [f \circ g \mid m]$.
\end{theorem}

\begin{proof}\mathlibok
\end{proof}

%%FILE blueprint/src/sections/deficiency.tex
\chapter{Deficiency}

TODO: this part is not used elsewhere for now, and should perhaps be deleted.

\begin{definition}[Le Cam deficiency]
  \label{def:deficiency}
  %\lean{}
  %\leanok
  \uses{def:TV}
  The Le Cam deficiency of $\kappa : \mathcal X \rightsquigarrow \mathcal Y$ with respect to $\eta : \mathcal X \rightsquigarrow \mathcal Z$ is
  \begin{align*}
  \delta_{\TV}(\kappa, \eta) = \inf_{\xi} \sup_{x \in \mathcal X} \TV(\xi \circ \kappa (x), \eta(x)) \: ,
  \end{align*}
  in which the infimum is taken over all Markov kernels $\xi : \mathcal Y \rightsquigarrow \mathcal Z$~.
\end{definition}

TODO: from any divergence on measures $D$, \cite{perrone2023markov} defines a divergence between kernels $\kappa, \kappa' : \mathcal X \rightsquigarrow \mathcal Y$ by $D(\kappa, \kappa') = \sup_{x \in \mathcal X} D(\kappa(x), \eta(x))$.
If we see measures as kernels from a point space, the divergence between measures-as-kernels coincides with the divergence between the measures.
With that notation, $\delta_{\TV}(\kappa, \eta) = \inf_{\xi} \TV(\xi \circ \kappa, \eta)$.
We could also proceed as \cite{raginsky2011shannon} and extend the notion of deficiency to deficiency with respect to a divergence: $\delta_D(\kappa, \eta) = \inf_{\xi} D(\eta, \xi\circ\kappa)$ (note the reversed arguments, which does not matter for the case of $\TV$ since it is symmetric).

\begin{lemma}
  \label{lem:deficiency_nonneg}
  %\lean{}
  %\leanok
  \uses{def:deficiency}
  Let $\kappa: \mathcal X \rightsquigarrow \mathcal Y$ and $\eta : \mathcal X \rightsquigarrow \mathcal Z$. Then $\delta_{\TV}(\kappa, \eta) \ge 0$~.
\end{lemma}

\begin{proof}%\leanok
\uses{lem:tv_nonneg}
This is a consequence of Lemma~\ref{lem:tv_nonneg}.
\end{proof}

\begin{lemma}
  \label{lem:deficiency_self}
  %\lean{}
  %\leanok
  \uses{def:deficiency}
  Let $\kappa: \mathcal X \rightsquigarrow \mathcal Y$. Then $\delta_{\TV}(\kappa, \kappa) = 0$~.
\end{lemma}

\begin{proof}%\leanok
\uses{lem:deficiency_nonneg}
By Lemma~\ref{lem:deficiency_nonneg}, $0 \le \delta_{\TV}(\kappa, \kappa)$.
In the infimum that defines $\delta_{\TV}(\kappa, \kappa)$, take $\xi = \textup{id}$ to get an upper bound of 0 on $\delta_{\TV}(\kappa, \kappa)$.
\end{proof}

\begin{lemma}
  \label{lem:deficiency_comp}
  %\lean{}
  %\leanok
  \uses{def:deficiency, def:blackwellOrder}
  Let $\kappa: \mathcal X \rightsquigarrow \mathcal Y$ and $\eta : \mathcal X \rightsquigarrow \mathcal Z$.
  If $\eta \le_B \kappa$ (that is, there exists a Markov kernel $\chi : \mathcal Y \rightsquigarrow \mathcal Z$ such that $\eta = \chi \circ \kappa$), then $\delta_{\TV}(\kappa, \eta) = 0$~.
\end{lemma}

\begin{proof}%\leanok
\uses{lem:deficiency_nonneg}
In the infimum that defines $\delta_{\TV}(\kappa, \eta)$, take $\xi = \chi$ to get an upper bound of 0 on $\delta_{\TV}(\kappa, \eta)$. The lower bound is Lemma~\ref{lem:deficiency_nonneg}.
\end{proof}

\begin{lemma}
  \label{lem:deficiency_antimono_left}
  %\lean{}
  %\leanok
  \uses{def:deficiency, def:blackwellOrder}
  Let $\kappa: \mathcal X \rightsquigarrow \mathcal Y$, $\eta : \mathcal X \rightsquigarrow \mathcal Z$ and $\zeta : \mathcal X \rightsquigarrow \mathcal W$ be Markov kernels.
  If $\eta \le_B \kappa$ then $\delta_{\TV}(\kappa, \zeta) \le \delta_{\TV}(\eta, \zeta)$~.
\end{lemma}

\begin{proof}%\leanok
\uses{}
Let $\chi$ be a Markov kernel such that $\eta = \chi \circ \kappa$~.
\begin{align*}
\delta_{\TV}(\eta, \zeta)
&= \inf_{\xi} \sup_{x \in \mathcal X} \TV(\xi \circ \eta (x), \zeta(x))
\\
&= \inf_\xi \sup_{x \in \mathcal X} \TV(\xi \circ \chi \circ \kappa (x), \zeta(x))
\: .
\end{align*}
Since $\xi \circ \chi$ is a Markov kernel, the infimum over Markov kernels of the form $\xi \circ \chi$ is larger than the infimum over all Markov kernels.
\end{proof}

\begin{lemma}
  \label{lem:deficiency_mono_right}
  %\lean{}
  %\leanok
  \uses{def:deficiency, def:blackwellOrder}
  Let $\kappa: \mathcal X \rightsquigarrow \mathcal Y$, $\eta : \mathcal X \rightsquigarrow \mathcal Z$ and $\zeta : \mathcal X \rightsquigarrow \mathcal W$ be Markov kernels.
  If $\zeta \le_B \eta$ then $\delta_{\TV}(\kappa, \zeta) \le \delta_{\TV}(\kappa, \eta)$~.
\end{lemma}

\begin{proof}%\leanok
\uses{thm:tv_data_proc}
Let $\chi$ be a Markov kernel such that $\zeta = \chi \circ \eta$~.
\begin{align*}
\delta_{\TV}(\kappa, \zeta)
&= \inf_{\xi : \mathcal Y \rightsquigarrow \mathcal W} \sup_{x \in \mathcal X} \TV(\xi \circ \kappa (x), \zeta(x))
\\
&= \inf_{\xi : \mathcal Y \rightsquigarrow \mathcal W} \sup_{x \in \mathcal X} \TV(\xi \circ \kappa (x), \chi \circ \eta(x))
\\
&\le \inf_{\xi' : \mathcal Y \rightsquigarrow \mathcal Z} \sup_{x \in \mathcal X} \TV(\chi \circ \xi' \circ \kappa (x), \chi \circ \eta(x))
\: .
\end{align*}
By the data-processing inequality for $\TV$ (Theorem \ref{thm:tv_data_proc}), for all $\xi' : \mathcal Y \rightsquigarrow \mathcal Z$, $\TV(\chi \circ \xi' \circ \kappa (x), \chi \circ \eta(x)) \le \TV(\xi' \circ \kappa (x), \eta(x))$.
This concludes the proof.
\end{proof}


%%FILE blueprint/src/sections/dpi_div.tex
A divergence between measures is a function which maps two measures on any common space $\mathcal X$ to a value in $[-\infty, +\infty]$.
For two probability measures $\mu$ and $\nu$, we will want the divergence $D(\mu, \nu)$ to be non-negative and to be zero if $\mu = \nu$~: divergences are meant to be notions of dissimilarity between measures.

We will pay particular attention to how the divergence behaves with respect to transformations of the measures.
For a Markov kernel $\kappa : \mathcal X \rightsquigarrow \mathcal Y$~, a fundamental property of all divergences we will consider is the \emph{data-processing inequality}, which states that $D(\kappa \circ \mu, \kappa \circ \nu) \le D(\mu, \nu)$.
That is, processing data with a (possibly random) common transformation $\kappa$ can only lose information about the difference between the two measures.

Another important concept is the \emph{conditional divergence} of two kernels $\kappa, \eta : \mathcal X \rightsquigarrow \mathcal Y$ given a measure $\mu \in \mathcal M(\mathcal X)$. This is the value $D(\mu \otimes \kappa, \mu \otimes \eta)$.

TODO: two possible concepts of conditional divergence are $D(\mu \otimes \kappa, \mu \otimes \eta)$ and $\mu\left[x \mapsto D(\kappa(x), \eta(x))\right]$. They are equal for $f$-divergences (which gives also convexity of the divergence) but not in general. 

\section{Generic divergences}

\begin{definition}[Divergence]
  \label{def:div}
  %\lean{}
  %\leanok
  \uses{}
  A divergence between measures is a function $D$ which for any measurable space $\mathcal X$ and any two measures $\mu, \nu \in \mathcal M(\mathcal X)$, returns a value $D(\mu, \nu) \in \mathbb{R} \cup \{+\infty\}$.
\end{definition}


\begin{leanremark}
In Lean, such a definition for a divergence would have to restrict the possible measurable spaces to be in a common universe.
We don't implement that definition, but will only implement particular divergences.
\end{leanremark}


\begin{definition}[Data-processing inequality]
  \label{def:dpi}
  %\lean{}
  %\leanok
  \uses{def:div}
  A divergence $D$ is said to satisfy the data-processing inequality (DPI) if for all measurable spaces $\mathcal X, \mathcal Y$, all $\mu, \nu \in \mathcal M(\mathcal X)$ and all Markov kernels $\kappa : \mathcal X \rightsquigarrow \mathcal Y$,
  \begin{align*}
  D(\kappa \circ \mu, \kappa \circ \nu) \le D(\mu, \nu) \: .
  \end{align*}
\end{definition}


\begin{lemma}[Second marginal]
  \label{lem:div_comp_le_div_compProd}
  %\lean{}
  %\leanok
  \uses{def:div, def:dpi}
  Let $D$ be a divergence that satisfies the DPI. Let $\mathcal \mu, \nu \in \mathcal M(\mathcal X)$ and let $\kappa, \eta : \mathcal X \rightsquigarrow \mathcal Y$. Then
  \begin{align*}
  D(\kappa \circ \mu, \eta \circ \nu) \le D(\mu \otimes \kappa, \nu \otimes \eta) \: .
  \end{align*}
\end{lemma}

\begin{proof}%\leanok
\uses{def:dpi}
Use the DPI for the deterministic kernel of the function $(x,y) \mapsto y$.
\end{proof}


\begin{lemma}[First marginal]
  \label{lem:div_le_div_compProd}
  %\lean{}
  %\leanok
  \uses{def:div, def:dpi}
  Let $D$ be a divergence that satisfies the DPI. Let $\mathcal \mu, \nu \in \mathcal M(\mathcal X)$ and let $\kappa, \eta : \mathcal X \rightsquigarrow \mathcal Y$ be Markov kernels. Then
  \begin{align*}
  D(\mu, \nu) \le D(\mu \otimes \kappa, \nu \otimes \eta) \: .
  \end{align*}
\end{lemma}

\begin{proof}%\leanok
\uses{def:dpi}
Use the DPI for the deterministic kernel of the function $(x,y) \mapsto x$.
\end{proof}


\begin{lemma}
  \label{lem:div_compProd_right}
  %\lean{}
  %\leanok
  \uses{def:div, def:dpi}
  Let $D$ be a divergence that satisfies the DPI. Let $\mathcal \mu, \nu \in \mathcal M(\mathcal X)$ and let $\kappa : \mathcal X \rightsquigarrow \mathcal Y$ be a Markov kernel. Then
  \begin{align*}
  D(\mu \otimes \kappa, \nu \otimes \kappa) = D(\mu, \nu) \: .
  \end{align*}
\end{lemma}

\begin{proof}%\leanok
\uses{def:dpi, lem:div_le_div_compProd}
Lemma~\ref{lem:div_le_div_compProd} gives one inequality. The other inequality is the DPI: $\mu \otimes \kappa = (\mathrm{id} \times \kappa) \circ \mu$.
\end{proof}


\begin{definition}[Conditional divergence]
  \label{def:condDiv}
  %\lean{}
  %\leanok
  \uses{def:div}
  Let $D$ be a divergence.
  The conditional divergence of kernels $\kappa, \eta : \mathcal X \rightsquigarrow \mathcal Y$ with respect to a measure $\mu \in \mathcal M(\mathcal X)$ is
  $\mu[x \mapsto D(\kappa(x), \eta(x)]$.
  It is denoted by $D(\kappa, \eta \mid \mu)$.
\end{definition}

For the important class of $f$-divergences that we will introduce later, $D(\kappa, \eta \mid \mu) = D(\mu \otimes \kappa, \mu \otimes \eta)$.

\begin{lemma}[Conditioning increases divergence]
  \label{lem:div_comp_le_div_compProd_right}
  %\lean{}
  %\leanok
  \uses{def:div, def:condDiv, def:dpi}
  Let $D$ be a divergence that satisfies the DPI and for which $D(\kappa, \eta \mid \mu) = D(\mu \otimes \kappa, \mu \otimes \eta)$.
  Let $\mathcal \mu \in \mathcal M(\mathcal X)$ and let $\kappa, \eta : \mathcal X \rightsquigarrow \mathcal Y$ be Markov kernels. Then
  \begin{align*}
  D(\kappa \circ \mu, \eta \circ \mu) \le D(\kappa, \eta \mid \mu) \: .
  \end{align*}
\end{lemma}

\begin{proof}%\leanok
\uses{lem:div_comp_le_div_compProd}
This is a special case of Lemma~\ref{lem:div_comp_le_div_compProd}.
\end{proof}

%%FILE blueprint/src/sections/f_divergence.tex
\section{f-divergences}

\subsection{Definition and basic properties}

Everywhere in this document, the functions used in divergences are assumed measurable, continuous and convex on $[0, +\infty)$.
For such a function $f : \mathbb{R} \to \mathbb{R}$, we write $f'(\infty) = \lim_{x \to +\infty} f(x)/x$ (which can be infinite).

For $\mu, \nu$ two measures, we denote by $\mu_{\perp \nu}$ the singular part of $\mu$ with respect to $\nu$.

\begin{definition}[f-divergence]
  \label{def:fDiv}
  \lean{ProbabilityTheory.fDiv}
  \leanok
  \uses{def:derivAtTop}
  Let $f : \mathbb{R} \to \mathbb{R}$ and let $\mu, \nu$ be two measures on a measurable space $\mathcal X$. The f-divergence between $\mu$ and $\nu$ is
  \begin{align*}
  D_f(\mu, \nu) = \nu\left[x \mapsto f\left(\frac{d \mu}{d \nu}(x)\right)\right] + f'(\infty) \mu_{\perp \nu}(\mathcal X)
  \end{align*}
  if $x \mapsto f\left(\frac{d \mu}{d \nu}(x)\right)$ is $\nu$-integrable and $+\infty$ otherwise.
\end{definition}


\begin{lemma}
  \label{lem:fDiv_ne_top_iff}
  \lean{ProbabilityTheory.fDiv_ne_top_iff}
  \leanok
  \uses{def:fDiv}
  For $\mu$ and $\nu$ two finite measures, $D_f(\mu, \nu)$ is finite if and only if $x \mapsto f\left(\frac{d \mu}{d \nu}(x)\right)$ is $\nu$-integrable and either $f'(\infty) < \infty$ or $\mu \ll \nu$.
\end{lemma}

\begin{proof}\leanok
\end{proof}



\paragraph{Divergences associated with basic (transformations of) functions}

\begin{lemma}
  \label{lem:fDiv_const}
  %\lean{}
  %\leanok
  \uses{def:fDiv}
  For $\nu$ a finite measure, for all $a \in \mathbb{R}$, $D_{x \mapsto a}(\mu, \nu) = a \nu(\mathcal X)$.
\end{lemma}

\begin{proof} %\leanok
Compute the integral.
\end{proof}


\begin{lemma}
  \label{lem:fDiv_self}
  \lean{ProbabilityTheory.fDiv_self}
  \leanok
  \uses{def:fDiv}
  If $f(1) = 0$ then $D_{f}(\mu, \mu) = 0$.
\end{lemma}

\begin{proof} \leanok
$\frac{d \mu}{d \mu}(x) = 1$ almost everywhere and $f(1) = 0$.
\end{proof}


\begin{lemma}
  \label{lem:fDiv_eq_zero_iff}
  \lean{ProbabilityTheory.fDiv_eq_zero_iff}
  \leanok
  \uses{def:fDiv}
  Let $\mu, \nu$ be two finite measures. Assume that $f$ is strictly convex on $(0, +\infty)$ and $f(1) = 0$. Then $D_f(\mu, \nu) = 0$ if and only if $\mu = \nu$.
\end{lemma}

\begin{proof}\leanok
\uses{lem:fDiv_self}
If $\mu = \nu$ then $D_f(\mu, \nu) = 0$ by Lemma~\ref{lem:fDiv_self}.
For the other direction, since $f \ge 0$ is strictly convex with $f(1) = 0$, $f(x) = 0$ only for $x = 1$, and $f'(\infty) > 0$ (since $f(2) \le 2 f'(\infty)$ and $f(2) > 0$).
$D_f(\mu, \nu) = 0$ then implies $\mu_{\perp \nu}(\mathcal X) = 0$, hence $\mu \ll \nu$, and $f(\frac{d\mu}{d\nu}) = 0$ $\nu$-almost everywhere, hence $\frac{d\mu}{d\nu} = 1$ $\nu$-almost everywhere.
\end{proof}


\begin{lemma}
  \label{lem:fDiv_id}
  %\lean{}
  %\leanok
  \uses{def:fDiv}
  $D_{x \mapsto x}(\mu, \nu) = \mu(\mathcal X)$.
\end{lemma}

\begin{proof} %\leanok
Compute the integral: its value is $(\frac{d\mu}{d\nu}\cdot \nu)(\mathcal X)$. Then
$D_{x\mapsto x}(\mu, \nu) = (\frac{d\mu}{d\nu}\cdot \nu)(\mathcal X) + \mu_{\perp \nu}(\mathcal X) = \mu (\mathcal X)$.
\end{proof}


\begin{lemma}
  \label{lem:fDiv_mul}
  \lean{ProbabilityTheory.fDiv_smul}
  \leanok
  \uses{def:fDiv}
  For all $a \ge 0$, $D_{a f}(\mu, \nu) = a D_{f}(\mu, \nu)$.
\end{lemma}

\begin{proof}\leanok
Linearity of the integral.
\end{proof}


\begin{lemma}
  \label{lem:fDiv_add}
  \lean{ProbabilityTheory.fDiv_add}
  \leanok
  \uses{def:fDiv}
  $D_{f + g}(\mu, \nu) = D_f(\mu, \nu) + D_g(\mu, \nu)$.
\end{lemma}

\begin{proof}\leanok
Linearity of the integral.
\end{proof}


\begin{lemma}
  \label{lem:fDiv_add_linear}
  %\lean{}
  %\leanok
  \uses{def:fDiv}
  For finite measures $\mu$ and $\nu$ with $\mu(\mathcal X) = \nu(\mathcal X)$, for all $a \in \mathbb{R}$, $D_{f + a(x - 1)}(\mu, \nu) = D_{f}(\mu, \nu)$.
\end{lemma}

\begin{proof}%\leanok
\uses{lem:fDiv_add, lem:fDiv_mul, lem:fDiv_const, lem:fDiv_id}
Linearity (Lemmas~\ref{lem:fDiv_add} and~\ref{lem:fDiv_mul}), then Lemma~\ref{lem:fDiv_const} and~\ref{lem:fDiv_id}.
\end{proof}


\begin{lemma}
  \label{lem:fDiv_map_measurableEmbedding}
  \lean{ProbabilityTheory.fDiv_map_measurableEmbedding}
  \leanok
  \uses{def:fDiv}
  Let $\mu$ and $\nu$ be two measures on $\mathcal X$ and let $g : \mathcal X \to \mathcal Y$ be a measurable embedding. Then $D_f(g_* \mu, g_* \nu) = D_f(\mu, \nu)$.
\end{lemma}

\begin{proof}\leanok
\end{proof}


\begin{lemma}
  \label{lem:fDiv_symm}
  \lean{ProbabilityTheory.fDiv_conj}
  \leanok
  \uses{def:fDiv}
  $D_f(\mu, \nu) = D_{x \mapsto xf(1/x)}(\nu, \mu)$~.
\end{lemma}

\begin{proof}\leanok
\uses{}

\end{proof}


\begin{lemma}
  \label{lem:fDiv_add_smul_left}
  \lean{ProbabilityTheory.fDiv_smul_add_smul_left}
  \leanok
  \uses{def:fDiv}
  For $a, b \ge 0$ with $a + b = 1$, $(\mu, \nu) \mapsto D_f(a \mu + b \nu, \nu)$ is an $f$-divergence for the function $x \mapsto f(ax + b)$
\end{lemma}

\begin{proof}\leanok
\uses{}
\end{proof}


\begin{lemma}
  \label{lem:fDiv_add_smul_right}
  \lean{ProbabilityTheory.fDiv_smul_add_smul_right}
  \leanok
  \uses{def:fDiv}
  For $a, b \ge 0$ with $a + b = 1$, $(\mu, \nu) \mapsto D_f(\mu, a \mu + b \nu)$ is an $f$-divergence for the function $x \mapsto (ax+b)f\left(\frac{x}{ax+b}\right)$~.
\end{lemma}

\begin{proof}\leanok
\uses{lem:fDiv_symm, lem:fDiv_add_smul_left}
We use Lemma~\ref{lem:fDiv_symm}, then Lemma~\ref{lem:fDiv_add_smul_left} and finally Lemma~\ref{lem:fDiv_symm} again:
\begin{align*}
D_f(\mu, a \mu + b \nu)
&= D_{xf(1/x)}(a \mu + b \nu, \mu)
\\
&= D_{(bx+a)f(1/(bx+a))}(\nu, \mu)
\\
&= D_{(ax+b)f(x/(ax+b))}(\mu, \nu)
\: .
\end{align*}

\end{proof}


\begin{lemma}
  \label{lem:fDiv_add_smul_right'}
  \lean{ProbabilityTheory.fDiv_smul_add_smul_right'}
  \leanok
  \uses{def:fDiv}
  For $a, b \ge 0$ with $a + b = 1$, $(\mu, \nu) \mapsto D_f(\nu, a \mu + b \nu)$ is an $f$-divergence for the function $x \mapsto (ax+b)f\left(\frac{1}{ax+b}\right)$~.
\end{lemma}

\begin{proof}\leanok
\uses{lem:fDiv_add_smul_right, lem:fDiv_symm}
We use Lemma~\ref{lem:fDiv_add_smul_right}, then Lemma~\ref{lem:fDiv_symm}.
\begin{align*}
D_f(\nu, a \mu + b \nu)
&= D_{(bx+a)f\left(\frac{x}{bx+a}\right)}(\nu, \mu)
= D_{(ax+b)f\left(\frac{1}{ax+b}\right)}(\mu, \nu)
\: .
\end{align*}
\end{proof}



\paragraph{Upper and lower bounds on $D_f$}

\begin{lemma}
  \label{lem:fDiv_absolutelyContinuous_add_mutuallySingular}
  \lean{ProbabilityTheory.fDiv_absolutelyContinuous_add_mutuallySingular}
  \leanok
  \uses{def:fDiv}
  Let $\mu_1, \mu_2$ and $\nu$ be finite measures on $\mathcal X$, with $\mu_1 \ll \nu$ and $\mu_2 \perp \nu$.
  Then $D_f(\mu_1 + \mu_2, \nu) = D_f(\mu_1, \nu) + \mu_2(\mathcal X) f'(\infty)$.
\end{lemma}

\begin{proof}\leanok
$\frac{d(\mu_1 + \mu_2)}{d \nu} = \frac{d \mu_1}{d \nu}$ a.e. and $(\mu_1 + \mu_2)_{\perp \nu} = \mu_2$.
\end{proof}


\begin{lemma}[Superseded by Lemma~\ref{lem:fDiv_add_measure_le}]
  \label{lem:fDiv_add_measure_le_of_ac}
  \lean{ProbabilityTheory.fDiv_add_measure_le_of_ac}
  \leanok
  \uses{def:fDiv}
  Let $\mu_1, \mu_2, \nu$ be three finite measures on $\mathcal X$ with $\mu_1 \ll \nu$ and $\mu_2 \ll \nu$. Then
  $D_f(\mu_1 + \mu_2, \nu) \le D_f(\mu_1, \nu) + \mu_2(\mathcal X) f'(\infty)$.
\end{lemma}

\begin{proof}\leanok
\begin{align*}
D_f(\mu_1 + \mu_2, \nu)
&= \int_x f \left( \frac{d \mu_1}{d\nu}(x) + \frac{d\mu_2}{d\nu}(x) \right) \partial \nu
\\
&\le \int_x f \left( \frac{d \mu_1}{d\nu}(x) \right) + \frac{d\mu_2}{d\nu}(x) f'(\infty) \partial \nu
\\
&= D_f(\mu_1, \nu) + \mu_2(\mathcal X) f'(\infty)
\: .
\end{align*}
\end{proof}


\begin{lemma}
  \label{lem:fDiv_add_measure_le}
  \lean{ProbabilityTheory.fDiv_add_measure_le}
  \leanok
  \uses{def:fDiv}
  Let $\mu_1, \mu_2, \nu$ be three finite measures on $\mathcal X$. Then
  $D_f(\mu_1 + \mu_2, \nu) \le D_f(\mu_1, \nu) + \mu_2(\mathcal X) f'(\infty)$.
\end{lemma}

\begin{proof}\leanok
\uses{lem:fDiv_absolutelyContinuous_add_mutuallySingular, lem:fDiv_add_measure_le_of_ac}
From Lemma~\ref{lem:fDiv_absolutelyContinuous_add_mutuallySingular}, then Lemma~\ref{lem:fDiv_add_measure_le_of_ac},
\begin{align*}
D_f(\mu_1 + \mu_2, \nu)
&= D_f(\frac{d\mu_1}{d \nu}\cdot \nu + \frac{d\mu_2}{d \nu}\cdot \nu, \nu) + (\mu_1)_{\perp\nu}(\mathcal X) f'(\infty) + (\mu_2)_{\perp\nu}(\mathcal X) f'(\infty)
\\
&\le D_f(\frac{d\mu_1}{d \nu}\cdot \nu, \nu) + (\frac{d\mu_2}{d \nu}\cdot \nu)(\mathcal X) f'(\infty) + (\mu_1)_{\perp\nu}(\mathcal X) f'(\infty) + (\mu_2)_{\perp\nu}(\mathcal X) f'(\infty)
\\
&= D_f(\mu_1, \nu) + \mu_2(\mathcal X) f'(\infty)
\: .
\end{align*}
\end{proof}


\begin{lemma}
  \label{lem:fDiv_eq_add_withDensity_derivAtTop}
  \lean{ProbabilityTheory.fDiv_eq_add_withDensity_derivAtTop}
  \leanok
  \uses{def:fDiv}
  Let $\mu$ and $\nu$ be two finite measures on $\mathcal X$.
  Then $D_f(\mu, \nu) = D_f(\frac{d\mu}{d\nu}\cdot \nu, \nu) + f'(\infty) \mu_{\perp \nu}(\mathcal X)$.
\end{lemma}

\begin{proof}\leanok
\uses{lem:fDiv_absolutelyContinuous_add_mutuallySingular}
Apply Lemma~\ref{lem:fDiv_absolutelyContinuous_add_mutuallySingular} to $\frac{d\mu}{d\nu}\cdot \nu$ and $\mu_{\perp \nu}$.
\end{proof}


\begin{lemma}[Superseded by Lemma~\ref{lem:le_fDiv}]
  \label{lem:le_fDiv_of_ac}
  \lean{ProbabilityTheory.le_fDiv_of_ac}
  \leanok
  \uses{def:fDiv}
  Let $\mu, \nu$ be two finite measures on the same space $\mathcal X$, such that $\mu \ll \nu$. Then $\nu(\mathcal X) f(\mu(\mathcal X)/\nu(\mathcal X)) \le D_f(\mu, \nu)$.
\end{lemma}

\begin{proof}\leanok
Since $\mu \ll \nu$ the $f$-divergence is only the integral part. If $\nu = 0$ there is nothing to prove. Otherwise, by Jensen's inequality for the probability measure $\nu / \nu(\mathcal X)$,
\begin{align*}
D_f(\mu, \nu)
&= \nu(\mathcal X) \frac{1}{\nu(\mathcal X)}\nu\left[ f\left( \frac{d\mu}{d\nu} \right) \right]
\ge \nu(\mathcal X) f\left( \frac{1}{\nu(\mathcal X)}\nu\left[\frac{d\mu}{d\nu} \right] \right)
= \nu(\mathcal X) f(\mu(\mathcal X)/\nu(\mathcal X))
\: .
\end{align*}

\end{proof}


\begin{lemma}
  \label{lem:le_fDiv}
  \lean{ProbabilityTheory.le_fDiv}
  \leanok
  \uses{def:fDiv}
  Let $\mu, \nu$ be two finite measures on the same space $\mathcal X$. Then $\nu(\mathcal X) f(\mu(\mathcal X)/\nu(\mathcal X)) \le D_f(\mu, \nu)$. In particular, if $\nu$ is a probability measure, $f(\mu(\mathcal X)) \le D_f(\mu, \nu)$.
\end{lemma}

\begin{proof}\leanok
\uses{lem:fDiv_eq_add_withDensity_derivAtTop, lem:le_fDiv_of_ac}
By convexity, then Lemma~\ref{lem:le_fDiv_of_ac} and finally Lemma~\ref{lem:fDiv_eq_add_withDensity_derivAtTop},
\begin{align*}
\nu(\mathcal X) f(\mu(\mathcal X)/\nu(\mathcal X))
&\le \nu(\mathcal X) f(\frac{d\mu}{d\nu}\cdot \nu (\mathcal X)/\nu(\mathcal X)) + f'(\infty)\mu_{\perp}(\mathcal X)
\\
&\le D_f(\frac{d\mu}{d\nu}\cdot \nu , \nu) + f'(\infty)\mu_{\perp}(\mathcal X)
\\
&= D_f(\mu, \nu)
\: .
\end{align*}
\end{proof}


\begin{lemma}
  \label{lem:fDiv_nonneg}
  %\lean{}
  \leanok
  \uses{def:fDiv}
  Let $\mu, \nu$ be two probability measures. If $f(1) = 0$ then $D_f(\mu, \nu) \ge 0$.
\end{lemma}

\begin{proof}\leanok
\uses{lem:le_fDiv}
Apply Lemma~\ref{lem:le_fDiv} and use $f(\mu(\mathcal X)) = f(1) = 0$.
\end{proof}



TODO: find a way to hide the following dummy lemma
\begin{lemma}[Dummy lemma: fDiv properties]
  \label{lem:fDiv_properties}
  %\lean{}
  \leanok
  \uses{def:fDiv}
  Dummy node to summarize properties of $f$-divergences.
\end{lemma}

\begin{proof}\leanok
\uses{
  lem:fDiv_ne_top_iff,
  lem:fDiv_const,
  lem:fDiv_self,
  lem:fDiv_eq_zero_iff,
  lem:fDiv_id,
  lem:fDiv_mul,
  lem:fDiv_add,
  lem:fDiv_add_linear,
  lem:fDiv_map_measurableEmbedding,
  lem:fDiv_symm,
  lem:fDiv_add_smul_left,
  lem:fDiv_add_smul_right,
  lem:fDiv_add_smul_right',
  lem:fDiv_absolutelyContinuous_add_mutuallySingular,
  lem:fDiv_add_measure_le,
  lem:fDiv_eq_add_withDensity_derivAtTop,
  lem:le_fDiv,
  lem:fDiv_nonneg,
}
\end{proof}


\subsection{Conditional f-divergence}

TODO: replace the definition by $D(\mu \otimes \kappa, \mu \otimes \eta)$, which allows for a more general proof of the ``conditioning increases divergence'' inequality.

\begin{definition}[Conditional f-divergence]
  \label{def:condFDiv}
  \lean{ProbabilityTheory.condFDiv}
  \leanok
  \uses{lem:fDiv_properties}
  Let $f : \mathbb{R} \to \mathbb{R}$, $\mu$ a measure on $\mathcal X$ and $\kappa, \eta : \mathcal X \rightsquigarrow \mathcal Y$ two Markov kernels from $\mathcal X$ to $\mathcal Y$. The conditional f-divergence between $\kappa$ and $\eta$ with respect to $\mu$ is
  \begin{align*}
  D_f(\kappa, \eta \mid \mu) = \mu\left[x \mapsto D_f(\kappa(x), \eta(x))\right]
  \end{align*}
  if $x \mapsto D_f(\kappa(x), \eta(x))$ is $\mu$-integrable and $+\infty$ otherwise.
\end{definition}


\begin{lemma}
  \label{lem:condFDiv_nonneg}
  %\lean{}
  \leanok
  \uses{def:condFDiv}
  Let $\mu$ be a measure on $\mathcal X$ and $\kappa, \eta : \mathcal X \rightsquigarrow \mathcal Y$ two Markov kernels. If $f(1) = 0$ then $D_f(\kappa, \eta \mid \mu) \ge 0$.
\end{lemma}

\begin{proof}\leanok
\uses{lem:fDiv_nonneg}
Apply Lemma~\ref{lem:fDiv_nonneg}.
\end{proof}


\begin{lemma}
  \label{lem:condFDiv_const}
  \lean{ProbabilityTheory.condFDiv_const}
  \leanok
  \uses{def:condFDiv, def:fDiv}
  Let $\mu, \nu$ be measures on $\mathcal X$, where $\mu$ is finite, and let $\xi$ be a finite measure on $\mathcal Y$.
  Then $D_f(x \mapsto \mu, x \mapsto \nu \mid \xi) = D_f(\mu, \nu) \xi (\mathcal X)$.
\end{lemma}

\begin{proof}\leanok
\uses{}
$$D_f(x \mapsto \mu, x \mapsto \nu \mid \xi) 
= \xi\left[x \mapsto D_f(\mu, \nu)\right] 
= D_f(\mu, \nu) \xi\left[1\right] 
= D_f(\mu, \nu) \xi (\mathcal X)$$
\end{proof}


% TODO : this statement is not the same as the lean code, in particular here there is only one measure mentioned, while in the lean we have 2, also the rest of the statement is not the same
\begin{lemma}
  \label{lem:integrable_fDiv_compProd_iff}
  \lean{ProbabilityTheory.integrable_f_rnDeriv_compProd_iff}
  \leanok
  \uses{def:fDiv}
  Let $\mu$ be a finite measure on $\mathcal X$ and $\kappa, \eta : \mathcal X \rightsquigarrow \mathcal Y$ be two finite kernels from $\mathcal X$ to $\mathcal Y$.
  Then $p \mapsto f \left(\frac{d(\mu \otimes \kappa)}{d(\mu \otimes \eta)}(p)\right)$ is $(\mu \otimes \eta)$-integrable iff
  \begin{itemize}
    \item $x \mapsto D_f(\kappa(x), \eta(x))$ is $\mu$-integrable and
    \item for $\mu$-almost all $x$, $y \mapsto f \left( \frac{d\kappa(x)}{d\eta(x)}(y) \right)$ is $\eta(x)$-integrable. 
  \end{itemize}
\end{lemma}

\begin{proof}
Since $x \mapsto f \left(\frac{d(\mu \otimes \kappa)}{d(\mu \otimes \eta)} x \right)$ is measurable, the integrability condition w.r.t. $\mu \otimes \eta$ is equivalent to
\begin{itemize}
    \item $x \mapsto \int_y \left\Vert f \left( \frac{d\kappa(x)}{d\eta(x)}(y) \right) \right\Vert \partial \eta(x)$ is $\mu$-integrable and
    \item for $\mu$-almost all $x$, $y \mapsto f \left( \frac{d\kappa(x)}{d\eta(x)}(y) \right)$ is $\eta(x)$-integrable. 
  \end{itemize}
It suffices to use convexity to show that integrability of $x \mapsto \int_y f \left( \frac{d\kappa(x)}{d\eta(x)}(y) \right) \partial \eta(x)$ implies integrability of $x \mapsto \int_y \left\Vert f \left( \frac{d\kappa(x)}{d\eta(x)}(y) \right) \right\Vert \partial \eta(x)$.

TODO
\end{proof}


\begin{lemma}
  \label{lem:condFDiv_ne_top_iff}
  \lean{ProbabilityTheory.condFDiv_ne_top_iff}
  \leanok
  \uses{def:fDiv, def:condFDiv}
  Let $\mu$ be a finite measure on $\mathcal X$ and let $\kappa, \eta : \mathcal X \rightsquigarrow \mathcal Y$ be two finite kernels, where $\kappa$ is a Markov kernel.
  Then $D_f(\kappa, \eta \mid \mu) \ne \infty$ if and only if
  \begin{itemize}
    \item for $\mu$-almost all $x$, $y \mapsto f \left( \frac{d\kappa(x)}{d\eta(x)}(y) \right)$ is $\eta(x)$-integrable,
    \item $x \mapsto \int_y f \left( \frac{d\kappa(x)}{d\eta(x)}(y) \right) \partial \eta(x)$ is $\mu$-integrable,
    \item either $f'(\infty) < \infty$ or for $\mu$-almost all $x$, $\kappa(x) \ll \eta(x)$.
  \end{itemize}
\end{lemma}

\begin{proof} \leanok
\uses{lem:integrable_fDiv_compProd_iff}
\end{proof}


\begin{lemma}
  \label{lem:fDiv_compProd_ne_top_iff}
  \lean{ProbabilityTheory.condFDiv_ne_top_iff_fDiv_compProd_ne_top}
  \leanok
  \uses{def:fDiv, def:condFDiv}
  Let $\mu$ be a finite measure on $\mathcal X$ and let $\kappa, \eta : \mathcal X \rightsquigarrow \mathcal Y$ be two finite kernels.
  Then $D_f(\mu \otimes \kappa, \mu \otimes \eta) \ne \infty \iff D_f(\kappa, \eta \mid \mu) \ne \infty$.
\end{lemma}

\begin{proof} \leanok
\uses{lem:fDiv_compProd_right}
Immediate from Lemma~\ref{lem:fDiv_compProd_right}.
\end{proof}


\begin{lemma}
  \label{lem:fDiv_compProd_right}
  \lean{ProbabilityTheory.fDiv_compProd_right}
  \leanok
  \uses{def:fDiv, def:condFDiv}
  Let $\mu$ be a finite measure on $\mathcal X$ and let $\kappa, \eta : \mathcal X \rightsquigarrow \mathcal Y$ be two finite kernels.
  Then $D_f(\mu \otimes \kappa, \mu \otimes \eta) = D_f(\kappa, \eta \mid \mu)$.
\end{lemma}

\begin{proof} \leanok
\uses{cor:rnDeriv_compProd_right, cor:rnDeriv_value}
The singular part of $\mu \otimes \kappa$ with respect to $\mu \otimes \eta$ is $\mu \otimes \kappa_{\perp \eta}$, hence $(\mu \otimes \kappa)_{\perp (\mu \otimes \eta)}(\mathcal X \times \mathcal Y) = \mu\left[x \mapsto \kappa(x)_{\perp \eta(x)}(\mathcal Y)\right]$.
For the integral part, by Lemma~\ref{cor:rnDeriv_compProd_right} and Corollary~\ref{cor:rnDeriv_value},
\begin{align*}
\int_{p} f\left(\frac{d (\mu \otimes \kappa)}{d (\mu \otimes \eta)}(p)\right) \partial(\mu \otimes \eta)
&= \int_{p} f\left(\frac{d \kappa}{d \eta}(p)\right) \partial(\mu \otimes \eta)
\\
&= \int_x \int_y f\left(\frac{d \kappa}{d \eta}(x, y)\right) \partial \eta(x) \partial \mu
\\
&= \int_x \int_y f\left(\frac{d \kappa(x)}{d \eta(x)}(y)\right) \partial \eta(x) \partial \mu
\: .
\end{align*}
Adding $f'(\infty)$ times the singular parts gives $D_f(\mu \otimes \kappa, \mu \otimes \eta) = \mu\left[D_f(\kappa(x), \eta(x))\right] = D_f(\kappa, \eta \mid \mu)$.
\end{proof}




TODO: find a way to hide the following dummy lemma
\begin{lemma}[Dummy lemma: condFDiv properties]
  \label{lem:condFDiv_properties}
  %\lean{}
  \leanok
  \uses{def:condFDiv}
  Dummy node to summarize properties of conditional $f$-divergences.
\end{lemma}

\begin{proof}\leanok
\uses{
  lem:condFDiv_nonneg,
  lem:condFDiv_const,
  lem:condFDiv_ne_top_iff,
  lem:fDiv_compProd_ne_top_iff,
  lem:fDiv_compProd_right
}
\end{proof}









\subsection{Integral representation of $f$-divergences}

We prove that any $f$-divergence is an integral over priors of statistical informations, in which the measure depends on $f$.
The main references for this part are \cite{liese2006divergences} and \cite{liese2012phi}.


\begin{definition}[Curvature measure]
  \label{def:curvatureMeasure}
  \lean{ProbabilityTheory.DivFunction.curvatureMeasure}
  \leanok
  \uses{}
  Let $f: \mathbb{R} \to \mathbb{R}$ be a convex function. Then its right derivative $f'_+(x) \coloneqq \lim_{y \downarrow x}\frac{f(y) - f(x)}{y - x}$ is a Stieltjes function (a monotone right continuous function) and it defines a measure $\gamma_f$ on $\mathbb{R}$ by $\gamma_f((x,y]) \coloneqq f'_+(y) - f'_+(x)$~. \cite{liese2012phi} calls $\gamma_f$ the curvature measure of $f$.
\end{definition}


\begin{lemma}
  \label{lem:curvatureMeasure_mul}
  \lean{ProbabilityTheory.DivFunction.curvatureMeasure_smul}
  \leanok
  \uses{def:curvatureMeasure}
  For $a > 0$ and $f$ a convex function, the curvature measure of $af$ is $\gamma_{af} = a \gamma_f$~.
\end{lemma}

\begin{proof}\leanok
\uses{}
The right derivative of $af$ is $a f'_+$ (with values in $[-\infty, +\infty]$), hence $\gamma_{af}((x,y]) = a f'_+(y) - a f'_+(x) = a \gamma_f((x,y])$. The equality of the measures is obtained at the level of the outer measures, which are built from these interval lengths.

For $f: \mathbb{R} \to \mathbb{R}$, the case $a = 0$ is trivial. It fails in the Lean formalization, where the right derivative is $-\infty$ to the left of the domain of $f$: $0 \cdot f$ is the zero function on $[0, +\infty)$, whose right derivative jumps from $-\infty$ to $0$ at $0$, so that $\gamma_{0 \cdot f}$ has an infinite atom at $0$.
\end{proof}


\begin{lemma}
  \label{lem:curvatureMeasure_add}
  \lean{ProbabilityTheory.DivFunction.curvatureMeasure_add}
  \leanok
  \uses{def:curvatureMeasure}
  Let $f,g$ be two convex functions whose effective domains have the same endpoints $a \le b$, and such that the right derivatives $f'_+$ and $g'_+$ either both tend to $-\infty$ at $a$ or both have a finite limit at $a$. Then the curvature measure of $f+g$ is $\gamma_{f+g} = \gamma_f + \gamma_g$~.
  In particular, this holds for $f,g: \mathbb{R} \to \mathbb{R}$ two convex functions.
\end{lemma}

\begin{proof}\leanok
\uses{}
The right derivative of $f+g$ is $f'_+ + g'_+$, seen as functions with values in $[-\infty, +\infty]$ ($-\infty$ below the domain, $+\infty$ above it). By the assumptions, $f'_+$ and $g'_+$ are equal to $-\infty$ at the same points and equal to $+\infty$ at the same points. On the region where both are finite, $\gamma_{f+g}((x,y]) = \gamma_f((x,y]) + \gamma_g((x,y])$. The region where $f'_+$ and $g'_+$ are $-\infty$, and the region above the first point where they are $+\infty$, are null for the three measures. At $a$ and $b$, the three measures have the same atoms (equal to $0$ or $+\infty$).

The assumption at $a$ is needed: if $f'_+(x) \to -\infty$ as $x \downarrow a$ but $g'_+$ has a finite limit at $a$, then $\gamma_g$ has an infinite atom at $a$ but $\gamma_{f+g}$ does not.
\end{proof}


\begin{theorem}
  \label{thm:integration_by_parts}
  %\lean{}
  %\leanok
  \notready
  \uses{}
  If $f$ and $g$ are two Stieltjes functions with associated measures $\mu_f$ and $\mu_g$ and $f$ is continuous on $[a, b]$, then
  \begin{align*}
  \int_x f(x) \partial\mu_g = f(b)g(b) - f(a)g(a) - \int_x g(x) \partial\mu_f \: .
  \end{align*}
\end{theorem}

\begin{proof} \notready
\uses{}

\end{proof}


\begin{lemma}
  \label{lem:convex_taylor}
  \lean{ProbabilityTheory.DivFunction.convex_taylor_one_right', ProbabilityTheory.DivFunction.convex_taylor_one_left'}
  \leanok
  \uses{def:curvatureMeasure}
  (In Lean, this is proved for a \texttt{DivFunction} and $x = 1$, by the fundamental theorem of calculus and Tonelli's theorem, without integration by parts.)
  For $f: \mathbb{R} \to \mathbb{R}$ a convex function and $x,y \in \mathbb{R}$,
  \begin{align*}
  f(y) - f(x) - (y - x)f'_+(x) &= \int_{z \in (x,y]} (y - z) \partial \gamma_f & \text{ if } x \le y \: ,
  \\
  f(y) - f(x) - (y - x)f'_+(x) &= \int_{z \in (y,x]} (z - y) \partial \gamma_f & \text{ if } y \le x \: .
  \end{align*}
\end{lemma}

\begin{proof}\leanok
\uses{thm:integration_by_parts}
Let $\Lambda$ be the Lebesgue measure and let $x < y$. Since $f$ has right derivative $f'_+$ in $(x,y)$ and $f'_+$ is integrable on that interval,
\begin{align*}
f(y) - f(x) = \int_{z \in (x, y]} f'_+(z) \partial \Lambda
\: .
\end{align*}
We now integrate by parts, using Theorem~\ref{thm:integration_by_parts}. $\Lambda$ is the measure obtained from the Stieltjes function $z \mapsto z - y$.
\begin{align*}
\int_{z \in (x, y]} f'_+(z) \partial \Lambda
&= - \int_{z \in (x,y]} (z - y)\partial \gamma_f + f'_+(y)(y - y) - f'_+(x)(x - y)
\\
&= \int_{z \in (x,y]} (y - z)\partial \gamma_f + f'_+(x)(y - x)
\: .
\end{align*}
Putting the two equalities together, we get the result for $x < y$. The proof for $x > y$ is similar, and for $x = y$ both sides of both equations are zero.
\end{proof}


\begin{definition}
  \label{def:statInfoFun}
  \lean{ProbabilityTheory.statInfoFun}
  \leanok
  \uses{}
  For $a,b \in (0, +\infty)$ let $\phi_{a,b} : \mathbb{R} \to \mathbb{R}$ be the function defined by
  \begin{align*}
  \phi_{a,b}(x) &= \max\left\{0, a x - b \right\} & \text{ for } a \le b \: ,
  \\
  \phi_{a,b}(x) &= \max\left\{0, b - a x \right\} & \text{ for } a > b \: .
  \end{align*}
\end{definition}


\begin{corollary}
  \label{cor:convex_taylor_statInfoFun}
  %\lean{}
  %\leanok
  \uses{def:curvatureMeasure, def:statInfoFun}
  For $f: \mathbb{R} \to \mathbb{R}$ a convex function, for all $x \in \mathbb{R}$~,
  \begin{align*}
  f(x) &= f(1) + f'_+(1) (x - 1) + \int_{y} \phi_{1,y}(x) \partial\gamma_f \: .
  \end{align*}
\end{corollary}

\begin{proof}%\leanok
\uses{lem:convex_taylor}
By Lemma~\ref{lem:convex_taylor}, for $x > 1$,
\begin{align*}
f(x) - f(1) - f'_+(1) (x - 1)
&= \int_{y \in (1, x]} (x - y) \partial\gamma_f
\\
&= \int_{y \in (1, +\infty)} \max\{0, x - y\} \partial\gamma_f
\\
&= \int_{y \in (1, +\infty)} \phi_{1,y}(x) \partial\gamma_f
\: .
\end{align*}
That final integral is equal to $\int_y \phi_{1,y}(x) \partial\gamma_f$ (without the restriction to $(1, +\infty)$) because $\phi_{1,y}(x)$ is zero outside of $(1, +\infty)$ for $x \ge 1$.

The case of $x < 1$ is similar.
\end{proof}


\begin{lemma}
  \label{lem:curvatureMeasure_statInfoFun}
  %\lean{}
  %\leanok
  \uses{def:curvatureMeasure, def:statInfoFun}
  The curvature measure of the function $\phi_{a,b}$ is $\gamma_{\phi_{a,b}} = a\delta_{b/a}$~, where $\delta_x$ is the Dirac measure at $x$.
\end{lemma}

\begin{proof}%\leanok
\uses{}

\end{proof}

The following theorem is inspired from \cite{liese2006divergences,liese2012phi} but gives a different integral representation.

\begin{theorem}
  \label{thm:fDiv_eq_integral_eGamma}
  \lean{ProbabilityTheory.fDiv_eq_lintegral_fDiv_statInfoFun}
  \leanok
  \uses{def:fDiv, def:curvatureMeasure, def:eGamma}
  For two finite measures $\mu, \nu \in \mathcal M(\mathcal X)$,
  \begin{align*}
  D_f(\mu, \nu) = f(1) \nu(\mathcal X) + f'_+(1)(\mu(\mathcal X) - \nu(\mathcal X)) + \int_y D_{\phi_{1,y}}(\mu, \nu) \partial\gamma_f \: .
  \end{align*}
\end{theorem}

For probability measures, $D_{\phi_{1,y}}(\mu, \nu) = E_y(\mu, \nu)$~, but these functions can differ if $\mu$ and $\nu$ are finite measures with different masses.

\begin{proof}\leanok
\uses{cor:convex_taylor_statInfoFun}
We prove the result for $f(1) = 0$ and $f'_+(1) = 0$ for simplicity of exposition.

From Corollary~\ref{cor:convex_taylor_statInfoFun} we get that $f(x) = \int_y \phi_{1,y}(x) \partial\gamma_f$~.
From this and the theorem of Fubini,
\begin{align*}
\int_x f\left( \frac{d \mu}{d \nu}(x) \right) \partial\nu
= \int_{x} \int_{y} \phi_{1,y} \left(\frac{d \mu}{d \nu}(x) \right) \partial\gamma_f \partial\nu
= \int_{y} \int_{x} \phi_{1,y} \left(\frac{d \mu}{d \nu}(x) \right) \partial\nu \partial\gamma_f
\: .
\end{align*}
For the mutually singular part of the divergence $D_f$,
\begin{align*}
f'(\infty)
= f'_+(\infty) - f'_+(1) 
= \gamma_f((1,+\infty))
= \int_{y > 1} 1 \partial \gamma_f
= \int_{y > 1} \phi_{1,y}'(\infty) \partial \gamma_f
\: .
\end{align*}
The last integral is equal to $\int_{y \in \mathbb{R}} \phi_{1,y}'(\infty) \partial \gamma_f$ since $\phi_{1,y}'(\infty) = 0$ for $y \le 1$.
Finally,
\begin{align*}
D_f(\mu, \nu)
&= \int_x f\left( \frac{d \mu}{d \nu}(x) \right) \partial\nu + f'(\infty) \mu_{\perp \nu}(\mathcal X)
\\
&= \int_{y} \left(\int_{x} \phi_{1,y} \left(\frac{d \mu}{d \nu}(x) \right) \partial\nu + \phi_{1,y}'(\infty) \mu(\mathcal X) \right) \partial\gamma_f
\\
&= \int_y D_{\phi_{1,y}}(\mu, \nu) \partial \gamma_f
\: .
\end{align*}
\end{proof}

TODO: define $\Gamma_f$.

TODO: in the lemma below, $I_\pi(\mu, \nu)$ is the correct value for probability measures, but in general it should be replaced by $D_{\phi_{\pi, 1 - \pi}}(\mu, \nu)$~.
These functions can differ if $\mu$ and $\nu$ are finite measures with different masses.

\begin{theorem}[\cite{liese2006divergences,liese2012phi}]
  \label{thm:fDiv_eq_integral}
  %\lean{}
  %\leanok
  \uses{def:fDiv, def:curvatureMeasure, def:deGrootInfo}
  \begin{align*}
  D_f(\mu, \nu) = f(1) \nu(\mathcal X) + f'_+(1)(\mu(\mathcal X) - \nu(\mathcal X)) + \int_\pi I_\pi(\mu, \nu) \partial\Gamma_f \: .
  \end{align*}
  
\end{theorem}

\begin{proof}%\leanok
\uses{thm:fDiv_eq_integral_eGamma}

\end{proof}




\subsection{Data-processing inequality}

This section contains three proofs of the data-processing inequality for $f$-divergences.
That inequality states that for $\mu, \nu$ two measures on $\mathcal X$ and $\kappa : \mathcal X \rightsquigarrow \mathcal Y$ a Markov kernel, $D_f(\kappa \circ \mu, \kappa \circ \nu) \le D_f(\mu, \nu)$.
\begin{enumerate}
  \item A first proof uses the convexity of $f$ through Jensen's inequality for conditional expectations. We first prove the DPI for measurable functions (or deterministic kernels) then use it to extend to the general case.
  \item A second proof deals with the general case and uses Radon-Nikodym derivatives of kernels. This restrict its applicability to measurable spaces that satisfy some particular countability assumptions.
  \item The last proof uses the representation of $f$-divergence as an integral over statistical (risk-based) divergences, then uses the data-processing inequality for those divergences.
\end{enumerate}

\begin{leanremark}
Currently we don't have Lean proofs of Jensen's inequality for conditional expectations (used in proof 1) or of an integration by part lemma that is used in proof 3.
Proof 2 gives the weakest results since it is restricted to spaces that satisfy countability assumptions, but it's the only one for which our Lean proof is \texttt{sorry}-free.
\end{leanremark}

We first remark that it is sufficient to prove the DPI under an absolute continuity assumption on the measures. 

\begin{lemma}
  \label{lem:DPI_of_DPI_ac}
  \lean{ProbabilityTheory.fDiv_comp_le_of_comp_le_of_ac}
  \leanok
  \uses{def:fDiv}
  Let $\kappa : \mathcal X \rightsquigarrow \mathcal Y$ be a Markov kernel and $\nu \in \mathcal M(\mathcal X)$ be a finite measure.
  Suppose that for all finite measures $\mu \in \mathcal M(\mathcal X)$ with $\mu \ll \nu$, $D_f(\kappa \circ \mu, \kappa \circ \nu) \le D_f(\mu, \nu)$.
  Then the same is true without the absolute continuity hypothesis.
\end{lemma}

\begin{proof}\leanok
\uses{lem:fDiv_add_measure_le}
We use $\kappa \circ \mu = \kappa \circ (\frac{d\mu}{d\nu}\cdot \nu) + \kappa \circ \mu_{\perp \nu}$~.
By Lemma~\ref{lem:fDiv_add_measure_le},
\begin{align*}
D_f(\kappa \circ \mu, \kappa \circ \nu)
\le D_f\left(\kappa \circ (\frac{d\mu}{d\nu}\cdot \nu), \kappa \circ \nu\right)
  + (\kappa \circ \mu_{\perp \nu})(\mathcal Y) f'(\infty)
\: .
\end{align*}
Since $\kappa$ is a Markov kernel, $(\kappa \circ \mu_{\perp \nu})(\mathcal Y) = \mu_{\perp \nu}(\mathcal Y)$.

By hypothesis, since $\frac{d\mu}{d\nu}\cdot \nu \ll \nu$, $D_f\left(\kappa \circ (\frac{d\mu}{d\nu}\cdot \nu), \kappa \circ \nu\right) \le D_f\left(\frac{d\mu}{d\nu}\cdot \nu, \nu\right)$.
We get
\begin{align*}
D_f(\kappa \circ \mu, \kappa \circ \nu)
&\le D_f\left(\frac{d\mu}{d\nu}\cdot \nu,\nu\right)
  + \mu_{\perp \nu}(\mathcal Y) f'(\infty)
\\
&= D_f(\mu, \nu)
\: .
\end{align*}
\end{proof}


The next Lemma implies that in order to prove the DPI, it is sufficient to prove $D_f(\kappa \circ \mu, \kappa \circ \nu) \le D_f(\mu \otimes \kappa, \nu \otimes \kappa)$.


\begin{lemma}[Composition-product with a kernel]
  \label{thm:fDiv_compProd_left_1}
  \lean{ProbabilityTheory.fDiv_compProd_left}
  \leanok
  \uses{def:fDiv}
  Let $\mu, \nu$ be two measures on $\mathcal X$ and let $\kappa : \mathcal X \rightsquigarrow \mathcal Y$ be a Markov kernel.
  Then $D_f(\mu \otimes \kappa, \nu \otimes \kappa) = D_f(\mu, \nu)$.
\end{lemma}

\begin{proof}\leanok
\uses{cor:rnDeriv_compProd_left}
By Lemma~\ref{cor:rnDeriv_compProd_left} and the fact that $(\mu \otimes \kappa)_{\perp \nu \otimes \kappa} = \mu_{\perp \nu} \otimes \kappa$ (Lemma TODO),
\begin{align*}
D_f(\mu \otimes \kappa, \nu \otimes \kappa)
&= \int_{p} f \left(\frac{d (\mu \otimes \kappa)}{d (\nu \otimes \kappa)}(p)\right) \partial(\nu \otimes \kappa)
  + f'(\infty) (\mu \otimes \kappa)_{\perp \nu \otimes \kappa}(\mathcal X)
\\
&= \int_{p} f \left(\frac{d \mu}{d \nu}(p_X)\right) \partial(\nu \otimes \kappa)
  + f'(\infty) (\mu_{\perp \nu} \otimes \kappa)(\mathcal X)
\\
&= \int_x \int_y f \left(\frac{d \mu}{d \nu}(x)\right) \partial \kappa(x) \partial \nu
  + f'(\infty) \mu_{\perp \nu}(\mathcal X)
\\
&= \int_x f \left(\frac{d \mu}{d \nu}(x)\right) \partial \nu
  + f'(\infty) \mu_{\perp \nu}(\mathcal X)
\\
&= D_f(\mu, \nu)
\: .
\end{align*}
\end{proof}


\begin{corollary}
  \label{cor:fDiv_prod_left_1}
  \lean{ProbabilityTheory.fDiv_prod_left}
  \leanok
  \uses{def:fDiv}
  Let $\mu, \nu$ be two finite measures on $\mathcal X$ and let $\xi$ be a probability measure on $\mathcal Y$.
  Then $D_f(\mu \times \xi, \nu \times \xi) = D_f(\mu, \nu)$.
\end{corollary}

\begin{proof}\leanok
\uses{thm:fDiv_compProd_left_1}
Apply Lemma~\ref{thm:fDiv_compProd_left_1} with $\kappa$ the constant kernel with value $\xi$.
\end{proof}


\subsubsection{DPI: proof using Jensen's inequality for the conditional expectation}

This proof lifts the DPI on deterministic kernels to a DPI on all kernels.
We first prove the DPI for the special cases of a composition with a deterministic kernel, which is the same as taking the map of the measures by a function.

\begin{theorem}
  \label{thm:fDiv_map_le}
  \lean{ProbabilityTheory.fDiv_map_le}
  \leanok
  \uses{def:fDiv}
  Let $\mu, \nu \in \mathcal M (\mathcal X)$ be two finite measures and let $g : \mathcal X \to \mathcal Y$ be a measurable function. Then
  $D_f(g_* \mu, g_* \nu) \le D_f(\mu, \nu)$.
\end{theorem}

\begin{proof}\leanok
\uses{lem:DPI_of_DPI_ac,lem:rnDeriv_map_eq_condexp,thm:condexp_jensen}
By Lemma~\ref{lem:DPI_of_DPI_ac}, it is sufficient to prove the inequality with the additional hypothesis $\mu \ll \nu$ (hence also $g_* \mu \ll g_*\nu$). In that case,
\begin{align*}
D_f(g_* \mu, g_* \nu)
&= \int_y f \left( \frac{d g_* \mu}{d g_* \nu}(y) \right) \partial(g_* \nu)
\\
&= \int_{x} f \left( \frac{d g_* \mu}{d g_* \nu}(g(x)) \right) \partial\nu
\: .
\end{align*}
By Lemma~\ref{lem:rnDeriv_map_eq_condexp}, $\nu$-almost everywhere,
\begin{align*}
\frac{d g_* \mu}{d g_* \nu}(g(x))
&= \nu\left[\frac{d \mu}{d \nu} \mid g^* \mathcal Y \right](x)
\: .
\end{align*}
Using that expression, then Jensen's inequality for the conditional expectation(Theorem~\ref{thm:condexp_jensen}), we get
\begin{align*}
D_f(g_* \mu, g_* \nu)
&= \int_{x} f \left( \nu\left[\frac{d \mu}{d \nu} \mid g^* \mathcal Y \right](x) \right) \partial\nu
\\
&\le \int_{x} \nu\left[ f \left( \frac{d \mu}{d \nu} \right) \mid g^* \mathcal Y \right](x) \partial\nu
\: .
\end{align*}
The integral of a conditional expectation is equal to the integral of the function, hence
\begin{align*}
D_f(\kappa \circ \mu, \kappa \circ \nu)
&\le \int_{x} f \left( \frac{d \mu}{d \nu}(x) \right) \partial\nu
\: .
\end{align*}
This is equal to $D_f(\mu, \nu)$.
\end{proof}


\begin{theorem}[Data-processing]
  \label{thm:fDiv_data_proc_3}
  \lean{ProbabilityTheory.fDiv_comp_right_le''}
  \leanok
  \uses{def:fDiv}
  Let $\mu, \nu$ be two finite measures on $\mathcal X$ and let $\kappa : \mathcal X \rightsquigarrow \mathcal Y$ be a Markov kernel.
  Then $D_f(\kappa \circ \mu, \kappa \circ \nu) \le D_f(\mu, \nu)$.
\end{theorem}

\begin{proof}\leanok
\uses{thm:fDiv_compProd_left_1, thm:fDiv_map_le}
Let $\pi_Y : \mathcal X \times \mathcal Y \to \mathcal Y$ be the projection $\pi_Y((x,y)) = y$.
Then $\kappa \circ \mu = \pi_{Y*}(\mu \otimes \kappa)$ and similarly for $\nu$.
We can use the DPI for that function (Theorem~\ref{thm:fDiv_map_le}), then Lemma~\ref{thm:fDiv_compProd_left_1},
\begin{align*}
D_f(\kappa \circ \mu, \kappa \circ \nu)
&= D_f(\pi_{Y*}(\mu \otimes \kappa), \pi_{Y*}(\nu \otimes \kappa))
\\
&\le D_f(\mu \otimes \kappa, \nu \otimes \kappa)
\\
&= D_f(\mu, \nu)
\: .
\end{align*}
\end{proof}



\subsubsection{DPI: proof using derivatives of kernels}

This proof does not use Jensen's inequality for the conditional expectation, but only for the ordinary expectation. 
It starts by establishing $D_f(\mu, \nu) \le D_f(\mu \otimes \kappa, \nu \otimes \eta)$ for Markov kernels $\kappa$ and $\eta$, then deduces the DPI.
Indeed, if the measurable spaces are standard Borel, we can write $\mu \otimes \kappa = ((\kappa \circ \mu) \otimes \kappa^\dagger_\mu)_\leftrightarrow$ for a Markov kernel $\kappa^\dagger_\mu$ (see Definition~\ref{def:bayesInv}; $(...)_\leftrightarrow$ is notation for composition with the map that swaps the coordinates) and then use the inequality to get
\begin{align*}
D(\kappa \circ \mu, \kappa \circ \nu)
&\le D((\kappa \circ \mu) \otimes \kappa^\dagger_\mu, (\kappa \circ \nu) \otimes \kappa^\dagger_\nu)
\\
&= D(((\kappa \circ \mu) \otimes \kappa^\dagger_\mu)_{\leftrightarrow}, ((\kappa \circ \nu) \otimes \kappa^\dagger_\nu)_{\leftrightarrow})
\\
&= D(\mu \otimes \kappa, \nu \otimes \kappa)
\\
&= D(\mu, \nu)
\: .
\end{align*}
The main drawback of this approach is that it requires that the Bayesian inverse exists, which is not always the case, unless for example the measurable spaces are standard Borel.

The inequality $D_f(\mu, \nu) \le D_f(\mu \otimes \kappa, \nu \otimes \eta)$ is obtained by describing the Radon-Nikodym derivative $\frac{d(\mu \otimes \kappa)}{d(\nu \otimes \eta)}$ and using convexity properties of the function $f$.

\begin{theorem}
  \label{thm:fDiv_le_compProd_1}
  \lean{ProbabilityTheory.le_fDiv_compProd}
  \leanok
  \uses{def:fDiv}
  Let $\mu, \nu$ be two finite measures on $\mathcal X$ and let $\kappa, \eta : \mathcal X \rightsquigarrow \mathcal Y$ be two Markov kernels, with either $\mathcal X$ countable or $\mathcal{Y}$ countably generated.
  Then $D_f(\mu, \nu) \le D_f(\mu \otimes \kappa, \nu \otimes \eta)$.
\end{theorem}

\begin{proof}\leanok
\uses{lem:rnDeriv_compProd, cor:rnDeriv_value}
TODO: the following proof assumes $\mu \ll \nu$ and $\kappa(x) \ll \eta(x)$ for $\nu$-almost all $x$. Describe the general case.

Using Lemma~\ref{lem:rnDeriv_compProd} and \ref{cor:rnDeriv_value},
\begin{align*}
D(\mu \otimes \kappa, \nu \otimes \eta)
&= \int_x \int_y f \left( \frac{\partial \mu}{\partial\nu}(x) \frac{\partial \kappa(x)}{\partial\eta(x)}(y) \right) \partial \eta(x) \partial \nu
\: .
\end{align*}
Since $f$ is convex, by Jensen's inequality,
\begin{align*}
\int_y f \left( \frac{\partial \mu}{\partial\nu}(x) \frac{\partial \kappa(x)}{\partial\eta(x)}(y) \right) \partial \eta(x)
&\ge f \left( \int_y \frac{\partial \mu}{\partial\nu}(x) \frac{\partial \kappa(x)}{\partial\eta(x)}(y) \partial \eta(x) \right)
\\
&= f \left( \frac{\partial \mu}{\partial\nu}(x) \int_y \frac{\partial \kappa(x)}{\partial\eta(x)}(y) \partial \eta(x) \right)
\: .
\end{align*}

Since $\kappa(x) \ll \eta(x)$ for $\nu$-almost all $x$, $\int_y \frac{\partial \kappa(x)}{\partial\eta(x)}(y) \partial \eta(x) = \int_y 1 \partial \kappa(x) = 1$ a.e.. We have obtained
\begin{align*}
D_f(\mu \otimes \kappa, \nu \otimes \eta)
\ge \int_x f \left( \frac{\partial \mu}{\partial\nu}(x)\right) \partial \nu
= D_f(\mu, \nu)
\: .
\end{align*}
\end{proof}


\begin{theorem}[Marginals]
  \label{thm:fDiv_fst_le_1}
  \lean{ProbabilityTheory.fDiv_fst_le}
  \leanok
  \uses{def:fDiv}
  Let $\mu$ and $\nu$ be two measures on $\mathcal X \times \mathcal Y$ where $\mathcal Y$ is standard Borel, and let $\mu_X, \nu_X$ be their marginals on $\mathcal X$.
  Then $D_f(\mu_X, \nu_X) \le D_f(\mu, \nu)$.
  Similarly, for $\mathcal X$ standard Borel and $\mu_Y, \nu_Y$ the marginals on $\mathcal Y$, $D_f(\mu_Y, \nu_Y) \le D_f(\mu, \nu)$.
\end{theorem}

\begin{proof}\leanok
\uses{thm:fDiv_le_compProd_1, lem:fDiv_map_measurableEmbedding}
We introduce conditional kernels and write $D(\mu, \nu) = D(\mu_X \otimes \mu_{Y|X}, \nu_X \otimes \nu_{Y|X})$. Then apply Theorem~\ref{thm:fDiv_le_compProd_1}.
For the second statement, we need Lemma~\ref{lem:fDiv_map_measurableEmbedding} to swap the role of the two coordinates.
\end{proof}


\begin{lemma}
  \label{lem:fDiv_comp_le_compProd_1}
  \lean{ProbabilityTheory.fDiv_comp_le_compProd}
  \leanok
  \uses{def:fDiv}
  Let $\mu, \nu$ be two finite measures on a standard Borel space $\mathcal X$ and let $\kappa, \eta : \mathcal X \rightsquigarrow \mathcal Y$ be two finite kernels.
  $D_f(\kappa \circ \mu, \eta \circ \nu) \le D_f(\mu \otimes \kappa, \nu \otimes \eta)$
\end{lemma}

\begin{proof}\leanok
\uses{thm:fDiv_fst_le_1}
By definition, $\kappa \circ \mu$ is the marginal of $\mu \otimes \kappa$ (a measure on $\mathcal X \times \mathcal Y$) on $\mathcal Y$, and similarly for the other measure. Hence by Theorem~\ref{thm:fDiv_fst_le_1}, $D_f(\kappa \circ \mu, \eta \circ \nu) \le D_f(\mu \otimes \kappa, \nu \otimes \eta)$. 
\end{proof}


\begin{theorem}[Conditioning increases f-divergence]
  \label{thm:fDiv_comp_le_condFDiv_1}
  \lean{ProbabilityTheory.fDiv_comp_left_le}
  \leanok
  \uses{def:fDiv, def:condFDiv}
  Let $\mu$ be a finite measure on $\mathcal X$ and let $\kappa, \eta : \mathcal X \rightsquigarrow \mathcal Y$ be two finite kernels.
  Then $D_f(\kappa \circ \mu, \eta \circ \mu) \le D_f(\mu \otimes \kappa, \mu \otimes \eta) = D_f(\kappa, \eta \mid \mu)$
\end{theorem}

\begin{proof}\leanok
\uses{lem:fDiv_comp_le_compProd_1, lem:fDiv_compProd_right}
By Lemma~\ref{lem:fDiv_comp_le_compProd_1}, $D_f(\kappa \circ \mu, \eta \circ \mu) \le D_f(\mu \otimes \kappa, \mu \otimes \eta)$. This is equal to $D_f(\kappa, \eta \mid \mu)$ by Lemma~\ref{lem:fDiv_compProd_right}.
\end{proof}


\begin{theorem}[Data-processing]
  \label{thm:fDiv_data_proc_1}
  \lean{ProbabilityTheory.fDiv_comp_right_le}
  \leanok
  \uses{def:fDiv}
  Let $\mu, \nu$ be two measures on $\mathcal X$ and let $\kappa : \mathcal X \rightsquigarrow \mathcal Y$ be a Markov kernel, where both $\mathcal X$ and $\mathcal Y$ are standard Borel.
  Then $D_f(\kappa \circ \mu, \kappa \circ \nu) \le D_f(\mu, \nu)$.
\end{theorem}

\begin{proof}\leanok
\uses{lem:fDiv_comp_le_compProd_1, thm:fDiv_compProd_left_1}
By Lemma~\ref{lem:fDiv_comp_le_compProd_1}, $D_f(\kappa \circ \mu, \kappa \circ \nu) \le D_f(\mu \otimes \kappa, \nu \otimes \kappa)$. Then the latter is equal to $D_f(\mu, \nu)$ by Lemma~\ref{thm:fDiv_compProd_left_1}.
\end{proof}


\subsubsection{DPI: proof using the integral representation}

This proof uses that every $f$-divergence is an integral over $\gamma \ge 0$ of $f$-divergences for the function $\phi_{1,\gamma}$ of Definition~\ref{def:statInfoFun}. It suffices then to prove the DPI for those particular $f$-divergences. Those are equal to the $E_\gamma$ divergence, up to a term that is invariant by composition with a Markov kernel. Finally, we already know that the $E_\gamma$ divergence satisfies the DPI, and that completes the proof.

\begin{lemma}
  \label{lem:fDiv_statInfoFun_eq}
  \lean{ProbabilityTheory.fDiv_statInfoFun_eq_StatInfo_of_nonneg_of_le, ProbabilityTheory.fDiv_statInfoFun_eq_StatInfo_of_nonneg_of_gt}
  \leanok
  \uses{def:fDiv, def:statInfoFun}
  Let $a,b \in [0, +\infty)$ and let $\mu, \nu$ be two measures on $\mathcal X$.
  \begin{align*}
  D_{\phi_{a,b}}(\mu, \nu)
  &= \text{sign}(b-a)\frac{1}{2}(a \mu(\mathcal X) - b \nu(\mathcal X)) + \frac{1}{2}\nu\left[ \left\vert a \frac{d\mu}{d\nu} - b \right\vert \right] + \frac{1}{2}a \mu_{\perp \nu}(\mathcal X)
  \: ,
  \end{align*}
  in which $\text{sign}(b-a)$ is $1$ if $b-a > 0$ and $-1$ if $b-a \le 0$.
\end{lemma}

\begin{proof}\leanok
\uses{}
If $a \le b$,
\begin{align*}
D_{\phi_{a,b}}(\mu, \nu)
&= \nu\left[ \max\{0, a \frac{d\mu}{d\nu} - b\} \right] + a \mu_{\perp \nu}(\mathcal X)
\\
&= \frac{1}{2}(a \mu_{\parallel \nu}(\mathcal X) - b \nu(\mathcal X)) + \frac{1}{2}\nu\left[ \left\vert a \frac{d\mu}{d\nu} - b \right\vert \right] + a \mu_{\perp \nu}(\mathcal X)
\\
&= \frac{1}{2}(a \mu(\mathcal X) - b \nu(\mathcal X)) + \frac{1}{2}\nu\left[ \left\vert a \frac{d\mu}{d\nu} - b \right\vert \right] + \frac{1}{2}a \mu_{\perp \nu}(\mathcal X)
\: . 
\end{align*}
Similar computations lead to the result for $b \le a$.
\end{proof}


\begin{corollary}
  \label{cor:fDiv_statInfoFun_eq_statInfo}
  \lean{ProbabilityTheory.fDiv_statInfoFun_eq_StatInfo_of_nonneg}
  \leanok
  \uses{def:fDiv, def:statInfoFun, def:statInfo}
  Let $a,b \in [0, +\infty)$ and let $\mu, \nu$ be two measures on $\mathcal X$.
  \begin{align*}
  D_{\phi_{a,b}}(\mu, \nu) = \mathcal I_{(a,b)}(\mu, \nu) + \frac{1}{2} \left\vert a \mu(\mathcal X) - b \nu(\mathcal X) \right\vert + \text{sign}(b-a)\frac{1}{2}(a \mu(\mathcal X) - b \nu(\mathcal X))
  \: .
  \end{align*}
\end{corollary}

\begin{proof}\leanok
\uses{lem:fDiv_statInfoFun_eq, cor:statInfo_eq_integral_abs}
Combine Lemma~\ref{lem:fDiv_statInfoFun_eq} and Corollary~\ref{cor:statInfo_eq_integral_abs}.
\end{proof}


\begin{lemma}
  \label{lem:fDiv_phi_data_proc}
  \lean{ProbabilityTheory.fDiv_statInfoFun_comp_right_le}
  \leanok
  \uses{def:fDiv, def:statInfoFun}
  Let $a,b \in [0, +\infty)$. Let $\mu, \nu$ be two finite measures on $\mathcal X$ and let $\kappa : \mathcal X \rightsquigarrow \mathcal Y$ be a Markov kernel.
  Then $D_{\phi_{a,b}}(\kappa \circ \mu, \kappa \circ \nu) \le D_{\phi_{a,b}}(\mu, \nu)$.
\end{lemma}

\begin{proof}\leanok
\uses{thm:data_proc_statInfo, cor:fDiv_statInfoFun_eq_statInfo}
By Corollary~\ref{cor:fDiv_statInfoFun_eq_statInfo},
\begin{align*}
D_{\phi_{a,b}}(\mu, \nu) = \mathcal I_{(a,b)}(\mu, \nu) + \frac{1}{2} \left\vert a \mu(\mathcal X) - b \nu(\mathcal X) \right\vert + \text{sign}(b-a)\frac{1}{2}(a \mu(\mathcal X) - b \nu(\mathcal X))
\: .
\end{align*}
The statistical information satisfies the data-processing inequality (Theorem~\ref{thm:data_proc_statInfo}) and the other terms are invariant by composition with a Markov kernel, hence $D_{\phi_{a,b}}$ also satisfies the data-processing inequality.
\end{proof}


\begin{theorem}[Data-processing]
  \label{thm:fDiv_data_proc_2}
  \lean{ProbabilityTheory.fDiv_comp_right_le'}
  \leanok
  \uses{def:fDiv}
  Let $\mu, \nu$ be two finite measures on $\mathcal X$ and let $\kappa : \mathcal X \rightsquigarrow \mathcal Y$ be a Markov kernel.
  Then $D_f(\kappa \circ \mu, \kappa \circ \nu) \le D_f(\mu, \nu)$.
\end{theorem}

\begin{proof}\leanok
\uses{thm:fDiv_eq_integral_eGamma, lem:fDiv_phi_data_proc}
By Theorem~\ref{thm:fDiv_eq_integral_eGamma},
\begin{align*}
D_f(\mu, \nu) = f(1) \nu(\mathcal X) + f'_+(1) (\mu(\mathcal X) - \nu (\mathcal X)) + \int_y D_{\phi_{1,y}}(\mu, \nu) \partial\gamma_f \: .
\end{align*}
The first two terms are unchanged by composition with a Markov kernel. It thus suffices to prove the data processing inequality for the divergences $D_{\phi_{1, \gamma}}$. This was done in Lemma~\ref{lem:fDiv_phi_data_proc}.
\end{proof}


\subsection{Sigma-algebras}

For $\mathcal A$ a sub-$\sigma$-algebra and $\mu$ a measure, we write $\mathcal \mu_{| \mathcal A}$ for the measure restricted to the $\sigma$-algebra.
The measure $\mathcal \mu_{| \mathcal A}$ is equal to the map of $\mu$ by the identity, seen as a function to $\mathcal X$ with the $\sigma$-algebra $\mathcal A$.
We can thus get a DPI for sub-sigma-algebras by using the DPI for maps.

\begin{theorem}
  \label{thm:fDiv_trim_le}
  \lean{ProbabilityTheory.fDiv_trim_le}
  \leanok
  \uses{def:fDiv}
  Let $\mu, \nu$ be two finite measures on $\mathcal X$ and let $\mathcal A$ be a sub-$\sigma$-algebra of $\mathcal X$. Then
  $D_f(\mu_{| \mathcal A}, \nu_{| \mathcal A}) \le D_f(\mu, \nu)$.
\end{theorem}

\begin{proof}\leanok
\uses{thm:fDiv_map_le}
The measure $\mathcal \mu_{| \mathcal A}$ is equal to the map of $\mu$ by the identity, seen as a function to $\mathcal X$ with the $\sigma$-algebra $\mathcal A$.
That function is measurable since $\mathcal A$ is a sub-sigma-algebra.
We can apply Theorem~\ref{thm:fDiv_map_le}.
\end{proof}


\begin{theorem}
  \label{thm:iSup_fDiv_trim}
  %\lean{}
  %\leanok
  \uses{def:fDiv}
  Let $\mu, \nu$ be two finite measures on $\mathcal X$. Then
  $\sup_{\mathcal A \text{ finite}} D_f(\mu_{| \mathcal A}, \nu_{| \mathcal A}) = D_f(\mu, \nu)$.
\end{theorem}

\begin{proof}
\uses{thm:fDiv_trim_le}
\end{proof}


\begin{lemma}
  \label{lem:fDiv_map_eq_fDiv_trim_of_ac}
  \lean{ProbabilityTheory.fDiv_map_eq_fDiv_trim_of_ac}
  \leanok
  \uses{def:fDiv}
  Let $\mu, \nu \in \mathcal M(\mathcal X)$ be finite measures with $\mu \ll \nu$ and let $g : \mathcal X \to \mathcal Y$ be a measurable function. Denote by $g^* \mathcal Y$ the comap of the $\sigma$-algebra on $\mathcal Y$ by $g$. Then
  $D_f(g_* \mu, g_* \nu) = D_f(\mu_{| g^* \mathcal Y}, \nu_{| g^* \mathcal Y})$~.
\end{lemma}

\begin{proof}\leanok
\uses{lem:rnDeriv_map_eq_rnDeriv_trim}
Using Lemma~\ref{lem:rnDeriv_map_eq_rnDeriv_trim},
\begin{align*}
D_f(g_* \mu, g_* \nu)
&= \int_y f \left( \frac{d g_*\mu}{d g_*\nu}(y)\right) \partial(g_*\nu)
\\
&= \int_x f \left( \frac{d g_*\mu}{d g_*\nu}(g(x))\right) \partial\nu
\\
&= \int_x f \left(\frac{d \mu_{| g^* \mathcal Y}}{d \nu_{| g^* \mathcal Y}}(x) \right) \partial\nu
\\
&= \int_x f \left(\frac{d \mu_{| g^* \mathcal Y}}{d \nu_{| g^* \mathcal Y}}(x) \right) \partial\nu_{| g^* \mathcal Y}
\\
&= D_f(\mu_{| g^* \mathcal Y}, \nu_{| g^* \mathcal Y})
\: .
\end{align*}
\end{proof}



\subsection{Consequences of the data-processing inequality}

\subsubsection{Direct applications}

\begin{theorem}
  \label{thm:fDiv_le_compProd_2}
  \lean{ProbabilityTheory.le_fDiv_compProd'}
  \leanok
  \uses{def:fDiv}
  Let $\mu, \nu$ be two finite measures on $\mathcal X$ and let $\kappa, \eta : \mathcal X \rightsquigarrow \mathcal Y$ be two Markov kernels.
  Then $D_f(\mu, \nu) \le D_f(\mu \otimes \kappa, \nu \otimes \eta)$.
\end{theorem}

\begin{proof}\leanok
\uses{thm:fDiv_data_proc_2}
Since $\kappa$ is a Markov kernel, $\mu$ is the composition of $\mu \otimes \kappa$ and the deterministic kernel for the function $(x,y) \mapsto x$. The same applies for $\nu$ and $\nu \otimes \eta$.
The result is then an application of the data-processing inequality, Theorem~\ref{thm:fDiv_data_proc_2}.
\end{proof}


\begin{theorem}
  \label{thm:fDiv_comp_le_compProd_2}
  \lean{ProbabilityTheory.fDiv_comp_le_compProd'}
  \leanok
  \uses{def:fDiv}
  Let $\mu, \nu$ be two finite measures on $\mathcal X$ and let $\kappa, \eta : \mathcal X \rightsquigarrow \mathcal Y$ be two finite kernels.
  Then $D_f(\kappa \circ \mu, \eta \circ \nu) \le D_f(\mu \otimes \kappa, \nu \otimes \eta)$.
\end{theorem}

\begin{proof}\leanok
\uses{thm:fDiv_data_proc_2}
$\kappa \circ \mu$ is the composition of $\mu \otimes \kappa$ and the deterministic kernel for the function $(x,y) \mapsto y$. The same applies for $\eta \circ \nu$ and $\nu \otimes \eta$.
The result is then an application of the data-processing inequality, Theorem~\ref{thm:fDiv_data_proc_2}.
\end{proof}


% Commented out because identical to something above
%
% \begin{lemma}[Composition-product with a kernel]
%   \label{thm:fDiv_compProd_left_2}
%   \lean{ProbabilityTheory.fDiv_compProd_left'}
%   \leanok
%   \uses{def:fDiv}
%   Let $\mu, \nu$ be two finite measures on $\mathcal X$ and let $\kappa : \mathcal X \rightsquigarrow \mathcal Y$ be a Markov kernel.
%   Then $D_f(\mu \otimes \kappa, \nu \otimes \kappa) = D_f(\mu, \nu)$.
% \end{lemma}

% \begin{proof}\leanok
% \uses{thm:fDiv_le_compProd_2, thm:fDiv_data_proc_2}
% By Theorem~\ref{thm:fDiv_le_compProd_2}, $D_f(\mu, \nu) \le D_f(\mu \otimes \kappa, \nu \otimes \kappa)$.
% For the reverse inequality, remark that $\mu \otimes \kappa = (\text{id} \times \kappa) \circ \mu$ where $\text{id}$ is the identity kernel and apply the data-processing inequality (Theorem~\ref{thm:fDiv_data_proc_2}). 
% \end{proof}


\begin{theorem}[Conditioning increases f-divergence]
  \label{thm:fDiv_comp_le_condFDiv_2}
  \lean{ProbabilityTheory.fDiv_comp_le_compProd_right'}
  \leanok
  \uses{def:fDiv, def:condFDiv}
  Let $\mu$ be a finite measure $\mathcal X$ and let $\kappa, \eta : \mathcal X \rightsquigarrow \mathcal Y$ be two finite kernels.
  Then $D_f(\kappa \circ \mu, \eta \circ \mu) \le D_f(\mu \otimes \kappa, \mu \otimes \eta) = D_f(\kappa, \eta \mid \mu)$~.
\end{theorem}

\begin{proof}\leanok
\uses{thm:fDiv_comp_le_compProd_2}
This is a particular case of Theorem~\ref{thm:fDiv_comp_le_compProd_2}.
\end{proof}


\begin{theorem}[Marginals]
  \label{thm:fDiv_fst_le_2}
  \lean{ProbabilityTheory.fDiv_fst_le', ProbabilityTheory.fDiv_snd_le'}
  \leanok
  \uses{def:fDiv}
  Let $\mu$ and $\nu$ be two measures on $\mathcal X \times \mathcal Y$, and let $\mu_X, \nu_X$ be their marginals on $\mathcal X$.
  Then $D_f(\mu_X, \nu_X) \le D_f(\mu, \nu)$.
  Similarly, for $\mu_Y, \nu_Y$ the marginals on $\mathcal Y$, $D_f(\mu_Y, \nu_Y) \le D_f(\mu, \nu)$.
\end{theorem}

\begin{proof}\leanok
\uses{thm:fDiv_data_proc_2}
The measure $\mu_X$ is the composition of $\mu$ and the deterministic kernel for the function $(x,y) \mapsto x$, and similarly $\mu_Y$ is the composition of $\mu$ and a deterministic kernel.
The results are both applications of the data-processing inequality, Theorem~\ref{thm:fDiv_data_proc_2}.
\end{proof}


\begin{lemma}
  \label{lem:fDiv_compProd_prod_eq}
  \lean{ProbabilityTheory.fDiv_comp_eq_of_fst}
  \leanok
  \uses{def:fDiv}
  Let $\mu, \nu$ be two finite measures on $\mathcal X$ and let $\kappa : \mathcal X \rightsquigarrow (\mathcal X \times \mathcal Y)$ be a Markov kernel such that for all $x$, $(\kappa(x))_X = \delta_x$. Then $D_f(\kappa \circ \mu, \kappa \circ \nu) = D_f(\mu, \nu)$.
\end{lemma}

\begin{proof}\leanok
\uses{thm:fDiv_data_proc_2, thm:fDiv_fst_le_2}
$D_f(\kappa \circ \mu, \kappa \circ \nu) \le D_f(\mu, \nu)$ by Theorem~\ref{thm:fDiv_data_proc_2}.
For the other inequality, remark that $\mu = (\kappa \circ \mu)_X$ (and similarly for $\nu$). Hence by Theorem~\ref{thm:fDiv_fst_le_2} $D_f(\mu, \nu) \le D_f(\kappa \circ \mu, \kappa \circ \nu)$.
\end{proof}



\subsubsection{Data-processing and mean values}

In this section $d_f(p, q)$ denotes the $f$-divergence between two Bernoulli distributions with means $p$ and $q$.

\begin{corollary}
  \label{cor:data_proc_event}
  \lean{ProbabilityTheory.fDiv_boolMeasure_le}
  \leanok
  \uses{def:fDiv}
  Let $\mu, \nu$ be two measures on $\mathcal X$ and let $E$ be an event. Then $D_f(\mu, \nu) \ge d_f(\mu(E), \nu(E))$.
\end{corollary}

\begin{proof}\leanok
\uses{thm:fDiv_data_proc_2}
Use the deterministic kernel $\kappa : \mathcal X \rightsquigarrow \{0, 1\}$ with $\kappa(x) = \delta_1 \mathbb{I}\{x \in E\} + \delta_0 \mathbb{I}\{x \notin E\}$ in Theorem~\ref{thm:fDiv_data_proc_2}.
\end{proof}


\begin{lemma}
  \label{lem:fDiv_bernoulli_convex}
  %\lean{}
  %\leanok
  \uses{def:fDiv}
  For all $y \in [0,1]$, $x \mapsto d_f(x, y)$ is convex and attains a minimum at $x = y$.
\end{lemma}

\begin{proof}%\leanok
\uses{}

\begin{align*}
d_f(x, y)
&= y f(\frac{x}{y}) + (1 - y) f(\frac{1 - x}{1 - y})
\: , \\
(x \mapsto d_f(x, y))'_+(x)
&= f'_+(\frac{x}{y}) - f'_+(\frac{1 - x}{1 - y})
\: .
\end{align*}
$f'_+$ is non-decreasing, hence we get that $(x \mapsto d_f(x, y))'_+$ is also non-decreasing and $x \mapsto d_f(x, y)$ is convex.

For $x \le y$, $x/y \le 1 \le (1-x)/(1-y)$, hence $f'_+(\frac{x}{y}) \le f'_+(\frac{1 - x}{1 - y})$ and $(x \mapsto d_f(x, y))'_+(x) \le 0$. The opposite holds for $x \ge y$.

\end{proof}


\begin{lemma}
  \label{lem:fDiv_bounded_ge_fDiv_mean}
  %\lean{}
  %\leanok
  \uses{def:fDiv}
  Let $\mu, \nu \in \mathcal P([0,1])$. Then
  \begin{align*}
  D_f(\mu, \nu) \ge d_f(\mu[X], \nu[X]) \: .
  \end{align*}
\end{lemma}

\begin{proof}%\leanok
\uses{thm:fDiv_data_proc_2,lem:fDiv_compProd_prod_eq}
Let $u$ be the uniform distribution on $[0,1]$. Then $D_f(\mu, \nu) = D_f(\mu \times u, \nu \times u)$ (by Lemma~\ref{lem:fDiv_compProd_prod_eq}).
Let $g : [0,1] \times [0,1] \to \{0,1\}$ be the measurable function with $g(x, y) = \mathbb{I}\{y \le x\}$. We also write $g$ for the corresponding deterministic kernel. By the data-processing inequality (Theorem~\ref{thm:fDiv_data_proc_2}),
\begin{align*}
D_f(\mu \times u, \nu \times u)
\ge D_f(g \circ (\mu \times u), g \circ (\nu \times u))
\: .
\end{align*}
Those two last measures are Bernoulli distributions with means $\mu[X]$ and $\nu[X]$.
\end{proof}


\begin{corollary}
  \label{cor:fDiv_ge_fDiv_mean_comp}
  %\lean{}
  %\leanok
  \uses{def:fDiv}
  Let $\mu, \nu \in \mathcal P(\mathcal X)$ and let $\kappa : \mathcal X \rightsquigarrow [0,1]$. Then
  \begin{align*}
  D_f(\mu, \nu) \ge d_f((\kappa \circ \mu)[X], (\kappa \circ \nu)[X]) \: .
  \end{align*}
\end{corollary}

\begin{proof}%\leanok
\uses{lem:fDiv_bounded_ge_fDiv_mean}
First, by the data-processing inequality (Theorem~\ref{thm:fDiv_data_proc_2}),
$D_f(\mu, \nu) \ge D_f(\kappa \circ \mu, \kappa \circ \nu)$~.
Then use Lemma~\ref{lem:fDiv_bounded_ge_fDiv_mean}.
\end{proof}


\begin{lemma}
  \label{lem:fDiv_estimation_ge}
  %\lean{}
  %\leanok
  \uses{def:fDiv, def:bayesRisk}
  Let $\pi, \xi \in \mathcal P(\Theta)$ and $P, Q : \Theta \rightsquigarrow \mathcal X$. Suppose that the loss $\ell'$ takes values in $[0,1]$. Then
  \begin{align*}
  D_f(\pi \otimes Q, \xi \otimes P)
  &\ge d_f(\mathcal R_\pi^Q, \mathcal R_\xi^P)
  \: .
  \end{align*}
\end{lemma}

\begin{proof}%\leanok
\uses{cor:fDiv_ge_fDiv_mean_comp, lem:fDiv_bernoulli_convex}
Suppose first that $\mathcal R_\pi^Q \ge \mathcal R_\xi^P$.
Let $\hat{y}_B$ be a Bayes estimator for the estimation task $(P, y, \ell')$ with prior $\xi$. If it does not exist, the proof can be adapted by taking an estimator with risk $\varepsilon$-close to the Bayes risk for any $\varepsilon > 0$.
By Corollary~\ref{cor:fDiv_ge_fDiv_mean_comp},
\begin{align*}
D_f(\pi \otimes Q, \xi \otimes P)
&\ge d_f(\ell' \circ (y \parallel \hat{y}_B) \circ (\pi \otimes Q) [X], \ell' \circ (y \parallel \hat{y}_B) \circ (\xi \otimes P) [X])
\\
&= d_f(\ell' \circ (y \parallel \hat{y}_B) \circ (\pi \otimes Q) [X], \mathcal R_\xi^P)
\: .
\end{align*}
By definition of the Bayes risk as an infimum,
\begin{align*}
\ell' \circ (y \parallel \hat{y}_B) \circ (\pi \otimes Q) [X]
\ge \mathcal R_\pi^Q
\ge \mathcal R_\xi^P
\: .
\end{align*}
Since $x \mapsto d_f(x, y)$ is monotone on $[y, 1]$,
\begin{align*}
D_f(\pi \otimes Q, \xi \otimes P)
&\ge d_f(\ell' \circ (y \parallel \hat{y}_B) \circ (\pi \otimes Q) [X], \mathcal R_\xi^P)
\\
&\ge d_f(\mathcal R_\pi^Q, \mathcal R_\xi^P)
\: .
\end{align*}
This concludes the proof for $\mathcal R_\pi^Q \ge \mathcal R_\xi^P$. If that inequality is reversed, we proceed similarly but take the Bayes estimator for $Q$.
\end{proof}





\subsection{Convexity}

\begin{theorem}[Joint convexity]
  \label{thm:fDiv_convex}
  \lean{ProbabilityTheory.fDiv_convex}
  \leanok
  \uses{def:fDiv}
  The function $(\mu, \nu) \mapsto D_f(\mu, \nu)$ is convex.
\end{theorem}

\begin{proof}\leanok
\uses{thm:fDiv_comp_le_condFDiv_2}
Let $\mu_0, \mu_1, \nu_0, \nu_1$ be four measures. Let $\lambda \in [0,1]$. Let $\xi$ be the probability measure on $\{0,1\}$ with $\xi(\{1\}) = \lambda$.
Let $\kappa$ be the kernel $\{0,1\} \rightsquigarrow \mathcal X$ defined by $\kappa(0) = \mu_0$ and $\kappa(1) = \mu_1$.
Let $\eta$ be the kernel $\{0,1\} \rightsquigarrow \mathcal X$ defined by $\eta(0) = \nu_0$ and $\eta(1) = \nu_1$.
\begin{align*}
D_f(\xi \otimes \kappa, \xi \otimes \eta)
&= D_f(\kappa, \eta \mid \xi)
= (1 - \lambda) D_f(\mu_0, \nu_0) + \lambda D_f(\mu_1, \nu_1)
\: , \\
D_f(\kappa \circ \xi, \eta \circ \xi)
&= D_f((1 - \lambda)\mu_0 + \lambda \mu_1, (1 - \lambda)\nu_0 + \lambda \nu_1)
\: .
\end{align*}
By Theorem~\ref{thm:fDiv_comp_le_condFDiv_2}, $D_f(\kappa \circ \xi, \eta \circ \xi) \le D_f(\xi \otimes \kappa, \xi \otimes \eta)$.
\end{proof}



\subsection{Variational representations}

TODO




\subsection{Comparisons between divergences}

We collate here general tools that can be used to compare $D_f$ and $D_g$ for two functions $f$ and $g$.
Specific divergences are discussed in the next sections.

\begin{lemma}
  \label{lem:fDiv_mono_fun}
  \lean{ProbabilityTheory.fDiv_mono}
  \leanok
  \uses{def:fDiv}
  If $f \le g$, then $D_f(\mu, \nu) \le D_g(\mu, \nu)$.
\end{lemma}

\begin{proof}\leanok
Monotonicity of the integral and of $f \mapsto f'(\infty)$.
\end{proof}


\begin{lemma}
  \label{lem:fDiv_le_of_deriv2_le}
  %\lean{}
  %\leanok
  \uses{def:fDiv}
  If $f(1) = 0$, $g(1) = 0$, $f'(1) = 0$, $g'(1) = 0$, and both $f$ and $g$ have a second derivative, then
  \begin{align*}
  D_f(\mu, \nu) \le \sup_x \frac{f''(x)}{g''(x)} D_g(\mu, \nu)
  \end{align*}
\end{lemma}

\begin{proof}%\leanok
\uses{lem:fDiv_mono_fun}
By Taylor with integral remainder, if $f'' \le \beta g''$,
\begin{align*}
f(x) = \int_1^x f''(t) (x - t) dt
\le \beta \int_1^x g''(t) (x - t) dt
= \beta g(x)
\: .
\end{align*}
Then use Lemma~\ref{lem:fDiv_mono_fun}.
\end{proof}


\begin{theorem}[\cite{sason2016f}]
  \label{thm:fDiv_eq_sup_mul_fDiv}
  %\lean{}
  %\leanok
  \uses{def:fDiv}
  Suppose that $f(1) = g(1) = 0$ and that $f'(1) = g'(1) = 0$, and that $g>0$ on $(0,1) \cup (1, +\infty)$. Suppose that the space $\mathcal X$ has at least two disjoint non-empty measurable sets. Then
  \begin{align*}
  \sup_{\mu, \nu \in \mathcal P(\mathcal X)} \frac{D_f(\mu, \nu)}{D_g(\mu, \nu)}
  = \sup_{t > 0, t \ne 1} \frac{f(t)}{g(t)}
  \: .
  \end{align*}
\end{theorem}

Remark: if $\mathcal X$ does not have two disjoint non-empty measurable sets then its measurable sets are only $\{\emptyset, \mathcal X\}$ and the ratio on the left is $0/0$.

\begin{proof}%\leanok
\uses{lem:fDiv_mono_fun}
TODO: proof done for measurable singletons?

By Lemma~\ref{lem:fDiv_mono_fun}, $\sup_{\mu, \nu \in \mathcal P(\mathcal X)} \frac{D_f(\mu, \nu)}{D_g(\mu, \nu)} \le \sup_{t > 0, t \ne 1} \frac{f(t)}{g(t)}$. Let's prove the other inequality.

Let $x$ and $y$ be two points in disjoint non-empty measurable sets. For $t \in (0,1) \cup (1, +\infty)$ and $\varepsilon \in (0, 1/t]$ let $\mu_{t, \varepsilon} = t \varepsilon \delta_x + (1 - t \varepsilon) \delta_y$ and $\nu_\varepsilon = \varepsilon \delta_x + (1 - \varepsilon) \delta_y$~.
\begin{align*}
D_f(\mu_{t, \varepsilon}, \nu_\varepsilon)
&= \varepsilon f(t) + (1 - \varepsilon) f(\frac{1 - t \varepsilon}{1 - \varepsilon})
\\
&= \varepsilon f(t) + \varepsilon (1 - t) \frac{1 - \varepsilon}{\varepsilon (1 - t)} f(1 + \frac{\varepsilon(1 - t)}{1 - \varepsilon})
\end{align*}
Let $\alpha = \frac{\varepsilon (1 - t)}{1 - \varepsilon}$. Then since $f(1) = g(1) = 0$,
\begin{align*}
\frac{D_f(\mu_{t, \varepsilon}, \nu_\varepsilon)}{D_g(\mu_{t, \varepsilon}, \nu_\varepsilon)}
&= \frac{f(t) + (1 - t)\frac{f(1 + \alpha) - f(1)}{\alpha}}{g(t) + (1 - t)\frac{g(1 + \alpha) - g(1)}{\alpha}} \: .
\end{align*}
As $\varepsilon \to 0$, $\alpha$ also tends to 0 and
\begin{align*}
\frac{D_f(\mu_{t, \varepsilon}, \nu_\varepsilon)}{D_g(\mu_{t, \varepsilon}, \nu_\varepsilon)}
\to \frac{f(t) + (1 - t)f'(1)}{g(t) + (1 - t)g'(1)}
= \frac{f(t)}{g(t)}
\: .
\end{align*}
\end{proof}

\paragraph{Bretagnolle-Huber inequalities}

\begin{theorem}
  \label{thm:bh_fDiv}
  \lean{ProbabilityTheory.apply_one_add_tv_add_apply_one_sub_tv_le_fDiv, ProbabilityTheory.conj_one_add_tv_add_conj_one_sub_tv_le_fDiv}
  \leanok
  \uses{def:fDiv, def:TV}
  Let $\mu, \nu \in \mathcal P(\mathcal X)$. If $f(1) = 0$,
  \begin{align*}
  D_f(\mu, \nu)
  &\ge f(1 + \TV(\mu, \nu)) + f(1 - \TV(\mu, \nu))
  \: , \\
  \text{ and }
  D_f(\mu, \nu)
  &\ge (1 + \TV(\mu, \nu))f\left(\frac{1}{1 + \TV(\mu, \nu)}\right) + (1 - \TV(\mu, \nu))f\left(\frac{1}{1 - \TV(\mu, \nu)}\right)
  \: .
  \end{align*}
\end{theorem}

\begin{proof}\leanok
\uses{lem:tv_eq_integral}
Let $V: x \mapsto \max\left\{0, \frac{d\mu}{d\nu}(x) - 1\right\}$ and $W: x \mapsto \max\left\{0, 1 - \frac{d\mu}{d\nu}(x)\right\}$~.
Then one of $1+V$ and $1-W$ is equal to $\frac{d\mu}{d\nu}$ and the other is equal to 1.
Since $f(1) = 0$, we obtain $f(\frac{d\mu}{d\nu}) = f(1+V) + f(1-W)$~.
The $f$-divergence is then
\begin{align*}
D_f(\mu, \nu)
&= \nu[f(\frac{d\mu}{d\nu})] + f'(\infty)\mu_{\perp \nu}(\mathcal X)
\\
&= \nu\left[f(1+V) + f'(\infty)\mu_{\perp \nu}(\mathcal X)\right] + \nu\left[f(1 - W)\right]
\: .
\end{align*}
By convexity of $f$, $f(1+V) + f'(\infty)\mu_{\perp \nu}(\mathcal X) \ge f(1+V + \mu_{\perp \nu}(\mathcal X))$~. Using convexity again, we get
\begin{align*}
D_f(\mu, \nu)
&\ge f(1+ \nu[V] + \mu_{\perp \nu}(\mathcal X)) + f(1 - \nu[W])
\: .
\end{align*}
By Lemma~\ref{lem:tv_eq_integral}, $\nu[V] + \mu_{\perp \nu}(\mathcal X) = \nu[W] = \TV(\mu, \nu)$. We obtain the first equality.

To get the second equality, remark that $D_f(\mu, \nu) = D_{x\mapsto xf(1/x)}(\nu, \mu)$ (Lemma~\ref{lem:fDiv_symm}), apply the first inequality to that function and use the symmetry of $\TV$.
\end{proof}

\begin{corollary}[Bretagnolle-Huber inequality]
  \label{cor:bh_kl}
  \lean{ProbabilityTheory.neg_log_one_sub_sq_tv_le_klDiv}
  \leanok
  \uses{def:KL, def:TV}
  TODO: move this somewhere after the definition of $KL$.

  For $\mu, \nu \in \mathcal P(\mathcal X)$,
  \begin{align*}
  \KL(\mu, \nu)
  &\ge - \log(1 - \TV^2(\mu, \nu))
  \: , \\
  \text{or equivalently, }
  \TV(\mu, \nu)
  &\le \sqrt{1 - \exp\left(-\KL(\mu, \nu) \right)}
  \: .
  \end{align*}
\end{corollary}

\begin{proof}\leanok
\uses{thm:bh_fDiv}
Use the second inequality of Theorem~\ref{thm:bh_fDiv}, for $f: x \mapsto x \log x$~.
\end{proof}

\begin{corollary}
  \label{cor:cor:bh_hellingerAlpha}
  %\lean{}
  %\leanok
  \uses{def:hellingerAlpha, def:TV}
  TODO: move this somewhere after the definition of $\He_\alpha$.

  Let $\mu, \nu \in \mathcal P(\mathcal X)$. For $\alpha > 0$,
  \begin{align*}
  \He_{1+\alpha}(\mu, \nu)
  &\ge \frac{1}{\alpha} \left( (1 - \TV(\mu, \nu))^{-\alpha} - 2 \right)
  \: , \\
  \text{or equivalently, }
  1 - \TV(\mu, \nu)
  &\ge \exp\left(-\frac{1}{\alpha} \log \left( 2 + \alpha \He_{1+\alpha}(\mu, \nu) \right)\right)
  \: .
  \end{align*}
\end{corollary}

\begin{proof}%\leanok
\uses{thm:bh_fDiv}
Use the second inequality of Theorem~\ref{thm:bh_fDiv}, for $f: x \mapsto \frac{x^{1+\alpha} - 1}{\alpha}$ to get
\begin{align*}
\He_{1+\alpha}(\mu, \nu) \ge \frac{1}{\alpha} \left((1 + \TV(\mu, \nu))^{-\alpha} + (1 - \TV(\mu, \nu))^{-\alpha} - 2 \right) \: .
\end{align*}
We then neglect one term: $(1 + \TV(\mu, \nu))^{-\alpha} \ge 0$.
\end{proof}


\paragraph{Joint range}

TODO \cite{harremoes2011pairs}


\subsection{Statistical divergences}

\begin{lemma}
  \label{lem:statInfo_eq_fDiv}
  \lean{ProbabilityTheory.statInfo_eq_fDiv}
  \leanok
  \uses{def:statInfo, def:fDiv}
  On probability measures, the statistical information $\mathcal I_\xi$ is an $f$-divergence for the function $\phi_{\xi_0, \xi_1}$ of Definition~\ref{def:statInfoFun}.
\end{lemma}

\begin{proof}\leanok
\uses{thm:statInfo_eq_integral}
This is a reformulation of Theorem~\ref{thm:statInfo_eq_integral}.
\end{proof}

% \begin{lemma}
%   \label{lem:deGrootInfo_eq_fDiv}
%   %\lean{}
%   %\leanok
%   \uses{def:deGrootInfo, def:statInfoFun, def:fDiv}
%   On probability measures, the DeGroot statistical information $I_\pi$ is an $f$-divergence for the function $\phi_{\pi, 1-\pi}$ of Definition~\ref{def:statInfoFun}.
% \end{lemma}

% \begin{proof}%\leanok
% \uses{lem:deGrootInfo_eq_integral}
% This is a reformulation of Lemma~\ref{lem:deGrootInfo_eq_integral}.
% \end{proof}

% \begin{lemma}
%   \label{lem:eGamma_eq_fDiv}
%   %\lean{}
%   %\leanok
%   \uses{def:eGamma, def:statInfoFun, def:fDiv}
%   On probability measures, the hockey-stick divergence $E_\gamma$ is an $f$-divergence for the function $\phi_{1,\gamma}$ of Definition~\ref{def:statInfoFun}.
% \end{lemma}

% \begin{proof}%\leanok
% \uses{lem:eGamma_eq_integral}
% This is a reformulation of Lemma~\ref{lem:eGamma_eq_integral}.
% \end{proof}

\begin{lemma}
  \label{lem:tv_eq_fDiv}
  \lean{ProbabilityTheory.fDiv_tvDivFun, ProbabilityTheory.tv_eq_fDiv_tvDivFun}
  \leanok
  \uses{def:TV, def:fDiv}
  On probability measures, the total variation distance $\TV$ is an $f$-divergence for the function $x \mapsto \frac{1}{2}\vert x - 1 \vert$.
\end{lemma}

\begin{proof}\leanok
\uses{lem:tv_eq_integral_abs}
This is a reformulation of Lemma~\ref{lem:tv_eq_integral_abs}.
\end{proof}

%%FILE blueprint/src/sections/hellinger.tex
\section{Hellinger distance}

TODO: refactor this section to link it to the Hellinger $\alpha$-divergence.

\begin{definition}[Squared Hellinger distance]
  \label{def:Hellinger}
  %\lean{}
  %\leanok
  \uses{def:fDiv, lem:fDiv_properties}
  Let $\mu, \nu$ be two measures. The squared Hellinger distance between $\mu$ and $\nu$ is
  \begin{align*}
    \Hsq(\mu, \nu) = D_f(\mu, \nu) \quad \text{with } f: x \mapsto \frac{1}{2}\left( 1 - \sqrt{x} \right)^2 \: .
  \end{align*}
\end{definition}

\begin{lemma}[Hellinger and Rényi]
  \label{lem:renyi_half_eq_log_hellinger}
  %\lean{}
  %\leanok
  \uses{def:Hellinger, def:Renyi}
  Let $\mu, \nu$ be two probability measures. Then $R_{1/2}(\mu, \nu) = -2\log(1 - \Hsq(\mu, \nu))$.
\end{lemma}

\begin{proof}
\uses{lem:fDiv_mul, lem:fDiv_add_linear}
$R_{1/2}(\mu, \nu) = -2 \log (1 - \frac{1}{2} D_f(\nu, \mu))$ for $f : x \mapsto -2 (\sqrt{x} - 1)$. Using Lemma~\ref{lem:fDiv_mul}, $R_{1/2}(\mu, \nu) = -2 \log (1 - D_g(\nu, \mu))$ for $g(x) = 1 - \sqrt{x}$.

It suffices then to show that $\Hsq(\mu, \nu) = D_g(\mu, \nu)$, which is true by an application of Lemma~\ref{lem:fDiv_add_linear}.
\end{proof}

\begin{lemma}
  \label{lem:hellinger_le_renyi}
  %\lean{}
  %\leanok
  \uses{def:Renyi, def:Hellinger}
  Let $\mu, \nu$ be two probability measures. Then $2 \Hsq(\mu, \nu) \le R_{1/2}(\mu, \nu)$.
\end{lemma}

\begin{proof}
\uses{lem:renyi_half_eq_log_hellinger}
Use Lemma~\ref{lem:renyi_half_eq_log_hellinger}, then $-\log(1 - x) \ge x$.
\end{proof}

\begin{lemma}
  \label{lem:hellinger_le_tv}
  %\lean{}
  %\leanok
  \uses{def:Hellinger, def:TV}
  Let $\mu, \nu$ be two probability measures. Then $\Hsq(\mu, \nu) \le \TV(\mu, \nu)$.
\end{lemma}

\begin{proof}
\end{proof}

\begin{lemma}
  \label{lem:tv_le_hellinger}
  %\lean{}
  %\leanok
  \uses{def:Hellinger, def:TV}
  Let $\mu, \nu$ be two probability measures. Then $\TV(\mu, \nu) \le \sqrt{\Hsq(\mu, \nu)(2 - \Hsq(\mu, \nu))}$.
\end{lemma}

\begin{proof}
\end{proof}

\begin{corollary}
  \label{cor:one_sub_hellinger_squared_le_one_sub_tv}
  %\lean{}
  %\leanok
  \uses{def:Hellinger, def:TV}
  Let $\mu, \nu$ be two probability measures. Then
  \begin{align*}
  \frac{1}{2}(1 - \Hsq(\mu, \nu))^2 \le 1 - \TV(\mu, \nu) - \frac{1}{2}(1 - \TV(\mu, \nu))^2
  \: .
  \end{align*}
\end{corollary}

\begin{proof}
\uses{lem:tv_le_hellinger}
We use Lemma~\ref{lem:tv_le_hellinger}.
\begin{align*}
1 - \TV(\mu, \nu) - \frac{1}{2}(1 - \Hsq(\mu, \nu))^2
&= \frac{1}{2} - \TV(\mu, \nu) + \frac{1}{2}(\Hsq(\mu, \nu)(2 - \Hsq(\mu, \nu)))
\\
&\ge \frac{1}{2} - \TV(\mu, \nu) + \frac{1}{2}\TV^2(\mu, \nu)
\\
&= \frac{1}{2}(1 - \TV(\mu, \nu))^2
\: .
\end{align*}
\end{proof}

\begin{corollary}
  \label{cor:one_sub_hellinger_le_one_sub_tv}
  %\lean{}
  %\leanok
  \uses{def:Hellinger, def:TV}
  Let $\mu, \nu$ be two probability measures. Then
  \begin{align*}
  1 - \sqrt{1 - (1 - \Hsq(\mu, \nu))^2} \le 1 - \TV(\mu, \nu) \le 1 - \Hsq(\mu, \nu)
  \: .
  \end{align*}
\end{corollary}

\begin{proof}
\uses{lem:tv_le_hellinger, lem:hellinger_le_tv}
Use Lemma~\ref{lem:tv_le_hellinger} and~\ref{lem:hellinger_le_tv}.
\end{proof}

\begin{corollary}
  \label{cor:one_sub_tv_bound_renyi}
  %\lean{}
  %\leanok
  \uses{def:Renyi, def:TV}
  Let $\mu, \nu$ be two probability measures on $\mathcal X$ and let $n \in \mathbb{N}$, and $\mu^{\otimes n}, \nu^{\otimes n}$ be their product measures on $\mathcal X^n$. Then
  \begin{align*}
  \frac{1}{2}e^{-n R_{1/2}(\mu, \nu)} \le 1 - \sqrt{1 - e^{-n R_{1/2}(\mu, \nu)}} \le 1 - \TV(\mu^{\otimes n}, \nu^{\otimes n}) \le e^{-\frac{1}{2} n R_{1/2}(\mu, \nu)}
  \: .
  \end{align*}
\end{corollary}

\begin{proof}
\uses{cor:one_sub_hellinger_le_one_sub_tv, lem:renyi_half_eq_log_hellinger, lem:renyi_prod_n}
Use Corollary~\ref{cor:one_sub_hellinger_le_one_sub_tv} and Lemma~\ref{lem:renyi_half_eq_log_hellinger} to get to a comparison of $\TV(\mu, \nu)$ and $R_{1/2}(\mu^{\otimes n}, \nu^{\otimes n})$. Finally $R_{1/2}(\mu^{\otimes n}, \nu^{\otimes n}) = n R_{1/2}(\mu, \nu)$ by Lemma~\ref{lem:renyi_prod_n}.
\end{proof}


%%FILE blueprint/src/sections/hellinger_alpha.tex
\section{Hellinger alpha-divergences}

\begin{definition}[Hellinger $\alpha$-divergence]
  \label{def:hellingerAlpha}
  \lean{ProbabilityTheory.hellingerDiv}
  \leanok
  \uses{def:KL, def:fDiv, lem:fDiv_properties, lem:kl_properties}
  Let $\mu, \nu$ be two measures on $\mathcal X$. The Hellinger divergence of order $\alpha \in [0,+\infty)$ between $\mu$ and $\nu$ is
  \begin{align*}
  \He_\alpha(\mu, \nu) = \left\{
  \begin{array}{ll}
    \nu \{x \mid \frac{d\mu}{d\nu} (x) = 0\} & \text{for } \alpha = 0 \\
    \KL(\mu, \nu) & \text{for } \alpha = 1
    \\
    D_{f_\alpha}(\mu, \nu) & \text{for } \alpha \in (0,+\infty) \backslash \{1\}
  \end{array}\right.
  \end{align*}
  with $f_\alpha : x \mapsto \frac{x^{\alpha} - 1}{\alpha - 1}$.
\end{definition}


\begin{lemma}
  \label{lem:hellingerAlpha_ne_top_of_lt_one}
  \lean{ProbabilityTheory.hellingerDiv_ne_top_of_lt_one}
  \leanok
  \uses{def:hellingerAlpha}
  For $\alpha \in [0, 1)$ and finite measures $\mu, \nu$, $\He_\alpha(\mu, \nu) < \infty$.
\end{lemma}

\begin{proof}\leanok
\uses{}
\end{proof}


\begin{lemma}
  \label{lem:hellingerAlpha_eq_integral}
  %\lean{}
  %\leanok
  \uses{def:hellingerAlpha}
  For $\alpha \in (0,1)\cup(1, \infty)$, $\mu$ a finite measure and $\nu$ a probability measure, if $\He_\alpha(\mu, \nu) < \infty$ then
  \begin{align*}
  \He_\alpha(\mu, \nu) = \frac{1}{\alpha - 1} \left( \int_x \left(\frac{d \mu}{d \nu}(x)\right)^\alpha \partial \nu - 1 \right)
  \: .
  \end{align*}
\end{lemma}

\begin{proof}%\leanok
\uses{}
\end{proof}


\begin{lemma}
  \label{lem:integral_rpow_rnDeriv_symm}
  \lean{ProbabilityTheory.integral_rpow_rnDeriv_symm}
  \leanok
  \uses{}
  For $\alpha \in (0,1)\cup(1, \infty)$, $\mu, \nu$ two sigma-finite measures,
  \begin{align*}
  \int_x \left(\frac{d \mu}{d \nu}(x)\right)^\alpha \partial \nu
  = \int_x \left(\frac{d \nu}{d \mu}(x)\right)^{1 - \alpha} \partial \mu
  \: .
  \end{align*}
\end{lemma}

\begin{proof}\leanok
\uses{}
\begin{align*}
  \int_x \left(\frac{d \mu}{d \nu}(x)\right)^\alpha \partial \nu
  &= \int_x \left(\frac{d \mu}{d (\mu + \nu)}(x) \left(\frac{d \nu}{d (\mu + \nu)}(x)\right)^{-1} \right)^\alpha \partial \nu \\
  &= \int_x \left(\frac{d \mu}{d (\mu + \nu)}(x) \left(\frac{d \nu}{d (\mu + \nu)}(x)\right)^{-1} \right)^\alpha \frac{d \nu}{d (\mu + \nu)}(x) \partial (\mu + \nu) \\
  &= \int_x \left(\frac{d \mu}{d (\mu + \nu)}(x) \left(\frac{d \nu}{d (\mu + \nu)}(x)\right)^{-1} \right)^{\alpha - 1} \frac{d \mu}{d (\mu + \nu)}(x) \partial (\mu + \nu) \\
  &= \int_x \left(\frac{d \mu}{d (\mu + \nu)}(x) \left(\frac{d \nu}{d (\mu + \nu)}(x)\right)^{-1} \right)^{\alpha - 1} \partial \mu \\
  &= \int_x \left(\frac{d \nu}{d \mu}(x)\right)^{1 - \alpha} \partial \mu
  \: .
  \end{align*}
\end{proof}


\begin{lemma}
  \label{lem:hellingerAlpha_symm}
  \lean{ProbabilityTheory.hellingerDiv_symm}
  \leanok
  \uses{def:hellingerAlpha}
  For $\alpha \in (0, 1)$ and finite measures $\mu, \nu$, $(1 - \alpha) \He_\alpha(\mu, \nu) = \alpha \He_{1 - \alpha}(\nu, \mu)$.
\end{lemma}

\begin{proof}\leanok
\uses{lem:integral_rpow_rnDeriv_symm}
Use Lemma~\ref{lem:integral_rpow_rnDeriv_symm}.
\end{proof}


\begin{definition}[Conditional Hellinger $\alpha$-divergence]
  \label{def:condHellingerAlpha}
  \lean{ProbabilityTheory.condHellingerDiv}
  \leanok
  \uses{def:condFDiv, lem:condFDiv_properties, lem:fDiv_properties}
  Let $\mu$ be a measure on $\mathcal X$ and $\kappa, \eta : \mathcal X \rightsquigarrow \mathcal Y$ be two kernels. The conditional Hellinger divergence of order $\alpha \in (0,+\infty) \backslash \{1\}$ between $\kappa$ and $\eta$ conditionally to $\mu$ is
  \begin{align*}
  \He_\alpha(\kappa, \eta \mid \mu) = D_{f_\alpha}(\kappa, \eta \mid \mu) \: ,
  \end{align*}
  for $f_\alpha : x \mapsto \frac{x^{\alpha} - 1}{\alpha - 1}$.
\end{definition}


\subsection{Properties inherited from f-divergences}

Since $\He_\alpha$ is an f-divergence, every inequality for f-divergences can be translated to $\He_\alpha$.

\begin{theorem}[Data-processing]
  \label{thm:hellingerAlpha_data_proc}
  \lean{ProbabilityTheory.hellingerDiv_comp_right_le}
  \leanok
  \uses{def:hellingerAlpha}
  Let $\alpha > 0$, $\mu, \nu$ be two finite measures on $\mathcal X$ and let $\kappa : \mathcal X \rightsquigarrow \mathcal Y$ be a Markov kernel.
  Then $\He_\alpha(\kappa \circ \mu, \kappa \circ \nu) \le \He_\alpha(\mu, \nu)$.
\end{theorem}

\begin{proof} \leanok
\uses{thm:fDiv_data_proc_2}
Apply Theorem~\ref{thm:fDiv_data_proc_2}.
\end{proof}


\begin{lemma}
  \label{lem:hellingerAlpha_nonneg}
  %\lean{}
  \leanok
  \uses{def:hellingerAlpha}
  Let $\mu, \nu$ be two probability measures. Then $\He_\alpha(\mu, \nu) \ge 0$.
\end{lemma}

\begin{proof} \leanok
\uses{lem:fDiv_nonneg}
Apply Lemma~\ref{lem:fDiv_nonneg}.
\end{proof}


\begin{lemma}
  \label{lem:hellingerAlpha_convex}
  %\lean{}
  %\leanok
  \uses{def:hellingerAlpha}
  $(\mu, \nu) \mapsto \He_\alpha(\mu, \nu)$ is convex.
\end{lemma}

\begin{proof}
\uses{thm:fDiv_convex}
Apply Theorem~\ref{thm:fDiv_convex}
\end{proof}



TODO: find a way to hide the following dummy lemma
\begin{lemma}[Dummy lemma: hellingerAlpha properties]
  \label{lem:hellingerAlpha_properties}
  %\lean{}
  \leanok
  \uses{def:hellingerAlpha, def:condHellingerAlpha}
  Dummy node to summarize properties of the Hellinger $\alpha$-divergence.
\end{lemma}

\begin{proof}\leanok
\uses{
  lem:hellingerAlpha_ne_top_of_lt_one,
  lem:hellingerAlpha_eq_integral,
  lem:hellingerAlpha_symm,
  thm:hellingerAlpha_data_proc,
  lem:hellingerAlpha_nonneg,
  lem:hellingerAlpha_convex
}
\end{proof}

%%FILE blueprint/src/sections/introduction.tex
\chapter*{Introduction}

The goal of this project is to formalize information divergences between probability measures, as well as results about error bounds for (sequential) hypothesis testing.
This is not an introductory text and we will not explain why the objects described here are useful, where they come from or what applications the results might have. For all those considerations, we refer the reader to the recent book~\cite{polyanskiy2024information}.

Our goal is to compile properties of information divergences and their uses in hypothesis testing. Most software libraries have tutorials, to learn how to use them, as well as a documentation listing all the available objects and methods of the library: in this project, we write the ``documentation'' of a part of the ``information theory library''.



\section*{This document and the Lean implementation}

This document is generated from a latex file written by the authors, thanks to the \texttt{leanblueprint} library developed by Patrick Massot (\url{https://github.com/PatrickMassot/leanblueprint}). It contains links to a Lean implementation of the results, but there is no technical guarantee that the results written here and the ones in the code match in any way.
In particular, nothing in this document is automatically generated from the code.
In summary, although the authors strive to make this document as close to the Lean implementation as possible, the only formally verified statements are those written in Lean.



\section*{Assumptions}

Information divergences are usually considered for probability measures, but we will be more granular about the assumptions we impose on the measures.
Since Radon-Nikodym derivatives don't have good properties unless the measures are sigma-finite and since those derivatives are the basic building block of all the information divergences we consider, every measure we consider will be supposed sigma-finite.
Many results additionally need that the measures are finite, and some will be about probability measures specifically.

In most results, the measures can be defined on a general measurable space, without additional assumptions.
However some statements, notably those involving conditional divergences, will require the sigma-algebra to be countably generated or the space to be standard Borel.
The countably generated assumption means that there is a countable set of sets that generates the sigma-algebra.
This is the case for example in any standard Borel space, which covers the vast majority of applications.



\section*{Notation}

$\mathcal X$, $\mathcal Y$ will denote measurable spaces. Additional assumptions such as $\mathcal X$ being standard Borel will be specified where necessary.

$\mu$, $\nu$ denote measures on a space $\mathcal X$. A measure $\mu$ is said to be finite if $\mu(\mathcal X) < \infty$, to be a probability measure if $\mu(\mathcal X) = 1$. A measure is s-finite if it can be written as a countable sum of finite measures, and is $\sigma$-finite if it can be written as a countable sum of finite measures which are pairwise mutually singular.
The set of measures on $\mathcal X$ is $\mathcal M(\mathcal X)$. The set of probability measures is denoted by $\mathcal P(\mathcal X)$.

For a function $g : \mathcal X \to \mathcal Y$ which is $a.e.$-measurable with respect to $\mu$, $g_* \mu$ denotes the pushforward or map of $\mu$ by $g$. This is a measure on $\mathcal Y$ such that $g_* \mu (A) = \mu(g^{-1} A)$.
For $g : \mathcal X \to \mathbb{R}_{+,\infty}$, $g \cdot \mu$ denotes the measure on $\mathcal X$ with density $g$ with respect to $\mu$. That is, $g \cdot \mu (A) = \int_{x \in A} g(x) \partial\mu$.
The Radon-Nikodym of a measure $\mu$ with respect to another measure $\mu$ is denoted by $\frac{d\mu}{d\nu}$. The singular part of $\mu$ with respect to $\mu$ is $\mu_{\perp \nu}$.

$\kappa, \eta$ are used to refer to transition kernels, and we write $\kappa : \mathcal X \rightsquigarrow \mathcal Y$ to indicate that $\kappa$ is a transition kernel from $\mathcal X$ to $\mathcal Y$. That is, $\kappa$ is a measurable map from $\mathcal X$ to the measures on $\mathcal Y$.
A kernel is finite if there exists $C < + \infty$ such that $\kappa(x, \mathcal Y) \le C$ for all $x \in \mathcal X$ (note that this is stronger than the condition that all $\kappa(x)$ are finite since it requires a uniform bound). It is s-finite if it can be written as a countable sum of finite kernels, and is called a Markov kernel if $\kappa(x)$ is a probability measure for all $x \in \mathcal X$.
See Appendix~\ref{app:kernels} for more details about kernels and Appendix~\ref{app:rnDeriv} for results about Radon-Nikodym derivatives of kernels.



\section*{Some references}

A very useful book (draft) about information theory and hypothesis testing: \cite{polyanskiy2024information} 

Main reference for the properties of the Rényi divergence: \cite{van2014renyi}


%%FILE blueprint/src/sections/jensen_shannon.tex
\section{Jensen-Shannon divergence}

\begin{definition}[Jensen-Shannon divergence]
  \label{def:jensenShannon}
  %\lean{}
  %\leanok
  \uses{def:KL, lem:kl_properties}
  The Jensen-Shannon divergence indexed by $\alpha \in (0,1)$ between two measures $\mu$ and $\nu$ is
  \begin{align*}
    \JS_\alpha(\mu, \nu) = \alpha \KL(\mu, \alpha \mu + (1 - \alpha)\nu) + (1 - \alpha) \KL(\nu, \alpha \mu + (1 - \alpha)\nu) \: .
  \end{align*}
\end{definition}


\begin{lemma}
  \label{lem:jensenShannon_eq_kl}
  %\lean{}
  %\leanok
  \uses{def:jensenShannon, def:KL}
  Let $\mu, \nu \in \mathcal P(\mathcal X)$ and let $\alpha \in (0, 1)$. Let $\pi_\alpha = (\alpha, 1 - \alpha) \in \mathcal P(\{0,1\})$ and let $P : \{0,1\} \rightsquigarrow \mathcal X$ be the kernel with $P(0) = \mu$ and $P(1) = \nu$. Then
  \begin{align*}
  \JS_\alpha(\mu, \nu) = \KL(\pi_\alpha \otimes P, \pi_\alpha \times (P \circ \pi_\alpha)) \: .
  \end{align*}
\end{lemma}

\begin{proof}%\leanok
\uses{}

\end{proof}


\begin{lemma}
  \label{lem:jensenShannon_eq_fDiv}
  %\lean{}
  %\leanok
  \uses{def:jensenShannon, def:KL}
  $\JS_\alpha$ is an $f$-divergence for $f(x) = \alpha x \log(x) - (\alpha x + 1 - \alpha) \log (\alpha x + 1 - \alpha)$~.
\end{lemma}

\begin{proof}%\leanok
\uses{lem:fDiv_add_smul_right, lem:fDiv_add_smul_right', lem:fDiv_add, lem:fDiv_mul, lem:fDiv_properties}
$(\mu, \nu) \mapsto \KL(\mu, \alpha \mu + (1 - \alpha) \nu)$ is an $f$-divergence for the function $x \mapsto x\log\left(\frac{x}{\alpha x + 1 - \alpha}\right)$ (Lemma~\ref{lem:fDiv_add_smul_right}).

$(\mu, \nu) \mapsto \KL(\nu, \alpha \mu + (1 - \alpha) \nu)$ is an $f$-divergence for the function $x \mapsto x\log\left(\frac{1}{\alpha x + 1 - \alpha}\right)$ (Lemma~\ref{lem:fDiv_add_smul_right'}).

By Lemmas~\ref{lem:fDiv_mul} and \ref{lem:fDiv_add}, we obtain the result.
\end{proof}


\begin{lemma}
  \label{lem:jensenShannon_symm}
  %\lean{}
  %\leanok
  \uses{def:jensenShannon}
  For $\alpha \in (0,1)$ and $\mu, \nu \in \mathcal M(\mathcal X)$,
  \begin{align*}
  \JS_\alpha(\mu, \nu) = \JS_{1 - \alpha}(\nu, \mu) \: .
  \end{align*}
\end{lemma}

\begin{proof}%\leanok
\uses{}
Immediate from the definition.
\end{proof}


\begin{lemma}
  \label{lem:jensenShannon_eq_inf_add_kl}
  %\lean{}
  %\leanok
  \uses{def:jensenShannon, def:KL}
  Let $\mu, \nu \in \mathcal P(\mathcal X)$ and let $\alpha \in (0, 1)$. Then
  \begin{align*}
  \JS_\alpha(\mu, \nu) = \inf_{\xi \in \mathcal P(\mathcal X)}\left( \alpha \KL(\mu, \xi) + (1 - \alpha)\KL(\nu, \xi) \right) \: .
  \end{align*}
  The infimum is attained at $\xi = \alpha \mu + (1 - \alpha) \nu$.
\end{lemma}

\begin{proof}%\leanok
\uses{}

\end{proof}


\begin{lemma}
  \label{lem:jensenShannon_eq_inf_kl}
  %\lean{}
  %\leanok
  \uses{def:jensenShannon, def:KL}
  Let $\mu, \nu \in \mathcal P(\mathcal X)$ and let $\alpha \in (0, 1)$. Let $\pi_\alpha = (\alpha, 1 - \alpha) \in \mathcal P(\{0,1\})$ and let $P : \{0,1\} \rightsquigarrow \mathcal X$ be the kernel with $P(0) = \mu$ and $P(1) = \nu$. Then
  \begin{align*}
  \JS_\alpha(\mu, \nu) = \inf_{\xi \in \mathcal P(\mathcal X)} \KL\left( \pi_\alpha \otimes P, \pi_\alpha \times \xi \right) \: .
  \end{align*}
  The infimum is attained at $\xi = \alpha \mu + (1 - \alpha) \nu$.
\end{lemma}

Compare with Lemma~\ref{lem:renyi_eq_inf_kl} on the Rényi divergence.

\begin{proof}%\leanok
\uses{lem:jensenShannon_eq_inf_add_kl}

\end{proof}


\begin{lemma}
  \label{lem:jensenShannon_prod_n}
  %\lean{}
  %\leanok
  \uses{def:jensenShannon}
  Let $\mu, \nu$ be two probability measures on $\mathcal X$. Let $n \in \mathbb{N}$ and write $\mu^{\otimes n}$ for the product measure on $\mathcal X^n$ of $n$ times $\mu$.
  Then $\JS_\alpha(\mu^{\otimes n}, \nu^{\otimes n}) \le n \JS_\alpha(\mu, \nu)$.
\end{lemma}

\begin{proof}%\leanok
\uses{lem:jensenShannon_eq_inf_add_kl, thm:kl_pi}
By Lemma~\ref{lem:jensenShannon_eq_inf_add_kl},
\begin{align*}
\JS_\alpha(\mu^{\otimes n}, \nu^{\otimes n})
&= \inf_{\xi \in \mathcal P(\mathcal X^n)}\left( \alpha \KL(\mu^{\otimes n}, \xi) + (1 - \alpha)\KL(\nu^{\otimes n}, \xi) \right)
\\
&\le \inf_{\zeta \in \mathcal P(\mathcal X)}\left( \alpha \KL(\mu^{\otimes n}, \zeta^{\otimes n}) + (1 - \alpha)\KL(\nu^{\otimes n}, \zeta^{\otimes n}) \right)
\\
&= n \inf_{\zeta \in \mathcal P(\mathcal X)}\left( \alpha \KL(\mu, \zeta) + (1 - \alpha)\KL(\nu, \zeta) \right)
\\
&= n \JS_\alpha(\mu, \nu)
\: .
\end{align*}

\end{proof}



\subsection{Comparisons with other divergences}

\begin{lemma}
  \label{lem:jensenShannon_le_hellingerAlpha}
  %\lean{}
  %\leanok
  \uses{def:jensenShannon, def:hellingerAlpha}
  For $\mu, \nu \in \mathcal P(\mathcal X)$ and $\alpha, \lambda \in (0,1)$~,
  \begin{align*}
  \JS_\alpha(\mu, \nu) \le \frac{(1 - \alpha)^{1 - \lambda} \alpha^\lambda}{\lambda^{1 - \lambda} (1 - \lambda)^\lambda} (1 - \lambda)\He_\lambda(\mu, \nu) \: .
  \end{align*}
  In particular,
  \begin{align*}
  \JS_\alpha(\mu, \nu)
  &\le \left(\frac{1 - \alpha}{\alpha}\right)^{1 - 2 \alpha} (1 - \alpha)\He_\alpha(\mu, \nu)
  \: , \\
  \JS_\alpha(\mu, \nu)
  &\le \alpha \He_{1 - \alpha}(\mu, \nu)
  \: .
  \end{align*}
\end{lemma}

\begin{proof}%\leanok
\uses{lem:fDiv_le_of_deriv2_le, lem:jensenShannon_eq_fDiv}
The main tool is Lemma~\ref{lem:fDiv_le_of_deriv2_le}.
$\JS_\alpha$ is an $f$-divergence with function $f(x) = \alpha x \log(x) - (\alpha x + 1 - \alpha) \log (\alpha x + 1 - \alpha)$. On probability measures, $(1 - \lambda)\He_\lambda$ is an $f$-divergence with function $g(x) = 1 - x^\lambda + \lambda(x-1)$.
These two functions satisfy $f(1) = 0$, $g(1) = 0$, $f'(1) = 0$, $g'(1) = 0$ and have second derivatives, hence Lemma~\ref{lem:fDiv_le_of_deriv2_le} applies and to obtain the inequality it suffices to compute $\sup_x\frac{f''(x)}{g''(x)}$~.
\begin{align*}
f''(x)
&= \frac{\alpha(1 - \alpha)}{x(\alpha x + 1 - \alpha)}
\: , \\
g''(x)
&= \lambda (1 - \lambda) x^{\lambda - 2}
\: .
\end{align*}
Let $h(x) = \frac{f''(x)}{g''(x)} = \frac{\alpha (1 - \alpha)}{\lambda (1 - \lambda)} \frac{x^{1 - \lambda}}{\alpha x + 1 - \alpha}$~. Its maximum is attained at a point where $h'(x) = 0$.
\begin{align*}
h'(x) = \frac{\alpha (1 - \alpha)}{\lambda (1 - \lambda)} \frac{(1 - \alpha)(1 - \lambda)x^{-\lambda} - \alpha \lambda x^{1 - \lambda}}{(\alpha x + 1 - \alpha)^2}
\: .
\end{align*}
The optimum for $h$ is attained at $x = \frac{(1 - \alpha)(1 - \lambda)}{\alpha \lambda}$ and the maximal value of $h$ is $\frac{(1 - \alpha)^{1 - \lambda} \alpha^\lambda}{\lambda^{1 - \lambda} (1 - \lambda)^\lambda}$~.
\end{proof}


\begin{corollary}
  \label{cor:jensenShannon_le_hellingerAlpha_of_le_half}
  %\lean{}
  %\leanok
  \uses{def:jensenShannon, def:hellingerAlpha}
  For $\mu, \nu \in \mathcal P(\mathcal X)$ and $\alpha \in (0,1)$, $\lambda \le 1/2$~,
  \begin{align*}
  \JS_\alpha(\mu, \nu)
  &\le \alpha^{1 - \lambda} \He_{1 - \lambda}(\mu, \nu)
  \: .
  \end{align*}
\end{corollary}

\begin{proof}%\leanok
\uses{lem:jensenShannon_le_hellingerAlpha}
By Lemma~\ref{lem:jensenShannon_le_hellingerAlpha} for $\alpha$ and $1 - \lambda$,
\begin{align*}
\JS_\alpha(\mu, \nu)
&\le \frac{(1 - \alpha)^{\lambda} \alpha^{1 - \lambda}}{\lambda^{1 - \lambda} (1 - \lambda)^\lambda} \lambda \He_{1 - \lambda}(\mu, \nu)
\\
&= \alpha^{1 - \lambda} (1 - \alpha)^\lambda \left(\frac{\lambda}{1 - \lambda}\right)^\lambda \He_{1 - \lambda}(\mu, \nu)
\: .
\end{align*}
To get the result, use $1 - \alpha \le 1$ and $\frac{\lambda}{1 - \lambda} \le 1$ since it is increasing in $\lambda$ with value 1 at $\lambda = 1/2$.

\end{proof}

%%FILE blueprint/src/sections/kernels.tex
\chapter{Probability transition kernels}
\label{app:kernels}

Consider, informally, a process where we first observe a random state (for example the result of a coin flip), and then, depending on the value of that first state, we produce another random value.
For example, if the coin lands on heads, we roll two dice and observe the sum. If on the other hand the coin lands on tails, we roll only one die and observe the result.
The law of the second observation is a function of the first observation.
This function from one space to measures on an other one is a transition kernel.
To get an usable definition for general measurable spaces, we need to enforce that the function is measurable.

\section{Definition}

\begin{definition}[Measurable space of measures]
  \label{def:measure_measurableSpace}
  \lean{MeasureTheory.Measure.instMeasurableSpace}
  \mathlibok
  \uses{}
  Let $\mathcal B$ be the Borel $\sigma$-algebra on $\mathbb{R}_{+,\infty}$. Let $\mathcal X$ be a measurable space.
  For a measurable set $s$ of $\mathcal X$, let $(\mu \mapsto \mu(s))^* \mathcal B$ be the $\sigma$-algebra on $\mathcal M(\mathcal X)$ defined by the comap of the evaluation function at $s$.
  Then $\mathcal M(\mathcal X)$ is a measurable space with $\sigma$-algebra $\bigsqcup_{s} (\mu \mapsto \mu(s))^* \mathcal B$ where the supremum is over all measurable sets $s$.
\end{definition}

It then makes sense to talk about a measurable function from a measurable space $\mathcal X$ to the measures on another measurable space $\mathcal Y$.
We now write more explicitly what it means for such a function to be measurable.

\begin{lemma}
  \label{lem:measurable_measure_fun}
  \lean{MeasureTheory.Measure.measurable_of_measurable_coe}
  \mathlibok
  \uses{def:measure_measurableSpace}
  Let $\mathcal X, \mathcal Y$ be measurable spaces, and let $f : \mathcal X \to \mathcal M(\mathcal Y)$ such that for all measurable sets $s$ of $\mathcal X$, the function $x \mapsto f(x)(s)$ is measurable.
  Then $f$ is measurable.
\end{lemma}

\begin{proof}\mathlibok
\uses{}
\end{proof}

We now have the tools needed to define probability transition kernels.

\begin{definition}[Kernel]
  \label{def:kernel}
  \lean{ProbabilityTheory.Kernel}
  \mathlibok
  \uses{def:measure_measurableSpace}
  Let $\mathcal X, \mathcal Y$ be two measurable spaces.
  A probability transition kernel (or simply kernel) from $\mathcal X$ to $\mathcal Y$ is a measurable map from $\mathcal X$ to $\mathcal M (\mathcal Y)$, the measurable space of measures on $\mathcal Y$.
  We write $\kappa : \mathcal X \rightsquigarrow \mathcal Y$ for a kernel $\kappa$ from $\mathcal X$ to $\mathcal Y$.
\end{definition}

Measures and measurable functions are special cases of kernels. A measure on $\mathcal Y$ can be seen as a constant kernel from any space to $\mathcal Y$ (for example from a space with 1 element to $\mathcal Y$).
A measurable function can be represented by a \emph{deterministic} kernel.

\begin{definition}[Deterministic kernel]
  \label{def:deterministic_kernel}
  \lean{ProbabilityTheory.Kernel.deterministic}
  \mathlibok
  \uses{def:kernel}
  The deterministic kernel defined by a measurable function $f : \mathcal X \to \mathcal Y$ is the kernel $d_f: \mathcal X \rightsquigarrow \mathcal Y$ defined by $d_f(x) = \delta_{f(x)}$, where for any $y \in \mathcal Y$, $\delta_y$ is the Dirac probability measure at $y$.
\end{definition}

A few deterministic kernels are particularly useful.
\begin{itemize}
  \item The identity kernel, deterministic kernel given by the identity function.
  \item The copy kernel $c_{\mathcal X} : \mathcal X \rightsquigarrow \mathcal X \times \mathcal X$, from the function $x \mapsto (x, x)$.
  \item The discard kernel $d_{\mathcal X}$ (TODO: conflicts with deterministic kernel notation). This is a kernel from $\mathcal X$ to $*$, the space with one point, corresponding to the constant function.
\end{itemize}

Kernels can be described by the integrals of measurable functions, as demonstrated in the following extensionality theorem.

\begin{theorem}
  \label{thm:kernel_ext}
  \lean{ProbabilityTheory.Kernel.ext_fun_iff}
  \mathlibok
  \uses{def:kernel}
  Two kernels $\kappa, \eta : \mathcal X \rightsquigarrow \mathcal Y$ are equal iff for all measurable functions $f : \mathcal Y \to \mathbb{R}_{+,\infty}$ and all $x \in \mathcal X$,
  \begin{align*}
  \int_y f(y) \partial \kappa(x) = \int_y f(y) \partial \eta(x) \: .
  \end{align*}
\end{theorem}

\begin{proof}\mathlibok
\uses{}

\end{proof}

\section{Classes of kernels}

\begin{definition}[Finite kernel]
  \label{def:finite_kernel}
  \lean{ProbabilityTheory.IsFiniteKernel}
  \mathlibok
  \uses{def:kernel}
  A kernel $\kappa : \mathcal X \rightsquigarrow \mathcal Y$ is said to be finite if there exists $C < +\infty$ such that for all $x \in \mathcal X$, $\kappa(x)(\mathcal Y) \le C$.
\end{definition}

Note the uniform bound: it is not enough to have $\kappa(x)(\mathcal Y) < \infty$ for all $x \in \mathcal X$. TODO: why? This is required for composition of finite kernels to still be finite (iirc): be precise about that.

\begin{definition}[Markov kernel]
  \label{def:markov_kernel}
  \lean{ProbabilityTheory.IsMarkovKernel}
  \mathlibok
  \uses{def:kernel}
  A kernel $\kappa : \mathcal X \rightsquigarrow \mathcal Y$ is said to be a Markov kernel if for all $x \in \mathcal X$, $\kappa(x)$ is a probability measure.
\end{definition}

A Markov kernel is finite, with constant $C = 1$.

\begin{definition}[s-finite kernel]
  \label{def:sFinite_kernel}
  \lean{ProbabilityTheory.IsSFiniteKernel}
  \mathlibok
  \uses{def:finite_kernel}
  A kernel $\kappa : \mathcal X \rightsquigarrow \mathcal Y$ is said to be a s-finite if it is equal to a countable sum of finite kernels.
\end{definition}

That last, most general class of kernels, is the right assumption to define operations on kernels like their product.



\section{Composition and product}

\begin{definition}[Composition-product]
  \label{def:kernel_compProd}
  \lean{ProbabilityTheory.Kernel.compProd}
  \mathlibok
  \uses{def:kernel, def:sFinite_kernel}
  Let $\kappa : \mathcal X \rightsquigarrow \mathcal Y$ and $\eta : (\mathcal X \times \mathcal Y) \rightsquigarrow \mathcal Z$ be two s-finite kernels.
  The composition-product of $\kappa$ and $\eta$ is a kernel $\kappa \otimes \eta : \mathcal X \rightsquigarrow (\mathcal Y \times \mathcal Z)$ such that for all measurable functions $f : \mathcal Y \times \mathcal Z \to \mathbb{R}_{+,\infty}$ and $x \in \mathcal X$~,
  \begin{align*}
  \int_{y,z} f(y,z) \partial(\kappa \otimes \eta)(x) = \int_y \int_z f(y,z) \partial\eta(x,y) \partial\kappa(x) \: .
  \end{align*}
\end{definition}

As for all other kernel definitions in this section, we omit the proof that this indeed defines a kernel and refer the reader to the Mathlib implementation.

TODO: explanation of that composition-product on an example.

The name \emph{composition-product} is not standard in the literature, where it can be called either composition or product depending on the source.
We reserve the name \emph{composition} for an operation that from kernels $\mathcal X \rightsquigarrow \mathcal Y$ and $\mathcal Y \rightsquigarrow \mathcal Z$ produces a kernel $\mathcal X \rightsquigarrow \mathcal Z$, and the name product for the one that takes kernels $\mathcal X \rightsquigarrow \mathcal Y$ and $\mathcal X \rightsquigarrow \mathcal Z$ and produces a kernel $\mathcal X \rightsquigarrow (\mathcal Y \times \mathcal Z)$.
Both are special cases of the composition-product, which justifies that name.

The composition-product is defined for $\kappa : \mathcal X \rightsquigarrow \mathcal Y$ and $\eta : \mathcal X \times \mathcal Y \rightsquigarrow \mathcal Z$, but since a kernel $\mathcal Y \rightsquigarrow \mathcal Z$ can be seen as a kernel $\mathcal X \times \mathcal Y \rightsquigarrow \mathcal Z$ independent of the first coordinate, we abuse notation and also write $\kappa \otimes \xi$ for $\xi : \mathcal Y \rightsquigarrow \mathcal Z$.

Since a measure is a constant kernel from the space with one element, we can also define the composition-product of a measure and a kernel. Note that this definition uses the abuse of notation we just introduced.

\begin{definition}[Composition-product of a measure and a kernel]
  \label{def:measure_compProd}
  \lean{MeasureTheory.Measure.compProd}
  \mathlibok
  \uses{def:kernel_compProd}
  Let $\mu \in \mathcal M(\mathcal X)$ be an s-finite measure and $\kappa : \mathcal X \rightsquigarrow \mathcal Y$ be an s-finite kernel.
  Let $\mathcal U$ be a measurable space with a unique element $u$. Let $\mu_k : \mathcal U \rightsquigarrow \mathcal X$ be the constant kernel with value $\mu$.
  The composition-product of $\mu$ and $\kappa$ is the measure on $\mathcal M(\mathcal X \times \mathcal Y)$ defined by $(\mu_k \otimes \kappa) (u)$~.
\end{definition}

\begin{definition}[Composition]
  \label{def:kernel_comp}
  \lean{ProbabilityTheory.Kernel.comp}
  \mathlibok
  \uses{def:kernel}
  Let $\kappa : \mathcal X \rightsquigarrow \mathcal Y$ and $\eta : \mathcal Y \rightsquigarrow \mathcal Z$ be two kernels.
  The composition of $\kappa$ and $\eta$ is the kernel $\eta \circ \kappa : \mathcal X \rightsquigarrow \mathcal Z$ such that for all measurable functions $f : \mathcal Z \to \mathbb{R}_{+,\infty}$ and all $x \in \mathcal X$,
  \begin{align*}
  \int_z f(z) \partial(\eta \circ \kappa)(x) = \int_y \int_z f(z) \partial\eta(y) \partial\kappa(x) \: .
  \end{align*}
\end{definition}

We can also define the composition of a measure and a kernel, by $\kappa \circ \mu = (\kappa \circ \mu_k) (u)$, where $\mathcal U$ is a measurable space with a unique element $u$ and $\mu_k : \mathcal U \rightsquigarrow \mathcal X$ is the constant kernel with value $\mu$.

\begin{definition}[Product]
  \label{def:kernel_prod}
  \lean{ProbabilityTheory.Kernel.prod}
  \mathlibok
  \uses{def:kernel, def:sFinite_kernel}
  Let $\kappa : \mathcal X \rightsquigarrow \mathcal Y$ and $\eta : \mathcal X \rightsquigarrow \mathcal Z$ be two s-finite kernels.
  The product of $\kappa$ and $\eta$ is the kernel $\kappa \times \eta : \mathcal X \rightsquigarrow \mathcal Y \times \mathcal Z$ such that for all measurable functions $f : \mathcal Y \times \mathcal Z \to \mathbb{R}_{+,\infty}$ and all $x \in \mathcal X$,
  \begin{align*}
  \int_{y,z} f(y,z) \partial(\kappa \times \eta)(x) = \int_y \int_z f(y,z) \partial\eta(x) \partial\kappa(x) \: .
  \end{align*}
\end{definition}

For $\kappa : \mathcal X \rightsquigarrow \mathcal Y$ and $\eta : \mathcal Y \rightsquigarrow \mathcal Z$ (independent of $\mathcal X$), the composition-product can be written as a composition: $\kappa \otimes \eta = (d_{\text{id}} \times \eta) \circ \kappa$ where $d_{\text{id}}$ is the deterministic kernel for the identity.
In particular, the composition-product of a measure and a kernel $\kappa$ is the composition of that measure with the kernel $d_{\text{id}} \times \kappa$.

\begin{definition}[Parallel product]
  \label{def:kernel_parallel_prod}
  \lean{ProbabilityTheory.Kernel.parallelComp}
  \mathlibok
  \uses{def:kernel, def:sFinite_kernel}
  Let $\kappa : \mathcal X \rightsquigarrow \mathcal Y$ and $\eta : \mathcal X' \rightsquigarrow \mathcal Y'$ be two s-finite kernels.
  The parallel product of $\kappa$ and $\eta$ is the kernel $\kappa \parallel \eta : \mathcal X \times \mathcal X' \rightsquigarrow \mathcal Y \times \mathcal Y'$ such that for all measurable functions $f : \mathcal Y \times \mathcal Y' \to \mathbb{R}_{+,\infty}$ and all $x \in \mathcal X \times \mathcal X'$,
  \begin{align*}
  \int_{y,y'} f(y,y') \partial(\kappa \parallel \eta)(x,x') = \int_y \int_{y'} f(y,y') \partial\eta(x') \partial\kappa(x) \: .
  \end{align*}
\end{definition}



\section{Copy-discard and Markov categories}

Measurable spaces and s-finite kernels are a copy-discard category (and in particular a symmetric monoidal category) with the following data:
\begin{itemize}
  \item Objects: measurable spaces
  \item Morphisms: s-finite kernels, with identity kernels and composition of kernels
  \item Tensor product of objects: product of measurable spaces $\mathcal X \times \mathcal Y$
  \item Tensor product of morphisms: parallel product $\kappa \parallel \eta$ 
  \item Unit: measurable space with one element (denoted by $*$)
  \item The associator, left and right unitors and braiding are given by the deterministic kernels of the corresponding equivalences between measurable spaces
  \item Copy morphism: $c_{\mathcal X}: \mathcal X \rightsquigarrow \mathcal X \times \mathcal X$
  \item Delete morphism: $d_{\mathcal X} : \mathcal X \rightsquigarrow *$
\end{itemize}

If we restrict the morphisms to Markov kernels, we get a Markov category (TODO ref), which is a copy-discard category in which for every morphism $\kappa$, $d_{\mathcal Y} \circ \kappa = d_{\mathcal X}$.

In that view of transition kernels, the following operations are not primitive, but derived from others:
\begin{itemize}
  \item Product: for $\kappa : \mathcal X \rightsquigarrow \mathcal Y$ and $\eta : \mathcal X \rightsquigarrow \mathcal Z$, $\kappa \times \eta = (\kappa \parallel \eta) \circ c_{\mathcal X}$
  \item Composition-product: for $\kappa : \mathcal X \rightsquigarrow \mathcal Y$ and $\eta : \mathcal X \times \mathcal Y \rightsquigarrow \mathcal Z$, $\kappa \otimes \eta = \mathrm{swap_{\mathcal Z, \mathcal Y}} \circ (\eta \parallel \mathrm{id}) \circ \alpha_{\mathcal X, \mathcal Y, \mathcal Y} \circ (\mathrm{id} \parallel c_{\mathcal Y}) \circ (\mathrm{id} \parallel \kappa) \circ c_{\mathcal X}$ (where $\alpha$ is the associator and $\mathrm{swap}$ is the braiding).
\end{itemize}



\section{Ordering}

\begin{definition}[Blackwell sufficiency]
  \label{def:blackwellOrder}
  %\lean{}
  %\leanok
  \uses{def:kernel, def:markov_kernel}
  We define a partial order on kernels by the following. Let $\kappa : \mathcal X \rightsquigarrow \mathcal Y$ and $\eta : \mathcal X \rightsquigarrow \mathcal Z$ (with same domain as $\kappa$).
  Then $\kappa$ is Blackwell sufficient for $\eta$, denoted by $\eta \le_B \kappa$, if there exists a Markov kernel $\xi : \mathcal Y \rightsquigarrow \mathcal Z$ such that $\eta = \xi \circ \kappa$.
\end{definition}



\section{Disintegration}

A kernel $\eta : \mathcal X \times \mathcal Y \rightsquigarrow \mathcal Z$ is said to disintegrate a kernel $\kappa : \mathcal X \rightsquigarrow \mathcal Y \times \mathcal Z$ if $\kappa = \kappa_{Y} \otimes \eta$.

TODO: if a disintegrating kernel $\eta$ exits for $\kappa$, the measure $\eta(x,y)$ is unique for all $x \in \mathcal X$ and $\kappa_{Y}(x)$-almost every $y \in \mathcal Y$.

TODO: if a disintegrating kernel $\eta$ exits for $\kappa$, the measure $\eta(x,y)$ is a probability measure for all $x \in \mathcal X$ and $\kappa_{Y}(x)$-almost every $y \in \mathcal Y$.

\begin{theorem}[Disintegration in standard Borel spaces]
  \label{thm:disintegration}
  \lean{ProbabilityTheory.Kernel.condKernel, ProbabilityTheory.Kernel.disintegrate}
  \mathlibok
  \uses{def:kernel, def:markov_kernel}
  If either $\mathcal X$ is countable of $\mathcal Y$ has a countably generated $\sigma$-algebra, and $\mathcal Z$ is standard Borel, then for every s-finite kernel $\kappa : \mathcal X \rightsquigarrow \mathcal Y \times \mathcal Z$, there exists a Markov kernel $\kappa_{Z \mid Y} : \mathcal X \times \mathcal Y \rightsquigarrow \mathcal Z$ that disintegrates $\kappa$.
\end{theorem}

TODO: the result in Mathlib is about finite kernels, but s-finite should be enough?

\begin{proof}\mathlibok
\uses{}
\end{proof}

TODO: as a consequence, in standard Borel spaces we also have disintegration of measures.



TODO: find a way to hide the following dummy lemma
\begin{lemma}[Dummy lemma: kernel properties]
  \label{lem:kernel_properties}
  %\lean{}
  \leanok
  \uses{
    def:measure_measurableSpace,
    def:kernel,
    def:deterministic_kernel,
    def:finite_kernel,
    def:markov_kernel,
    def:sFinite_kernel,
    def:kernel_compProd,
    def:kernel_comp,
    def:kernel_prod,
    def:kernel_parallel_prod,
    def:measure_compProd,
  }
  Dummy node to summarize kernel properties.
\end{lemma}

\begin{proof}\leanok
\uses{
  lem:measurable_measure_fun,
  thm:kernel_ext,
  thm:disintegration,
}
\end{proof}




\section{Bayesian inverse of a kernel}

\begin{definition}
  \label{def:bayesInv}
  \lean{ProbabilityTheory.posterior}
  \mathlibok
  \uses{lem:kernel_properties}
  For $\mu \in \mathcal M(\mathcal X)$ and $\kappa : \mathcal X \rightsquigarrow \mathcal Y$, a Bayesian inverse of $\kappa$ is a Markov kernel $\kappa_\mu^\dagger : \mathcal Y \rightsquigarrow \mathcal X$ such that $\mu \otimes \kappa = ((\kappa \circ \mu) \otimes \kappa_\mu^\dagger)_\leftrightarrow$ in which $(\cdot)_\leftrightarrow$ denotes swapping the two coordinates.
  If such an inverse exists it is unique up to a $(\kappa \circ \mu)$-null set, and we talk about \emph{the} Bayesian inverse of $\kappa$ with respect to $\mu$.
\end{definition}

The bayesian interpretation of this definition is that $\mu$ represents a prior over a parameter and $\kappa$ gives the distribution of data given the parameter. $\kappa \circ \mu$ is the distribution of the data, and $\kappa_\mu^\dagger$ gives the posterior distribution of the parameter given the data.

The measure $\rho = \mu \otimes \kappa$ is a joint probability distribution over $\mathcal X \times \mathcal Y$ with marginals $\rho_X = \mu$ and $\rho_Y = \kappa \circ \mu$ (if $\kappa$ is Markov). By construction, that joint probability distribution can be written as the composition-product of the first marginal $\rho_X$ and a kernel $\kappa$. The Bayesian inverse, if it exists, gives a way to write the joint as a composition-product of the second marginal $\rho_Y$ and a kernel. 

If the joint space is such that every product distribution can be disintegrated (see Theorem~\ref{thm:disintegration}), then the Bayesian inverse exists.

\begin{lemma}[Existence of the Bayesian inverse]
  \label{lem:exists_bayesInv}
  \lean{ProbabilityTheory.posterior, ProbabilityTheory.compProd_posterior_eq_map_swap}
  \mathlibok
  \uses{def:bayesInv}
  For $\mathcal X$ standard Borel, $\mu$ and $\kappa$ s-finite, the Bayesian inverse of $\kappa$ with respect to $\mu$ exists and is obtained by disintegration of the measure $\mu \otimes \kappa$ on $\mathcal X \times \mathcal Y$ into a measure $\kappa \circ \mu \in \mathcal M(\mathcal Y)$ and a Markov kernel $\kappa_\mu^\dagger : \mathcal Y \rightsquigarrow \mathcal X$.
\end{lemma}

\begin{proof}\mathlibok
\uses{thm:disintegration}

\end{proof}


See \cite{clerc2017pointless} and \cite{dahlqvist2018borel} for a category theory point of view on Bayesian inversion.

\begin{lemma}[a.e.-uniqueness of the Bayesian inverse]
  \label{lem:eq_bayesInv_of_compProd_eq}
  \lean{ProbabilityTheory.ae_eq_posterior_of_compProd_eq}
  \mathlibok
  \uses{def:bayesInv}
  Let $\mu \in \mathcal M(\mathcal X)$ be a finite measure and let $\kappa : \mathcal X \rightsquigarrow \mathcal Y$ be a finite kernel.
  If $\eta : \mathcal Y \rightsquigarrow \mathcal X$ is such that $\mu \otimes \kappa = ((\kappa \circ \mu) \otimes \eta)_\leftrightarrow$, then $\eta(y) = \kappa_\mu^\dagger(y)$ for $(\kappa \circ \mu)$-almost all $y$.
\end{lemma}

\begin{proof}\mathlibok
\uses{}
TODO: uniqueness of the conditional kernel.
\end{proof}


\begin{lemma}
  \label{lem:bayesInv_comp_self}
  \lean{ProbabilityTheory.posterior_comp_self}
  \mathlibok
  \uses{def:bayesInv}
  For $\mu \in \mathcal M(\mathcal X)$ s-finite, $\kappa : \mathcal X \rightsquigarrow \mathcal Y$ a Markov kernel and $\kappa_\mu^\dagger$ the Bayesian inverse of $\kappa$ with respect to $\mu$, these objects satisfy the equality $\kappa_\mu^\dagger \circ \kappa \circ \mu = \mu$.
\end{lemma}

\begin{proof}\mathlibok
\uses{}
The measure $\kappa_\mu^\dagger \circ (\kappa \circ \mu)$ is the projection on $\mathcal X$ of $(\kappa \circ \mu) \otimes \kappa_\mu^\dagger$. But $(\kappa \circ \mu) \otimes \kappa_\mu^\dagger = (\mu \otimes \kappa)_\leftrightarrow$ by definition of the Bayesian inverse. Since $\kappa$ is a Markov kernel, the projection on $\mathcal X$ of that measure is simply $\mu$.
\end{proof}


\begin{lemma}
  \label{lem:bayesInv_self}
  \lean{ProbabilityTheory.posterior_posterior}
  \mathlibok
  \uses{def:bayesInv}
  For $\mathcal X$ and $\mathcal Y$ two standard Borel spaces, $\mu \in \mathcal M(\mathcal X)$ s-finite and $\kappa : \mathcal X \rightsquigarrow \mathcal Y$ a Markov kernel, $(\kappa_\mu^\dagger)_{\kappa \circ \mu}^\dagger = \kappa$ ($\mu$-a.e.).
\end{lemma}

\begin{proof}\mathlibok
\uses{lem:bayesInv_comp_self, lem:eq_bayesInv_of_compProd_eq}
By uniqueness of the disintegration, it suffices to show that $(\kappa \circ \mu) \otimes \kappa_\mu^\dagger = ((\kappa_\mu^\dagger \circ \kappa \circ \mu) \otimes \kappa)_\leftrightarrow$~.
By Lemma~\ref{lem:bayesInv_comp_self}, $\kappa_\mu^\dagger \circ \kappa \circ \mu = \mu$ and we need to prove $(\kappa \circ \mu) \otimes \kappa_\mu^\dagger = (\mu \otimes \kappa)_\leftrightarrow$~.
This is true by definition of $\kappa_\mu^\dagger$~.
\end{proof}


\begin{lemma}
  \label{lem:bayesInv_id}
  \lean{ProbabilityTheory.posterior_id}
  \mathlibok
  \uses{def:bayesInv}
  Let $\mu \in \mathcal M (\mathcal X)$ and let $\textup{id} : \mathcal X \rightsquigarrow \mathcal X$ be the identity kernel. Then $\textup{id}_\mu^\dagger = \textup{id}$.
\end{lemma}

\begin{proof}\mathlibok
\uses{lem:eq_bayesInv_of_compProd_eq}
It suffices to show that $\mu \otimes \textup{id} = ((\textup{id} \circ \mu) \otimes \textup{id})_\leftrightarrow$~, which is clear.
\end{proof}


\begin{lemma}
  \label{lem:bayesInv_comp}
  \lean{ProbabilityTheory.posterior_comp}
  \mathlibok
  \uses{def:bayesInv,def:kernel_comp}
  Let $\mu \in \mathcal M(\mathcal X)$, $\kappa : \mathcal X \rightsquigarrow \mathcal Y$ and $\eta : \mathcal Y \rightsquigarrow \mathcal Z$. Then $(\eta \circ \kappa \circ \mu)$-a.e.,
  \begin{align*}
  (\eta \circ \kappa)_\mu^\dagger = \kappa_{\mu}^\dagger \circ \eta_{\kappa \circ \mu}^\dagger
  \: .
  \end{align*}
\end{lemma}

\begin{proof}\mathlibok
\uses{lem:eq_bayesInv_of_compProd_eq}
It suffices to show that $\mu \otimes (\eta \circ \kappa) = ((\eta \circ \kappa \circ \mu) \otimes (\kappa_{\mu}^\dagger \circ \eta_{\kappa \circ \mu}^\dagger))_\leftrightarrow$.

TODO: transcribe the code.
\end{proof}


\begin{lemma}
  \label{lem:rnDeriv_bayesInv}
  \lean{ProbabilityTheory.rnDeriv_posterior}
  \mathlibok
  \uses{def:bayesInv}
  Let $\mu \in \mathcal M(\mathcal X)$ and $\kappa : \mathcal X \rightsquigarrow \mathcal Y$ such that $\kappa_\mu^\dagger$ exists and such that $\kappa(x) \ll \kappa \circ \mu$ almost surely.
  Then for $\mu$-almost all $x$ and $(\kappa \circ \mu)$-almost all $y$,
  \begin{align*}
  \frac{d \kappa_\mu^\dagger(y)}{d \mu}(x) = \frac{d \kappa(x)}{d(\kappa \circ \mu)}(y)
  \: .
  \end{align*}
\end{lemma}

If $\mathcal{X}$ is countable, then we do have $\kappa(x) \ll \kappa \circ \mu$ almost surely.

\begin{proof}\mathlibok
\uses{}
$\mu$-almost all $x$ and $(\kappa \circ \mu)$-almost all $y$ is the same as $(\mu \times (\kappa \circ \mu))$-almost all $(x, y)$.
It suffices to show that for all measurable sets of $\mathcal X \times \mathcal Y$, the integrals with respect to that product measure of the two sides of the equality on the set are equal. Furthermore, it is sufficient to show it for sets of the form $s \times t$ since those are a $\pi$-system.

Since $\kappa(x) \ll \kappa \circ \mu$ almost surely, we have $\mu \otimes \kappa \ll \mu \times (\kappa \circ \mu)$, hence $(\kappa \circ \mu) \otimes \kappa_\mu^\dagger \ll (\kappa \circ \mu) \times \mu$.
We get that $\kappa_\mu^\dagger(y) \ll \mu$ almost surely.
\begin{align*}
\int_{x \in s, y \in t} \frac{d \kappa_\mu^\dagger(y)}{d \mu}(x) \partial (\mu \times (\kappa \circ \mu))
&= \int_{y \in t} \int_{x \in s} \frac{d \kappa_\mu^\dagger(y)}{d \mu}(x) \partial \mu \partial (\kappa \circ \mu)
\\
&= \int_{y \in t} \kappa_\mu^\dagger(y)(s) \partial (\kappa \circ \mu)
\\
&= ((\kappa \circ \mu) \otimes \kappa_\mu^\dagger) (t \times s)
\: .
\end{align*}

\begin{align*}
\int_{x \in s, y \in t} \frac{d \kappa(x)}{d(\kappa \circ \mu)}(y) \partial (\mu \times (\kappa \circ \mu))
&= \int_{x \in s} \int_{y \in t} \frac{d \kappa(x)}{d(\kappa \circ \mu)}(y) \partial (\kappa \circ \mu) \partial \mu
\\
&= \int_{x \in s} \kappa(x)(t) \partial \mu
\\
&= (\mu \otimes \kappa) (s \times t)
\: .
\end{align*}

By the defining property of the Bayesian inverse, those two quantities are equal.
\end{proof}



TODO: find a way to hide the following dummy lemma
\begin{lemma}[Dummy lemma: bayesInv properties]
  \label{lem:bayesInv_properties}
  %\lean{}
  \leanok
  \uses{
    def:bayesInv
  }
  Dummy node to summarize properties of the Bayesian inverse.
\end{lemma}

\begin{proof}\leanok
\uses{
  lem:exists_bayesInv,
  lem:eq_bayesInv_of_compProd_eq,
  lem:bayesInv_comp_self,
  lem:bayesInv_self,
  lem:bayesInv_id,
  lem:bayesInv_comp,
  lem:rnDeriv_bayesInv
}
\end{proof}

%%FILE blueprint/src/sections/kl_divergence.tex
\section{Kullback-Leibler divergence}

\begin{definition}[Kullback-Leibler divergence]
  \label{def:KL}
  \lean{InformationTheory.klDiv}
  \mathlibok
  %\uses{}
  Let $\mu, \nu$ be two measures on $\mathcal X$. The Kullback-Leibler divergence between $\mu$ and $\nu$ is
  \begin{align*}
  \KL(\mu, \nu) = \mu\left[\log \frac{d \mu}{d \nu}\right]
  \end{align*}
  if $\mu \ll \nu$ and $x \mapsto \log \frac{d \mu}{d \nu}(x)$ is $\mu$-integrable and $+\infty$ otherwise.
\end{definition}


\begin{lemma}
  \label{lem:kl_eq_fDiv}
  \lean{ProbabilityTheory.klDiv_eq_fDiv}
  \leanok
  \uses{def:KL, def:fDiv, lem:fDiv_properties}
  $\KL(\mu, \nu) = D_f(\mu, \nu)$ for $f: x \mapsto x \log x$.
\end{lemma}

\begin{proof}\leanok
Simple computation.
\end{proof}


\begin{definition}[Conditional Kullback-Leibler divergence]
  \label{def:condKL}
  \lean{ProbabilityTheory.condKL}
  \leanok
  \uses{def:KL}
  Let $\mu$ be a measure on $\mathcal X$ and $\kappa, \eta : \mathcal X \rightsquigarrow \mathcal Y$ be two kernels. The conditional Kullback-Leibler divergence between $\kappa$ and $\eta$ with respect to $\mu$ is
  \begin{align*}
  \KL(\kappa, \eta \mid \mu) = \mu\left[x \mapsto \KL(\kappa(x), \eta(x))\right]
  \end{align*}
  if $x \mapsto \KL(\kappa(x), \eta(x))$ is $\mu$-integrable and $+\infty$ otherwise.
\end{definition}


\begin{lemma}
  \label{lem:condKL_eq_condFDiv}
  \lean{ProbabilityTheory.klDiv_eq_fDiv}
  \leanok
  \uses{def:condKL, def:condFDiv, lem:condFDiv_properties}
  $\KL(\kappa, \eta \mid \mu) = D_f(\kappa, \eta \mid \mu)$ for $f: x \mapsto x \log x$.
\end{lemma}

\begin{proof}\leanok
Simple computation.
\end{proof}



\subsection{Properties inherited from f-divergences}

Since $\KL$ is an f-divergence, every inequality for f-divergences can be translated to $\KL$.

\begin{lemma}
  \label{lem:kl_self}
  \lean{InformationTheory.klDiv_self}
  \mathlibok
  \uses{def:KL}
  $\KL(\mu, \mu) = 0$.
\end{lemma}

\begin{proof}\mathlibok
  \uses{lem:kl_eq_fDiv, lem:fDiv_self}
  KL is an f-divergence thanks to Lemma~\ref{lem:kl_eq_fDiv}, then apply Lemma~\ref{lem:fDiv_self}.
\end{proof}

\begin{lemma}
  \label{lem:condKL_const}
  \lean{ProbabilityTheory.condKL_const}
  \leanok
  \uses{def:KL, def:condKL}
  Let $\mu, \nu$ be finite measures on $\mathcal X$, $\xi$ be a finite measure on $\mathcal Y$. 
  Then $\KL(x \mapsto \mu, x \mapsto \nu \mid \xi) = \KL(\mu, \nu) \xi (\mathcal X)$.
\end{lemma}

\begin{proof} \leanok
  \uses{lem:kl_eq_fDiv, lem:condFDiv_const, lem:condKL_eq_condFDiv}
  KL is an f-divergence thanks to Lemma~\ref{lem:kl_eq_fDiv} and Lemma~\ref{lem:condKL_eq_condFDiv}, then apply Lemma~\ref{lem:condFDiv_const}.
\end{proof}


\begin{theorem}[Marginals]
  \label{thm:kl_fst_le}
  \lean{ProbabilityTheory.klDiv_fst_le, ProbabilityTheory.klDiv_snd_le}
  \leanok
  \uses{def:KL}
  Let $\mu$ and $\nu$ be two measures on $\mathcal X \times \mathcal Y$, and let $\mu_X, \nu_X$ be their marginals on $\mathcal X$.
  Then $\KL(\mu_X, \nu_X) \le \KL(\mu, \nu)$.
  Similarly, for $\mu_Y, \nu_Y$ the marginals on $\mathcal Y$, $\KL(\mu_Y, \nu_Y) \le \KL(\mu, \nu)$.
\end{theorem}

\begin{proof} \leanok
\uses{thm:fDiv_fst_le_2, lem:kl_eq_fDiv}
Apply Theorem~\ref{thm:fDiv_fst_le_2}.
\end{proof}


\begin{theorem}[Data-processing]
  \label{thm:kl_data_proc}
  \lean{InformationTheory.klDiv_comp_right_le}
  \mathlibok
  \uses{def:KL}
  Let $\mu, \nu$ be two measures on $\mathcal X$ and let $\kappa : \mathcal X \rightsquigarrow \mathcal Y$ be a Markov kernel.
  Then $\KL(\kappa \circ \mu, \kappa \circ \nu) \le \KL(\mu, \nu)$.
\end{theorem}

\begin{proof}\mathlibok
\uses{thm:fDiv_data_proc_2, lem:kl_eq_fDiv}
Apply Theorem~\ref{thm:fDiv_data_proc_2}.
\end{proof}


\begin{lemma}
  \label{lem:kl_ge}
  \lean{InformationTheory.mul_log_le_klDiv}
  \mathlibok
  \uses{def:KL}
  Let $\mu, \nu$ be two finite measures on $\mathcal X$. Then $\KL(\mu, \nu) \ge \mu(\mathcal X) \log \frac{\mu(\mathcal X)}{\nu(\mathcal X)}$.
\end{lemma}

\begin{proof}\mathlibok
\uses{lem:le_fDiv, lem:kl_eq_fDiv}
\end{proof}


\begin{lemma}[Gibbs' inequality]
  \label{lem:kl_nonneg}
  %\lean{}
  \leanok
  \uses{def:KL}
  Let $\mu, \nu$ be two probability measures. Then $\KL(\mu, \nu) \ge 0$.
\end{lemma}

\begin{proof}\leanok
\uses{lem:kl_ge}
Apply Lemma~\ref{lem:kl_ge} and use $\mu(\mathcal X) = \nu (\mathcal X)$.
\end{proof}


\begin{lemma}[Converse Gibbs' inequality]
  \label{lem:kl_eq_zero_iff}
  \lean{InformationTheory.klDiv_eq_zero_iff}
  \mathlibok
  \uses{def:KL}
  Let $\mu, \nu$ be two probability measures. Then $\KL(\mu, \nu) = 0$ if and only if $\mu = \nu$.
\end{lemma}

\begin{proof}\mathlibok
\uses{lem:kl_eq_fDiv, lem:fDiv_eq_zero_iff}
KL is an f-divergence thanks to Lemma~\ref{lem:kl_eq_fDiv}, then apply Lemma~\ref{lem:fDiv_eq_zero_iff}.
\end{proof}


\begin{lemma}
  \label{lem:kl_convex}
  %\lean{}
  %\leanok
  \uses{def:KL}
  $(\mu, \nu) \mapsto \KL(\mu, \nu)$ is convex.
\end{lemma}

\begin{proof}
\uses{thm:fDiv_convex, lem:kl_eq_fDiv}
Apply Theorem~\ref{thm:fDiv_convex}
\end{proof}


%TODO : this may be moved to another section about integrability
\begin{lemma}
  \label{lem:integrable_llr_compProd_iff}
  \lean{ProbabilityTheory.integrable_llr_compProd_iff}
  \leanok
  \uses{}
  Let $\mu, \nu$ be two finite measures on $\mathcal X$ and $\kappa, \eta : \mathcal X \rightsquigarrow \mathcal Y$ be two Markov kernels such that $\mu \otimes \kappa \ll \nu \otimes \eta$. Then $p \mapsto \log \frac{d \mu \otimes \kappa}{d \nu \otimes \eta}(p)$ is $\mu \otimes \kappa$-integrable if and only if the following hold:
  \begin{enumerate}
    \item $x \mapsto \log \frac{d \mu}{d \nu}(x)$ is $\mu$-integrable \label{:1}
    \item $x \mapsto \int_y  \log \frac{d \kappa(x)}{d \eta(x)}(y) \partial \kappa(x)$ is $\mu$-integrable \label{:2}
    \item for $\mu$-almost all $x$, $y \mapsto \log \frac{d \kappa(x)}{d \eta(x)}(y)$ is $\kappa(x)$-integrable \label{:3}
  \end{enumerate}
\end{lemma}

\begin{proof}
\uses{def:KL, lem:integrable_fDiv_compProd_iff}

Note that from the hypothesis it easily follows that $\mu \ll \nu$ and $\mu$-a.e.\ $\kappa(x) \ll \eta(x)$.

Applying Lemma~\ref{lem:integrable_fDiv_compProd_iff} we know that $\log \frac{d \mu \otimes \kappa}{d \nu \otimes \eta}$ is $\nu \otimes \eta$-integrable iff the following hold:

\begin{enumerate}
    \item[i.] for $\nu$-almost all $x$, $y \mapsto \frac{d \mu}{d \nu}(x) \frac{d \kappa(x)}{d \eta(x)}(y) \log \left(\frac{d \mu}{d \nu}(x) \frac{d \kappa(x)}{d \eta(x)}(y) \right) $ is $\eta(x)$-integrable \label{:i}
    \item[ii.] $x \mapsto  \int_y \frac{d \mu}{d \nu}(x) \frac{d \kappa(x)}{d \eta(x)}(y) \log \left(\frac{d \mu}{d \nu}(x) \frac{d \kappa(x)}{d \eta(x)}(y) \right) \partial \eta (x)$ is $\nu$-integrable \label{:ii}
\end{enumerate}

{\bfseries(i.$\Leftarrow$)} 3.\ implies that $\nu$-a.e.\ $y \mapsto \frac{d \mu}{d \nu}(x) \log \frac{d \kappa(x)}{d \eta(x)}(y)$ is $\kappa(x)$-integrable, hence $y \mapsto \frac{d \mu}{d \nu}(x) \frac{d \kappa(x)}{d \eta(x)}(y) \log \frac{d \kappa(x)}{d \eta(x)}(y)$ is $\eta(x)$-integrable $\nu$-a.e.\ 
Moreover $y \mapsto \frac{d \kappa(x)}{d \eta(x)}(y)$ is $\eta(x)$-integrable, then, since $\frac{d \mu}{d \nu}(x) \log \frac{d \mu}{d \nu}(x)$ is constant in $y$, $y \mapsto \frac{d \mu}{d \nu}(x) \frac{d \kappa(x)}{d \eta(x)}(y) \log \frac{d \mu}{d \nu}(x)$ is $\eta(x)$-integrable. 
This allows us to conclude that $y \mapsto \frac{d \mu}{d \nu}(x) \frac{d \kappa(x)}{d \eta(x)}(y) \log \frac{d \kappa(x)}{d \eta(x)}(y) + \frac{d \mu}{d \nu}(x) \frac{d \kappa(x)}{d \eta(x)}(y) \log \frac{d \mu}{d \nu}(x) = \frac{d \mu}{d \nu}(x) \frac{d \kappa(x)}{d \eta(x)}(y) \log \left(\frac{d \mu}{d \nu}(x) \frac{d \kappa(x)}{d \eta(x)}(y) \right)$ is $\eta(x)$-integrable $\nu$-a.e.

{\bfseries(ii.$\Leftarrow$)} 1.\ implies that $x \mapsto \frac{d \mu}{d \nu}(x) \log \frac{d \mu}{d \nu}(x)$ is $\nu$-integrable. 
2.\ implies that $x \mapsto \frac{d \mu}{d \nu}(x) \int_y \frac{d \kappa(x)}{d \eta(x)}(y) \log \frac{d \kappa(x)}{d \eta(x)}(y) \partial \eta(x)$ is $\nu$-integrable. 
Therefore $$x \mapsto \frac{d \mu}{d \nu}(x) \log \frac{d \mu}{d \nu}(x) + \frac{d \mu}{d \nu}(x) \int_y \frac{d \kappa(x)}{d \eta(x)}(y) \log \frac{d \kappa(x)}{d \eta(x)}(y) \partial \eta(x) = \int_y \frac{d \mu}{d \nu}(x) \frac{d \kappa(x)}{d \eta(x)}(y) \log \left(\frac{d \mu}{d \nu}(x) \frac{d \kappa(x)}{d \eta(x)}(y) \right) \partial \eta (x)$$ is $\nu$-integrable.

{\bfseries($\Rightarrow$1.)} This follows from the hypothesis applying a more general Lemma for f divergences, see ProbabilityTheory.integrable\_f\_rnDeriv\_of\_integrable\_compProd' in the lean code.
% TODO : add the lemma ProbabilityTheory.integrable_f_rnDeriv_of_integrable_compProd' to the blueprint and reference it here

{\bfseries($\Rightarrow$3.)} i.\ implies that $\mu$-a.e.\ $y \mapsto \frac{d \kappa(x)}{d \eta(x)}(y) \log \left(\frac{d \mu}{d \nu}(x) \frac{d \kappa(x)}{d \eta(x)}(y) \right)$ is $\eta(x)$-integrable, hence $y \mapsto \log \left(\frac{d \mu}{d \nu}(x) \frac{d \kappa(x)}{d \eta(x)}(y) \right) = \log \frac{d \mu}{d \nu}(x) + \log \frac{d \kappa(x)}{d \eta(x)}(y)$ is $\kappa(x)$-integrable. Then, since $\log \frac{d \mu}{d \nu}(x)$ is constant in $y$ and $\kappa(x)$ is a finite measure, $y \mapsto \log \frac{d \kappa(x)}{d \eta(x)}(y)$ is $\kappa(x)$-integrable.

{\bfseries($\Rightarrow$2.)} ii.\ implies that $x \mapsto  \int_y \frac{d \kappa(x)}{d \eta(x)}(y) \log \left(\frac{d \mu}{d \nu}(x) \frac{d \kappa(x)}{d \eta(x)}(y) \right) \partial \eta (x) = \int_y \log \left(\frac{d \mu}{d \nu}(x) \frac{d \kappa(x)}{d \eta(x)}(y) \right) \partial \kappa (x)$ is $\mu$-integrable. 
Thanks to 3.\ we can split the integral:
$$\int_y \log \left(\frac{d \mu}{d \nu}(x) \frac{d \kappa(x)}{d \eta(x)}(y) \right) \partial \kappa (x)
= \int_y \log \frac{d \mu}{d \nu}(x) \partial \kappa (x) + \int_y \log \frac{d \kappa(x)}{d \eta(x)}(y) \partial \kappa (x) = \log \frac{d \mu}{d \nu}(x) + \int_y \log \frac{d \kappa(x)}{d \eta(x)}(y) \partial \kappa (x)$$
Then using 1.\ we can conclude that $x \mapsto \int_y \log \frac{d \kappa(x)}{d \eta(x)}(y) \partial \kappa (x)$ is $\mu$-integrable.

\end{proof}



\subsection{Chain rule and tensorization}

\begin{theorem}
  \label{thm:kl_compProd_aux}
  %\lean{}
  %\leanok
  \uses{def:KL}
  Let $\mu, \nu$ be two finite measures on $\mathcal X$ and $\kappa, \eta : \mathcal X \rightsquigarrow \mathcal Y$ two Markov kernels.
  Then $\KL(\mu \otimes \kappa, \nu \otimes \eta) = \KL(\mu, \nu) + \KL(\mu \otimes \kappa, \mu \otimes \eta)$.
\end{theorem}

\begin{proof}%\leanok
\uses{thm:rnDeriv_chain_compProd, lem:integrable_llr_compProd_iff, lem:ac_compProd_iff}
By Lemma~\ref{lem:ac_compProd_iff}, $\mu \otimes \kappa \ll \nu \otimes \eta \iff \left( \mu \ll \nu \ \wedge \ \mu \otimes \kappa \ll \mu \otimes \eta \right)$.
Hence if any of the two sides of the equality is infinite due to a lack of absolute continuity, the other is as well.
We now suppose that $\mu \otimes \kappa \ll \nu \otimes \eta$, $\mu \ll \nu$ and $\mu \otimes \kappa \ll \mu \otimes \eta$.

TODO: deal with the non-integrable case, without using Lemma~\ref{lem:integrable_llr_compProd_iff} that requires countably generated sigma-algebras.

Suppose then that all three log-likelihood ratios are integrable.
We use the chain rule for Radon-Nikodym derivatives (Theorem~\ref{thm:rnDeriv_chain_compProd}).
\begin{align*}
\KL(\mu \otimes \kappa, \nu \otimes \eta)
&= \int_p \log \frac{d(\mu \otimes \kappa)}{d(\nu \otimes \eta)}(p) \partial (\mu \otimes \kappa)
\\
&= \int_x \int_y\log \left(\frac{d\mu}{d \nu}(x)\frac{d(\mu \otimes \kappa)}{d (\mu \otimes \eta)}(x,y)\right) \partial \kappa(x) \partial\mu
\\
&= \int_x \int_y\log \left(\frac{d\mu}{d \nu}(x)\right) + \log \left(\frac{d(\mu \otimes \kappa)}{d (\mu \otimes \eta)}(x,y)\right) \partial \kappa(x) \partial\mu
\\
&= \int_x \log \left(\frac{d\mu}{d \nu}(x)\right)\partial\mu + \int_x \int_y\log \left(\frac{d(\mu \otimes \kappa)}{d (\mu \otimes \eta)}(x,y)\right) \partial \kappa(x) \partial\mu
\\
&= \int_x \log \left(\frac{d\mu}{d \nu}(x)\right)\partial\mu + \int_p\log \left(\frac{d(\mu \otimes \kappa)}{d (\mu \otimes \eta)}(p)\right) \partial (\mu \otimes \kappa)
\\
&= \KL(\mu, \nu) + \KL(\mu \otimes \kappa, \mu \otimes \eta)
\: .
\end{align*}
\end{proof}


\begin{theorem}[Chain rule, kernel version]
  \label{thm:kl_compProd}
  \lean{ProbabilityTheory.klDiv_compProd_eq_add_condKL}
  \leanok
  \uses{def:KL, def:condKL}
  Let $\mu, \nu$ be two finite measures on $\mathcal X$ and $\kappa, \eta : \mathcal X \rightsquigarrow \mathcal Y$ two Markov kernels, with either $\mathcal X$ countable or $\mathcal{Y}$ countably generated.
  Then $\KL(\mu \otimes \kappa, \nu \otimes \eta) = \KL(\mu, \nu) + \KL(\kappa, \eta \mid \mu)$.
\end{theorem}

\begin{proof} \leanok
\uses{lem:rnDeriv_compProd, cor:rnDeriv_value, lem:integrable_llr_compProd_iff}
Handle the case where $\mu \otimes \kappa$ is not absolutely continuous with respect to  $\nu \otimes \eta$ first, then the case where $p \mapsto \log \frac{d \mu \otimes \kappa}{d \nu \otimes \eta}(p)$ is not $\mu \otimes \kappa$-integrable using Lemma~\ref{lem:integrable_llr_compProd_iff}.

For the case where $\mu \otimes \kappa \ll \nu \otimes \eta$ and $p \mapsto \log \frac{d \mu \otimes \kappa}{d \nu \otimes \eta}(p)$ is $\mu \otimes \kappa$-integrable use Lemma~\ref{lem:rnDeriv_compProd} and Corollary~\ref{cor:rnDeriv_value} in a computation:
\begin{align*}
\KL(\mu \otimes \kappa, \nu \otimes \eta)
&= \int_p \log \frac{d(\mu \otimes \kappa)}{d(\nu \otimes \eta)}(p) \partial (\mu \otimes \kappa)
\\
&= \int_x \int_y\log \left(\frac{d\mu}{d \nu}(x)\frac{d\kappa}{d \eta}(x,y)\right) \partial \kappa(x) \partial\mu
\\
&= \int_x \int_y\log \left(\frac{d\mu}{d \nu}(x)\right) + \log \left(\frac{d\kappa}{d \eta}(x,y)\right) \partial \kappa(x) \partial\mu
\\
&= \int_x \log \left(\frac{d\mu}{d \nu}(x)\right)\partial\mu + \int_x \int_y\log \left(\frac{d\kappa}{d \eta}(x,y)\right) \partial \kappa(x) \partial\mu
\\
&= \int_x \log \left(\frac{d\mu}{d \nu}(x)\right)\partial\mu + \int_x \int_y\log \left(\frac{d\kappa(x)}{d \eta(x)}(y)\right) \partial \kappa(x) \partial\mu
\\
&= \KL(\mu, \nu) + \KL(\kappa, \eta \mid \mu)
\: .
\end{align*}

\end{proof}


\begin{theorem}[Chain rule, with Bayesian inverse]
  \label{thm:kl_compProd_bayesInv}
  %\lean{}
  %\leanok
  \uses{def:KL, def:condKL, def:bayesInv}
  Let $\mu, \nu$ be two finite measures on $\mathcal X$ and $\kappa, \eta : \mathcal X \rightsquigarrow \mathcal Y$ two Markov kernels with Bayesian inverses $\kappa_\mu^\dagger$ and $\eta_\nu^\dagger$.
  Then 
  \begin{align*}
  \KL(\mu \otimes \kappa, \nu \otimes \eta) = \KL(\kappa \circ \mu, \eta \circ \nu) + \KL(\kappa_\mu^\dagger, \eta_\nu^\dagger \mid \kappa \circ \mu)
  \: .
  \end{align*}
\end{theorem}

\begin{proof}%\leanok
\uses{thm:kl_compProd}

\end{proof}


\begin{theorem}[Chain rule, product version]
  \label{thm:kl_fst_add_condKL}
  \lean{ProbabilityTheory.klDiv_fst_add_condKL}
  \leanok
  \uses{def:KL, def:condKL}
  Let $\mu, \nu$ be two finite measures on $\mathcal X \times \mathcal Y$, where $\mathcal Y$ is standard Borel.
  Then $\KL(\mu, \nu) = \KL(\mu_X, \nu_X) + \KL(\mu_{Y|X}, \nu_{Y|X} \mid \mu_X)$.
\end{theorem}

\begin{proof} \leanok
\uses{thm:kl_compProd}
Write $\mu = \mu_X \otimes \mu_{Y|X}$ and $\nu = \nu_X \otimes \nu_{Y|X}$, then use Theorem~\ref{thm:kl_compProd}.
\end{proof}


\begin{lemma}
  \label{lem:kl_prod_two'}
  \lean{ProbabilityTheory.klDiv_prod_two'}
  \leanok
  \uses{def:KL}
  Let $\mu_1, \nu_1$ be finite measures on $\mathcal X$ and $\mu_2, \nu_2$ probability measures on $\mathcal{Y}$. Then
  $$\KL(\mu_1\times \mu_2, \nu_1 \times \nu_2) = \KL(\mu_1, \nu_1) + \KL(\mu_2, \nu_2) \mu_1 (\mathcal X) \: .$$
\end{lemma}

\begin{proof} \leanok
\uses{thm:kl_compProd, lem:condKL_const}
Write $\mu_1 \times \mu_2$ and $\nu_1 \times \nu_2$ as composition products of a measure and a constant kernel, then apply the chain rule (Theorem~\ref{thm:kl_compProd}) and Lemma~\ref{lem:condKL_const}.
\end{proof}
 

\begin{theorem}[Tensorization]
  \label{thm:kl_prod_two}
  \lean{ProbabilityTheory.klDiv_prod_two}
  \leanok
  \uses{def:KL}
  For $\mu_1$ a probability measure on $\mathcal X$, $\nu_1$ a finite measure on $\mathcal{X}$ and $\mu_2, \nu_2$ two probability measures on $\mathcal Y$,
  \begin{align*}
  \KL(\mu_1 \times \mu_2, \nu_1 \times \nu_2) = \KL(\mu_1, \nu_1) + \KL(\mu_2 \times \nu_2) \: .
  \end{align*}
\end{theorem}

\begin{proof} \leanok
\uses{lem:kl_prod_two'}
This is a particular case of Lemma~\ref{lem:kl_prod_two'}.
\end{proof}


\begin{theorem}[Tensorization - finite product]
  \label{thm:kl_pi}
  \lean{ProbabilityTheory.klDiv_pi}
  \leanok
  \uses{def:KL}
  Let $I$ be a finite index set. Let $(\mu_i)_{i \in I}, (\nu_i)_{i \in I}$ be probability measures on spaces $(\mathcal X_i)_{i \in I}$. Then
  \begin{align*}
  \KL(\prod_{i \in I} \mu_i, \prod_{i \in I} \nu_i) = \sum_{i \in I} \KL(\mu_i, \nu_i)
  \: .
  \end{align*}
\end{theorem}

\begin{proof} \leanok
\uses{thm:kl_prod_two}
Induction on the size of $I$, using Theorem~\ref{thm:kl_prod_two}.
\end{proof}


\begin{lemma}
  \label{lem:kl_chain_rule_cond_event}
  %\lean{}
  %\leanok
  \uses{def:KL}
  Let $\mu, \nu$ be two measures and $E$ an event. Let $\mu_{|E}$ be the measure defined by $\mu_{|E}(A) = \frac{\mu(A \cap E)}{\mu(E)}$ and define $\nu_{|E}$, $\mu_{| E^c}$ and $\nu_{| E^c}$ similarly. Let $\kl(p,q) = p\log\frac{p}{q} + (1-p)\log\frac{1-p}{1-q}$ be the Kullback-Leibler divergence between Bernoulli distributions with means $p$ and $q$. Then
  \begin{align*}
  \KL(\mu, \nu) \ge \kl(\mu(E), \nu(E)) + \mu(E) \KL(\mu_{|E}, \nu_{|E}) + \mu(E^c) \KL(\mu_{|E^c}, \nu_{|E^c}) \: .
  \end{align*}
\end{lemma}

\begin{proof}
\uses{thm:kl_compProd, lem:fDiv_compProd_prod_eq}
Let $\mu_E$ be the Bernoulli distribution with $\mu_E(\{1\}) = \mu(E)$, and define $\nu_E$ similarly.
Let $\kappa : \{0,1\} \rightsquigarrow \mathcal X$ be the kernel defined by $\kappa(1) = \mu_{|E}$ and $\kappa(0) = \mu_{|E^c}$. Define a kernel $\eta$ similarly for $\nu$.
Then $\mu = \kappa \circ \mu_E$ and $\nu = \eta \circ \nu_E$.

$\mu_E \otimes \kappa$ is equal to the composition of $\mu$ and the deterministic kernel given by the function $f : x \mapsto (\mathbb{I}_E(x), x)$. By Lemma~\ref{lem:fDiv_compProd_prod_eq}, that gives $\KL(\mu_E \otimes \kappa, \nu_E \otimes \eta) = \KL(\mu, \nu)$.

Using the chain rule (Theorem~\ref{thm:kl_compProd}),
\begin{align*}
\KL(\mu_E \otimes \kappa, \nu_E \otimes \eta)
&= \KL(\mu_E, \nu_E) + \KL(\kappa, \eta \mid \mu_E)
\\
&= \KL(\mu_E, \nu_E) + \mu(E) \KL(\mu_{|E}, \nu_{|E}) + \mu(E^c) \KL(\mu_{|E^c}, \nu_{|E^c})
\: .
\end{align*}

\end{proof}


\begin{lemma}
  \label{lem:expectation_llr_event}
  %\lean{}
  %\leanok
  \uses{def:KL}
  Let $\mu, \nu$ be two measures and $E$ an event. Then
  $\mu(E)\log\frac{\mu(E)}{\nu(E)} \le \mu\left[\mathbb{I}(E)\log \frac{d \mu}{d \nu}\right]$ .
\end{lemma}

\begin{proof}
\uses{lem:kl_ge}
Apply Lemma~\ref{lem:kl_ge} to the measures $\mu$ and $\nu$ restricted to $E$. $\mu\left[\mathbb{I}(E)\log \frac{d \mu}{d \nu}\right]$ is the KL between those two distributions.
\end{proof}




TODO: find a way to hide the following dummy lemma
\begin{lemma}[Dummy lemma: KL properties]
  \label{lem:kl_properties}
  %\lean{}
  \leanok
  \uses{def:KL, def:condKL}
  Dummy node to summarize properties of the Kullback-Leibler divergence.
\end{lemma}

\begin{proof}\leanok
\uses{
  lem:kl_eq_fDiv,
  lem:condKL_eq_condFDiv,
  lem:kl_self,
  lem:condKL_const,
  thm:kl_fst_le,
  thm:kl_data_proc,
  lem:kl_ge,
  lem:kl_nonneg,
  lem:kl_eq_zero_iff,
  lem:kl_convex,
  thm:kl_compProd,
  thm:kl_compProd_bayesInv,
  thm:kl_fst_add_condKL,
  lem:kl_prod_two',
  thm:kl_prod_two,
  thm:kl_pi,
  lem:kl_chain_rule_cond_event
}
\end{proof}

%%FILE blueprint/src/sections/mutual_info.tex

\section{Mutual information}

\paragraph{Definitions}

\begin{definition}
  \label{def:mutualInfo}
  %\lean{}
  %\leanok
  \uses{def:KL, lem:kl_properties}
  The mutual information is, for $\rho \in \mathcal M(\mathcal X \times \mathcal Y)$~,
  \begin{align*}
  I(\rho) = \KL(\rho, \rho_X \times \rho_Y) \: .
  \end{align*}
\end{definition}


\begin{definition}
  \label{def:condMutualInfo}
  %\lean{}
  %\leanok
  \uses{def:condKL, lem:kl_properties}
  Let $\kappa : \mathcal Z \rightsquigarrow \mathcal X \times \mathcal Y$. The conditional mutual information of $\kappa$ with respect to $\nu \in \mathcal M(\mathcal Z)$ is
  \begin{align*}
  I(\kappa \mid \nu) = \KL(\kappa, \kappa_X \times \kappa_Y \mid \nu) \: .
  \end{align*}
\end{definition}


\paragraph{Properties}

\begin{lemma}
  \label{lem:mutualInfo_eq_condKL}
  %\lean{}
  %\leanok
  \uses{def:mutualInfo, def:KL, def:condKL}
  For $\mu \in \mathcal M(\mathcal X)$ and $\kappa : \mathcal X \rightsquigarrow \mathcal Y$ a Markov kernel,
  \begin{align*}
  I(\mu \otimes \kappa)
  = KL(\mu \otimes \kappa, \mu \times (\kappa \circ \mu))
  = KL(\kappa, \kappa \circ \mu \mid \mu)
  \: ,
  \end{align*}
  where in the conditional divergence the measure $\kappa \circ \mu$ should be understood as the constant kernel from $\mathcal X$ to $\mathcal Y$ with that value.
\end{lemma}

\begin{proof}%\leanok
\uses{}
The first equality is the definition of $I$, together with $(\mu \otimes \kappa)_X = \mu$ (since $\kappa$ is Markov) and $(\mu \otimes \kappa)_Y = \kappa \circ \mu$.
The second equality is from $KL(\kappa, \eta \mid \mu) = KL(\mu \otimes \kappa, \mu \otimes \eta)$. (TODO: now from a lemma, soon by definition).
\end{proof}


\begin{lemma}
  \label{lem:mutualInfo_symm}
  %\lean{}
  %\leanok
  \uses{def:mutualInfo}
  $I(\rho_\leftrightarrow) = I(\rho)$~.
\end{lemma}

\begin{proof}%\leanok
\uses{lem:fDiv_map_measurableEmbedding}

\end{proof}


\begin{lemma}
  \label{lem:mutualInfo_eq_jensenShannon}
  %\lean{}
  %\leanok
  \uses{def:mutualInfo, def:jensenShannon}
  Let $\mu, \nu \in \mathcal P(\mathcal X)$ and let $\alpha \in (0, 1)$. Let $\pi_\alpha = (\alpha, 1 - \alpha) \in \mathcal P(\{0,1\})$ and let $P : \{0,1\} \rightsquigarrow \mathcal X$ be the kernel with $P(0) = \mu$ and $P(1) = \nu$. Then
  \begin{align*}
  I(\pi_\alpha \otimes P) = \JS_\alpha(\mu, \nu) \: .
  \end{align*}
\end{lemma}

\begin{proof}%\leanok
\uses{lem:jensenShannon_eq_kl, lem:mutualInfo_eq_condKL}
Use Lemma~\ref{lem:mutualInfo_eq_condKL} and Lemma~\ref{lem:jensenShannon_eq_kl}.
\end{proof}


\begin{theorem}[Data-processing]
  \label{thm:mutualInfo_data_proc}
  %\lean{}
  %\leanok
  \uses{def:mutualInfo}
  For $\rho \in \mathcal M(\mathcal X \times \mathcal Y)$, $\kappa : \mathcal X \rightsquigarrow \mathcal X'$ and $\eta : \mathcal Y \rightsquigarrow \mathcal Y'$ two Markov kernels,
  \begin{align*}
  I((\kappa \parallel \eta) \circ \rho) \le I(\rho) \: .
  \end{align*}
\end{theorem}

\begin{proof}%\leanok
\uses{thm:kl_data_proc}
$((\kappa \parallel \eta) \circ \rho)_{X'} = \kappa \circ \rho_X$ and $((\kappa \parallel \eta) \circ \rho)_{Y'} = \eta \circ \rho_Y$ hence
\begin{align*}
((\kappa \parallel \eta) \circ \rho)_{X'} \times ((\kappa \parallel \eta) \circ \rho)_{Y'}
&= (\kappa \circ \rho_X) \times (\eta \circ \rho_Y)
= (\kappa \parallel \eta) \circ (\rho_X \times \rho_Y)
\: .
\end{align*}

By the data-processing inequality for $\KL$ (Theorem~\ref{thm:kl_data_proc}),
\begin{align*}
I((\kappa \parallel \eta) \circ \rho)
&= \KL((\kappa \parallel \eta) \circ \rho, (\kappa \parallel \eta) \circ (\rho_X \times \rho_Y))
\\
&\le \KL(\rho, \rho_X \times \rho_Y)
\\
&= I(\rho)
\: .
\end{align*}
\end{proof}

\begin{corollary}
  \label{cor:mutualInfo_compProd_le}
  %\lean{}
  %\leanok
  \uses{def:mutualInfo}
  For $\mu \in \mathcal X$, $\kappa : \mathcal X \rightsquigarrow \mathcal Y$ and $\eta : \mathcal Y \rightsquigarrow \mathcal Z$ two Markov kernels,
  \begin{align*}
  I(\mu \otimes (\eta \circ \kappa)) \le I(\mu \otimes \kappa)
  \: .
  \end{align*}
\end{corollary}

That corollary corresponds to the following situation: 3 random variables $X, Y, Z$ have the dependency structure $X \to Y \to Z$ and $\mu$ is the law of $X$, $\mu \otimes \kappa$ is the law of $(X, Y)$ and $\mu \otimes (\eta \circ \kappa)$ is the law of $(X, Z)$.
With the usual $I(X, Y)$ notation, the corollary statement is $I(X, Z) \le I(X, Y)$~.

\begin{proof}%\leanok
\uses{thm:mutualInfo_data_proc}
First, rewrite $\mu \otimes (\eta \circ \kappa) = (\mathrm{id} \parallel \eta) \circ (\mu \otimes \kappa)$. Then apply Theorem~\ref{thm:mutualInfo_data_proc}.
\end{proof}




\section{New definitions related to the mutual information}

\begin{definition}
  \label{def:mutualInfoLeft}
  %\lean{}
  %\leanok
  \uses{def:fDiv, lem:fDiv_properties}
  Let $D$ be a divergence between measures.
  The left $D$-mutual information for a measure $\mu \in \mathcal M(\mathcal X)$ and a kernel $\kappa : \mathcal X \rightsquigarrow \mathcal Y$ is
  \begin{align*}
  I_D^L(\mu, \kappa) = \inf_{\xi \in \mathcal P(\mathcal Y)} D(\mu \times \xi, \mu \otimes \kappa) \: .
  \end{align*}
\end{definition}


\begin{definition}
  \label{def:mutualInfoRight}
  %\lean{}
  %\leanok
  \uses{def:fDiv, lem:fDiv_properties}
  Let $D$ be a divergence between measures.
  The right $D$-mutual information for a measure $\mu \in \mathcal M(\mathcal X)$ and a kernel $\kappa : \mathcal X \rightsquigarrow \mathcal Y$ is
  \begin{align*}
  I_D^R(\mu, \kappa) = \inf_{\xi \in \mathcal P(\mathcal Y)} D(\mu \otimes \kappa, \mu \times \xi) \: .
  \end{align*}
\end{definition}


\begin{lemma}
  \label{lem:mutualInfoLeft_eq_renyi}
  %\lean{}
  %\leanok
  \uses{def:mutualInfoLeft, def:Renyi}
  For $\mu \in \mathcal P(\{0,1\})$ and $\kappa : \{0,1\} \rightsquigarrow \mathcal Y$,
  \begin{align*}
  I_{\KL}^L(\mu, \kappa) = (1 - \mu_0) R_{\mu_0}(\kappa_0, \kappa_1) \: .
  \end{align*}
\end{lemma}

\begin{proof}%\leanok
\uses{lem:renyi_eq_inf_kl}

\end{proof}


\begin{lemma}
  \label{lem:mutualInfoRight_eq_mutualInfo}
  %\lean{}
  %\leanok
  \uses{def:mutualInfoRight, def:mutualInfo}
  For $\mu \in \mathcal P(\mathcal X)$ and $\kappa : \mathcal X \rightsquigarrow \mathcal Y$,
  \begin{align*}
  I_{\KL}^R(\mu, \kappa) = I(\mu \otimes \kappa) \: .
  \end{align*}
\end{lemma}

\begin{proof}%\leanok
\uses{}

\end{proof}


\begin{corollary}
  \label{cor:mutualInfoRight_eq_jensenShannon}
  %\lean{}
  %\leanok
  \uses{def:mutualInfoRight, def:jensenShannon}
  For $\mu \in \mathcal P(\{0,1\})$ and $\kappa : \{0,1\} \rightsquigarrow \mathcal Y$,
  \begin{align*}
  I_{\KL}^R(\mu, \kappa) = \JS_{\mu_0}(\kappa_0, \kappa_1) \: .
  \end{align*}
\end{corollary}

\begin{proof}%\leanok
\uses{lem:mutualInfoRight_eq_mutualInfo, lem:mutualInfo_eq_jensenShannon}

\end{proof}



\paragraph{Data-processing inequality}


\begin{theorem}
  \label{thm:mutualInfoLeft_comp_le}
  %\lean{}
  %\leanok
  \uses{def:mutualInfoLeft, def:mutualInfoRight}
  If the divergence $D$ satisfies the data-processing inequality, then for all $\mu \in \mathcal M(\mathcal X)$, $\kappa : \mathcal X \rightsquigarrow \mathcal Y$ and all Markov kernels $\eta : \mathcal Y \rightsquigarrow \mathcal Z$ then
  \begin{align*}
  I_D^L(\mu, \eta \circ \kappa) \le I_D^L(\mu, \kappa)
  \: , \\
  I_D^R(\mu, \eta \circ \kappa) \le I_D^R(\mu, \kappa)
  \: .
  \end{align*}
  
\end{theorem}

\begin{proof}%\leanok
\uses{}

\end{proof}


\begin{lemma}
  \label{lem:mutualInfoLeft_estimation_ge}
  %\lean{}
  %\leanok
  \uses{def:mutualInfoLeft, def:mutualInfoRight, def:bayesRisk}
  Let $\pi \in \mathcal P(\Theta)$ and $P : \Theta \rightsquigarrow \mathcal X$. Suppose that the loss $\ell'$ of an estimation task with kernel $P$ takes values in $[0,1]$. Then
  \begin{align*}
  I_{D_f}^L(\pi, P)
  &\ge d_f(\mathcal R_\pi^{d_\Theta}, \mathcal R_\pi^P)
  \: , \\
  I_{D_f}^R(\pi, P)
  &\ge d_f(\mathcal R_\pi^P, \mathcal R_\pi^{d_\Theta})
  \: .
  \end{align*}
\end{lemma}

\begin{proof}%\leanok
\uses{lem:fDiv_estimation_ge}
The left mutual information $I_D^L(\pi, P)$ is an infimum: $I_D^L(\pi, P) = \inf_{\xi \in \mathcal P(\mathcal X)} D(\pi \times \xi, \pi \otimes P)$.
For all $\xi \in \mathcal P(\mathcal X)$, by Lemma~\ref{lem:fDiv_estimation_ge},
\begin{align*}
D_f(\pi \times \xi, \pi \otimes P) &\ge d_f(\mathcal R_\pi^{\theta \mapsto \xi}, \mathcal R_\pi^P) \: .
\end{align*}
The risk of a constant Markov kernel is equal to the risk of the discard kernel $d_\Theta$ (TODO ref).

The proof for the right mutual information is similar.
\end{proof}



%%FILE blueprint/src/sections/renyi_divergence.tex
\section{Rényi divergences}

\begin{definition}[Rényi divergence]
  \label{def:Renyi}
  \lean{ProbabilityTheory.renyiDiv}
  \leanok
  \uses{def:KL, def:hellingerAlpha, lem:hellingerAlpha_properties}
  Let $\mu, \nu$ be two finite measures on $\mathcal X$. The Rényi divergence of order $\alpha \in \mathbb{R}$ between $\mu$ and $\nu$ is
  \begin{align*}
  R_\alpha(\mu, \nu) = \left\{
 \begin{array}{ll}
  - \log \left( \nu\{x \mid \frac{d \mu}{d \nu}(x) > 0\} / \nu(\mathcal X) \right) & \text{for } \alpha = 0
  \\
  \KL(\mu / \mu(\mathcal X), \nu / \nu(\mathcal X)) & \text{for } \alpha = 1
  \\
  \frac{1}{\alpha - 1} \log \left( \frac{(1 - \alpha) \nu(\mathcal X) + \alpha \mu(\mathcal X) + (\alpha - 1) \He_\alpha(\mu, \nu)}{\mu(\mathcal X)^\alpha \nu(\mathcal X)^{1 - \alpha}} \right) & \text{for } \alpha \notin \{0, 1\}
  \end{array}\right.
  \end{align*}
  It is the Rényi divergence of the normalized measures $\mu / \mu(\mathcal X)$ and $\nu / \nu(\mathcal X)$. For $\alpha \notin \{0, 1\}$, the numerator inside the logarithm is $\int \left(\frac{d \mu}{d \nu}\right)^\alpha d\nu$.
\end{definition}


\begin{lemma}
  \label{lem:renyiDiv_zero}
  \lean{ProbabilityTheory.renyiDiv_zero}
  \leanok
  \uses{def:Renyi}
  For $\mu$ a sigma-finite measure and $\nu$ a finite measure 
  \[R_0(\mu, \nu) = - \log(\nu\{x \mid \frac{d \mu}{d \nu}(x) > 0\}) \: .\]
\end{lemma}

\begin{proof}\leanok
\uses{def:hellingerAlpha}
\end{proof}


\begin{lemma}
  \label{lem:renyiDiv_eq_top_iff_mutuallySingular_of_lt_one}
  \lean{ProbabilityTheory.renyiDiv_eq_top_iff_mutuallySingular_of_lt_one}
  \leanok
  \uses{def:Renyi}
  For $\alpha \in [0, 1)$ and finite measures $\mu, \nu$, 
  \[R_\alpha(\mu, \nu) = \infty \iff \mu \perp \nu \: .\]
\end{lemma}

\begin{proof}\leanok
\uses{lem:hellingerAlpha_ne_top_of_lt_one}
\end{proof}


\begin{lemma}
  \label{lem:renyi_eq_log_integral}
  %\lean{}
  %\leanok
  \uses{def:Renyi}
  For $\alpha \in (0,1)\cup(1, \infty)$ and finite measures $\mu, \nu$, if 
  $\left(\frac{d \mu}{d \nu}\right)^\alpha$ is integrable with respect to $\nu$ and $\mu \ll \nu$ then
  \begin{align*}
  R_\alpha(\mu, \nu) = \frac{1}{\alpha - 1} \log \int_x \left(\frac{d \mu}{d \nu}(x)\right)^\alpha \partial \nu
  \: .
  \end{align*}
\end{lemma}

\begin{proof}%\leanok
\uses{lem:hellingerAlpha_eq_integral}
\end{proof}


\begin{lemma}
  \label{lem:renyi_eq_log_integral'}
  %\lean{}
  %\leanok
  \uses{def:Renyi}
  For $\alpha \in (0,1)\cup(1, \infty)$ and finite measures $\mu, \nu$, if 
  $\left(\frac{d \mu}{d \nu}\right)^\alpha$ is integrable with respect to $\nu$ and $\mu \ll \nu$ then
  \begin{align*}
  R_\alpha(\mu, \nu) = \frac{1}{\alpha - 1} \log \int_x \left(\frac{d \mu}{d \nu}(x)\right)^{\alpha - 1} \partial \mu
  \: .
  \end{align*}
\end{lemma}

\begin{proof}%\leanok
\uses{lem:renyi_eq_log_integral, lem:integral_rpow_rnDeriv_symm}
\end{proof}


\begin{lemma}
  \label{lem:renyi_symm}
  %\lean{}
  %\leanok
  \uses{def:Renyi}
  For $\alpha \in (0, 1)$ and finite measures $\mu, \nu$ with $\mu(\mathcal X) = \nu(\mathcal X)$, 
  \[(1 - \alpha) R_\alpha(\mu, \nu) = \alpha R_{1 - \alpha}(\nu, \mu) \: .\]
\end{lemma}

\begin{proof}%\leanok
\uses{lem:hellingerAlpha_symm}
Use Lemma~\ref{lem:hellingerAlpha_symm}.
\end{proof}


\begin{lemma}
  \label{lem:renyi_cgf}
  %\lean{}
  %\leanok
  \uses{def:Renyi}
  The cumulant generating function of $\log\frac{d\mu}{d\nu}$ under $\nu$ is $\alpha \mapsto (\alpha - 1) R_\alpha(\mu, \nu)$.
\end{lemma}

\begin{proof}%\leanok
\uses{lem:renyi_eq_log_integral}
Unfold the definitions, using Lemma~\ref{lem:renyi_eq_log_integral} for the Rényi divergence.
\end{proof}


\begin{lemma}
  \label{lem:renyi_cgf_2}
  %\lean{}
  %\leanok
  \uses{def:Renyi}
  Set $\alpha > 0$. If $R_{1+\alpha}(\mu, \nu) < \infty$, the cumulant generating function of $\log\frac{d\mu}{d\nu}$ under $\mu$ has value $\alpha R_{1+\alpha}(\mu, \nu)$ at $\alpha$.
\end{lemma}

\begin{proof}%\leanok
\uses{lem:renyi_eq_log_integral'}
Unfold the definitions, using Lemma~\ref{lem:renyi_eq_log_integral'} for the Rényi divergence.
\end{proof}


\begin{definition}[Conditional Rényi divergence]
  \label{def:condRenyi}
  \lean{ProbabilityTheory.condRenyiDiv}
  \leanok
  \uses{def:condHellingerAlpha, lem:hellingerAlpha_properties}
  Let $\mu$ be a measure on $\mathcal X$ and $\kappa, \eta : \mathcal X \rightsquigarrow \mathcal Y$ be two kernels. The conditional Rényi divergence of order $\alpha \in (0,+\infty) \backslash \{1\}$ between $\kappa$ and $\eta$ conditionally to $\mu$ is
  \begin{align*}
  R_\alpha(\kappa, \eta \mid \mu) =\frac{1}{\alpha - 1} \log (1 + (\alpha - 1) \He_\alpha(\kappa, \eta \mid \mu)) \: .
  \end{align*}
\end{definition}

TODO: the reasons for this definition is that it is equal to $R_\alpha(\mu \otimes \kappa, \mu \otimes \eta)$, which is the generic conditional divergence definition, and also that this is the definition for which we have a chain rule.




\subsection{Properties inherited from f-divergences}

Since Rényi divergences are monotone transfomations of f-divergences, every inequality for f-divergences can be translated to Rényi divergences.

\begin{theorem}[Data-processing]
  \label{thm:renyi_data_proc}
  \lean{ProbabilityTheory.renyiDiv_comp_right_le}
  \leanok
  \uses{def:Renyi}
  Let $\mu, \nu$ be two finite measures on $\mathcal X$ and let $\kappa : \mathcal X \rightsquigarrow \mathcal Y$ be a Markov kernel.
  Then $R_\alpha(\kappa \circ \mu, \kappa \circ \nu) \le R_\alpha(\mu, \nu)$.
\end{theorem}

\begin{proof} \leanok
\uses{thm:fDiv_data_proc_2}
The function $x \mapsto \frac{1}{\alpha - 1}\log (1 + (\alpha - 1)x)$ is non-decreasing and $D_f$ satisfies the DPI (Theorem~\ref{thm:fDiv_data_proc_2}), hence we get the DPI for $R_\alpha$.
\end{proof}


\begin{lemma}
  \label{lem:renyi_data_proc_event}
  %\lean{}
  %\leanok
  \uses{def:Renyi}
  Let $\mu, \nu$ be two measures on $\mathcal X$ and let $E$ be an event. Let $\mu_E$ and $\nu_E$ be the two Bernoulli distributions with respective means $\mu(E)$ and $\nu(E)$.
  Then $R_\alpha(\mu, \nu) \ge R_\alpha(\mu_E, \nu_E)$.
\end{lemma}

\begin{proof}
\uses{cor:data_proc_event}
By Corollary~\ref{cor:data_proc_event}, $D_f(\mu, \nu) \ge D_f(\mu_E, \nu_E)$, hence $R_\alpha(\mu, \nu) \ge R_\alpha(\mu_E, \nu_E)$.
\end{proof}




\subsection{Comparisons between orders}

\begin{lemma}
  \label{lem:renyi_tendsto_renyi_zero}
  %\lean{}
  %\leanok
  \uses{def:Renyi}
  Let $\mu, \nu$ be two finite measures. $R_0(\mu, \nu) = \lim_{\alpha \downarrow 0} R_\alpha(\mu, \nu)$.
\end{lemma}

\begin{proof}
\end{proof}

\begin{lemma}
  \label{lem:renyi_tendsto_renyi_one}
  %\lean{}
  %\leanok
  \uses{def:Renyi}
  Let $\mu, \nu$ be two finite measures. $R_1(\mu, \nu) = \lim_{\alpha \uparrow 1} R_\alpha(\mu, \nu)$.
\end{lemma}

\begin{proof}
\end{proof}

\begin{lemma}
  \label{lem:renyi_tendsto_renyi_one_above}
  %\lean{}
  %\leanok
  \uses{def:Renyi}
  Let $\mu, \nu$ be two finite measures such that there exists $\alpha > 1$ with $R_\alpha(\mu, \nu)$ finite. Then $R_1(\mu, \nu) = \lim_{\alpha \downarrow 1} R_\alpha(\mu, \nu)$.
\end{lemma}

\begin{proof}
\end{proof}

\begin{lemma}
  \label{lem:renyi_monotone}
  %\lean{}
  %\leanok
  \uses{def:Renyi}
  Let $\mu, \nu$ be two finite measures. Then $\alpha \mapsto R_\alpha(\mu, \nu)$ is nondecreasing on $[0, + \infty)$.
\end{lemma}

\begin{proof}
\end{proof}

\begin{lemma}
  \label{lem:renyi_continuous}
  %\lean{}
  %\leanok
  \uses{def:Renyi}
  Let $\mu, \nu$ be two finite measures. Then $\alpha \mapsto R_\alpha(\mu, \nu)$ is continuous on $[0, 1]$ and on $[0, \sup \{\alpha \mid R_\alpha(\mu, \nu) < \infty\})$.
\end{lemma}

\begin{proof}
\end{proof}




\subsection{The tilted measure}

\begin{definition}
  \label{def:renyiMeasure}
  %\lean{}
  %\leanok
  \uses{def:Renyi}
  Let $\mu, \nu$ be two measures on $\mathcal X$ and let $\alpha \in (0, +\infty) \backslash \{1\}$. Let $p = \frac{d \mu}{d (\mu + \nu)}$ and $q = \frac{d \nu}{d (\mu + \nu)}$. We define a measure $\mu^{(\alpha, \nu)}$, absolutely continuous with respect to $\mu + \nu$ with density
  \begin{align*}
  \frac{d \mu^{(\alpha, \nu)}}{d (\mu + \nu)} = p^\alpha q^{1 - \alpha} e^{-(\alpha - 1)R_\alpha(\mu, \nu)} \: .
  \end{align*}
\end{definition}


\begin{lemma}
  \label{lem:renyiMeasure_ac}
  %\lean{}
  %\leanok
  \uses{def:renyiMeasure}
  For $\alpha \in (0,1)$, $\mu^{(\alpha, \nu)} \ll \mu$ and $\mu^{(\alpha, \nu)} \ll \nu$.
\end{lemma}

\begin{proof}%\leanok
\uses{}

\end{proof}


\begin{lemma}
  \label{lem:rnDeriv_renyiMeasure}
  %\lean{}
  %\leanok
  \uses{def:renyiMeasure}
  $\frac{d \mu^{(\alpha, \nu)}}{d \nu} = \left(\frac{d\mu}{d\nu}\right)^\alpha e^{-(\alpha - 1) R_\alpha(\mu, \nu)}$~, $\nu$-a.e.,
  and
  $\frac{d \mu^{(\alpha, \nu)}}{d \mu} = \left(\frac{d\nu}{d\mu}\right)^{1 - \alpha} e^{-(\alpha - 1) R_\alpha(\mu, \nu)}$~, $\mu$-a.e..
\end{lemma}

\begin{proof}%\leanok
\uses{lem:renyiMeasure_ac}

\end{proof}


\begin{lemma}
  \label{lem:kl_renyiMeasure_eq}
  %\lean{}
  %\leanok
  \uses{def:renyiMeasure, def:Renyi, def:KL}
  Let $\mu, \nu, \xi$ be three measures on $\mathcal X$ and let $\alpha \in (0, 1)$. Then
  \begin{align*}
  \KL(\xi, \mu^{(\alpha, \nu)}) = \alpha \KL(\xi, \mu) + (1 - \alpha)\KL(\xi, \nu) - (1 - \alpha) R_\alpha(\mu, \nu) \: .
  \end{align*}
\end{lemma}

\begin{proof}
\uses{def:renyiMeasure, def:KL}
Unfold definitions and compute?
\end{proof}


\begin{corollary}
  \label{cor:renyi_eq_add_kl}
  %\lean{}
  %\leanok
  \uses{def:renyiMeasure, def:Renyi, def:KL}
  Let $\mu, \nu, \xi$ be three measures on $\mathcal X$ and let $\alpha \in (0, 1)$. Then
  \begin{align*}
  (1 - \alpha) R_\alpha(\mu, \nu) = \alpha \KL(\mu^{(\alpha, \nu)}, \mu) + (1 - \alpha)\KL(\mu^{(\alpha, \nu)}, \nu) \: .
  \end{align*}
\end{corollary}

\begin{proof}
\uses{lem:kl_renyiMeasure_eq}
Apply Lemma~\ref{lem:kl_renyiMeasure_eq} to $\xi = \mu^{(\alpha, \nu)}$.
\end{proof}


\begin{lemma}
  \label{lem:renyi_eq_inf_add_kl}
  %\lean{}
  %\leanok
  \uses{def:Renyi, def:KL}
  Let $\mu, \nu$ be two probability measures on $\mathcal X$ and let $\alpha \in (0, 1)$. Then
  \begin{align*}
  (1 - \alpha) R_\alpha(\mu, \nu) = \inf_{\xi \in \mathcal P(\mathcal X)}\left( \alpha \KL(\xi, \mu) + (1 - \alpha)\KL(\xi, \nu) \right) \: .
  \end{align*}
  The infimum is attained at $\xi = \mu^{(\alpha, \nu)}$.
\end{lemma}

\begin{proof}
\uses{lem:kl_renyiMeasure_eq, lem:kl_nonneg, cor:renyi_eq_add_kl}
By Lemma~\ref{lem:kl_renyiMeasure_eq} and Lemma~\ref{lem:kl_nonneg}, for all $\xi$,
\begin{align*}
\alpha \KL(\xi, \mu) + (1 - \alpha)\KL(\xi, \nu) - (1 - \alpha) R_\alpha(\mu, \nu)
= \KL(\xi, \mu^{(\alpha, \nu)})
\ge 0 \: .
\end{align*}
We have then $(1 - \alpha) R_\alpha(\mu, \nu) \le \inf_{\xi \in \mathcal P(\mathcal X)}\left( \alpha \KL(\xi, \mu) + (1 - \alpha)\KL(\xi, \nu) \right)$.
Corollary~\ref{cor:renyi_eq_add_kl} proves the other inequality.
\end{proof}


\begin{lemma}
  \label{lem:renyi_eq_inf_kl}
  %\lean{}
  %\leanok
  \uses{def:Renyi, def:KL}
  Let $\mu, \nu \in \mathcal P(\mathcal X)$ and let $\alpha \in (0, 1)$. Let $\pi_\alpha = (\alpha, 1 - \alpha) \in \mathcal P(\{0,1\})$ and let $P : \{0,1\} \rightsquigarrow \mathcal X$ be the kernel with $P(0) = \mu$ and $P(1) = \nu$. Then
  \begin{align*}
  (1 - \alpha) R_\alpha(\mu, \nu) = \inf_{\xi \in \mathcal P(\mathcal X)} \KL\left( \pi_\alpha \times \xi, \pi_\alpha \otimes P \right) \: .
  \end{align*}
  The infimum is attained at $\xi = \mu^{(\alpha, \nu)}$.
\end{lemma}

Compare with Lemma~\ref{lem:jensenShannon_eq_inf_kl} on the Jensen-Shannon divergence.

\begin{proof}%\leanok
\uses{lem:renyi_eq_inf_add_kl}
\begin{align*}
\KL\left( \pi_\alpha \times \xi, \pi_\alpha \otimes P \right)
&= \alpha \KL(\xi, \mu) + (1 - \alpha) \KL(\xi, \nu)
\: .
\end{align*}
The result is a rephrasing of Lemma~\ref{lem:renyi_eq_inf_add_kl}.
\end{proof}



\subsection{Chain rule and tensorization}

\begin{theorem}[Chain rule]
  \label{thm:renyi_chain_rule}
  %\lean{}
  %\leanok
  \uses{def:renyiMeasure, def:condRenyi}
  Let $\mu, \nu$ be two measures on $\mathcal X$ and let $\kappa, \eta : \mathcal X \rightsquigarrow \mathcal Y$ be two Markov kernels.
  Then $R_\alpha(\mu \otimes \kappa, \nu \otimes \eta) = R_\alpha(\mu, \nu) + R_\alpha(\kappa, \eta \mid \mu^{(\alpha, \nu)})$.
\end{theorem}

\begin{proof}%\leanok
\uses{thm:kl_compProd}
\end{proof}


\begin{theorem}[Chain rule, with Bayesian inverse]
  \label{thm:renyi_compProd_bayesInv}
  %\lean{}
  %\leanok
  \uses{def:Renyi, def:condRenyi, def:bayesInv}
  Let $\mu, \nu$ be two finite measures on $\mathcal X$ and $\kappa, \eta : \mathcal X \rightsquigarrow \mathcal Y$ two Markov kernels with Bayesian inverses $\kappa_\mu^\dagger$ and $\eta_\nu^\dagger$.
  Then 
  \begin{align*}
  R_\alpha(\mu \otimes \kappa, \nu \otimes \eta)
  = R_\alpha(\kappa \circ \mu, \eta \circ \nu) + R_\alpha(\kappa_\mu^\dagger, \eta_\nu^\dagger \mid (\kappa \circ \mu)^{(\alpha, \eta \circ \nu)})
  \: .
  \end{align*}
\end{theorem}

\begin{proof}%\leanok
\uses{thm:renyi_chain_rule}

\end{proof}


\begin{corollary}
  \label{cor:renyi_prod_two}
  %\lean{}
  %\leanok
  \uses{def:Renyi}
  Let $\mu, \nu \in \mathcal P(\mathcal X)$ and $\xi, \lambda \in \mathcal P(\mathcal Y)$.
  Then $R_\alpha(\mu \times \xi, \nu \times \lambda) = R_\alpha(\mu, \nu) + R_\alpha(\xi, \lambda)$.
\end{corollary}

\begin{proof}%\leanok
\uses{thm:renyi_chain_rule}
Apply Theorem~\ref{thm:renyi_chain_rule} (the product of measures is a special case of $\otimes$). TODO: more details.
\end{proof}


\begin{theorem}[Tensorization - finite product]
  \label{thm:renyi_prod}
  %\lean{}
  %\leanok
  \uses{def:Renyi}
  Let $I$ be a finite index set. Let $(\mu_i)_{i \in I}, (\nu_i)_{i \in I}$ be probability measures on measurable spaces $(\mathcal X_i)_{i \in I}$.
  Then $R_\alpha (\prod_{i \in I} \mu_i, \prod_{i \in I} \nu_i) = \sum_{i \in I} R_\alpha(\mu_i, \nu_i)$.
\end{theorem}

\begin{proof}%\leanok
\uses{cor:renyi_prod_two}
Induction over the finite index set, using Corollary~\ref{cor:renyi_prod_two}.
\end{proof}


\begin{corollary}
  \label{lem:renyi_prod_n}
  %\lean{}
  %\leanok
  \uses{def:Renyi}
  Let $\mu, \nu$ be two probability measures on $\mathcal X$. Let $n \in \mathbb{N}$ and write $\mu^{\otimes n}$ for the product measure on $\mathcal X^n$ of $n$ times $\mu$.
  Then $R_\alpha(\mu^{\otimes n}, \nu^{\otimes n}) = n R_\alpha(\mu, \nu)$.
\end{corollary}

\begin{proof}%\leanok
\uses{thm:renyi_prod}
Apply Theorem~\ref{thm:renyi_prod}.
\end{proof}


\begin{theorem}[Tensorization - countable product]
  \label{thm:renyi_prod_countable}
  %\lean{}
  %\leanok
  \uses{def:Renyi}
  Let $I$ be a countable index set. Let $(\mu_i)_{i \in I}, (\nu_i)_{i \in I}$ be probability measures on measurable spaces $(\mathcal X_i)_{i \in I}$.
  Then $R_\alpha (\prod_{i \in I} \mu_i, \prod_{i \in I} \nu_i) = \sum_{i \in I} R_\alpha(\mu_i, \nu_i)$.
\end{theorem}

\begin{proof}%\leanok
\uses{thm:renyi_prod}
\end{proof}


TODO: find a way to hide the following dummy lemma
\begin{lemma}[Dummy lemma: Renyi properties]
  \label{lem:renyi_properties}
  %\lean{}
  \leanok
  \uses{def:Renyi, def:condRenyi, def:renyiMeasure}
  Dummy node to summarize properties of the Rényi divergence.
\end{lemma}

\begin{proof}\leanok
\uses{
  lem:renyiDiv_zero,
  lem:renyiDiv_eq_top_iff_mutuallySingular_of_lt_one,
  lem:renyi_eq_log_integral,
  lem:renyi_eq_log_integral',
  lem:renyi_symm,
  lem:renyi_cgf,
  lem:renyi_cgf_2,
  thm:renyi_data_proc,
  lem:renyi_data_proc_event,
  lem:renyiMeasure_ac,
  lem:rnDeriv_renyiMeasure,
  lem:kl_renyiMeasure_eq,
  cor:renyi_eq_add_kl,
  lem:renyi_eq_inf_add_kl,
  lem:renyi_eq_inf_kl,
  thm:renyi_chain_rule,
  thm:renyi_compProd_bayesInv,
  cor:renyi_prod_two,
  thm:renyi_prod,
  lem:renyi_prod_n,
  thm:renyi_prod_countable
}
\end{proof}

%%FILE blueprint/src/sections/risk.tex

In the estimation problem, we are given a family of measures and we want to estimate an unknown parameter for the measure from which we observe samples.
Let $\mathcal X$ be the measurable space on which the measures are defined.
The measure family is described by a kernel $P : \Theta \rightsquigarrow \mathcal X$. We write $P_\theta$ for $P(\theta)$.
In \emph{parametric} problems, $\Theta$ is a low-dimensional space. In \emph{nonparametric} problems $\Theta$ is large, for example $\Theta = \mathcal M (\mathcal X)$ or $\Theta = \mathcal P(\mathcal X)$.

The goal of the estimation task is to find, given the observation of samples of $P_\theta$ for an unknown $\theta$, the value of a function $y : \Theta \to \mathcal Y$~, which maps the parameter $\theta$ to the quantity we are interested in. For example, $y(\theta)$ might be the expected value of the distribution $P_\theta$.

An estimator is a Markov kernel $\hat{y} : \mathcal X \rightsquigarrow \mathcal Z$, where most often $\mathcal Z = \mathcal Y$.
Note that, while $y$ is a function of a parameter in $\Theta$, $\hat{y}$ has domain $\mathcal X$ as the estimation is done based on a sample of the measure, while the measure itself and its parameter remain unknown.
In general this is a randomized estimator, but often deterministic estimators (corresponding to deterministic kernels) are enough to attain optimal performance.
The quality of the estimation is quantified by a loss $\ell' : \mathcal Y \times \mathcal Z \to \mathbb{R}_{+, \infty}$. We write $\ell$ for a loss on $\mathcal Y \times \mathcal Y$ and $\ell'$ in the heterogeneous case.

In short, an estimation problem is specified by the data $(P, y, \ell')$.

% TODO: I can't make tikz-cd work with 'leanblueprint web'

%\begin{figure}[ht]
%\centering
% https://tikzcd.yichuanshen.de/#N4Igdg9gJgpgziAXAbVABwnAlgFyxMJZABgBpiBdUkANwEMAbAVxiRAB12AVACxhzogAvqXSZc+QigCM5KrUYs2nALZ0cPAMaMABAE1hokBmx4CRMtPn1mrRB3ZqN2hjoAahsaclFZV6jZK9qrqWroAWp7G4mZSyABMpP4KtsqOoQBGGcAASkIA+sAA1KScWGAAZjgAnkLC8jBQAObwRKAVAE4QKkhkIDgQSIkgDOV2IFB0cHxQIAGK4wAKUZ3dQ9QDSADM1KNg45PTjXMpQQ486sC1K109iH2biLKn49U3a08bg4gALNQZMDAs0QAFotn1plgqr1qHw6MCQBcOhhzPNUsF2DAGAwAOTvO47frfP4gAFAmEgSHQ+6wmDwthIlFgRg4epCIA
%\begin{tikzcd}
%\Theta \arrow[d, "P", dashed] \arrow[r, "y"] & \mathcal Y \arrow[rd, "\ell'", harpoon, bend right] &                         \\
%\mathcal X \arrow[r, "\hat{y}", dashed]      & \mathcal Z \arrow[r, harpoon']                      & {\mathbb{R}_{+,\infty}}
%\end{tikzcd}
%\caption{The estimation problem. Dashed arrow represent transition kernels. Plain arrows are measurable functions. }
%\label{fig:estimation}
%\end{figure}

\paragraph{Example: simple binary hypothesis testing (SBHT)}

We say that an estimation problem is a testing problem if $\mathcal Y = \mathcal Z$ is discrete and the loss takes values in $\{0, 1\}$.
A test is said to be binary if $\mathcal Y$ has two elements.
A test is simple if $\{\theta \mid y(\theta) = y_0\}$ is a singleton for all $y_0 \in \mathcal Y$, i.e. $y$ is bijective.

In summary, in simple binary hypothesis testing, $\Theta = \mathcal Y = \mathcal Z = \{0,1\}$, $y(0) = 0$, $y(1) = 1$ (that is, $y$ is the identity). For $y_0, z \in \{0,1\}$, $\ell(y_0, z) = \mathbb{I}\{y_0 \ne z\}$.

\section{Risk}

Let $\hat{\mu}_\theta$ be the law of $\hat{y}(\theta)$. That is, $\hat{\mu} : \mathcal \Theta \rightsquigarrow \mathcal X$ is the kernel $\hat{y} \circ P$.

\begin{definition}[Risk]
  \label{def:risk}
  \mathlibok
  \uses{lem:kernel_properties}
  The risk of an estimator $\hat{y}$ on the estimation problem $(P, y, \ell')$ at $\theta \in \Theta$ is $r^P_\theta(\hat{y}) = (\hat{y} \circ P)(\theta)\left[z \mapsto \ell'(y(\theta), z)\right]$~.
\end{definition}

\emph{Example (SBHT):} $r^P_\theta(\hat{y}) = \hat{\mu}_\theta(\hat{y}(\theta) \ne y(\theta))$.


\begin{definition}[Bayesian risk]
  \label{def:bayesianRisk}
  \lean{ProbabilityTheory.avgRisk}
  \mathlibok
  \uses{def:risk}
  The Bayesian risk of an estimator $\hat{y}$ on $(P, y, \ell')$ for a prior $\pi \in \mathcal M(\Theta)$ is $R^P_\pi(\hat{y}) = \pi\left[\theta \mapsto r^P_\theta(\hat{y})\right]$~. It can also be expanded as $R^P_\pi(\hat{y}) = (\pi \otimes (\hat{y} \circ P))\left[ (\theta, z) \mapsto \ell'(y(\theta), z) \right]$~.
\end{definition}

If we write $\ell'$ and $y$ for their deterministic kernels, the Bayesian risk of $\hat{y}$ is the mean of the measure $\ell' \circ (y \parallel \hat{y}) \circ (\pi \otimes P)$~.


\begin{definition}[Bayes risk]
  \label{def:bayesRisk}
  \lean{ProbabilityTheory.bayesRisk}
  \mathlibok
  \uses{def:bayesianRisk}
  The Bayes risk of $(P, y, \ell')$ for prior $\pi \in \mathcal M(\Theta)$ is $\mathcal R^P_\pi = \inf_{\hat{y} : \mathcal X \rightsquigarrow \mathcal Z} R^P_\pi(\hat{y})$~, where the infimum is over Markov kernels.

  The Bayes risk of $(P, y, \ell')$ is $\mathcal R^*_B = \sup_{\pi \in \mathcal P(\Theta)} \mathcal R^P_\pi \: .$
\end{definition}

\emph{Example (SBHT):} $\mathcal R^*_B = \sup_{\gamma \in [0,1]}\inf_{\hat{y}}\left(\gamma \hat{\mu}_0(\hat{y} \ne 0) + (1 - \gamma) \hat{\mu}_1(\hat{y} \ne 1)\right)$.

\begin{definition}[Bayes estimator]
  \label{def:bayesEstimator}
  \lean{ProbabilityTheory.IsBayesEstimator}
  \mathlibok
  \uses{def:bayesRisk}
  An estimator $\hat{y}$ is said to be a Bayes estimator for a prior $\pi \in \mathcal P(\Theta)$ if $R^P_\pi(\hat{y}) = \mathcal R^P_\pi$.
\end{definition}

\begin{definition}[Minimax risk]
  \label{def:minimaxRisk}
  \lean{ProbabilityTheory.minimaxRisk}
  \mathlibok
  \uses{def:risk}
  The minimax risk of $(P, y, \ell')$ is $\mathcal R^* = \inf_{\hat{y} : \mathcal X \rightsquigarrow \mathcal Z} \sup_{\theta \in \Theta} r^P_\theta(\hat{y})$~.
\end{definition}

\emph{Example (SBHT):} $\mathcal R^* = \inf_{\hat{y}} \max\{\hat{\mu}_0(\hat{y} \ne 0), \hat{\mu}_1(\hat{y} \ne 1)\}$.

\begin{lemma}
  \label{lem:bayesRisk_le_minimaxRisk}
  \lean{ProbabilityTheory.bayesRisk_le_minimaxRisk}
  \mathlibok
  \uses{def:bayesRisk, def:minimaxRisk}
  $\mathcal R_B^* \le \mathcal R^*$.
\end{lemma}

\begin{proof}\mathlibok
For any $\pi \in \mathcal P(\mathcal X)$ and any estimator, $\pi\left[\hat{\mu}_\theta\left[\ell'(y(\theta), \hat{y}(\theta))\right]\right] \le \sup_{\theta \in \Theta}\hat{\mu}_\theta\left[\ell'(y(\theta), \hat{y}(\theta))\right]$.
\end{proof}

TODO: ``often'', $\mathcal R^*_B = \mathcal R^*$.




\subsection{Properties of the Bayes risk of a prior}

\begin{lemma}
  \label{lem:bayesRisk_le_const}
  \lean{ProbabilityTheory.bayesRisk_le_iInf}
  \mathlibok
  \uses{def:bayesRisk}
  The Bayes risk of a prior $\pi \in \mathcal M(\Theta)$ on $(P, y, \ell')$ with $P$ a Markov kernel satisfies
  \begin{align*}
  \mathcal R^P_\pi \le \inf_{z \in \mathcal Z} \pi\left[ \theta \mapsto \ell'(y(\theta), z) \right]
  \: .
  \end{align*}
\end{lemma}

\begin{proof}\mathlibok
\uses{}
The infimum over all Markov kernels in the definition of the Bayes risk of $\pi$ is bounded from above by the infimum restricted to constant, deterministic Markov kernels.
\end{proof}


\begin{lemma}
  \label{lem:bayesRisk_const}
  \lean{ProbabilityTheory.bayesRisk_const}
  \mathlibok
  \uses{def:bayesRisk}
  The Bayes risk of a prior $\pi \in \mathcal M(\Theta)$ on $(P, y, \ell')$ with $P$ a constant Markov kernel is
  \begin{align*}
  \mathcal R^P_\pi
  = \inf_{z \in \mathcal Z} \pi\left[\theta \mapsto \ell'(y(\theta), z)\right]
  \: .
  \end{align*}
  In particular, it does not depend on $P$.
\end{lemma}

\begin{proof}\mathlibok
\uses{}
Let $\xi$ be the measure such that $P(\theta) = \xi$ for all $\theta$.
\begin{align*}
\mathcal R^P_\pi
&= \inf_{\hat{y}}(\pi \times (\hat{y} \circ \xi))\left[(\theta, z)\mapsto \ell'(y(\theta), z)\right]
\\
&= \inf_{\mu \in \mathcal P(\mathcal X)} (\pi \times \mu)\left[(\theta, z)\mapsto \ell'(y(\theta), z)\right]
\\
&= \inf_{\mu \in \mathcal P(\mathcal X)} \mu\left[ z \mapsto \pi\left[\theta \mapsto \ell'(y(\theta), z)\right]\right]
\\
&= \inf_{z \in \mathcal Z} \pi\left[\theta \mapsto \ell'(y(\theta), z)\right]
\: .
\end{align*}

\end{proof}



\paragraph{Data-processing inequality and consequences}

\begin{theorem}[Data-processing inequality]
  \label{thm:data_proc_bayesRisk}
  \lean{ProbabilityTheory.bayesRisk_le_bayesRisk_comp}
  \mathlibok
  \uses{def:bayesRisk}
  For $P : \Theta \rightsquigarrow \mathcal X$ and $\kappa : \mathcal X \rightsquigarrow \mathcal X'$ a Markov kernel, $\mathcal R^{\kappa \circ P}_\pi \ge \mathcal R^{P}_\pi$ (where the estimation problems differ only in the kernel).
\end{theorem}

\begin{proof}\mathlibok
\uses{}
The risk $\mathcal R^{\kappa \circ P}_\pi$ is
\begin{align*}
\mathcal R^{\kappa \circ P}_\pi
&= \inf_{\hat{y} : \mathcal X' \rightsquigarrow \mathcal Z} (\pi \otimes (\hat{y} \circ \kappa \circ P))\left[(\theta, z) \mapsto \ell'(y(\theta), z)\right]
\\
&\ge \inf_{\hat{y} : \mathcal X \rightsquigarrow \mathcal Z} (\pi \otimes (\hat{y} \circ P))\left[(\theta, z) \mapsto \ell'(y(\theta), z)\right]
\\
&= \mathcal R^{P}_\pi
\: .
\end{align*}
The inequality is due to taking an infimum over a larger set: $\{\hat{y} \circ \kappa \mid \hat{y} : \mathcal X' \rightsquigarrow \mathcal Z \text{ Markov}\} \subseteq \{\hat{y}' : \mathcal X \rightsquigarrow \mathcal Z \mid \hat{y} \text{ Markov}\}$~.
\end{proof}

\begin{lemma}
  \label{lem:bayesRisk_compProd_le_fst}
  \lean{ProbabilityTheory.bayesRisk_compProd_le_bayesRisk}
  \mathlibok
  \uses{def:bayesRisk}
  For $P : \Theta \rightsquigarrow \mathcal X$ and $\kappa : \Theta \times \mathcal X \rightsquigarrow \mathcal X'$ a Markov kernel, $\mathcal R^{P \otimes \kappa}_\pi \le \mathcal R^{P}_\pi$~.
\end{lemma}

\begin{proof}\mathlibok
\uses{thm:data_proc_bayesRisk}
Use Theorem~\ref{thm:data_proc_bayesRisk}: $P$ is the composition of $P \otimes \kappa$ and the deterministic kernel $\mathcal X \times \mathcal X' \rightsquigarrow \mathcal X$ that forgets the second value.
\end{proof}

\begin{lemma}
  \label{lem:bayesRisk_compProd_le_snd}
  %\lean{}
  %\leanok
  \uses{def:bayesRisk}
  For $P : \Theta \rightsquigarrow \mathcal X$ and $\kappa : \Theta \times \mathcal X \rightsquigarrow \mathcal X'$ a Markov kernel, $\mathcal R^{P \otimes \kappa}_\pi \le \mathcal R^{(P \otimes \kappa)_{\mathcal X'}}_\pi$~, in which $(P \otimes \kappa)_{\mathcal X'} : \mathcal \Theta \rightsquigarrow \mathcal X'$ is the kernel obtained by marginalizing over $\mathcal X$ in the output of $P \otimes \kappa$~.
\end{lemma}

\begin{proof}%\leanok
\uses{thm:data_proc_bayesRisk}
Use Theorem~\ref{thm:data_proc_bayesRisk}: $(P \otimes \kappa)_{\mathcal X'}$ is the composition of $P \otimes \kappa$ and the deterministic kernel $\mathcal X \times \mathcal X' \rightsquigarrow \mathcal X'$ that forgets the first value.
\end{proof}

\begin{lemma}
  \label{lem:bayesRisk_concave}
  %\lean{}
  %\leanok
  \uses{def:bayesRisk}
  The Bayes risk $\mathcal R_\pi^P$ is concave in $P : \Theta \rightsquigarrow \mathcal X$~.
\end{lemma}

\begin{proof}%\leanok
\uses{}
The infimum of a sum is larger than the sum of the infimums:
\begin{align*}
\mathcal R_\pi^{\lambda P_1 + (1 - \lambda)P_2}
&= \inf_{\hat{y} : \mathcal X \rightsquigarrow \mathcal Z} (\pi \otimes (\hat{y} \circ (\lambda P_1 + (1 - \lambda)P_2)))\left[(\theta, z) \mapsto \ell'(y(\theta), z)\right]
\\
&= \inf_{\hat{y} : \mathcal X \rightsquigarrow \mathcal Z}  \left( \lambda (\pi \otimes (\hat{y} \circ P_1))\left[(\theta, z) \mapsto \ell'(y(\theta), z)\right] \right.
  \\&\qquad \left. + (1 - \lambda) (\pi \otimes (\hat{y} \circ P_2))\left[(\theta, z) \mapsto \ell'(y(\theta), z)\right] \right)
\\
&\ge \lambda \inf_{\hat{y} : \mathcal X \rightsquigarrow \mathcal Z} (\pi \otimes (\hat{y} \circ P_1))\left[(\theta, z) \mapsto \ell'(y(\theta), z)\right]
  \\&\qquad + (1 - \lambda) \inf_{\hat{y} : \mathcal X \rightsquigarrow \mathcal Z} (\pi \otimes (\hat{y} \circ P_2))\left[(\theta, z) \mapsto \ell'(y(\theta), z)\right]
\\
&= \lambda \mathcal R_\pi^{P_1} + (1 - \lambda)\mathcal R_\pi^{P_2}
\: .
\end{align*}
TODO: can we build kernels such that this lemma is a consequence of one of the lemmas above?
\end{proof}



\paragraph{Generalized Bayes estimator}

\begin{lemma}
  \label{lem:bayesianRisk_bayesInv}
  \lean{ProbabilityTheory.avgRisk_eq_lintegral_posterior_prod}
  \mathlibok
  \uses{def:bayesianRisk, def:bayesInv}
  The Bayesian risk of a Markov kernel $\hat{y} : \mathcal X \rightsquigarrow \mathcal Z$ with respect to a prior $\pi \in \mathcal M(\Theta)$ on $(P, y, \ell')$ satisfies
  \begin{align*}
  R^P_\pi(\hat{y}) = ((P_\pi^\dagger \times \hat{y}) \circ P \circ \pi)\left[(\theta, z) \mapsto \ell'(y(\theta), z)\right] \: ,
  \end{align*}
  whenever the Bayesian inverse $P_\pi^\dagger$ of $P$ with respect to $\pi$ exists (Definition~\ref{def:bayesInv}).
\end{lemma}

\begin{proof}\mathlibok
\uses{lem:bayesInv_properties}
Use the main property of the Bayesian inverse.
\begin{align*}
  (P_\pi^\dagger \times \hat{y}) \circ P \circ \pi
  &= (\mathrm{id} \parallel \hat{y}) \circ (P_\pi^\dagger \times \mathrm{id}) \circ (P \circ \pi)
  \\
  &= (\mathrm{id} \parallel \hat{y}) \circ (\mathrm{id} \times P) \circ \pi
  \\
  &= (\mathrm{id} \times (\hat{y} \circ P)) \circ \pi
  \\
  &= \pi \otimes (\hat{y} \circ P)
  \: .
  \end{align*}
\end{proof}

\begin{lemma}
  \label{lem:bayesianRisk_ge_inf_bayesInv}
  \lean{ProbabilityTheory.lintegral_iInf_posterior_le_avgRisk}
  \mathlibok
  \uses{def:bayesianRisk, def:bayesInv}
  The Bayesian risk of a Markov kernel $\hat{y} : \mathcal X \rightsquigarrow \mathcal Z$ with respect to a prior $\pi \in \mathcal M(\Theta)$ on $(P, y, \ell')$ satisfies
  \begin{align*}
  R^P_\pi(\hat{y}) \ge (P \circ \pi)\left[x \mapsto \inf_{z \in \mathcal Z} P_\pi^\dagger(x) \left[\theta \mapsto \ell'(y(\theta), z)\right]\right] \: .
  \end{align*}
\end{lemma}

\begin{proof}\mathlibok
\uses{lem:bayesianRisk_bayesInv}
Starting from the equality of Lemma~\ref{lem:bayesianRisk_bayesInv}, we get
\begin{align*}
R^P_\pi(\hat{y})
&= ((P_\pi^\dagger \times \hat{y}) \circ P \circ \pi)\left[(\theta, z) \mapsto \ell'(y(\theta), z)\right]
\\
&= (P \circ \pi)\left[x \mapsto (P_\pi^\dagger(x) \times \hat{y}(x)) \left[(\theta, z) \mapsto \ell'(y(\theta), z)\right]\right]
\\
&= (P \circ \pi)\left[x \mapsto \hat{y}(x)\left[z \mapsto P_\pi^\dagger(x) \left[\theta \mapsto \ell'(y(\theta), z)\right]\right]\right]
\\
&\ge (P \circ \pi)\left[x \mapsto \inf_{z \in \mathcal Z} P_\pi^\dagger(x) \left[\theta \mapsto \ell'(y(\theta), z)\right]\right]
\: .
\end{align*}

\end{proof}

\begin{definition}[Generalized Bayes estimator]
  \label{def:genBayesEstimator}
  \lean{ProbabilityTheory.IsArgminEstimator}
  \mathlibok
  \uses{def:risk}
  The generalized Bayes estimator for prior $\pi \in \mathcal P(\Theta)$ on $(P, y, \ell')$ is the deterministic estimator $\mathcal X \to \mathcal Z$ given by
  \begin{align*}
  x \mapsto \arg\min_z P_\pi^\dagger(x)\left[\theta \mapsto \ell'(y(\theta), z)\right] \: ,
  \end{align*}
  if there exists such a measurable argmin.
\end{definition}


\begin{lemma}
  \label{lem:bayesianRisk_genBayesEstimator}
  \lean{ProbabilityTheory.IsArgminEstimator.avgRisk_eq_lintegral_iInf}
  \mathlibok
  \uses{def:bayesianRisk, def:genBayesEstimator}
  The Bayesian risk of the generalized Bayes estimator $\hat{y}_B$ is
  \begin{align*}
  R^P_\pi(\hat{y}_B) = (P \circ \pi)\left[x \mapsto \inf_{z \in \mathcal Z} P_\pi^\dagger(x) \left[\theta \mapsto \ell'(y(\theta), z)\right]\right]
  \: .
  \end{align*}
\end{lemma}

\begin{proof}\mathlibok
\uses{lem:bayesianRisk_bayesInv}
Start from the equality of Lemma~\ref{lem:bayesianRisk_bayesInv} and use the definition of the generalized Bayes estimator.
\end{proof}

\begin{theorem}
  \label{thm:isBayesEstimator_genBayesEstimator}
  \lean{ProbabilityTheory.IsArgminEstimator.isBayesEstimator, ProbabilityTheory.HasArgminEstimator.bayesRisk_eq}
  \mathlibok
  \uses{def:bayesEstimator, def:genBayesEstimator, def:bayesRisk}
  When the generalized Bayes estimator is well defined, it is a Bayes estimator. The value of the Bayes risk with respect to the prior $\pi \in \mathcal M(\Theta)$ is then
  \begin{align*}
  \mathcal R^P_\pi = (P \circ \pi)\left[x \mapsto \inf_{z \in \mathcal Z} P_\pi^\dagger(x) \left[\theta \mapsto \ell'(y(\theta), z)\right]\right]
  \: .
  \end{align*}
\end{theorem}

\begin{proof}\mathlibok
\uses{lem:bayesianRisk_ge_inf_bayesInv, lem:bayesianRisk_genBayesEstimator}
By Lemma~\ref{lem:bayesianRisk_ge_inf_bayesInv}, the Bayesian risk of the generalized Bayes estimator obtained in Lemma~\ref{lem:bayesianRisk_genBayesEstimator} is a lower bound for the Bayesian risk.
\end{proof}


\begin{lemma}
  \label{lem:bayesRisk_eq_rnDeriv}
  %\lean{}
  %\leanok
  \uses{def:genBayesEstimator, def:bayesRisk}
  When the generalized Bayes estimator is well defined, the Bayes risk with respect to the prior $\pi \in \mathcal M(\Theta)$ is
  \begin{align*}
  \mathcal R^P_\pi = (P \circ \pi)\left[x \mapsto \inf_{z \in \mathcal Z} \pi \left[\theta \mapsto \frac{d P_\theta}{d (P \circ \pi)}(x) \ell'(y(\theta), z)\right]\right]
  \: .
  \end{align*}
\end{lemma}

\begin{proof}%\leanok
\uses{lem:bayesInv_properties}

\end{proof}


\begin{lemma}
  \label{lem:bayesRisk_binary_eq_sub_bayesInv}
  %\lean{}
  %\leanok
  \uses{def:genBayesEstimator, def:bayesRisk}
  When the generalized Bayes estimator is well defined, the Bayes risk with respect to the prior $\pi \in \mathcal M(\Theta)$ for $\mathcal Y = \mathcal Z = \Theta$, $y = \mathrm{id}$ and $\ell' = \mathbb{I}\{\theta \ne z\}$ is
  \begin{align*}
  \mathcal R^P_\pi = (P \circ \pi)[1] - (P \circ \pi)\left[x \mapsto \sup_{z \in \mathcal Z} P_\pi^\dagger(x) \{z\}\right]
  \: .
  \end{align*}
\end{lemma}

\begin{proof}%\leanok
\uses{thm:isBayesEstimator_genBayesEstimator,lem:bayesInv_properties}

\end{proof}


\begin{lemma}
  \label{lem:bayesRisk_binary_eq_sub}
  %\lean{}
  %\leanok
  \uses{def:genBayesEstimator, def:bayesRisk}
  When the generalized Bayes estimator is well defined, the Bayes risk with respect to the prior $\pi \in \mathcal M(\Theta)$ for $\mathcal Y = \mathcal Z = \Theta$, $y = \mathrm{id}$ and $\ell' = \mathbb{I}\{\theta \ne z\}$ is
  \begin{align*}
  \mathcal R^P_\pi = (P \circ \pi)[1] - (P \circ \pi)\left[x \mapsto \sup_{z \in \mathcal Z} \frac{d P_z}{d (P \circ \pi)}(x) \pi\{z\} \right]
  \: .
  \end{align*}
  When $\pi$ is a probability measure and $P$ is a Markov kernel, $(P \circ \pi)[1] = 1$.
\end{lemma}

\begin{proof}%\leanok
\uses{lem:bayesRisk_binary_eq_sub_bayesInv}

\end{proof}


\begin{lemma}
  \label{lem:genBayesEstimator_binary}
  %\lean{}
  %\leanok
  \uses{def:genBayesEstimator}
  The generalized Bayes estimator for prior $\pi \in \mathcal P(\Theta)$ on the estimation problem defined by $\mathcal Y = \mathcal Z = \Theta$, $y = \mathrm{id}$ and $\ell' = \mathbb{I}\{\theta \ne z\}$ is
  \begin{align*}
  x \mapsto \arg\max_z \left( \pi\{z\} \frac{d P_z}{d (P \circ \pi)}(x) \right)
  \: .
  \end{align*}
\end{lemma}

\begin{proof}%\leanok
\uses{lem:bayesRisk_binary_eq_sub}

\end{proof}


\begin{lemma}
  \label{lem:bayesRisk_binary_le_sub_prod}
  %\lean{}
  %\leanok
  \uses{def:bayesRisk, def:genBayesEstimator}
  Suppose that $\Theta$ is finite and let $\xi \in \mathcal P(\Theta)$.
  The Bayes risk with respect to the prior $\pi \in \mathcal P(\Theta)$ for $\mathcal Y = \mathcal Z = \Theta$, $y = \mathrm{id}$, $P$ a Markov kernel and $\ell' = \mathbb{I}\{\theta \ne z\}$ satisfies
  \begin{align*}
  \mathcal R^P_\pi \le 1 - \left(\prod_{\theta \in \Theta} \pi_\theta^{\xi_\theta}\right) (P \circ \pi)\left[x \mapsto \prod_{\theta \in \Theta} \left(\frac{d P_\theta}{d (P \circ \pi)}(x)\right)^{\xi_\theta} \right]
  \: .
  \end{align*}
\end{lemma}

\begin{proof}%\leanok
\uses{lem:bayesRisk_binary_eq_sub}

\end{proof}


\subsection{Bayes risk increase}

TODO: remove this section? It's not used for anything now.

The following is a new definition, which is a natural generalization of the DeGroot statistical information.

\begin{definition}
  \label{def:riskIncrease}
  \lean{ProbabilityTheory.riskIncrease}
  \mathlibok
  \uses{def:bayesRisk}
  The Bayes risk increase $I^P_{\pi}(\kappa)$ of a kernel $\kappa : \mathcal X \rightsquigarrow \mathcal X'$ with respect to the estimation problem $(P, y, \ell')$ and the prior $\pi \in \mathcal M(\Theta)$ is the difference of the Bayes risk of $(\kappa \circ P, y, \ell')$ and that of $(P, y, \ell')$. That is,
  \begin{align*}
  I^P_{\pi}(\kappa)
  &= \mathcal R^{\kappa \circ P}_\pi - \mathcal R^P_\pi
  \\
  &= \inf_{\hat{y} : \mathcal X' \rightsquigarrow \mathcal Z} (\pi \otimes (\hat{y} \circ \kappa \circ P))\left[(\theta, z) \mapsto \ell'(y(\theta), z)\right]
    - \inf_{\hat{y} : \mathcal X \rightsquigarrow \mathcal Z} (\pi \otimes (\hat{y} \circ P))\left[(\theta, z) \mapsto \ell'(y(\theta), z)\right]
  \: .
  \end{align*}
\end{definition}

In particular, if we take $\kappa : \mathcal X \rightsquigarrow *$ equal to the Markov kernel to the point space (denoted by $d_{\mathcal X}$, for ``discard'' kernel), we get $I^P_{\pi}(d_{\mathcal X}) = \mathcal R^{d_{\mathcal X} \circ P}_\pi - \mathcal R^P_\pi$, where $d_{\mathcal X} \circ P = d_\Theta$ if $P$ is a Markov kernel. We recover a difference between the risk of the best constant kernel and the Bayes estimator, as in the DeGroot statistical information.

The risk increase quantifies how much risk we accrue by forgetting information due to the composition with $\kappa$.

\begin{lemma}
  \label{lem:riskIncrease_nonneg}
  %\lean{}
  %\leanok
  \uses{def:riskIncrease}
  For $\kappa$ a Markov kernel, $I^P_\pi(\kappa) \ge 0$~.
\end{lemma}

\begin{proof}%\leanok
\uses{thm:data_proc_bayesRisk}
Use Theorem~\ref{thm:data_proc_bayesRisk}.
\end{proof}

\begin{lemma}
  \label{lem:riskIncrease_comp}
  %\lean{}
  %\leanok
  \uses{def:riskIncrease}
  For $\kappa : \mathcal X \rightsquigarrow \mathcal X'$ and $\eta : \mathcal X' \rightsquigarrow \mathcal X''$ two Markov kernels,
  $I^P_\pi(\eta \circ \kappa) = I^P_\pi(\kappa) + I^{\kappa \circ P}_\pi(\eta)$~.
\end{lemma}

\begin{proof}%\leanok
\uses{}
\begin{align*}
I^P_\pi(\kappa) + I^{\kappa \circ P}_\pi(\eta)
&= \mathcal R^{\kappa \circ P}_\pi - \mathcal R^{P}_\pi + \mathcal R^{\eta \circ \kappa \circ P}_\pi - \mathcal R^{\kappa \circ P}_\pi
\\
&= \mathcal R^{\eta \circ \kappa \circ P}_\pi - \mathcal R^{P}_\pi
\\
&= I^P_\pi(\eta \circ \kappa)
\: .
\end{align*}
\end{proof}

\begin{lemma}[Data-processing inequality]
  \label{lem:riskIncrease_comp_del}
  \lean{ProbabilityTheory.riskIncrease_comp_le}
  \mathlibok
  \uses{def:riskIncrease}
  For any measurable space $\mathcal X$, let $d_{\mathcal X} : \mathcal X \rightsquigarrow *$ be the Markov kernel to the point space.
  For all Markov kernels $\kappa : \mathcal X \rightsquigarrow \mathcal X'$,
  \begin{align*}
  I_\pi^P(d_{\mathcal X}) \ge I_\pi^{\kappa \circ P}(d_{\mathcal X'}) \: .
  \end{align*}
\end{lemma}

\begin{proof}\mathlibok
\uses{lem:riskIncrease_comp, lem:riskIncrease_nonneg}
By Lemma~\ref{lem:riskIncrease_comp}, then Lemma~\ref{lem:riskIncrease_nonneg},
\begin{align*}
I_\pi^P(d_{\mathcal X'} \circ \kappa)
= I^P_\pi(\kappa) + I^{\kappa \circ P}_\pi(d_{\mathcal X'})
\ge I_\pi^{\kappa \circ P}(d_{\mathcal X'})
\: .
\end{align*}
Finally, $d_{\mathcal X'} \circ \kappa = d_{\mathcal X}$ since $\kappa$ is a Markov kernel.
\end{proof}




\section{Simple binary hypothesis testing}

In this section, for a measure $\xi$ on $\{0,1\}$, we write $\xi_0$ for $\xi(\{0\})$ and $\xi_1$ for $\xi(\{1\})$~. We sometimes write $(a,b)$ for the measure $\xi$ on $\{0,1\}$ with $\xi_0 = a$ and $\xi_1 = b$.

\begin{lemma}
  \label{lem:bayesInv_binary}
  \lean{ProbabilityTheory.posterior_boolKernel_apply_true, ProbabilityTheory.posterior_boolKernel_apply_false}
  \mathlibok
  \uses{def:bayesInv}
  The Bayesian inverse of a kernel $P : \{0,1\} \rightsquigarrow \mathcal X$ with respect to a prior $\xi \in \mathcal M(\{0,1\})$ is $P_\xi^\dagger(x) = \left(\xi_0\frac{d P(0)}{d(P \circ \xi)}(x), \xi_1\frac{d P(1)}{d(P \circ \xi)}(x)\right)$ (almost surely w.r.t. $P \circ \xi = \xi_0 P(0) + \xi_1 P(1)$).
\end{lemma}

\begin{proof}\mathlibok
\uses{lem:bayesInv_properties}

\end{proof}

\begin{lemma}
  \label{lem:bayesRisk_binary}
  \lean{ProbabilityTheory.bayesRisk_boolKernel_eq_lintegral_iInf}
  \leanok
  \uses{def:bayesRisk}
  For $\Theta = \{0,1\}$, the Bayes risk of a prior $\xi \in \mathcal M(\{0,1\})$ is
  \begin{align*}
  \mathcal R^P_\xi = (P \circ \xi)\left[x \mapsto \inf_{z \in \mathcal Z} \left( \xi_0\frac{d P(0)}{d(P \circ \xi)}(x)\ell'(y(0), z) + \xi_1\frac{d P(1)}{d(P \circ \xi)}(x)\ell'(y(1), z) \right) \right]
  \: .
  \end{align*}
\end{lemma}

\begin{proof}\leanok
\uses{lem:bayesInv_binary, thm:isBayesEstimator_genBayesEstimator}
Use Theorem~\ref{thm:isBayesEstimator_genBayesEstimator}, with the value of $P_\xi^\dagger$ given by Lemma~\ref{lem:bayesInv_binary}.
\end{proof}



\subsection{Bayes risk of a prior}

\begin{definition}
  \label{def:bayesBinaryRisk}
  \lean{ProbabilityTheory.bayesBinaryRisk}
  \leanok
  \uses{def:bayesRisk}
  The Bayes binary risk between measures $\mu$ and $\nu$ with respect to prior $\xi \in \mathcal M(\{0,1\})$, denoted by $\mathcal B_\xi(\mu, \nu)$, is the Bayes risk $\mathcal R^P_\xi$ for $\Theta = \mathcal Y = \mathcal Z = \{0,1\}$, $\ell(y,z) = \mathbb{I}\{y \ne z\}$, $P$ the kernel sending 0 to $\mu$ and 1 to $\nu$ and prior $\xi$.
  That is,
  \begin{align*}
  \mathcal B_\xi(\mu, \nu) = \inf_{\hat{y} : \mathcal X \rightsquigarrow \{0,1\}}\left(\xi_0 (\hat{y} \circ \mu)(\{1\}) + \xi_1 (\hat{y} \circ \nu)(\{0\})\right)
  \: ,
  \end{align*}
  in which the infimum is over Markov kernels.

  If the prior is a probability measure with weights $(\pi, 1 - \pi)$, we write $B_\pi(\mu, \nu) = \mathcal B_{(\pi, 1 - \pi)}(\mu, \nu)$~.
\end{definition}

Thanks to the following lemma, we can prove properties of the Bayes binary risk for non-zero finite measures by considering only probability measures.


\begin{lemma}
  \label{lem:bayesBinaryRisk_mul}
  \lean{ProbabilityTheory.bayesBinaryRisk_smul_smul}
  \leanok
  \uses{def:bayesBinaryRisk}
  For all $a, b > 0$, $\mathcal B_\xi(\mu, \nu) = \mathcal B_{(a \xi_0, b \xi_1)}(a^{-1} \mu, b^{-1} \nu)$~.
\end{lemma}

\begin{proof}\leanok
\uses{}

\end{proof}

Alternatively, we can reduce the study of the Bayes binary risk for any prior to the study of the prior $(1, 1)$, for which that risk is related to the total variation distance (defined later).


\begin{lemma}
  \label{lem:bayesBinaryRisk_one_one}
  \lean{ProbabilityTheory.bayesBinaryRisk_eq_bayesBinaryRisk_one_one}
  \leanok
  \uses{def:bayesRisk}
  $\mathcal B_\xi(\mu, \nu) = \mathcal B_{(1,1)}(\xi_0\mu, \xi_1\nu)$~.
\end{lemma}

\begin{proof}\leanok
\uses{lem:bayesBinaryRisk_mul}
\end{proof}


\begin{theorem}[Data-processing inequality]
  \label{thm:data_proc_bayesBinaryRisk}
  \lean{ProbabilityTheory.bayesBinaryRisk_le_bayesBinaryRisk_comp}
  \leanok
  \uses{def:bayesBinaryRisk}
  For $\mu, \nu \in \mathcal M(\mathcal X)$ and $\kappa : \mathcal X \rightsquigarrow \mathcal Y$ a Markov kernel, $\mathcal B_\xi(\kappa \circ \mu, \kappa \circ \nu) \ge \mathcal B_\xi(\mu, \nu)$~.
\end{theorem}

\begin{proof}\leanok
\uses{thm:data_proc_bayesRisk}
Apply Theorem~\ref{thm:data_proc_bayesRisk}.
\end{proof}


\begin{lemma}
  \label{lem:bayesBinaryRisk_self}
  \lean{ProbabilityTheory.bayesBinaryRisk_self}
  \leanok
  \uses{def:bayesBinaryRisk}
  For $\mu \in \mathcal M(\mathcal X)$, $\mathcal B_\xi(\mu, \mu) = \min\{\xi_0, \xi_1\} \mu(\mathcal X)$~.
\end{lemma}

\begin{proof}\leanok
\uses{}
\end{proof}


\begin{lemma}
  \label{lem:bayesBinaryRisk_nonneg}
  %\lean{}
  %\leanok
  \uses{def:bayesBinaryRisk}
  For all measures $\mu, \nu$, $\mathcal B_\xi(\mu, \nu) \ge 0$~.
\end{lemma}

\begin{proof}%\leanok
\uses{}
It is an infimum of non-negative values.
\end{proof}


\begin{lemma}
  \label{lem:bayesBinaryRisk_le}
  \lean{ProbabilityTheory.bayesBinaryRisk_le_min}
  \leanok
  \uses{def:bayesBinaryRisk}
  For all measures $\mu, \nu$, $\mathcal B_\xi(\mu, \nu) \le \min\{\xi_0 \mu(\mathcal X), \xi_1 \nu(\mathcal X)\}$~.
\end{lemma}

\begin{proof}\leanok
\uses{thm:data_proc_bayesBinaryRisk, lem:bayesBinaryRisk_self, lem:bayesBinaryRisk_mul}
Let $d_{\mathcal X} : \mathcal X \rightsquigarrow *$ be the discard kernel and $\delta_*$ be the only probability measure over $*$, that is a Dirac delta.
Then applying the DPI (Theorem~\ref{thm:data_proc_bayesBinaryRisk}), followed by Lemma~\ref{lem:bayesBinaryRisk_mul} and Lemma~\ref{lem:bayesBinaryRisk_self}, we obtain:
\begin{align*}
\mathcal B_\xi(\mu, \nu)
&\le \mathcal B_\xi(d_{\mathcal X} \circ \mu, d_{\mathcal X} \circ \nu)
\\
&= \mathcal B_\xi(\mu(\mathcal X) \delta_*, \nu(\mathcal X) \delta_*)
\\
&= \mathcal B_{(\mu(\mathcal X) \xi_0, \nu(\mathcal X) \xi_1)}(\delta_*, \delta_*)
\\
&= \min\{\mu(\mathcal X) \xi_0, \nu(\mathcal X) \xi_1\} \delta_*(\mathcal X)
\\
&= \min\{\mu(\mathcal X) \xi_0, \nu(\mathcal X) \xi_1\}
\: .
\end{align*}
\end{proof}


\begin{lemma}
  \label{lem:bayesBinaryRisk_symm}
  \lean{ProbabilityTheory.bayesBinaryRisk_symm}
  \leanok
  \uses{def:bayesBinaryRisk}
  For $\mu, \nu \in \mathcal M(\mathcal X)$ and $\xi \in \mathcal M(\{0,1\})$, $\mathcal B_\xi(\mu, \nu) = \mathcal B_{\xi_{\leftrightarrow}}(\nu, \mu)$ where $\xi_{\leftrightarrow} \in \mathcal M(\{0,1\})$ is such that $\xi_{\leftrightarrow}(\{0\}) = \xi_1$ and $\xi_{\leftrightarrow}(\{1\}) = \xi_0$.
  For $\pi \in [0,1]$, $B_\pi(\mu, \nu) = B_{1 - \pi}(\nu, \mu)$~.
\end{lemma}

\begin{proof}\leanok
\uses{}

\end{proof}


\begin{lemma}
  \label{lem:genBayesEstimator_bayesBinaryRisk}
  \lean{ProbabilityTheory.isArgminEstimator_binaryGenBayesEstimator}
  \leanok
  \uses{def:genBayesEstimator, def:bayesBinaryRisk}
  The generalized Bayes estimator for the Bayes binary risk with prior $\xi \in \mathcal M(\{0,1\})$ is $x \mapsto \text{if } \xi_1\frac{d \nu}{d(P \circ \xi)}(x) \le \xi_0\frac{d \mu}{d(P \circ \xi)}(x) \text{ then } 0 \text{ else } 1$, i.e. it is equal to $\mathbb{I}_E$ for $E = \{x \mid \xi_1\frac{d \nu}{d(P \circ \xi)}(x) > \xi_0\frac{d \mu}{d(P \circ \xi)}(x)\}$~.
\end{lemma}

\begin{proof}\leanok
\uses{lem:bayesInv_binary}
The generalized Bayes estimator is defined by
\begin{align*}
x \mapsto \arg\min_z P_\xi^\dagger(x)\left[\theta \mapsto \ell'(y(\theta), z)\right] \: .
\end{align*}
For the simple binary hypothesis testing problem,
\begin{align*}
P_\xi^\dagger(x)\left[\theta \mapsto \ell'(y(\theta), z)\right]
&= P_\xi^\dagger(x)(\theta \ne z)
\: .
\end{align*}
For $z = 0$, this is $P_\xi^\dagger(x)(\{1\})$ and for $z = 1$, this is $P_\xi^\dagger(x)(\{0\})$.
The generalized Bayes estimator is $x \mapsto \text{if } P_\xi^\dagger(x)(\{1\}) \le P_\xi^\dagger(x)(\{0\}) \text{ then } 0 \text{ else } 1$.

The expression of the lemma statement is obtained by replacing $P_\xi^\dagger$ by its value (Lemma~\ref{lem:bayesInv_binary}).
\end{proof}


\begin{lemma}
  \label{lem:bayesBinaryRisk_eq_event}
  \lean{ProbabilityTheory.bayesBinaryRisk_eq_iInf_measurableSet}
  \leanok
  \uses{def:bayesRisk}
  \begin{align*}
  \mathcal B_\xi(\mu, \nu) = \inf_{E \text{ event}} \left( \xi_0 \mu(E) + \xi_1 \nu(E^c) \right) \: .
  \end{align*}
\end{lemma}

\begin{proof}\leanok
\uses{lem:genBayesEstimator_bayesBinaryRisk}
By definition,
\begin{align*}
\mathcal B_\xi(\mu, \nu)
= \inf_{\hat{y} : \mathcal X \rightsquigarrow \{0,1\}}\left(\xi_0 (\hat{y} \circ \mu)(\{1\}) + \xi_1 (\hat{y} \circ \nu)(\{0\})\right)
\: .
\end{align*}
If we restrict the infimum to deterministic kernels corresponding to functions of the form $x \mapsto \mathbb{I}\{x \in E\}$, we get that $\mathcal B_\xi(\mu, \nu)$ is bounded from above by the expression for which we want to show equality.

Since the Bayes estimator has that shape (Lemma~\ref{lem:genBayesEstimator_bayesBinaryRisk}), $\mathcal B_\xi(\mu, \nu)$ is also greater than the infimum and we have equality.
\end{proof}


\begin{theorem}
  \label{thm:bayesBinaryRisk_eq}
  \lean{ProbabilityTheory.bayesBinaryRisk_eq_lintegral_min, ProbabilityTheory.toReal_bayesBinaryRisk_eq_integral_min}
  \leanok
  \uses{def:bayesBinaryRisk}
  The Bayes risk of simple binary hypothesis testing for prior $\xi \in \mathcal M(\{0,1\})$ is
  \begin{align*}
  \mathcal B_\xi(\mu, \nu) = (P \circ \xi)\left[x \mapsto \min \left\{\xi_0\frac{d \mu}{d(P \circ \xi)}(x), \xi_1\frac{d \nu}{d(P \circ \xi)}(x)\right\}\right]
  \: .
  \end{align*}
\end{theorem}

\begin{proof}\leanok
\uses{lem:bayesianRisk_genBayesEstimator, thm:isBayesEstimator_genBayesEstimator, lem:bayesInv_binary}
By Theorem~\ref{thm:isBayesEstimator_genBayesEstimator}, the generalized Bayes estimator is a Bayes estimator, hence it suffices to compute its Bayesian risk.
By Lemma~\ref{lem:bayesianRisk_genBayesEstimator},
\begin{align*}
R_\xi(\hat{y}_B)
&= (P \circ \xi)\left[x \mapsto \inf_{z \in \mathcal Z} P_\xi^\dagger(x) \left[\theta \mapsto \ell'(y(\theta), z)\right]\right]
\\
&= (P \circ \xi)\left[x \mapsto \min \left\{P_\xi^\dagger(x)(\{0\}), P_\xi^\dagger(x)(\{1\})\right\}\right]
\\
&= (P \circ \xi)\left[x \mapsto \min \left\{\xi_0\frac{d \mu}{d(P \circ \xi)}(x), \xi_1\frac{d \nu}{d(P \circ \xi)}(x)\right\}\right]
\: .
\end{align*}
The last line is obtained by replacing $P_\xi^\dagger$ by its value (Lemma~\ref{lem:bayesInv_binary}).
\end{proof}


\begin{corollary}
  \label{cor:bayesBinaryRisk_eq_abs}
  \lean{ProbabilityTheory.toReal_bayesBinaryRisk_eq_integral_abs}
  \leanok
  \uses{def:bayesBinaryRisk}
  \begin{align*}
  \mathcal B_\xi(\mu, \nu) = \frac{1}{2}\left((P \circ \xi)(\mathcal X) -  (P \circ \xi)\left[x \mapsto \left\vert \xi_0\frac{d \mu}{d(P \circ \xi)}(x) - \xi_1\frac{d \nu}{d(P \circ \xi)}(x)\right\vert\right] \right)
  \: .
  \end{align*}
\end{corollary}

\begin{proof}\leanok
\uses{thm:bayesBinaryRisk_eq}
Use Theorem~\ref{thm:bayesBinaryRisk_eq} and the equality $\min\{a,b\} = \frac{1}{2}(a + b - \vert a - b \vert)$.
\end{proof}


\begin{lemma}
  \label{lem:bayesBinaryRisk_bernoulli}
  %\lean{}
  %\leanok
  \uses{def:bayesBinaryRisk}
  Let $\hat{y}_B$ be the generalized Bayes estimator for simple binary hypothesis testing.
  The distribution $\ell \circ (\mathrm{id} \parallel \hat{y}_B) \circ (\pi \otimes P)$ (in which $\ell$ stands for the associated deterministic kernel) is a Bernoulli with mean $\mathcal B_\pi(\mu, \nu)$.
\end{lemma}

\begin{proof}%\leanok
\uses{thm:isBayesEstimator_genBayesEstimator}
$\ell$ takes values in $\{0,1\}$ so the law has to be Bernoulli. It has mean $\mathcal B_\pi(\mu, \nu)$ because the expected risk of the generalized Bayes estimator is $\mathcal B_\pi(\mu, \nu)$.
\end{proof}



TODO: find a way to hide the following dummy lemma
\begin{lemma}[Dummy lemma: bayesBinaryRisk properties]
  \label{lem:bayesBinaryRisk_properties}
  %\lean{}
  \leanok
  \uses{def:bayesBinaryRisk}
  Dummy node to summarize properties of the Bayes binary risk.
\end{lemma}

\begin{proof}\leanok
\uses{
  lem:bayesBinaryRisk_mul,
  lem:bayesBinaryRisk_one_one,
  thm:data_proc_bayesBinaryRisk,
  lem:bayesBinaryRisk_self,
  lem:bayesBinaryRisk_nonneg,
  lem:bayesBinaryRisk_le,
  lem:bayesBinaryRisk_symm,
  lem:genBayesEstimator_bayesBinaryRisk,
  lem:bayesBinaryRisk_eq_event,
  thm:bayesBinaryRisk_eq,
  cor:bayesBinaryRisk_eq_abs,
  lem:bayesBinaryRisk_bernoulli
}
\end{proof}


\section{Reduction to testing}

TODO: lower bounds on estimation by reduction to testing problems.

%%FILE blueprint/src/sections/rn_deriv.tex
\chapter{Radon-Nikodym derivatives and disintegration}
\label{app:rnDeriv}

For $\mu$ and $\nu$ two $\sigma$-finite measures on a measurable space $\mathcal{X}$, we can decompose $\mu$ into $\frac{d \mu}{d \nu} \cdot \nu + \mu_{\perp \nu}$, where $\frac{d \mu}{d \nu}$ is a measurable function, unique $\nu$-a.e., and $\mu_{\perp \nu} \perp \nu$.

If we then introduce two kernels $\kappa, \eta : \mathcal X \rightsquigarrow \mathcal Y$, we will want to compare the joint measures $\mu \otimes \kappa$ and $\nu \otimes \eta$, as well as the compositions $\kappa \circ \mu$ and $\eta \circ \nu$.


\section{Kernel Radon-Nikodym derivative}

In standard Borel spaces, we can define a Radon-Nikodym derivative for finite kernels.
This is a measurable function $\frac{d \kappa}{d \eta} : \mathcal X \times \mathcal Y \to \mathbb{R}_{+, \infty}$ with $\kappa = \frac{d \kappa}{d \eta} \cdot \eta + \kappa_{\perp \eta}$, where for all $x$, $\kappa_{\perp \eta}(x) \perp \eta(x)$.
The hypothesis that the spaces are standard Borel is used to ensure the measurability of the function on the product space.
Taking for all $x$ the ratio $\frac{d \kappa(x)}{d \eta(x)}$ would give us measurable function $\mathcal Y \to \mathbb{R}_{+, \infty}$ for all $x \in \mathcal X$, but not necessarily a jointly measurable function.

\begin{definition}
  \label{def:kernel_rnDeriv}
  \lean{ProbabilityTheory.Kernel.rnDeriv}
  \mathlibok
  \uses{lem:kernel_properties}
  Let $\kappa, \eta : \mathcal X \rightsquigarrow \mathcal Y$ be two finite kernels, with either $\mathcal X$ countable or $\mathcal{Y}$ countably generated. The Radon-Nikodym derivative of $\kappa$ with respect to $\eta$, denoted by $\frac{d \kappa}{d \eta}$, is a measurable function $\mathcal X \times \mathcal Y \to \mathbb{R}_{+, \infty}$ with $\kappa = \frac{d \kappa}{d \eta} \cdot \eta + \kappa_{\perp \eta}$, where for all $x$, $\kappa_{\perp \eta}(x) \perp \eta(x)$.
\end{definition}


\begin{lemma}
  \label{lem:rnDeriv_unique}
  \lean{ProbabilityTheory.Kernel.eq_rnDeriv_measure}
  \mathlibok
  \uses{def:kernel_rnDeriv}
  Let $\kappa, \eta : \mathcal X \rightsquigarrow \mathcal Y$ be two finite kernels, with either $\mathcal X$ countable or $\mathcal{Y}$ countably generated. If for some $f$ and $\xi$, $\kappa = f \cdot \eta + \xi$ with $\xi(x) \perp \eta(x)$ for all $x$, then for all $x$, $f(x, y) = \frac{d \kappa(x)}{d \eta(x)}(y)$ for $\eta(x)$-almost all $y \in \mathcal Y$.
\end{lemma}

\begin{proof}\mathlibok
Let $f$ and $\xi$ be such that $\kappa = f \cdot \eta + \xi$ with $\xi(x) \perp \eta(x)$ for all $x$. Then for all $x \in \mathcal X$, $\kappa(x) = f(x, \cdot) \cdot \eta(x) + \xi(x)$ with $\xi(x) \perp \eta(x)$. By the uniqueness result for the Radon-Nikodym derivative of measures, $f(x, \cdot) = \frac{d \kappa(x)}{d \eta(x)}$ almost everywhere.
\end{proof}


\begin{corollary}
  \label{cor:rnDeriv_value}
  \lean{ProbabilityTheory.Kernel.rnDeriv_eq_rnDeriv_measure}
  \mathlibok
  \uses{def:kernel_rnDeriv}
  For all $x \in \mathcal X$, for $\eta(x)$-almost all $y \in \mathcal Y$, $\frac{d \kappa}{d \eta}(x, y) = \frac{d \kappa(x)}{d \eta(x)}(y)$.
\end{corollary}

\begin{proof}\mathlibok
\uses{lem:rnDeriv_unique}
Apply Lemma~\ref{lem:rnDeriv_unique}.
\end{proof}



\section{Composition-product}

\subsection{Absolute Continuity}

\begin{lemma}
  \label{lem:ac_compProd_iff}
  \lean{MeasureTheory.Measure.absolutelyContinuous_compProd_of_compProd, MeasureTheory.Measure.absolutelyContinuous_of_compProd}
  \mathlibok
  \uses{lem:kernel_properties}
  Let $\mu, \nu$ be two $\sigma$-finite measures on $\mathcal X$ and let $\kappa, \eta : \mathcal X \rightsquigarrow \mathcal Y$ be two s-finite kernels. Then
  \begin{itemize}
    \item[(1a)] if $\mu \otimes \kappa \ll \nu \otimes \eta$ then $\mu \otimes \kappa \ll \mu \otimes \eta$,
    \item[(1b)] if $\mu \otimes \kappa \ll \nu \otimes \eta$ and $\kappa(x) \ne 0$ for all $x$ then $\mu \ll \nu$,
    \item[(2)]  if $\mu \ll \nu$ and $\mu \otimes \kappa \ll \mu \otimes \eta$ then $\mu \otimes \kappa \ll \nu \otimes \eta$.
  \end{itemize}
  In particular,
  \begin{itemize}
    \item if $\kappa(x) \ne 0$ for all $x$ then
      $\mu \otimes \kappa \ll \nu \otimes \eta \iff \left( \mu \ll \nu \ \wedge \ \mu \otimes \kappa \ll \mu \otimes \eta \right)$~.
    \item If $\mu \ll \nu$ then $\mu \otimes \kappa \ll \nu \otimes \eta \iff \mu \otimes \kappa \ll \mu \otimes \eta$~.
  \end{itemize}
  
\end{lemma}

\begin{proof}\mathlibok
\uses{lem:kernel_properties}

\end{proof}


\begin{lemma}
  \label{lem:mutuallySingular_compProd}
  \lean{MeasureTheory.Measure.MutuallySingular.compProd_of_left, MeasureTheory.Measure.mutuallySingular_compProd_iff, MeasureTheory.Measure.mutuallySingular_of_mutuallySingular_compProd}
  \mathlibok
  \uses{lem:kernel_properties}
  Let $\mu, \nu$ be two $\sigma$-finite measures on $\mathcal X$ and let $\kappa, \eta : \mathcal X \rightsquigarrow \mathcal Y$ be two finite kernels.
  Let $\mu \sqcap \nu$ denote the infimum of $\mu$ and $\nu$.
  Then
  \begin{enumerate}
    \item if $\mu \perp \nu$ then $\mu \otimes \kappa \perp \nu \otimes \eta$,
    \item $\mu \otimes \kappa \perp \nu \otimes \eta \iff (\mu \sqcap \nu) \otimes \kappa \perp (\mu \sqcap \nu) \otimes \eta$, and the same holds for any measure which is equivalent to $\mu \sqcap \nu$, like $\frac{d \mu}{d \nu} \cdot \nu$~,
    \item if $\mu \otimes \kappa \perp \nu \otimes \eta$ then for $(\mu \sqcap \nu)$-almost every $x$, $\kappa(x) \perp \eta(x)$.
  \end{enumerate}
\end{lemma}

\begin{proof}\mathlibok
\uses{}
First, let's state two facts about mutually singular measures that we will use without proof:
\begin{itemize}
  \item $\mu \perp (\nu + \nu') \iff \mu \perp \nu \ \wedge \ \mu \perp \nu'$~,
  \item if $\mu \ll \mu'$ and  $\mu' \ll \mu$ then $\mu \perp \nu \iff \mu' \perp \nu$~.
\end{itemize}

\emph{1.} Let $s$ be a measurable set of $\mathcal X$ such that $\mu(s) = 0$ and $\nu(s^c) = 0$. One can check that the sets $s \times \mathcal Y$ and $s^c \times \mathcal Y$ demonstrate $\mu \otimes \kappa \perp \nu \otimes \eta$.

\emph{2.} Write $\mu = \frac{d \mu}{d \nu} \cdot \nu + \mu_{\perp \nu}$ and $\nu = \frac{d \nu}{d \mu} \cdot \mu + \nu_{\perp \mu}$.
Then
\begin{align*}
\mu \otimes \kappa \perp \nu \otimes \eta
&\iff \left(\frac{d \mu}{d \nu} \cdot \nu\right) \otimes \kappa \perp \left(\frac{d \nu}{d \mu} \cdot \mu\right) \otimes \eta
\\&\qquad \wedge \ \left(\frac{d \mu}{d \nu} \cdot \nu\right) \otimes \kappa \perp \nu_{\perp \mu} \otimes \eta
\\&\qquad \wedge \ \mu_{\perp \nu} \otimes \kappa \perp \left(\frac{d \nu}{d \mu} \cdot \mu\right) \otimes \eta
\\&\qquad \wedge \ \mu_{\perp \nu} \otimes \kappa \perp \nu_{\perp \mu} \otimes \eta
\end{align*}
The three last lines are true since $\mu \perp \nu_{\perp \mu}$ and $\mu_{\perp \nu} \perp \nu$. Only the first line remains.
It suffices then to prove that $\mu \sqcap \nu$, $\frac{d \mu}{d \nu} \cdot \nu$ and $\frac{d \nu}{d \mu} \cdot \mu$ are equivalent.

TODO

\emph{3.} By 2. it suffices to consider the case $\nu = \mu$ and show that for $\mu$-almost every $x$, $\kappa(x) \perp \eta(x)$.
Let $s$ be a measurable set of $\mathcal X \times \mathcal Y$ such that $(\mu \otimes \kappa)(s) = 0$ and $(\mu \otimes \eta)(s^c) = 0$.
Then
\begin{align*}
0 &= (\mu \otimes \kappa)(s)
\\
&= \int_{x} \kappa(x)(\{y \mid (x, y) \in s\}) \partial \mu
\: .
\end{align*}
Hence for $\mu$-almost all $x$, $\kappa(x)(\{y \mid (x, y) \in s\}) = 0$. Similarly, for $\mu$-almost all $x$, $\eta(x)(\{y \mid (x, y) \in s^c\}) = 0$. Since $\{y \mid (x, y) \in s^c\} = \{y \mid (x, y) \in s\}^c$, we have a measurable set witnessing $\kappa(x) \perp \eta(x)$ for $\mu$-almost all $x$.
\end{proof}



\subsection{Radon-Nikodym derivative and singular part}

\begin{lemma}
  \label{lem:singularPart_compProd}
  %\lean{}
  %\leanok
  \uses{lem:kernel_properties}
  Let $\mu, \nu$ be two $\sigma$-finite measures on $\mathcal X$ and let $\kappa, \eta : \mathcal X \rightsquigarrow \mathcal Y$ be two s-finite kernels.
  We denote $\frac{d\mu}{d\nu}\cdot \nu$ by $\mu_{\parallel \nu}$.
  Then
  \begin{align*}
  (\mu \otimes \kappa)_{\perp (\nu \otimes \eta)} = \mu_{\perp \nu} \otimes \kappa + (\mu_{\parallel \nu} \otimes \kappa)_{\perp (\mu_{\parallel \nu} \otimes \eta)}
  \: .
  \end{align*}
\end{lemma}

\begin{proof}%\leanok
\uses{}

\end{proof}


\begin{lemma}
  \label{lem:rnDeriv_eq_ac_left}
  \lean{ProbabilityTheory.rnDeriv_compProd_withDensity_rnDeriv}
  \mathlibok
  \uses{}
  Let $\mu, \nu$ be two finite measures on $\mathcal X$ and let $\kappa, \eta : \mathcal X \rightsquigarrow \mathcal Y$ be two finite kernels.
  Let $\mu_{\parallel \nu} = \left(\frac{\partial \mu}{\partial \nu}\right) \cdot \nu$.
  Then for $(\nu \otimes \eta)$-almost all $z$, $\frac{d (\mu_{\parallel \nu} \otimes \kappa)}{d (\nu \otimes \eta)}(z) = \frac{d (\mu \otimes \kappa)}{d (\nu \otimes \eta)}(z)$.
\end{lemma}

\begin{proof}\mathlibok
\uses{lem:kernel_properties}
\end{proof}


\begin{lemma}
  \label{cor:rnDeriv_compProd_left}
  \lean{ProbabilityTheory.rnDeriv_measure_compProd_left}
  \mathlibok
  \uses{}
  Let $\mu, \nu \in \mathcal M(\mathcal X)$ and let $\kappa : \mathcal X \rightsquigarrow \mathcal Y$ be a finite kernel. Then for $(\nu \otimes \kappa)$-almost all $(x, y)$, $\frac{d (\mu \otimes \kappa)}{d (\nu \otimes \kappa)}(x,y) = \frac{d\mu}{d\nu}(x)$.
\end{lemma}

\begin{proof}\mathlibok
\uses{lem:rnDeriv_eq_ac_left}
We can suppose $\mu \ll \nu$ without loss of generality (by Lemma~\ref{lem:rnDeriv_eq_ac_left}).
That implies $\mu \otimes \kappa \ll \nu \otimes \kappa$.
It suffices to show that the integrals of the two functions agree on all sets in the $\pi$-system of products of measurable sets. Let $s, t$ be two measurable sets of $\mathcal X$ and $\mathcal Y$ respectively.
\begin{align*}
\int_{p \in s \times t} \frac{d (\mu \otimes \kappa)}{d (\nu \otimes \kappa)}(p) \partial(\nu \otimes \kappa)
&= (\mu \otimes \kappa) (s \times t)
\: .
\end{align*}

\begin{align*}
\int_{p \in s \times t} \frac{d \mu}{d \nu}(p_1) \partial(\nu \otimes \kappa)
&= \int_{x \in s} \int_{y \in t} \frac{d \mu}{d \nu}(x) \partial \kappa(x) \partial \nu
\\
&= \int_{x \in s} \frac{d \mu}{d \nu}(x) \kappa(x)(t) \partial \nu
\\
&= \int_{x \in s} \kappa(x)(t) \partial \mu
\\
&= (\mu \otimes \kappa) (s \times t)
\: .
\end{align*}
\end{proof}


\begin{lemma}
  \label{lem:rnDeriv_chain}
  \lean{MeasureTheory.Measure.rnDeriv_mul_rnDeriv, MeasureTheory.Measure.rnDeriv_mul_rnDeriv'}
  \mathlibok
  \uses{}
  Let $\mu, \nu, \xi$ be $\sigma$-finite measures on $\mathcal X$.
  \begin{enumerate}
    \item If $\mu \ll \nu$ then $\xi$-almost surely, $\frac{d \mu}{d \xi} = \frac{d \mu}{d \nu} \frac{d \nu}{d \xi}$.
    \item If $\nu \ll \xi$ then $\nu$-almost surely, $\frac{d \mu}{d \xi} = \frac{d \mu}{d \nu} \frac{d \nu}{d \xi}$.
  \end{enumerate}
\end{lemma}

\begin{proof}\mathlibok
\uses{}
\end{proof}


\begin{theorem}[Chain rule for Radon-Nikodym derivatives]
  \label{thm:rnDeriv_chain_compProd}
  \lean{ProbabilityTheory.rnDeriv_compProd}
  \mathlibok
  \uses{}
  Let $\mu, \nu$ be two finite measures on $\mathcal X$ and let $\kappa, \eta : \mathcal X \rightsquigarrow \mathcal Y$ be two finite kernels with $\mu \otimes \kappa \ll \mu \otimes \eta$. Then for $(\nu \otimes \eta)$-almost all $(x,y)$,
  \begin{align*}
  \frac{d(\mu \otimes \kappa)}{d(\nu \otimes \eta)}(x, y)
  = \frac{d \mu}{d \nu}(x) \frac{d(\mu \otimes \kappa)}{d(\mu \otimes \eta)}(x, y)
  \: .
  \end{align*}
\end{theorem}

\begin{proof}\mathlibok
\uses{lem:rnDeriv_chain,cor:rnDeriv_compProd_left}
By the first point of Lemma~\ref{lem:rnDeriv_chain}, $(\nu \otimes \eta)$-almost surely,
\begin{align*}
\frac{d(\mu \otimes \kappa)}{d(\nu \otimes \eta)}
  = \frac{d(\mu \otimes \kappa)}{d(\mu \otimes \eta)} \frac{d (\mu \otimes \eta)}{d (\nu \otimes \eta)} 
\end{align*}
Then, by Lemma~\ref{cor:rnDeriv_compProd_left}, $\frac{d (\mu \otimes \eta)}{d (\nu \otimes \eta)}$ is almost everywhere equal to $\frac{d\mu}{d\nu}$.
\end{proof}


\begin{lemma}
  \label{lem:rnDeriv_eq_ac}
  %\lean{}
  %\leanok
  \uses{def:kernel_rnDeriv}
  Let $\mu, \nu$ be two measures on $\mathcal X$ and let $\kappa, \eta : \mathcal X \rightsquigarrow \mathcal Y$ be two finite kernels, with either $\mathcal X$ countable or $\mathcal{Y}$ countably generated.
  Let $\mu' = \left(\frac{\partial \mu}{\partial \nu}\right) \cdot \nu$ and $\kappa' = \left(\frac{\partial \kappa}{\partial \eta}\right) \cdot \eta$.
  Then for $(\nu \otimes \eta)$-almost all $z$, $\frac{\partial(\mu' \otimes \kappa')}{\partial (\nu \otimes \eta)}(z) = \frac{\partial(\mu \otimes \kappa)}{\partial (\nu \otimes \eta)}(z)$.
\end{lemma}

\begin{proof} %\leanok
\end{proof}


\begin{lemma}
  \label{cor:rnDeriv_compProd_right}
  \lean{ProbabilityTheory.rnDeriv_measure_compProd_right}
  \mathlibok
  \uses{def:kernel_rnDeriv}
  Let $\mu \in \mathcal M(\mathcal X)$ be a finite measure and $\kappa, \eta : \mathcal X \rightsquigarrow \mathcal Y$ be two finite kernels, with either $\mathcal X$ countable or $\mathcal{Y}$ countably generated.
  Then $(\mu \otimes \eta)$-almost surely,
  \begin{align*}
  \frac{d (\mu \otimes \kappa)}{d (\mu \otimes \eta)} = \frac{d \kappa}{d \eta}
  \: .
  \end{align*}
\end{lemma}

\begin{proof}\mathlibok
\uses{lem:rnDeriv_eq_ac}

\end{proof}

That last lemma is significant because it means that we can get a kind of kernel derivative without any assumptions on the spaces.
As long as we are happy with $\mu$-almost sure statements, we can use $\frac{d (\mu \otimes \kappa)}{d (\mu \otimes \eta)}$ (which is jointly measurable) instead of the possibly non-existent $\frac{d \kappa}{d \eta}$.


\begin{lemma}
  \label{lem:ae_rnDeriv_ne_zero}
  \lean{MeasureTheory.Measure.ae_rnDeriv_ne_zero_imp_of_ae}
  \mathlibok
  \uses{}
  Let $\mu, \nu \in \mathcal M(\mathcal X)$ be two $\sigma$-finite measures and let $p$ be a predicate on $\mathcal X$.
  If $p$ is true $\mu$-almost surely, then for $\nu$-almost all $x$, either $\frac{d\mu}{d\nu}(x) = 0$ or $p(x)$.
\end{lemma}

\begin{proof}\mathlibok
\uses{}

\end{proof}


\begin{lemma}
  \label{lem:rnDeriv_compProd}
  \lean{ProbabilityTheory.rnDeriv_measure_compProd}
  \mathlibok
  \uses{def:kernel_rnDeriv}
  Let $\mu, \nu$ be two finite measures on $\mathcal X$ and let $\kappa, \eta : \mathcal X \rightsquigarrow \mathcal Y$ be two finite kernels, with either $\mathcal X$ countable or $\mathcal{Y}$ countably generated.
  Then for $(\nu \otimes \eta)$-almost all $(x, y)$,
  \begin{align*}
  \frac{d (\mu \otimes \kappa)}{d (\nu \otimes \eta)}(x,y) = \frac{d\mu}{d\nu}(x)\frac{d \kappa}{d \eta}(x,y)
  \: .
  \end{align*}
  This implies that the equality is true for $\nu$-almost all $x$, for $\eta(x)$-almost all $y$.
\end{lemma}

\begin{proof}\mathlibok
\uses{thm:rnDeriv_chain_compProd, cor:rnDeriv_compProd_right, lem:ae_rnDeriv_ne_zero, cor:rnDeriv_compProd_left}
First, by Theorem~\ref{thm:rnDeriv_chain_compProd}, for $(\nu \otimes \eta)$-almost all $(x,y)$,
\begin{align*}
\frac{d(\mu \otimes \kappa)}{d(\nu \otimes \eta)}(x, y)
= \frac{d \mu}{d \nu}(x) \frac{d(\mu \otimes \kappa)}{d(\mu \otimes \eta)}(x, y)
\: .
\end{align*}
It remains to show that $\frac{d \mu}{d \nu}(x) \frac{d(\mu \otimes \kappa)}{d(\mu \otimes \eta)}(x, y)$ is $(\nu \otimes \eta)$-almost everywhere equal to $\frac{d\mu}{d\nu}(x)\frac{d \kappa}{d \eta}(x,y)$.
By Lemma~\ref{cor:rnDeriv_compProd_right}, $(\mu \otimes \eta)$-almost surely,
\begin{align*}
\frac{d \kappa}{d \eta} = \frac{d (\mu \otimes \kappa)}{d (\mu \otimes \eta)}
\: .
\end{align*}
We then use Lemma~\ref{lem:ae_rnDeriv_ne_zero} to turn that into a $(\nu \otimes \eta)$-almost sure equality:
\begin{align*}
\frac{d(\mu \otimes \eta)}{d(\nu \otimes \eta)}\frac{d \kappa}{d \eta} = \frac{d(\mu \otimes \eta)}{d(\nu \otimes \eta)}\frac{d (\mu \otimes \kappa)}{d (\mu \otimes \eta)}
\: .
\end{align*}
Finally, by Lemma~\ref{cor:rnDeriv_compProd_left}, $\frac{d(\mu \otimes \eta)}{d(\nu \otimes \eta)} = \frac{d\mu}{d\nu}$ almost surely.
\end{proof}


\section{Composition}


\begin{lemma}
  \label{lem:rnDeriv_map_eq_condexp}
  \lean{MeasureTheory.toReal_rnDeriv_map}
  \mathlibok
  \uses{}
  Let $\mu, \nu \in \mathcal M(\mathcal X)$ with $\mu \ll \nu$, $g : \mathcal X \to \mathcal Y$ a measurable function and denote by $g^* \mathcal Y$ the comap of the $\sigma$-algebra on $\mathcal Y$ by $g$.
  Then $\nu$-almost everywhere,
  \begin{align*}
  \frac{d g_*\mu}{d g_*\nu}(g(x)) = \nu\left[ \frac{d \mu}{d \nu} \mid g^* \mathcal Y\right](x)
  \: .
  \end{align*}
\end{lemma}

\begin{proof}\mathlibok
\uses{}
We show that the integrals of the two functions agree on all $g^* \mathcal Y$-measurable sets. It suffices to show equality on all sets oc $\mathcal X$ of the form $g^{-1}(t)$ for $t$ a measurable set of $\mathcal Y$.
TODO: also show measurability.
\begin{align*}
\int_{x \in g^{-1}(t)}\frac{d g_*\mu}{d g_*\nu}(g(x)) \partial\nu
&= \int_{y \in t}\frac{d g_*\mu}{d g_*\nu}(y) \partial(g_*\nu)
\\
&= g_*\mu (t)
\: .
\end{align*}

\begin{align*}
\int_{x \in g^{-1}(t)}\nu\left[ \frac{d \mu}{d \nu} \mid g^* \mathcal Y\right](x) \partial\nu
&= \int_{x \in g^{-1}(t)}\frac{d \mu}{d \nu}(x) \partial\nu
\\
&= \mu(g^{-1}(t))
\\
&= g_*\mu(t)
\: .
\end{align*}


\end{proof}


\begin{lemma}
  \label{lem:rnDeriv_comp_eq_condexp}
  %\lean{ProbabilityTheory.toReal_rnDeriv_comp_eq_condexp_compProd}
  %\leanok
  \uses{}
  Let $\mu, \nu \in \mathcal M(\mathcal X)$ with $\mu \ll \nu$ and let $\kappa, \eta : \mathcal X \rightsquigarrow \mathcal Y$ be finite kernels with $\kappa(x) \ll \eta(x)$ $\nu$-a.e..
  Let $\mathcal B$ be the sigma-algebra on $\mathcal X \times \mathcal Y$ obtained by taking the comap of the sigma-algebra of $\mathcal Y$ by the projection.
  Then for $(\nu \otimes \eta)$-almost every $(x,y)$,
  \begin{align*}
  \frac{d(\kappa \circ \mu)}{d(\eta \circ \nu)}(y)
  &= (\nu \otimes \eta)\left[ \frac{d (\mu \otimes \kappa)}{d (\nu \otimes \eta)} \mid \mathcal B \right](x,y)
  \: .
  \end{align*}
\end{lemma}

\begin{proof}\leanok
\uses{lem:rnDeriv_map_eq_condexp}
Let $\pi_Y : \mathcal X \times \mathcal Y \to \mathcal Y$ be the projection $\pi_Y(x,y) = y$.
Remark that $\kappa \circ \mu = \pi_{Y*}(\mu \otimes \kappa)$ and similarly for $\nu, \eta$.
By Lemma~\ref{lem:rnDeriv_map_eq_condexp}, $(\nu \otimes \eta)$-almost everywhere,
\begin{align*}
\frac{d(\kappa \circ \mu)}{d(\eta \circ \nu)}(y)
&= \frac{d \pi_{Y*}(\mu \otimes \kappa)}{d \pi_{Y*}(\nu \otimes \eta)}(\pi_Y((x,y)))
\\
&= (\nu \otimes \eta)\left[ \frac{d (\mu \otimes \kappa)}{d (\nu \otimes \eta)} \mid \mathcal B\right](x,y)
\: .
\end{align*}
\end{proof}


\begin{lemma}[\cite{csiszar1963informationstheoretische}]
  \label{lem:rnDeriv_comp_eq_condexp_right}
  %\lean{ProbabilityTheory.toReal_rnDeriv_comp}
  %\leanok
  \uses{}
  Let $\mu, \nu \in \mathcal M(\mathcal X)$ with $\mu \ll \nu$ and let $\kappa : \mathcal X \rightsquigarrow \mathcal Y$ be a finite kernel.
  Let $\mathcal B$ be the sigma-algebra on $\mathcal X \times \mathcal Y$ obtained by taking the comap of the sigma-algebra of $\mathcal Y$ by the projection.
  Then for $(\nu \otimes \kappa)$-almost every $(x,y)$,
  \begin{align*}
  \frac{d(\kappa \circ \mu)}{d(\kappa \circ \nu)}(y)
  &= (\nu \otimes \kappa)\left[ (x, y) \mapsto \frac{d \mu}{d \nu}(x) \mid \mathcal B \right](x,y)
  \: .
  \end{align*}
  
\end{lemma}

\begin{proof}\leanok
\uses{lem:rnDeriv_comp_eq_condexp,cor:rnDeriv_compProd_left}
Apply Lemma~\ref{lem:rnDeriv_comp_eq_condexp} to get, $(\nu \otimes \kappa)$-almost everywhere,
\begin{align*}
\frac{d(\kappa \circ \mu)}{d(\kappa \circ \nu)}(y)
&= (\nu \otimes \kappa)\left[ \frac{d (\mu \otimes \kappa)}{d (\nu \otimes \kappa)} \mid \mathcal B\right](x,y)
\: .
\end{align*}
Finally, by Corollary~\ref{cor:rnDeriv_compProd_left}, $\frac{d (\mu \otimes \kappa)}{d (\nu \otimes \kappa)} = \frac{d \mu}{d \nu}$ a.e.
\end{proof}


\section{Sigma-algebras}

For $\mathcal A$ a sub-$\sigma$-algebra and $\mu$ a measure, we write $\mathcal \mu_{| \mathcal A}$ for the measure restricted to the $\sigma$-algebra.

\begin{lemma}
  \label{lem:rnDeriv_trim_of_ac}
  \lean{MeasureTheory.toReal_rnDeriv_trim}
  \mathlibok
  %\uses{}
  Let $\mu, \nu$ be two finite measures on $\mathcal X$ with $\mu \ll \nu$ and let $\mathcal A$ be a sub-$\sigma$-algebra of $\mathcal X$.
  Then $\frac{d \mu_{| \mathcal A}}{d \nu_{| \mathcal A}}$ is $\nu_{| \mathcal A}$-almost everywhere (hence also $\nu$-a.e.) equal to $\nu\left[ \frac{d \mu}{d \nu} \mid \mathcal A\right]$.
\end{lemma}

\begin{proof}\mathlibok
\uses{lem:rnDeriv_map_eq_condexp}
The restriction $\mu_{| \mathcal A}$ is the map of $\mu$ by the identity, seen as a function from $\mathcal X$ with its $\sigma$-algebra to $\mathcal X$ with $\mathcal{A}$.
We can thus apply Lemma~\ref{lem:rnDeriv_map_eq_condexp}.
\end{proof}


\begin{lemma}
  \label{lem:rnDeriv_map_eq_rnDeriv_trim}
  %\lean{}
  %\leanok
  \uses{}
  Let $\mu, \nu \in \mathcal M(\mathcal X)$ with $\mu \ll \nu$, $g : \mathcal X \to \mathcal Y$ a measurable function and denote by $g^* \mathcal Y$ the comap of the $\sigma$-algebra on $\mathcal Y$ by $g$.
  Then $\nu$-almost everywhere,
  \begin{align*}
  \frac{d g_*\mu}{d g_*\nu}(g(x)) = \frac{d \mu_{| g^* \mathcal Y}}{d \nu_{| g^* \mathcal Y}}(x)
  \: .
  \end{align*}
\end{lemma}

\begin{proof}%\leanok
\uses{lem:rnDeriv_map_eq_condexp, lem:rnDeriv_trim_of_ac}
Combine Lemma~\ref{lem:rnDeriv_map_eq_condexp} and Lemma~\ref{lem:rnDeriv_trim_of_ac}.
\end{proof}

%%FILE blueprint/src/sections/sample_complexity.tex
\section{Sample complexity}

TODO: we could generalize this to a sequence of kernels instead of iid observations.

An estimation problem of the shape $(P^{\otimes n}, y, \ell')$ corresponds to estimating from $n \in \mathbb{N}$ i.i.d. samples of the distribution. From the data-processing inequality, we get that the risk (Bayes risk, minimax risk...) decreases with the number of samples (Lemma~\ref{lem:bayesRisk_mono_prod} for example).
Intuitively, having more information allows for better estimation.

In a sequence of estimation tasks $(P^{\otimes n}, y, \ell')$ for $n \in \mathbb{N}$, the \emph{sample complexity} is the minimal $n$ such that the risk is less than a desired value.

\begin{lemma}
  \label{lem:bayesRisk_mono_prod}
  %\lean{}
  %\leanok
  \uses{def:bayesRisk}
  If $n \le m$ then $\mathcal R_\pi^{P^{\otimes n}} \ge \mathcal R_\pi^{P^{\otimes m}}$.
\end{lemma}

\begin{proof}%\leanok
\uses{thm:data_proc_bayesRisk}
Use the data-processing inequality, Theorem~\ref{thm:data_proc_bayesRisk}, for the deterministic kernel that projects on the first $n$ coordinates.
\end{proof}


\begin{definition}
  \label{def:priorSampleComplexity}
  %\lean{}
  %\leanok
  \uses{def:bayesRisk}
  The sample complexity of Bayesian estimation with respect to a prior $\pi \in \mathcal M(\Theta)$ at risk level $\delta \in \mathbb{R}_{+,\infty}$ is
  \begin{align*}
  n_\pi^P(\delta) = \min \{n \in \mathbb{N} \mid \mathcal R_\pi^{P^{\otimes n}} \le \delta\} \: .
  \end{align*}
\end{definition}


%%FILE blueprint/src/sections/sequential.tex
\section{Sequential testing}

For $\mu \in \mathcal M(\mathcal X)$, let $\mu^{\otimes \mathbb{N}}$ be the corresponding product measure on $\mathcal X^{\mathbb{N}}$.
For $n \in \mathbb{N}$, let $\mathcal F_n$ be the $\sigma$-algebra of $\mathcal X^{\mathbb{N}}$ generated by the projection to the first $n$ coordinates. Denote that projection by $\pi_{[n]}$.
$\mathcal F = (\mathcal F_n)_{n \in \mathbb{N}}$ is a filtration.
Let $\tau$ be a stopping time with respect to $\mathcal F$. $\mu_\tau$ denotes the restriction of the measure $\mu^{\otimes \mathbb{N}}$ to $\mathcal F_\tau$, the sigma-algebra generated by the stopping time $\tau$.

In sequential testing, two components need to be designed for estimation: the stopping time $\tau$ and the $\mathcal F_\tau$-measurable estimator that returns an answer given the data available at $\tau$.
We will be interested in two performance metrics:
\begin{itemize}
  \item The risk, where the data generating kernel is $P_\tau$ which takes values $\mu_\tau$ and $\nu_\tau$.
  \item The stopping time $\tau$ (through its mean, other moments, or high probability bounds).
\end{itemize}
The idea is that a good estimator stops quickly with low risk.


\begin{lemma}
  \label{lem:llr_filtration_nat}
  %\lean{}
  %\leanok
  %\uses{}
  For $\mu, \nu \in \mathcal P(\mathcal X)$ with $\mu \ll \nu$ and $n \in \mathbb{N}$, $\nu^{\otimes \mathbb{N}}_{| \mathcal F_n}$-almost surely,
  $\frac{d \mu^{\otimes \mathbb{N}}_{| \mathcal F_n}}{d \nu^{\otimes \mathbb{N}}_{| \mathcal F_n}}(x) = \prod_{m=1}^n \frac{d \mu}{d \nu}(x_m)$.
\end{lemma}

\begin{proof}%\leanok
\uses{lem:rnDeriv_map_eq_rnDeriv_trim}
By Lemma~\ref{lem:rnDeriv_map_eq_rnDeriv_trim}, $\frac{d \mu^{\otimes \mathbb{N}}_{| \mathcal F_n}}{d \nu^{\otimes \mathbb{N}}_{| \mathcal F_n}}(x) = \frac{d \pi_{[n]*}\mu^{\otimes \mathbb{N}}}{d \pi_{[n]*}\nu^{\otimes \mathbb{N}}}(\pi_{[n]}(x))$ a.s. . The pushforwards simplify: $\pi_{[n]*}\nu^{\otimes \mathbb{N}} = \nu^{\otimes n}$.
\end{proof}


\begin{lemma}
  \label{lem:llr_stopping_time}
  %\lean{}
  %\leanok
  %\uses{}
  For $\mu, \nu \in \mathcal P(\mathcal X)$ with $\mu \ll \nu$, $\nu_\tau$-almost surely,
  $\frac{d \mu_\tau}{d \nu_\tau}(x) = \prod_{n=1}^\tau \frac{d \mu}{d \nu}(x_n)$.
\end{lemma}

\begin{proof}
\uses{lem:rnDeriv_trim_of_ac, lem:llr_filtration_nat}
It suffices to show the equality of their integrals on $\mathcal F_\tau$-measurable sets. Let $E$ be such a set.
Then $E = \sum_{n=1}^\infty E \cap \{\tau = n\}$~. For $n \in \mathbb{N}$, let $E_n = E \cap \{\tau = n\}$.

We show that $\nu_\tau \left[\mathbb{I}(E_n) \frac{d \mu_\tau}{d \nu_\tau}\right] = \nu_\tau\left[x \mapsto \mathbb{I}(E_n)\prod_{m=1}^n \frac{d \mu}{d \nu}(x_m)\right]$.
The result for $E$ then follows from the monotone convergence theorem.
The product $\prod_{m=1}^n \frac{d \mu}{d \nu}(x_m)$ is equal to $\frac{d \mu^{\otimes \mathbb{N}}_{| \mathcal F_n}}{d \nu^{\otimes \mathbb{N}}_{| \mathcal F_n}}$ by Lemma~\ref{lem:llr_filtration_nat}.

$E_n$ is $\mathcal F_n$-measurable, hence (using Lemma~\ref{lem:rnDeriv_trim_of_ac})
TODO

\end{proof}


\begin{lemma}
  \label{lem:wald_equation}
  %\lean{}
  \notready
  \uses{}
  Let $\mu \in \mathcal P(\mathcal X)$ and let $(X_n)_{n \in \mathbb{N}}$ be $\mu$-integrable real random variables. Let $\tau$ be a random variable with value in $\mathbb{N}$.
  Suppose that $\mu[X_n \mathbb{I}\{\tau \ge n\}] = \mu[X_n]\mu\{\tau \ge n\}$ for all $n \in \mathbb{N}$ and that $\sum_{n=1}^\infty \mu[|X_n| \mathbb{I}\{\tau \ge n\}] < \infty$. Then the sums $S_\tau = \sum_{n=1}^\tau X_n$ and $T_\tau = \sum_{n=1}^\tau \mu[X_n]$ are integrable and $\mu[S_\tau] = \mu[T_\tau]$. 
\end{lemma}

\begin{proof}%\leanok
\uses{}

\end{proof}


\begin{theorem}
  \label{thm:kl_stopping_time}
  %\lean{}
  %\leanok
  \uses{def:KL}
  For $\mu, \nu \in \mathcal P(\mathcal X)$, if $\tau$ has finite expectation then $\KL(\mu_\tau, \nu_\tau) = \mu[\tau] \KL(\mu, \nu)$.
\end{theorem}

\begin{proof}
\uses{lem:llr_stopping_time,lem:wald_equation}
TODO: need Wald's first identity.
\end{proof}


\begin{lemma}
  \label{lem:renyiMeasure_stopping_time}
  %\lean{}
  %\leanok
  \uses{def:renyiMeasure}
  For $\mu, \nu$ two probability measures on $\mathcal X$ and $\alpha \in (0,1)$, $(\mu^{(\alpha, \nu)})_\tau = (\mu_\tau)^{(\alpha, \nu_\tau)}$.
\end{lemma}

\begin{proof}
This is a guess, I did not check it. TODO
\end{proof}


\begin{theorem}
  \label{thm:renyi_stopping_time}
  %\lean{}
  %\leanok
  \uses{def:KL, def:Renyi, def:renyiMeasure}
  For $\mu, \nu$ two probability measures on $\mathcal X$ and $\alpha \in (0,1)$,
  $R_\alpha(\mu_\tau, \nu_\tau) = \mu^{\alpha, \nu}[\tau] R_\alpha(\mu, \nu)$.
\end{theorem}

\begin{proof}
\uses{thm:kl_stopping_time, cor:renyi_eq_add_kl, lem:renyiMeasure_stopping_time}
From Corollary~\ref{cor:renyi_eq_add_kl}, $(1 - \alpha) R_\alpha(\mu, \nu) = \alpha \KL(\mu^{(\alpha, \nu)}, \mu) + (1 - \alpha)\KL(\mu^{(\alpha, \nu)}, \nu)$, hence we can use Theorem~\ref{thm:kl_stopping_time} to write
\begin{align*}
\mu^{\alpha, \nu}[\tau] R_\alpha(\mu, \nu)
&= \frac{1}{1 - \alpha}\left(\alpha \mu^{\alpha, \nu}[\tau]\KL(\mu^{(\alpha, \nu)}, \mu) + (1 - \alpha)\mu^{\alpha, \nu}[\tau]\KL(\mu^{(\alpha, \nu)}, \nu) \right)
\\
&= \frac{1}{1 - \alpha}\left(\alpha \KL((\mu^{(\alpha, \nu)})_\tau, \mu_\tau) + (1 - \alpha)\KL((\mu^{(\alpha, \nu)})_\tau, \nu_\tau) \right)
\end{align*}
It then suffices to remark that $(\mu^{(\alpha, \nu)})_\tau = (\mu_\tau)^{(\alpha, \nu_\tau)}$, thanks to Lemma~\ref{lem:renyiMeasure_stopping_time}.

TODO: if that lemma does not hold, we still have the inequality $R_\alpha(\mu_\tau, \nu_\tau) \le \mu^{\alpha, \nu}[\tau] R_\alpha(\mu, \nu)$ by Lemma~\ref{lem:renyi_eq_inf_add_kl}.
\end{proof}

%%FILE blueprint/src/sections/stat_div.tex

\section{Statistical divergences}

We define a family of divergences based on the risk of a binary estimation task.

\subsection{Statistical information}

\begin{definition}
  \label{def:statInfo}
  \lean{ProbabilityTheory.statInfo}
  \leanok
  \uses{def:bayesBinaryRisk, def:riskIncrease, lem:bayesBinaryRisk_properties}
  The statistical information between measures $\mu$ and $\nu$ with respect to prior $\xi \in \mathcal M(\{0,1\})$ is
  $\mathcal I_\xi(\mu, \nu) = \min\{\xi_0 \mu(\mathcal X), \xi_1 \nu(\mathcal X)\} - \mathcal B_\xi(\mu, \nu)$.
  This is the risk increase $I_\xi^P(d_{\mathcal X})$ in the binary hypothesis testing problem for $d_{\mathcal X} : \mathcal X \rightsquigarrow *$ the Markov kernel to the point space.
\end{definition}

The statistical information is the difference between the minimal risk of an estimator that does not depend on the data and that of a Bayes estimator.
This is a simple generalization of both the DeGroot statistical information as well as the hockey-stick (or $E_\gamma$) divergence.

\begin{lemma}
  \label{lem:statInfo_self}
  \lean{ProbabilityTheory.statInfo_self}
  \leanok
  \uses{def:statInfo}
  For $\mu \in \mathcal M(\mathcal X)$, $\mathcal I_\xi(\mu, \mu) = 0$~.
\end{lemma}

\begin{proof}\leanok
\uses{lem:bayesBinaryRisk_self}
Use Lemma~\ref{lem:bayesBinaryRisk_self}.
\end{proof}

\begin{lemma}
  \label{lem:statInfo_nonneg}
  %\lean{}
  %\leanok
  \uses{def:statInfo}
  For $\mu, \nu \in \mathcal M(\mathcal X)$, $\mathcal I_\xi(\mu, \nu) \ge 0$~.
\end{lemma}

\begin{proof}%\leanok
\uses{lem:bayesBinaryRisk_le}
Use Lemma~\ref{lem:bayesBinaryRisk_le}.
\end{proof}


\begin{lemma}
  \label{lem:statInfo_le}
  \lean{ProbabilityTheory.statInfo_le_min}
  \leanok
  \uses{def:statInfo}
  For $\mu, \nu \in \mathcal M(\mathcal X)$, $\mathcal I_\xi(\mu, \nu) \le \min\{\xi_0 \mu(\mathcal X), \xi_1 \nu(\mathcal X)\}$~.
\end{lemma}

\begin{proof}\leanok
\uses{lem:bayesBinaryRisk_nonneg}
Use Lemma~\ref{lem:bayesBinaryRisk_nonneg}.
\end{proof}


\begin{lemma}
  \label{lem:statInfo_symm}
  \lean{ProbabilityTheory.statInfo_symm}
  \leanok
  \uses{def:statInfo}
  For $\mu, \nu \in \mathcal M(\mathcal X)$ and $\xi \in \mathcal M(\{0,1\})$, $\mathcal I_\xi(\mu, \nu) = \mathcal I_{\xi_\leftrightarrow}(\nu, \mu)$~.
\end{lemma}

\begin{proof}\leanok
\uses{lem:bayesBinaryRisk_symm}
Use Lemma~\ref{lem:bayesBinaryRisk_symm}.
\end{proof}


\begin{theorem}[Data-processing inequality]
  \label{thm:data_proc_statInfo}
  \lean{ProbabilityTheory.statInfo_comp_le}
  \leanok
  \uses{def:statInfo}
  For $\mu, \nu \in \mathcal M(\mathcal X)$, $\xi \in \mathcal M(\{0,1\})$ and $\kappa : \mathcal X \rightsquigarrow \mathcal Y$ a Markov kernel, $\mathcal I_\xi(\kappa \circ \mu, \kappa \circ \nu) \le \mathcal I_\xi(\mu, \nu)$~.
\end{theorem}

\begin{proof}\leanok
\uses{thm:data_proc_bayesBinaryRisk, lem:riskIncrease_comp_del}
Since $\kappa$ is a Markov kernel,
\begin{align*}
\mathcal I_\xi(\kappa \circ \mu, \kappa \circ \nu)
&= \min\{\xi_0(\kappa \circ \mu)(\mathcal X), \xi_1(\kappa \circ \nu)(\mathcal X)\} - \mathcal B_\xi(\kappa \circ \mu, \kappa \circ \nu)
\\
&= \min\{\xi_0 \mu(\mathcal X), \xi_1 \nu(\mathcal X)\} - \mathcal B_\xi(\kappa \circ \mu, \kappa \circ \nu)
\end{align*}
It suffices to prove $\mathcal B_\xi(\kappa \circ \mu, \kappa \circ \nu) \ge \mathcal B_\xi(\mu, \nu)$ since $\mathcal I_\xi(\mu, \nu) = \min\{\xi_0\mu(\mathcal X), \xi_1\nu(\mathcal X)\} - \mathcal B_\xi(\mu, \nu)$.
The inequality was proved in Theorem~\ref{thm:data_proc_bayesBinaryRisk}.

Alternatively, we can use Lemma~\ref{lem:riskIncrease_comp_del}.
\end{proof}


\begin{corollary}
  \label{cor:statInfo_data_proc_event}
  \lean{ProbabilityTheory.statInfo_boolMeasure_le_statInfo}
  \leanok
  \uses{def:statInfo}
  Let $\mu, \nu$ be two measures on $\mathcal X$, $\xi \in \mathcal M(\{0,1\})$ and let $E$ be an event on $\mathcal X$. Let $\mu_E$ and $\nu_E$ be the two Bernoulli distributions with respective means $\mu(E)$ and $\nu(E)$.
  Then $\mathcal I_\xi(\mu, \nu) \ge \mathcal I_\xi(\mu_E, \nu_E)$.
\end{corollary}

\begin{proof}\leanok
\uses{thm:data_proc_statInfo}
Apply Theorem~\ref{thm:data_proc_statInfo} to the deterministic kernel $\mathcal X \rightsquigarrow \{0,1\}$ given by the function $x \mapsto \mathbb{I}\{x \in E\}$.
\end{proof}


\begin{lemma}
  \label{lem:statInfo_eq_sub_min}
  \lean{ProbabilityTheory.statInfo_eq_min_sub_lintegral, ProbabilityTheory.statInfo_eq_min_sub_lintegral', ProbabilityTheory.toReal_statInfo_eq_min_sub_integral}
  \leanok
  \uses{def:statInfo}
  For finite measures $\mu, \nu$ and $\xi \in \mathcal M(\{0,1\})$, for any measure $\zeta$ with $\mu \ll \zeta$ and $\nu \ll \zeta$~,
  \begin{align*}
  \mathcal I_\xi(\mu, \nu)
  &= \min\{\xi_0\mu(\mathcal X), \xi_1\nu(\mathcal X)\} - \zeta\left[x \mapsto \min \left\{\xi_0\frac{d \mu}{d\zeta}(x), \xi_1\frac{d \nu}{d\zeta}(x)\right\}\right]
  \: .
  \end{align*}
  This holds in particular for $\zeta = P \circ \xi$.
\end{lemma}

\begin{proof}\leanok
\uses{thm:bayesBinaryRisk_eq}
Apply Theorem~\ref{thm:bayesBinaryRisk_eq} to get that result for $\zeta = P \circ \xi$. Then for other $\zeta$ with $\mu \ll \zeta$ and $\nu \ll \zeta$~, use $\frac{d \mu}{d\zeta} = \frac{d \mu}{d(P \circ \xi)} / \frac{d (P \circ \xi)}{d\zeta}$~.
\end{proof}


\begin{lemma}
  \label{lem:statInfo_eq_integral_abs_sub}
  \lean{ProbabilityTheory.toReal_statInfo_eq_integral_abs}
  \leanok
  \uses{def:statInfo}
  For finite measures $\mu, \nu$ and $\xi \in \mathcal M(\{0,1\})$, for any measure $\zeta$ with $\mu \ll \zeta$ and $\nu \ll \zeta$~,
  \begin{align*}
  \mathcal I_\xi(\mu, \nu)
  &= - \frac{1}{2} \left\vert \xi_0\mu(\mathcal X) - \xi_1\nu(\mathcal X) \right\vert + 
  \frac{1}{2} \zeta\left[x \mapsto \left\vert \xi_0\frac{d \mu}{d\zeta}(x) - \xi_1\frac{d \nu}{d\zeta}(x)\right\vert\right] 
  \: .
  \end{align*}
  This holds in particular for $\zeta = P \circ \xi$.
\end{lemma}

\begin{proof}%\leanok
\uses{cor:bayesBinaryRisk_eq_abs}
Apply Corollary~\ref{cor:bayesBinaryRisk_eq_abs} and write $\min\{\xi_0\mu(\mathcal X), \xi_1\nu(\mathcal X)\} = \frac{1}{2}\left((P \circ \xi)(\mathcal X) - \left\vert \xi_0\mu(\mathcal X) - \xi_1\nu(\mathcal X) \right\vert \right)$ to get the result for $\zeta = P \circ \xi$.
The result for other $\zeta$ follows from simple manipulations of Radon-Nikodym derivatives.
\end{proof}


\begin{theorem}[Integral form of the statistical information]
  \label{thm:statInfo_eq_integral}
  \lean{ProbabilityTheory.toReal_statInfo_eq_integral_max_of_le, ProbabilityTheory.toReal_statInfo_eq_integral_max_of_ge}
  \leanok
  \uses{def:statInfo}
  For finite measures $\mu, \nu$ and $\xi \in \mathcal M(\{0,1\})$,
  \begin{align*}
  \mathcal I_\xi(\mu, \nu)
  &= \nu\left[ x \mapsto \max \left\{0 , \xi_0\frac{d \mu}{d\nu}(x) - \xi_1 \right\} \right] + \xi_0 \mu_{\perp \nu}(\mathcal X) & \text{ if } \xi_0 \mu(\mathcal X) \le \xi_1 \nu(\mathcal X)
  \: , \\
  \mathcal I_\xi(\mu, \nu)
  &= \nu\left[ x \mapsto \max \left\{0 , \xi_1 - \xi_0\frac{d \mu}{d\nu}(x) \right\} \right] & \text{ if } \xi_0 \mu(\mathcal X) \ge \xi_1 \nu(\mathcal X)
  \: .
  \end{align*}
  
\end{theorem}

For probability measures, this theorem makes the statistical information an $f$-divergence (which is defined later in this document).

\begin{proof}\leanok
\uses{thm:bayesBinaryRisk_eq}
By Theorem~\ref{thm:bayesBinaryRisk_eq},
\begin{align*}
\mathcal I_\xi(\mu, \nu)
&= \min\{\xi_0\mu(\mathcal X), \xi_1\nu(\mathcal X)\} - (P \circ \xi)\left[x \mapsto \min \left\{\xi_0\frac{d \mu}{d(P \circ \xi)}(x), \xi_1\frac{d \nu}{d(P \circ \xi)}(x)\right\}\right]
\: .
\end{align*}
Suppose that $\xi_0\mu(\mathcal X) \le \xi_1\nu(\mathcal X)$~. Then
\begin{align*}
\mathcal I_\xi(\mu, \nu)
&= \xi_0\mu(\mathcal X) - (P \circ \xi)\left[x \mapsto \xi_0\frac{d \mu}{d(P \circ \xi)}(x) + \min \left\{0 , \xi_1\frac{d \nu}{d(P \circ \xi)}(x) - \xi_0\frac{d \mu}{d(P \circ \xi)}(x)\right\}\right]
\\
&= - (P \circ \xi)\left[x \mapsto \min \left\{0 , \xi_1\frac{d \nu}{d(P \circ \xi)}(x) - \xi_0\frac{d \mu}{d(P \circ \xi)}(x) \right\}\right]
\\
&= (P \circ \xi)\left[x \mapsto \max \left\{0 , \xi_0\frac{d \mu}{d(P \circ \xi)}(x) - \xi_1\frac{d \nu}{d(P \circ \xi)}(x) \right\}\right]
\\
&= \nu\left[ x \mapsto \max \left\{0 , \xi_0\frac{d \mu}{d\nu}(x) - \xi_1 \right\} \right] + \xi_0 \mu_{\perp \nu}(\mathcal X)
\: .
\end{align*}
If on the other hand $\xi_0\mu(\mathcal X) \ge \xi_1\nu(\mathcal X)$~, with similar computations,
\begin{align*}
\mathcal I_\xi(\mu, \nu)
&= (P \circ \xi)\left[x \mapsto \max \left\{0 , \xi_1\frac{d \nu}{d(P \circ \xi)}(x) - \xi_0\frac{d \mu}{d(P \circ \xi)}(x) \right\}\right]
\\
&= \nu\left[ x \mapsto \max \left\{0 , \xi_1 - \xi_0\frac{d \mu}{d\nu}(x) \right\} \right]
\: .
\end{align*}
\end{proof}


\begin{corollary}
  \label{cor:statInfo_eq_integral_abs}
  \lean{ProbabilityTheory.toReal_statInfo_eq_integral_abs}
  \leanok
  \uses{def:statInfo}
  For finite measures $\mu, \nu$ and $\xi \in \mathcal M(\{0,1\})$,
  \begin{align*}
  \mathcal I_\xi(\mu, \nu)
  &= -\frac{1}{2} \left\vert\xi_0 \mu(\mathcal X) - \xi_1 \nu(\mathcal X)\right\vert + \frac{1}{2}\left( \nu\left[ x \mapsto \left\vert\xi_0\frac{d \mu}{d\nu}(x) - \xi_1 \right\vert \right]  + \xi_0 \mu_{\perp \nu}(\mathcal X)\right)
  \: .
  \end{align*}
\end{corollary}

\begin{proof}\leanok
\uses{thm:statInfo_eq_integral}
Start from Theorem~\ref{thm:statInfo_eq_integral}, then use $\max\{a,b\} = \frac{1}{2}\left( a + b + \vert a - b \vert \right)$~.
The two cases of Theorem~\ref{thm:statInfo_eq_integral} lead to the same expression.
\end{proof}


\begin{lemma}
  \label{lem:statInfo_eq_sub_inf_event}
  \lean{ProbabilityTheory.statInfo_eq_min_sub_iInf_measurableSet}
  \leanok
  \uses{def:statInfo}
  For finite measures $\mu, \nu$ and $\xi \in \mathcal M(\{0,1\})$,
  \begin{align*}
  \mathcal I_\xi(\mu, \nu)
  &= \min\{\xi_0 \mu(\mathcal X), \xi_1 \nu(\mathcal X)\} - \inf_{E \text{ event}} \left( \xi_0 \mu(E) + \xi_1 \nu(E^c) \right)
  \: .
  \end{align*}
\end{lemma}

\begin{proof}\leanok
\uses{lem:bayesBinaryRisk_eq_event}
This is a direct application of Lemma~\ref{lem:bayesBinaryRisk_eq_event}.
\end{proof}


\begin{lemma}
  \label{lem:statInfo_eq_sup_event}
  %\lean{}
  %\leanok
  \uses{def:statInfo}
  For finite measures $\mu, \nu$ and $\xi \in \mathcal M(\{0,1\})$,
  \begin{align*}
  \mathcal I_\xi(\mu, \nu)
  &= - \max\{0, \xi_1 \nu(\mathcal X) - \xi_0 \mu(\mathcal X) \} + \sup_{E \text{ event}} \left( \xi_1 \nu(E) - \xi_0 \mu(E) \right)
  \\
  &= - \max\{0, \xi_0 \mu(\mathcal X) - \xi_1 \nu(\mathcal X) \} + \sup_{E \text{ event}} \left( \xi_0 \mu(E) - \xi_1 \nu(E) \right)
  \: .
  \end{align*}
\end{lemma}

\begin{proof}%\leanok
\uses{lem:statInfo_eq_sub_inf_event}

\end{proof}


TODO: find a way to hide the following dummy lemma
\begin{lemma}[Dummy lemma: statInfo properties]
  \label{lem:statInfo_properties}
  %\lean{}
  \leanok
  \uses{def:statInfo}
  Dummy node to summarize properties of the statistical information.
\end{lemma}

\begin{proof}\leanok
\uses{
  lem:statInfo_self,
  lem:statInfo_nonneg,
  lem:statInfo_le,
  lem:statInfo_symm,
  thm:data_proc_statInfo,
  cor:statInfo_data_proc_event,
  lem:statInfo_eq_sub_min,
  lem:statInfo_eq_integral_abs_sub,
  thm:statInfo_eq_integral,
  cor:statInfo_eq_integral_abs,
  lem:statInfo_eq_sub_inf_event,
  lem:statInfo_eq_sup_event
}
\end{proof}


\begin{definition}
  \label{def:deGrootInfo}
  \lean{ProbabilityTheory.deGrootInfo}
  \leanok
  \uses{def:statInfo, lem:statInfo_properties}
  The DeGroot statistical information between finite measures $\mu$ and $\nu$ for $\pi \in [0,1]$ is $I_\pi (\mu, \nu) = \mathcal I_{(\pi, 1 - \pi)}(\mu, \nu)$~.
\end{definition}


\begin{definition}
  \label{def:eGamma}
  \lean{ProbabilityTheory.eGamma}
  \leanok
  \uses{def:statInfo, lem:statInfo_properties}
  The $E_\gamma$ or hockey-stick divergence between finite measures $\mu$ and $\nu$ for $\gamma \in (0,+\infty)$ is $E_\gamma(\mu, \nu) = \mathcal I_{(1,\gamma)}(\mu, \nu)$~.
\end{definition}

Note that our definition of the $E_\gamma$ divergence extends the one from the literature: the divergence is usually defined for $\gamma \ge 1$ only.
The extension makes sense in light of Theorem~\ref{thm:fDiv_eq_integral_eGamma}.
It is also extended from probability measures to finite measures.


%\subsection{DeGroot statistical information}
%
%TODO: we could also call the more general statistical information defined above the DeGroot statistical information.
%
%\begin{lemma}
%   \label{lem:deGrootInfo_self}
%   \lean{ProbabilityTheory.deGrootInfo}
%   \leanok
%   \uses{def:deGrootInfo}
%   For $\mu \in \mathcal M(\mathcal X)$, $I_\pi(\mu, \mu) = 0$~.
% \end{lemma}

% \begin{proof}%\leanok
% \uses{lem:statInfo_self}
% Use Lemma~\ref{lem:statInfo_self}.
% \end{proof}

% \begin{lemma}
%   \label{lem:deGrootInfo_nonneg}
%   %\lean{}
%   %\leanok
%   \uses{def:deGrootInfo}
%   For $\mu, \nu \in \mathcal M(\mathcal X)$, $I_\pi(\mu, \nu) \ge 0$~.
% \end{lemma}

% \begin{proof}%\leanok
% \uses{lem:statInfo_nonneg}
% Use Lemma~\ref{lem:statInfo_nonneg}.
% \end{proof}

% \begin{lemma}
%   \label{lem:deGrootInfo_symm}
%   %\lean{}
%   %\leanok
%   \uses{def:deGrootInfo}
%   For $\mu, \nu \in \mathcal M(\mathcal X)$ and $\pi \in [0,1]$, $I_\pi(\mu, \nu) = I_{1 - \pi}(\nu, \mu)$~.
% \end{lemma}

% \begin{proof}%\leanok
% \uses{lem:statInfo_symm}
% Use Lemma~\ref{lem:statInfo_symm}.
% \end{proof}

% \begin{theorem}[Data-processing inequality]
%   \label{thm:data_proc_deGrootInfo}
%   \lean{ProbabilityTheory.deGrootInfo_comp_le}
%   \leanok
%   \uses{def:deGrootInfo}
%   For $\mu, \nu \in \mathcal M(\mathcal X)$ and $\kappa : \mathcal X \rightsquigarrow \mathcal Y$ a Markov kernel, $I_\pi(\kappa \circ \mu, \kappa \circ \nu) \le I_\pi(\mu, \nu)$~.
% \end{theorem}

% \begin{proof}\leanok
% \uses{thm:data_proc_statInfo}
% Apply Theorem~\ref{thm:data_proc_statInfo}.
% \end{proof}


% \begin{lemma}[Integral form of the DeGroot statistical information]
%   \label{lem:deGrootInfo_eq_integral}
%   %\lean{}
%   %\leanok
%   \uses{def:deGrootInfo}
%   For finite measures $\mu, \nu$ and $\pi \in [0,1]$,
%   \begin{align*}
%   I_\pi(\mu, \nu)
%   &= \nu\left[ x \mapsto \pi \max \left\{0 , \frac{d \mu}{d\nu}(x) - \frac{1 - \pi}{\pi} \right\} \right] + \pi \mu_{\perp \nu}(\mathcal X) & \text{ if } \pi \mu(\mathcal X) \le (1 - \pi) \nu(\mathcal X)
%   \: , \\
%   I_\pi(\mu, \nu)
%   &= \nu\left[ x \mapsto \pi \max \left\{0 , \frac{1 - \pi}{\pi} - \frac{d \mu}{d\nu}(x) \right\} \right] & \text{ if } \pi \mu(\mathcal X) \ge (1 - \pi) \nu(\mathcal X)
%   \: .
%   \end{align*}
% \end{lemma}

% \begin{proof}%\leanok
% \uses{thm:statInfo_eq_integral}
% Apply Theorem~\ref{thm:statInfo_eq_integral}.
% \end{proof}




% \subsection{Hockey-stick divergence}


% \begin{lemma}
%   \label{lem:eGamma_self}
%   %\lean{}
%   %\leanok
%   \uses{def:eGamma}
%   For $\mu \in \mathcal M(\mathcal X)$, $E_\gamma(\mu, \mu) = 0$~.
% \end{lemma}

% \begin{proof}%\leanok
% \uses{lem:statInfo_self}
% Use Lemma~\ref{lem:statInfo_self}.
% \end{proof}

% \begin{lemma}
%   \label{lem:eGamma_nonneg}
%   %\lean{}
%   %\leanok
%   \uses{def:eGamma}
%   For $\mu, \nu \in \mathcal M(\mathcal X)$, $E_\gamma(\mu, \nu) \ge 0$~.
% \end{lemma}

% \begin{proof}%\leanok
% \uses{lem:statInfo_nonneg}
% Use Lemma~\ref{lem:statInfo_nonneg}.
% \end{proof}

% \begin{theorem}[Data-processing inequality]
%   \label{thm:data_proc_eGamma}
%   \lean{ProbabilityTheory.eGamma_comp_le}
%   \leanok
%   \uses{def:eGamma}
%   For $\mu, \nu \in \mathcal M(\mathcal X)$ and $\kappa : \mathcal X \rightsquigarrow \mathcal Y$ a Markov kernel, $E_\gamma(\kappa \circ \mu, \kappa \circ \nu) \le E_\gamma(\mu, \nu)$~.
% \end{theorem}

% \begin{proof}\leanok
% \uses{thm:data_proc_statInfo}
% Apply Theorem~\ref{thm:data_proc_statInfo}.
% \end{proof}

% \begin{lemma}[Integral form of the hockey-stick divergence]
%   \label{lem:eGamma_eq_integral}
%   %\lean{}
%   %\leanok
%   \uses{def:eGamma}
%   For finite measures $\mu, \nu$ and $\gamma \in (0,+\infty)$,
%   \begin{align*}
%   E_\gamma(\mu, \nu)
%   &= \nu\left[ x \mapsto \max \left\{0 , \frac{d \mu}{d\nu}(x) - \gamma \right\} \right] + \mu_{\perp \nu}(\mathcal X) & \text{ if } \mu(\mathcal X) \le \gamma \nu(\mathcal X)
%   \: , \\
%   E_\gamma(\mu, \nu)
%   &= \nu\left[ x \mapsto \max \left\{0 , \gamma - \frac{d \mu}{d\nu}(x) \right\} \right] & \text{ if } \mu(\mathcal X) \ge \gamma \nu(\mathcal X)
%   \: .
%   \end{align*}
% \end{lemma}

% \begin{proof}%\leanok
% \uses{thm:statInfo_eq_integral}
% Apply Theorem~\ref{thm:statInfo_eq_integral}.
% \end{proof}



\subsection{Total variation distance}

\begin{definition}
  \label{def:TV}
  \lean{ProbabilityTheory.tv}
  \leanok
  \uses{def:statInfo, lem:statInfo_properties}
  The total variation distance between finite measures $\mu$ and $\nu$ is $\TV(\mu, \nu) = \mathcal I_{(1,1)}(\mu, \nu)$~.
\end{definition}

This is also equal to $2 I_{1/2}(\mu, \nu)$ and to $E_1(\mu, \nu)$.

\begin{lemma}
  \label{lem:tv_self}
  \lean{ProbabilityTheory.tv_self}
  \leanok
  \uses{def:TV}
  For $\mu \in \mathcal M(\mathcal X)$, $\TV(\mu, \mu) = 0$~.
\end{lemma}

\begin{proof}\leanok
\uses{lem:statInfo_self}
Use Lemma~\ref{lem:statInfo_self}.
\end{proof}

\begin{lemma}
  \label{lem:tv_nonneg}
  \lean{ProbabilityTheory.tv_nonneg}
  \leanok
  \uses{def:TV}
  For $\mu, \nu \in \mathcal M(\mathcal X)$, $\TV(\mu, \nu) \ge 0$~.
\end{lemma}

\begin{proof}\leanok
\uses{lem:statInfo_nonneg}
Use Lemma~\ref{lem:statInfo_nonneg}.
\end{proof}


\begin{lemma}
  \label{lem:tv_le}
  \lean{ProbabilityTheory.tv_le}
  \leanok
  \uses{def:TV}
  For $\mu, \nu \in \mathcal M(\mathcal X)$, $\TV(\mu, \nu) \le \min\{\mu(\mathcal X), \nu(\mathcal X)\}$~.
\end{lemma}

\begin{proof}\leanok
\uses{lem:statInfo_le}
Use Lemma~\ref{lem:statInfo_le}.
\end{proof}


\begin{theorem}[Data-processing inequality]
  \label{thm:tv_data_proc}
  \lean{ProbabilityTheory.tv_comp_le}
  \leanok
  \uses{def:TV}
  For $\mu, \nu \in \mathcal M(\mathcal X)$ and $\kappa : \mathcal X \rightsquigarrow \mathcal Y$ a Markov kernel, $\TV(\kappa \circ \mu, \kappa \circ \nu) \le \TV(\mu, \nu)$~.
\end{theorem}

\begin{proof}\leanok
\uses{thm:data_proc_statInfo}
Apply Theorem~\ref{thm:data_proc_statInfo}.
\end{proof}


TODO: find a way to hide the following dummy lemma
\begin{lemma}[Dummy lemma: TV properties]
  \label{lem:tv_properties}
  %\lean{}
  \leanok
  \uses{def:TV}
  Dummy node to summarize properties of the total variation distance.
\end{lemma}

\begin{proof}\leanok
\uses{
  lem:tv_self,
  lem:tv_nonneg,
  lem:tv_le,
  thm:tv_data_proc
}
\end{proof}


\begin{lemma}
  \label{lem:tv_eq_sub_min}
  %\lean{}
  %\leanok
  \uses{def:TV}
  For finite measures $\mu, \nu$, for any measure $\zeta$ with $\mu \ll \zeta$ and $\nu \ll \zeta$~,
  \begin{align*}
  \TV(\mu, \nu)
  &= \min\{\mu(\mathcal X), \nu(\mathcal X)\} - \zeta\left[x \mapsto \min \left\{\frac{d \mu}{d\zeta}(x), \frac{d \nu}{d\zeta}(x)\right\}\right]
  \: .
  \end{align*}
  This holds in particular for $\zeta = P \circ \xi$.
\end{lemma}

\begin{proof}%\leanok
\uses{lem:statInfo_eq_sub_min}
Apply Lemma~\ref{lem:statInfo_eq_sub_min}.
\end{proof}


\begin{lemma}
  \label{lem:tv_eq_integral_abs_sub}
  %\lean{}
  %\leanok
  \uses{def:TV}
  For finite measures $\mu, \nu$, for any measure $\zeta$ with $\mu \ll \zeta$ and $\nu \ll \zeta$~,
  \begin{align*}
  \TV(\mu, \nu)
  &= - \frac{1}{2} \left\vert \mu(\mathcal X) - \nu(\mathcal X) \right\vert + 
  \frac{1}{2} \zeta\left[x \mapsto \left\vert \frac{d \mu}{d\zeta}(x) - \frac{d \nu}{d\zeta}(x)\right\vert\right] 
  \: .
  \end{align*}
  This holds in particular for $\zeta = P \circ \xi$.
\end{lemma}

\begin{proof}%\leanok
\uses{lem:statInfo_eq_integral_abs_sub}
Apply Lemma~\ref{lem:statInfo_eq_integral_abs_sub}.
\end{proof}


\begin{lemma}[Integral form of the total variation distance]
  \label{lem:tv_eq_integral}
  %\lean{}
  %\leanok
  \uses{def:TV}
  For finite measures $\mu, \nu$,
  \begin{align*}
  \TV(\mu, \nu)
  &= \nu\left[ x \mapsto \max \left\{0 , \frac{d \mu}{d\nu}(x) - 1 \right\} \right] + \mu_{\perp \nu}(\mathcal X) & \text{ if } \mu(\mathcal X) \le \nu(\mathcal X)
  \: , \\
  \TV(\mu, \nu)
  &= \nu\left[ x \mapsto \max \left\{0 , 1 - \frac{d \mu}{d\nu}(x) \right\} \right] & \text{ if } \mu(\mathcal X) \ge \nu(\mathcal X)
  \: .
  \end{align*}
\end{lemma}

\begin{proof}%\leanok
\uses{thm:statInfo_eq_integral}
Apply Theorem~\ref{thm:statInfo_eq_integral}.
\end{proof}


\begin{lemma}
  \label{lem:tv_eq_integral_abs}
  %\lean{}
  %\leanok
  \uses{def:TV}
  For finite measures $\mu, \nu$,
  \begin{align*}
  \TV(\mu, \nu)
  &= -\frac{1}{2} \left\vert \mu(\mathcal X) - \nu(\mathcal X)\right\vert + \frac{1}{2}\left( \nu\left[ x \mapsto \left\vert \frac{d \mu}{d\nu}(x) - 1 \right\vert \right]  + \mu_{\perp \nu}(\mathcal X)\right)
  \: .
  \end{align*}
\end{lemma}

\begin{proof}%\leanok
\uses{cor:statInfo_eq_integral_abs}
Apply Corollary~\ref{cor:statInfo_eq_integral_abs}.
\end{proof}


\begin{lemma}
  \label{lem:tv_eq_sub_inf_event}
  %\lean{}
  %\leanok
  \uses{def:TV}
  For finite measures $\mu, \nu$,
  \begin{align*}
  \TV(\mu, \nu)
  &= \min\{\mu(\mathcal X), \nu(\mathcal X)\} - \inf_{E \text{ event}} \left( \mu(E) + \nu(E^c) \right)
  \: .
  \end{align*}
\end{lemma}

\begin{proof}%\leanok
\uses{lem:statInfo_eq_sub_inf_event}
Apply Lemma~\ref{lem:statInfo_eq_sub_inf_event}.
\end{proof}


\begin{theorem}
  \label{thm:tv_eq_sup_sub_measure}
  %\lean{}
  %\leanok
  \uses{}
  Let $\mu, \nu$ be two finite measures on $\mathcal X$ with $\mu(\mathcal X) \leq \nu(\mathcal X)$.\\
  Then $TV(\mu, \nu) = \sup_{E \text{ event}} \left( \mu(E) - \nu(E) \right)$.
\end{theorem}

\begin{proof}%\leanok
\uses{lem:tv_eq_sub_inf_event}
Starting from Lemma~\ref{lem:tv_eq_sub_inf_event},
\begin{align*}
\TV(\mu, \nu)
&= \min\{\mu(\mathcal X), \nu(\mathcal X)\} - \inf_{E \text{ event}} \left( \mu(E) + \nu(E^c) \right)
\\
&= \mu(\mathcal X) - \inf_{E \text{ event}} \left( \mu(E^c) + \nu(E) \right)
\\
&= \mu(\mathcal X) - \mu(\mathcal X) - \inf_{E \text{ event}} \left( -\mu(E) + \nu(E) \right)
\\
&= \sup_{E \text{ event}} \left( \mu(E) - \nu(E) \right)
\: .
\end{align*}
\end{proof}


\begin{lemma}
  \label{lem:tv_eq_sup_aux}
  %\lean{}
  %\leanok
  \uses{def:TV}
  Let $\mathcal F = \{f : \mathcal X \to \mathbb{R} \mid \Vert f \Vert_\infty \le 1\}$.
  Then for $\mu, \nu$ finite measures with $\mu(\mathcal X) = \nu(\mathcal X)$, $\frac{1}{2} \sup_{f \in \mathcal F} \left( \mu[f] - \nu[f] \right) \le \TV(\mu, \nu)$.
\end{lemma}

\begin{proof}
\uses{lem:tv_eq_integral_abs_sub}
Let $p,q$ be the respective densities of $\mu, \nu$ with respect to $\xi=\mu+\nu$.
For any $f \in \mathcal F$,
\begin{align*}
\mu[f] - \nu[f]
= \int_x f(x)(p(x) - q(x)) \partial\xi
&\le \int_x \Vert f(x) \Vert_\infty \vert p(x) - q(x) \vert \partial\xi
\\
&\le \int_x \vert p(x) - q(x) \vert \partial\xi
= 2 \TV(\mu, \nu)
\: .
\end{align*}
The last equality is Lemma~\ref{lem:tv_eq_integral_abs_sub}.
\end{proof}


\begin{theorem}
  \label{thm:tv_eq_sup_sub_integral}
  %\lean{}
  %\leanok
  \uses{def:TV}
  Let $\mathcal F = \{f : \mathcal X \to \mathbb{R} \mid \Vert f \Vert_\infty \le 1\}$.
  Then for $\mu, \nu$ finite measures with $\mu(\mathcal X) = \nu(\mathcal X)$, $\TV(\mu, \nu) = \frac{1}{2} \sup_{f \in \mathcal F} \left( \mu[f] - \nu[f] \right)$.
\end{theorem}

\begin{proof}
\uses{lem:tv_eq_sup_aux, thm:tv_eq_sup_sub_measure}
Lemma~\ref{lem:tv_eq_sup_aux} gives $\TV(\mu, \nu) \ge \frac{1}{2} \sup_{f \in \mathcal F} \left( \mu[f] - \nu[f] \right)$.

The bayes estimator is a function of the form $x \mapsto \mathbb{I}\{x \in E\}$ for some event $E$, hence it belongs to $\mathcal F$. This gives equality.
\end{proof}


\begin{lemma}
  \label{lem:tv_data_proc_event}
  %\lean{}
  %\leanok
  \uses{def:TV}
  Let $\mu, \nu$ be two measures on $\mathcal X$ and let $E$ be an event. Let $\mu_E$ and $\nu_E$ be the two Bernoulli distributions with respective means $\mu(E)$ and $\nu(E)$.
  Then $\TV(\mu, \nu) \ge \TV(\mu_E, \nu_E)$.
\end{lemma}

\begin{proof}
\uses{cor:statInfo_data_proc_event}
Apply Corollary~\ref{cor:statInfo_data_proc_event}.
\end{proof}


%%FILE blueprint/src/sections/testing.tex

The goal of this section is to obtain bounds on the risk of simple binary hypothesis testing $B_\pi(\mu^{\otimes n}, \nu^{\otimes n})$ or on $\inf_{E} \max\{\mu^{\otimes n}(E), \nu^{\otimes n}(E^c)\}$ that show how those quantities depends on the number of samples $n \in \mathbb{N}$.

The main idea is to relate these to divergences that tensorize, or on which we have tensorization inequalities. Those divergences are $\KL$, $R_\alpha$, $C_\alpha$, $\JS_\alpha$ (and more generally the mutual information if we want to consider more than two hypotheses).



\section{Bounding the Bayes risk with tensorizing divergences}

\subsection{Upper bound with the Chernoff information}

\begin{lemma}[\cite{zhou2018non}]
  \label{lem:bayesBinaryRisk_eq_exp_renyi_mul_integral}
  %\lean{}
  %\leanok
  \uses{def:bayesBinaryRisk, def:Renyi, def:renyiMeasure}
  Let $\zeta$ be a measure such that $\mu \ll \zeta$ and $\nu \ll \zeta$. Let $p = \frac{d \mu}{d\zeta}$ and $q = \frac{d \nu}{d\zeta}$.
  For $\alpha \in (0,1)$, for $g_\alpha(x) = \min\{(\alpha-1)x, \alpha x\}$,
  \begin{align*}
  \mathcal B_\xi(\mu, \nu)
  = e^{-(1 - \alpha) R_\alpha(\xi_0\mu, \xi_1\nu)} \int_x \exp \left(g_\alpha\left( \log \frac{\xi_1 q(x)}{\xi_0 p(x)} \right)\right) \partial(\xi_0\mu)^{(\alpha, \xi_1\nu)}
  \: .
  \end{align*}
\end{lemma}

\begin{proof}%\leanok
\uses{thm:bayesBinaryRisk_eq}
\begin{align*}
\mathcal B_\xi(\mu, \nu)
&= \int_x \min \left\{\xi_0 p(x), \xi_1 q(x)\right\} \partial\zeta
\\
&= \int_x (\xi_0 p(x))^\alpha (\xi_1 q(x))^{1-\alpha} \min \left\{ \left(\frac{\xi_0p(x)}{\xi_1q(x)}\right)^{1 - \alpha}, \left(\frac{\xi_1q(x)}{\xi_0p(x)}\right)^\alpha\right\} \partial\zeta
\\
&= e^{-(1 - \alpha) R_\alpha(\xi_0\mu, \xi_1\nu)} \int_x \min \left\{ \left(\frac{\xi_1q(x)}{\xi_0p(x)}\right)^{\alpha - 1}, \left(\frac{\xi_1q(x)}{\xi_0p(x)}\right)^\alpha\right\} \partial(\xi_0\mu)^{(\alpha, \xi_1\nu)}
\: .
\end{align*}
\end{proof}


\begin{corollary}
  \label{cor:bayesBinaryRisk_le_exp_renyi}
  %\lean{}
  %\leanok
  \uses{def:bayesBinaryRisk, def:Renyi}
  For $\alpha \in (0,1)$,
  \begin{align*}
  \mathcal B_\xi(\mu, \nu)
  \le e^{-(1 - \alpha) R_\alpha(\xi_0\mu, \xi_1\nu)}
  = \xi_0^\alpha \xi_1^{1-\alpha} e^{-(1 - \alpha) R_\alpha(\mu, \nu)}
  \: .
  \end{align*}
\end{corollary}

\begin{proof}%\leanok
\uses{lem:bayesBinaryRisk_eq_exp_renyi_mul_integral}
Use $g_\alpha(x) \le 0$ in Lemma~\ref{lem:bayesBinaryRisk_eq_exp_renyi_mul_integral}.
\end{proof}


\begin{lemma}
  \label{lem:bayesBinaryRisk_le_exp_chernoff}
  %\lean{}
  %\leanok
  \uses{def:bayesBinaryRisk, def:Chernoff}
  \begin{align*}
  \mathcal B_\xi(\mu, \nu)
  \le e^{- C_1(\xi_0\mu, \xi_1\nu)}
  \: .
  \end{align*}
\end{lemma}

\begin{proof}%\leanok
\uses{cor:bayesBinaryRisk_le_exp_renyi, lem:chernoff_eq_max_renyi}
Optimize over $\alpha$ in Corollary~\ref{cor:bayesBinaryRisk_le_exp_renyi}: the optimal value gives the Chernoff information thanks to Lemma~\ref{lem:chernoff_eq_max_renyi}.
\end{proof}


\begin{lemma}
  \label{lem:one_sub_tv_le_exp_chernoff}
  %\lean{}
  %\leanok
  \uses{def:TV, def:Chernoff}
  For probability measures,
  \begin{align*}
  1 - \TV(\mu, \nu) \le e^{- C_1(\mu, \nu)} \: .
  \end{align*}
\end{lemma}

\begin{proof}%\leanok
\uses{lem:bayesBinaryRisk_le_exp_chernoff}
By definition, since $\mu$ and $\nu$ are probability measures, $1 - \TV(\mu, \nu) = \mathcal B_{(1,1)}(\mu, \nu)$. Then apply Lemma~\ref{lem:bayesBinaryRisk_le_exp_chernoff}.
\end{proof}



\subsection{Lower bounds using the data-processing inequality}

The ideas behind the next two theorems are the same, up to the order of the arguments of a Kullback-Leibler divergence.
The hidden parameter and data in simple binary hypothesis testing are generated from a distribution $\pi \otimes P \in \mathcal P(\Theta \times \mathcal X)$.
The reason that an estimator can get low risk is that $P$ depends on the parameter: if it did not, the Bayes risk with prior $(\alpha, 1 - \alpha)$ would be $\min\{\alpha, 1 - \alpha\}$.
We will compare $\pi \otimes P$ to a product distribution $\pi' \times \xi$ for $\pi' \in \mathcal P(\Theta)$ and $\xi \in \mathcal P(\mathcal X)$.
In order to do that, we use the data-processing inequality for $\KL$, starting either from $\KL(\pi \otimes P, \pi' \times \xi)$ or from $\KL(\pi' \times \xi, \pi \otimes P)$ (the two theorems correspond to the two choices).
The data-processing inequality is used to say that those divergences are greater than the divergences between the losses of the two corresponding estimation tasks.

We start with some preparatory lemmas.


\begin{lemma}
  \label{lem:kl_estimation_ge}
  %\lean{}
  %\leanok
  \uses{def:KL, def:bayesRisk}
  Let $\pi, \xi \in \mathcal P(\Theta)$ and $P, Q : \Theta \rightsquigarrow \mathcal X$. Suppose that the loss $\ell'$ takes values in $[0,1]$. Then
  \begin{align*}
  \KL(\pi \otimes Q, \xi \otimes P)
  &\ge \kl(\mathcal R_\pi^Q, \mathcal R_\xi^P)
  \: .
  \end{align*}
\end{lemma}

\begin{proof}%\leanok
\uses{lem:fDiv_estimation_ge}
This is Lemma~\ref{lem:fDiv_estimation_ge} specialized to the Kullback-Leibler divergence.
\end{proof}


\begin{corollary}
  \label{cor:kl_estimation_ge_sub_kl'}
  %\lean{}
  %\leanok
  \uses{def:KL, def:bayesRisk}
  Let $\pi, \xi \in \mathcal P(\Theta)$ and $P, Q : \Theta \rightsquigarrow \mathcal X$. Suppose that the loss $\ell'$ takes values in $[0,1]$. Then
  \begin{align*}
  \KL(\pi \otimes Q, \pi \otimes P)
  &\ge \kl(\mathcal R_\pi^Q, \mathcal R_\xi^P) - \KL(\pi, \xi)
  \: .
  \end{align*}
\end{corollary}

\begin{proof}%\leanok
\uses{lem:kl_estimation_ge, thm:kl_compProd}
Use $\KL(\pi \otimes Q, \xi \otimes P) = \KL(\pi, \xi) + \KL(\pi \otimes Q, \pi \otimes P)$ in Lemma~\ref{lem:kl_estimation_ge}.
\end{proof}


\begin{corollary}
  \label{cor:kl_estimation_ge_binary}
  %\lean{}
  %\leanok
  \uses{def:KL, def:bayesBinaryRisk}
  Let $\alpha, \beta \in (0, 1)$. Let $P, Q : \{0,1\} \rightsquigarrow \mathcal X$. We write $\pi_\alpha$ for the probability measure on $\{0,1\}$ with $\pi_\alpha(\{0\}) = \alpha$. Then
  \begin{align*}
  \KL(\pi_\beta \otimes Q, \pi_\alpha \otimes P)
  &\ge \kl(B_\beta(Q(0), Q(1)), B_\alpha(P(0), P(1)))
  \: .
  \end{align*}
\end{corollary}

\begin{proof}%\leanok
\uses{lem:kl_estimation_ge}
Apply Lemma~\ref{lem:kl_estimation_ge}.
\end{proof}


\begin{corollary}
  \label{cor:kl_estimation_ge_sub_kl}
  %\lean{}
  %\leanok
  \uses{def:KL, def:bayesBinaryRisk}
  Let $\mu, \nu \in \mathcal P(\mathcal X)$ and let $\alpha, \beta \in (0, 1)$. Let $P, Q : \{0,1\} \rightsquigarrow \mathcal X$. We write $\pi_\alpha$ for the probability measure on $\{0,1\}$ with $\pi_\alpha(\{0\}) = \alpha$. Then
  \begin{align*}
  \KL(\pi_\beta \otimes Q, \pi_\beta \otimes P)
  &\ge \kl(B_\beta(Q(0), Q(1)), B_\alpha(P(0), P(1))) - \kl(\beta, \alpha)
  \: .
  \end{align*}
\end{corollary}

\begin{proof}%\leanok
Apply Corollary~\ref{cor:kl_estimation_ge_sub_kl'}.
\end{proof}


\begin{lemma}
  \label{lem:kl_estimation_ge_prod}
  %\lean{}
  %\leanok
  \uses{def:KL, def:bayesBinaryRisk}
  Let $\mu, \nu, \xi \in \mathcal P(\mathcal X)$ and let $\alpha, \beta \in (0, 1)$. Let $P : \{0,1\} \rightsquigarrow \mathcal X$ be the kernel with $P(0) = \mu$ and $P(1) = \nu$. We write $\pi_\alpha$ for the probability measure on $\{0,1\}$ with $\pi_\alpha(\{0\}) = \alpha$. Let $\bar{\beta} = \min\{\beta, 1 - \beta\}$.
  Then
  \begin{align*}
  \KL(\pi_\alpha \otimes P, \pi_\alpha \times \xi)
  &\ge \kl(B_\alpha(\mu, \nu), \bar{\beta}) - \kl(\alpha, \bar{\beta})
  \: , \\
  \KL(\pi_\beta \times \xi, \pi_\beta \otimes P)
  &\ge \kl(\bar{\beta}, B_\alpha(\mu, \nu)) - \kl(\bar{\beta}, \alpha)
  \: .
  \end{align*}
\end{lemma}

\begin{proof}%\leanok
\uses{cor:kl_estimation_ge_sub_kl, lem:bayesBinaryRisk_self}
We apply Corollary~\ref{cor:kl_estimation_ge_sub_kl} for the constant kernel with value $\xi$ and use Lemma~\ref{lem:bayesBinaryRisk_self}.
\end{proof}


\begin{theorem}[Fano's inequality, binary case]
  \label{thm:sub_bayesBinaryRisk_le_jensenShannon}
  %\lean{}
  %\leanok
  \uses{def:bayesBinaryRisk, def:jensenShannon}
  Let $\mu, \nu \in \mathcal P(\mathcal X)$ and let $\alpha \in (0, 1)$.
  \begin{align*}
  h_2(\alpha) - h_2(B_\alpha(\mu, \nu)) \le \JS_\alpha(\mu, \nu) \: ,
  \end{align*}
  in which $h_2: x \mapsto x\log\frac{1}{x} + (1 - x)\log\frac{1}{1 - x}$ is the binary entropy function.
\end{theorem}

\begin{proof}%\leanok
\uses{lem:kl_estimation_ge_prod, thm:kl_compProd, lem:jensenShannon_eq_inf_kl}
Apply Lemma~\ref{lem:jensenShannon_eq_inf_kl} and then Lemma~\ref{lem:kl_estimation_ge_prod} with $\beta = 1/2$ to get
\begin{align*}
\JS_\alpha(\mu, \nu)
&= \inf_{\xi \in \mathcal P(\mathcal X)}\KL(\pi_\alpha \otimes P, \pi_\alpha \times \xi)
\\
&\ge \kl(B_\alpha(\mu, \nu), 1/2) - \kl(\alpha, 1/2)
\\
&= h_2(\alpha) - h_2(B_\alpha(\mu, \nu)) \: .
\end{align*}
Note that $\beta = 1/2$ is the value that results in the best bound: the lower bound on $\JS_\alpha(\mu, \nu)$ from some $\beta \le 1/2$ would be
\begin{align*}
\kl(B_\alpha(\mu, \nu), \beta) - \kl(\alpha, \beta)
&= h_2(\alpha) - h_2(B_\alpha(\mu, \nu)) - (\alpha - B_\alpha(\mu, \nu)) \log \frac{1 - \beta}{\beta}
\\
&\le h_2(\alpha) - h_2(B_\alpha(\mu, \nu)) \: .
\end{align*}

\end{proof}


\begin{theorem}
  \label{thm:log_inv_bayesBinaryRisk_le_renyi}
  %\lean{}
  %\leanok
  \uses{def:bayesBinaryRisk, def:Renyi}
  For $\alpha, \beta \in (0, 1/2)$,
  \begin{align*}
  \beta \log\frac{\alpha}{B_\alpha(\mu, \nu)} + (1 - \beta) \log\frac{1 - \alpha}{1 - B_\alpha(\mu, \nu)}
  &\le (1 - \beta) R_{\beta}(\mu, \nu)
  \: .
  \end{align*}
  As a consequence,
  \begin{align*}
  \beta \log\frac{\alpha}{B_\alpha(\mu, \nu)}
  &\le (1 - \beta) (R_{\beta}(\mu, \nu) + \log 2)
  \: .
  \end{align*}
  In particular, $\log\frac{\alpha}{B_\alpha(\mu, \nu)} \le R_{1/2}(\mu, \nu) + \log 2$~.
\end{theorem}

\begin{proof}%\leanok
\uses{lem:kl_estimation_ge_prod, lem:renyi_eq_inf_kl}
By Lemma~\ref{lem:renyi_eq_inf_kl} and then Lemma~\ref{lem:kl_estimation_ge_prod},
\begin{align*}
(1 - \beta) R_{\beta}(\mu, \nu)
&= \inf_{\xi \in \mathcal P(\mathcal X)} \KL(\pi_\beta \times \xi, \pi_\beta \otimes P)
\\
&\ge \kl(\beta, B_\alpha(\mu, \nu)) - \kl(\beta, \alpha)
\\
&= \beta \log\frac{\alpha}{B_\alpha(\mu, \nu)} + (1 - \beta)\log\frac{1 - \alpha}{1 - B_\alpha(\mu, \nu)}
\: .
\end{align*}
The first particular case is obtained by using $\alpha \le 1/2$ and $B_\pi(\mu, \nu) \ge 0$.
The second one further sets $\beta = 1/2$.
\end{proof}


%\begin{theorem}[TODO]
%  \label{thm:renyi_bayesBinaryRisk_le_renyi_mutual}
%  %\lean{}
%  %\leanok
%  \uses{def:bayesBinaryRisk, def:Renyi}
%  Let $\mu, \nu \in \mathcal P(\mathcal X)$ and let $\alpha \in (0, 1)$, $\lambda \in [0, +\infty)$. Let $r_\lambda(a,b)$ be the Rényi divergence between Bernoulli distributions with means $a$ and $b$.
%  \begin{align*}
%  r_\lambda(B_\alpha(\mu, \nu), \alpha) \le \inf_{\xi \in \mathcal P(\mathcal X)} R_\lambda(\pi_\alpha \otimes P, \pi_\alpha \times \xi) \: .
%  \end{align*}
%\end{theorem}



\paragraph{TODO: move or remove}


\begin{lemma}
  \label{lem:testing_bound_renyi_mean}
  %\lean{}
  %\leanok
  \uses{def:Renyi}
  Let $\mu, \nu$ be two probability measures on $\mathcal X$ and $E$ an event. Let $\alpha \in (0,1)$. Then
  \begin{align*}
  \mu(E)^\alpha + \nu(E^c)^{1 - \alpha}
  \ge \exp\left(-(1 - \alpha) R_{\alpha}(\mu, \nu)\right)
  \: .
  \end{align*}
\end{lemma}

\begin{proof}
\uses{lem:renyi_data_proc_event}
Let $\mu_E$ and $\nu_E$ be the two Bernoulli distributions with respective means $\mu(E)$ and $\nu(E)$.
By Lemma~\ref{lem:renyi_data_proc_event}, $R_\alpha(\mu, \nu) \ge R_\alpha(\mu_E, \nu_E)$. That divergence is
\begin{align*}
R_\alpha(\mu_E, \nu_E)
&= \frac{1}{\alpha - 1}\log \left(\mu(E)^\alpha \nu(E)^{1 - \alpha}
  + \mu(E^c)^\alpha \nu(E^c)^{1 - \alpha}\right)
\\
&\ge \frac{1}{\alpha - 1}\log \left(\mu(E)^\alpha + \nu(E^c)^{1 - \alpha}\right)
\: .
\end{align*}
\end{proof}

\begin{corollary}
  \label{cor:testing_bound_hellinger}
  %\lean{}
  %\leanok
  \uses{def:Hellinger, def:Renyi}
  Let $\mu, \nu$ be two probability measures on $\mathcal X$ and $E$ an event. Then
  \begin{align*}
  \sqrt{\mu(E)} + \sqrt{\nu(E^c)}
  \ge \exp\left(-\frac{1}{2} R_{1/2}(\mu, \nu)\right)
  = 1 - \Hsq(\mu, \nu)
  \: .
  \end{align*}
\end{corollary}

\begin{proof}
\uses{lem:testing_bound_renyi_mean, lem:renyi_half_eq_log_hellinger}
The inequality is an application of Lemma~\ref{lem:testing_bound_renyi_mean} for $\alpha = 1/2$. The equality is Lemma~\ref{lem:renyi_half_eq_log_hellinger}.
\end{proof}



\subsection{Lower bounds using the change of measure lemma}

The main tool of this section is the next lemma, which relates the probability of an event under two measures and deviations of the log-likelihood ratio between the measures.

\begin{lemma}[Change of measure lemma]
  \label{lem:llr_change_measure}
  \lean{ProbabilityTheory.measure_sub_le_measure_mul_exp'}
  \leanok
  %\uses{}
  Let $\mu, \nu$ be two measures on $\mathcal X$ with $\mu \ll \nu$ and let $E$ be an event on $\mathcal X$. Let $\beta \in \mathbb{R}$. Then
  \begin{align*}
  \nu(E) e^{\beta} \ge \mu(E) - \mu\left\{ \log\frac{d \mu}{d \nu} > \beta \right\} \: .
  \end{align*}
\end{lemma}

\begin{proof}\leanok
\begin{align*}
\nu(E)
\ge \mu\left[\mathbb{I}(E) e^{- \log\frac{d \mu}{d \nu} }\right]
&\ge \mu\left[\mathbb{I}\left(E \cap \left\{\log\frac{d \mu}{d \nu} \le \beta\right\}\right) e^{- \log\frac{d \mu}{d \nu} }\right]
\\
&\ge e^{- \beta}\mu\left(E \cap \left\{\log\frac{d \mu}{d \nu} \le \beta\right\}\right)
\\
&\ge e^{- \beta}\left( \mu(E) - \mu\left\{ \log\frac{d \mu}{d \nu} > \beta \right\} \right)
\: .
\end{align*}
\end{proof}


\begin{lemma}[Change of measure - functions]
  \label{lem:llr_change_measure_fun}
  %\lean{}
  %\leanok
  \uses{}
  Let $\mu, \nu$ be two measures on $\mathcal X$ with $\mu \ll \nu$ and let $f : \mathcal X \to [0,1]$ be a measurable function. Let $\beta \in \mathbb{R}$. Then
  \begin{align*}
  \nu[f] e^{\beta} \ge \mu[f] - \mu\left\{ \log\frac{d \mu}{d \nu} > \beta \right\} \: .
  \end{align*}
\end{lemma}

\begin{proof}%\leanok
\uses{}
\begin{align*}
\nu[f]
\ge \mu\left[f e^{- \log\frac{d \mu}{d \nu} }\right]
&\ge \mu\left[f \mathbb{I}\left\{\log\frac{d \mu}{d \nu} \le \beta\right\} e^{- \log\frac{d \mu}{d \nu} }\right]
\\
&\ge e^{- \beta}\mu\left[f \mathbb{I}\left\{\log\frac{d \mu}{d \nu} \le \beta\right\} \right]
\\
&\ge e^{- \beta}\left( \mu[f] - \mu\left[f \mathbb{I}\left\{ \log\frac{d \mu}{d \nu} > \beta \right\}\right] \right)
\\
&\ge e^{- \beta}\left( \mu[f] - \mu \left\{ \log\frac{d \mu}{d \nu} > \beta \right\}\right)
\: .
\end{align*}
The first and second inequalities use $f \ge 0$ and the last inequality uses $f \le 1$~.
\end{proof}


\begin{lemma}
  \label{lem:change_measure_risk}
  %\lean{}
  %\leanok
  \uses{def:bayesRisk}
  Consider an estimation problem with loss $\ell' : \mathcal Y \times \mathcal Z \to [0,1]$. Let $\pi, \zeta \in \mathcal P(\Theta)$ and $P, Q : \Theta \rightsquigarrow \mathcal X$ be such that $\zeta \otimes Q \ll \pi \otimes P$. Then for all $\beta \in \mathbb{R}$,
  \begin{align*}
  \mathcal R_\pi^P e^{\beta}
  \ge \mathcal R_\zeta^Q - (\zeta \otimes Q)\left\{ \log\frac{d (\zeta \otimes Q)}{d (\pi \otimes P)} > \beta \right\}
  \: .
  \end{align*}
\end{lemma}

\begin{proof}%\leanok
\uses{lem:llr_change_measure_fun}
Let $\hat{y}_B$ be a Bayes estimator for $(P, y, \ell')$. If no such estimator exists, the proof can be adapted by taking an estimator with risk $\varepsilon$-close to the Bayes risk and then minimizing over $\varepsilon > 0$.
Apply Lemma~\ref{lem:llr_change_measure_fun} to the function $(\theta, x) \mapsto \hat{y}_B(x)\left[ z \mapsto \ell'(y(\theta), z) \right]$~, which takes values in $[0,1]$ and to the measures $\pi \otimes P$ and $\zeta \otimes Q$.
\begin{align*}
\mathcal R_\pi^P e^{\beta}
&\ge (\zeta \otimes Q)\left[ (\theta, x) \mapsto \hat{y}_B(x)\left[ z \mapsto \ell'(y(\theta), z) \right] \right]
  - (\zeta \otimes Q)\left\{ \log\frac{d (\zeta \otimes Q)}{d (\pi \otimes P)} > \beta \right\}
\\
&\ge \mathcal R_\zeta^Q - (\zeta \otimes Q)\left\{ \log\frac{d (\zeta \otimes Q)}{d (\pi \otimes P)} > \beta \right\}
\: .
\end{align*}
\end{proof}


\begin{lemma}
  \label{lem:change_measure_risk_inf}
  %\lean{}
  %\leanok
  \uses{def:bayesRisk}
  Consider an estimation problem with loss $\ell' : \mathcal Y \times \mathcal Z \to [0,1]$. Let $\pi, \zeta \in \mathcal P(\Theta)$ and $P : \Theta \rightsquigarrow \mathcal X$. Then for all $\beta \in \mathbb{R}$,
  \begin{align*}
  \mathcal R_\pi^P e^{\beta}
  \ge \mathcal R_\zeta^{d_{\mathcal X}} - \inf_{\xi \in \mathcal P(\mathcal X)}(\zeta \times \xi)\left\{ \log\frac{d (\zeta \times \xi)}{d (\pi \otimes P)} > \beta \right\}
  \: ,
  \end{align*}
  in which the infimum over $\xi$ is restricted to probability measures such that $\zeta \times \xi \ll \pi \otimes P$ and $d_{\mathcal X} : \Theta \rightsquigarrow *$ is the discard kernel.
\end{lemma}

\begin{proof}%\leanok
\uses{lem:change_measure_risk, lem:bayesRisk_const}
For any $\xi$, we apply Lemma~\ref{lem:change_measure_risk} to $\zeta \times \xi$ and $\pi \otimes P$ and remark that by Lemma~\ref{lem:bayesRisk_const} the risk for a constant kernel with value $\xi$ is independent of $\xi$ and equal to $\mathcal R_\zeta^{d_{\mathcal X}}$.
\end{proof}


\begin{lemma}[3 points change of measure]
  \label{lem:llr_change_measure_add}
  \lean{ProbabilityTheory.one_sub_le_add_measure_mul_exp}
  \leanok
  %\uses{}
  Let $\mu, \nu, \xi \in \mathcal P(\mathcal X)$ and let $E$ be an event on $\mathcal X$. Let $\beta_1, \beta_2 \in \mathbb{R}$. Then
  \begin{align*}
  \mu(E) e^{\beta_1} + \nu(E^c) e^{\beta_2} \ge 1 - \xi\left\{ \log\frac{d \xi}{d \mu} > \beta_1 \right\} - \xi\left\{ \log\frac{d \xi}{d \nu} > \beta_2 \right\} \: .
  \end{align*}
\end{lemma}

\begin{proof}\leanok
\uses{lem:llr_change_measure}
Two applications of Lemma~\ref{lem:llr_change_measure}, then sum them and use $\xi(E)+\xi(E^c) = 1$.
\end{proof}



\paragraph{Change of measure and the moments of the log-likelihood ratio}

\begin{corollary}[Change of measure - mean]
  \label{cor:kl_change_measure}
  %\lean{}
  %\leanok
  \uses{def:KL}
  Let $\mu, \nu$ be two measures on $\mathcal X$ and let $E$ be an event on $\mathcal X$. Let $\beta \in \mathbb{R}$. Then
  \begin{align*}
  \nu(E) e^{\KL(\mu, \nu) + \beta} \ge \mu(E) - \mu\left\{ \log\frac{d \mu}{d \nu} - \KL(\mu, \nu) > \beta \right\} \: .
  \end{align*}
\end{corollary}

\begin{proof}
\uses{lem:llr_change_measure}
Use Lemma~\ref{lem:llr_change_measure} with the choice $\KL(\mu, \nu) + \beta$ for $\beta$.
\end{proof}


\begin{lemma}
  \label{lem:change_measure_risk_mean}
  %\lean{}
  %\leanok
  \uses{def:bayesBinaryRisk, def:Renyi, def:KL, def:renyiMeasure}
  For $\alpha \in (0,1)$, let $\pi_\alpha \in \mathcal P(\{0,1\})$ be the measure $(\alpha, 1 - \alpha)$.
  Let $\alpha, \gamma \in (0,1)$, $\mu, \nu \in \mathcal P(\mathcal X)$ and let $P : \{0,1\} \rightsquigarrow \mathcal X$ be the kernel with $P(0) = \mu$ and $P(1) = \nu$~. Then for all $\beta \in \mathbb{R}$~,
  \begin{align*}
  &\mathcal B_{\pi_\alpha}(\mu, \nu) e^{\kl(\gamma, \alpha) + (1 - \gamma) R_\gamma(\mu, \nu) + \beta}
  \\
  &\ge \min\{\gamma, 1 - \gamma\} - (\pi_\gamma \times \mu^{(\gamma, \nu)})\left\{ \log\frac{d (\pi_\gamma \times \mu^{(\gamma, \nu)})}{d (\pi_\alpha \otimes P)} > \KL(\pi_\gamma \times \mu^{(\gamma, \nu)}, \pi_\alpha \otimes P) + \beta \right\}
  \: .
  \end{align*}
\end{lemma}

\begin{proof}%\leanok
\uses{lem:change_measure_risk_inf}
An application of Lemma~\ref{lem:change_measure_risk_inf} gives
\begin{align*}
\mathcal B_{\pi_\alpha}(\mu, \nu) e^{\beta}
\ge \min\{\gamma, 1 - \gamma\} - \inf_{\xi \in \mathcal P(\mathcal X)}(\pi_\gamma \times \xi)\left\{ \log\frac{d (\pi_\gamma \times \xi)}{d (\pi_\alpha \otimes P)} > \beta \right\}
\: .
\end{align*}

Take $\xi = \mu^{(\gamma, \nu)}$ and
\begin{align*}
\beta
&= \KL(\pi_\gamma \times \xi, \pi_\alpha \otimes P) + \beta'
\\
&= \kl(\gamma, \alpha) + \KL(\pi_\gamma \times \xi, \pi_\gamma \otimes P) + \beta'
\\
&= \kl(\gamma, \alpha) + (1 - \gamma) R_\gamma(\mu, \nu) + \beta'
\: .
\end{align*}
We obtain
\begin{align*}
&\mathcal B_{\pi_\alpha}(\mu, \nu) e^{\kl(\gamma, \alpha) + (1 - \gamma) R_\gamma(\mu, \nu) + \beta'}
\\
&\ge \min\{\gamma, 1 - \gamma\} - (\pi_\gamma \times \xi)\left\{ \log\frac{d (\pi_\gamma \times \xi)}{d (\pi_\alpha \otimes P)} > \KL(\pi_\gamma \times \xi, \pi_\alpha \otimes P) + \beta' \right\}
\: .
\end{align*}
\end{proof}


\begin{lemma}[Change of measure - variance]
  \label{lem:llr_change_measure_variance}
  %\lean{}
  %\leanok
  \uses{}
  Let $\mu, \nu$ be two measures on $\mathcal X$ such that $\mu\left[\left(\log\frac{d \mu}{d \nu}\right)^2\right] < \infty$. Let $E$ be an event on $\mathcal X$ and let $\beta > 0$. Then
  \begin{align*}
  \nu(E) e^{\KL(\mu, \nu) + \sqrt{\Var_\mu[\log\frac{d \mu}{d \nu}]\beta}} \ge \mu(E) - \frac{1}{\beta} \: .
  \end{align*}
\end{lemma}

\begin{proof}%\leanok
\uses{lem:llr_change_measure}
Use Lemma~\ref{lem:llr_change_measure} with the choice $\KL(\mu, \nu) + \sqrt{\Var_\mu[\log\frac{d \mu}{d \nu}]\beta}$ for $\beta$ and bound the probability of deviation of the log-likelihood ratio with Chebychev's inequality.
\end{proof}

\begin{lemma}[Change of measure - c.g.f.]
  \label{lem:renyi_chernoff_bound}
  \lean{ProbabilityTheory.measure_llr_gt_renyiDiv_le_exp}
  \leanok
  \uses{def:Renyi}
  For $\mu, \nu$ finite measures and $\alpha, \beta > 0$,
  \begin{align*}
  \mu\left\{ \log\frac{d \mu}{d \nu} > R_{1+\alpha}(\mu, \nu) + \beta \right\}
  \le e^{- \alpha \beta}
  \: .
  \end{align*}
\end{lemma}

\begin{proof}\leanok
\uses{lem:renyi_cgf_2}
This is a Chernoff bound, using that the cumulant generating function of $\log\frac{d\mu}{d\nu}$ under $\mu$ has value $\alpha R_{1+\alpha}(\mu, \nu)$ at $\alpha$ by Lemma~\ref{lem:renyi_cgf_2}.
\begin{align*}
\mu\left\{ \log\frac{d \mu}{d \nu} > R_{1+\alpha}(\mu, \nu) + \beta \right\}
&= \mu\left\{ \exp\left(\alpha\log\frac{d \mu}{d \nu}\right) > \exp\left(\alpha R_{1+\alpha}(\mu, \nu) + \alpha \beta\right) \right\}
\\
&\le e^{-\alpha R_{1+\alpha}(\mu, \nu) - \alpha \beta} \mu\left[\left(\frac{d \mu}{d \nu}\right)^\alpha \right]
\end{align*}
Then $\mu\left[\left(\frac{d \mu}{d \nu}\right)^\alpha \right] = \nu\left[\left(\frac{d \mu}{d \nu}\right)^{1+\alpha} \right] = e^{\alpha R_{1+\alpha}(\mu, \nu)}$.
\end{proof}


\begin{lemma}
  \label{lem:renyi_change_measure}
  \lean{ProbabilityTheory.measure_sub_le_measure_mul_exp_renyiDiv}
  \leanok
  \uses{def:Renyi}
  Let $\mu, \nu$ be two finite measures on $\mathcal X$ and let $E$ be an event on $\mathcal X$. Let $\alpha,\beta > 0$. Then
  \begin{align*}
  \nu(E) e^{R_{1+\alpha}(\mu, \nu) + \beta} \ge \mu(E) - e^{-\alpha \beta} \: .
  \end{align*}
\end{lemma}

\begin{proof}\leanok
\uses{lem:llr_change_measure, lem:renyi_chernoff_bound}
Use Lemma~\ref{lem:llr_change_measure} with the choice $R_{1+\alpha}(\mu, \nu) + \beta$ for $\beta$. Then apply Lemma~\ref{lem:renyi_chernoff_bound}.
\end{proof}


\begin{lemma}
  \label{lem:change_measure_risk_cgf}
  %\lean{}
  %\leanok
  \uses{def:bayesBinaryRisk, def:Renyi}
  For $\alpha \in (0,1)$, let $\pi_\alpha \in \mathcal P(\{0,1\})$ be the measure $(\alpha, 1 - \alpha)$.
  Let $\alpha, \gamma \in (0,1)$, $\mu, \nu \in \mathcal P(\mathcal X)$ and let $P : \{0,1\} \rightsquigarrow \mathcal X$ be the kernel with $P(0) = \mu$ and $P(1) = \nu$~. Then for all $\beta > 0$ and $\varepsilon>0$~,
  \begin{align*}
  \mathcal B_{\pi_\alpha}(\mu, \nu) e^{\inf_{\xi \in \mathcal P(\mathcal X)} R_{1 + \varepsilon}(\pi_\gamma \times \xi, \pi_\alpha \otimes P) + \beta}
  &\ge \min\{\gamma, 1 - \gamma\} - e^{-\varepsilon \beta}
  \: .
  \end{align*}
\end{lemma}

\begin{proof}
\uses{lem:change_measure_risk_inf, lem:renyi_chernoff_bound}
\end{proof}



\paragraph{Applications of the change of measure with 3 points}

\begin{lemma}
  \label{lem:change_measure_variance_add}
  % \lean{}
  % \leanok
  \uses{}
  Let $\mu, \nu, \xi$ be three probability measures on $\mathcal X$ and let $E$ be an event on $\mathcal X$. For $\beta > 0$~,
  \begin{align*}
  \mu(E) e^{\KL(\xi, \mu) + \sqrt{\beta \Var_{\xi}\left[\log\frac{d\xi}{d\mu}\right]}} + \nu(E^c) e^{\KL(\xi, \nu) + \sqrt{\beta \Var_{\xi}\left[\log\frac{d\xi}{d\nu}\right]}}
  \ge 1 - \frac{2}{\beta} \: .
  \end{align*}
\end{lemma}

\begin{proof} %\leanok
\uses{lem:llr_change_measure_add}
Use Lemma~\ref{lem:llr_change_measure_add} with the choices $\KL(\xi, \mu) + \sqrt{\beta \Var_{\xi}\left[\log\frac{d\xi}{d\mu}\right]}$ and $\KL(\xi, \nu) + \sqrt{\beta \Var_{\xi}\left[\log\frac{d\xi}{d\nu}\right]}$ for $\beta_1$ and $\beta_2$.
Then use Chebyshev's inequality to bound the probabilities of deviation of the log-likelihood ratios.
\end{proof}

\begin{lemma}
  \label{lem:renyi_change_measure_add}
  \lean{ProbabilityTheory.one_sub_exp_le_add_measure_mul_exp_max_renyiDiv}
  \leanok
  \uses{def:Renyi}
  Let $\mu, \nu, \xi$ be three probability measures on $\mathcal X$ and let $E$ be an event on $\mathcal X$. Let $\alpha, \beta \ge 0$. Then
  \begin{align*}
  \mu(E) e^{R_{1+\alpha}(\xi, \mu) + \beta} + \nu(E^c) e^{R_{1+\alpha}(\xi, \nu) + \beta} \ge 1 - 2 e^{-\alpha \beta} \: .
  \end{align*}
\end{lemma}

\begin{proof}\leanok
\uses{lem:llr_change_measure_add, lem:renyi_chernoff_bound}
Use Lemma~\ref{lem:llr_change_measure_add} with the choices $R_{1+\alpha}(\xi, \mu) + \beta$ and $R_{1+\alpha}(\xi, \nu) + \beta$ for $\beta_1$ and $\beta_2$.
Then apply Lemma~\ref{lem:renyi_chernoff_bound}.
\end{proof}


\begin{lemma}
  \label{lem:testing_bound_renyi_one_add}
  %\lean{}
  %\leanok
  \uses{def:Renyi, def:Chernoff}
  Let $\mu, \nu$ be two probability measures on $\mathcal X$ and let $E$ be an event on $\mathcal X$. Let $\alpha > 0$. Then
  \begin{align*}
  \mu(E) + \nu(E^c) \ge \frac{1}{2}\exp\left( - C_{1+\alpha}(\mu, \nu) - \frac{\log 4}{\alpha}\right) \: .
  \end{align*}
\end{lemma}

\begin{proof}%\leanok
\uses{lem:renyi_change_measure_add}
Use Lemma~\ref{lem:renyi_change_measure_add} with $\beta = \log(4)/\alpha$ and use that
\begin{align*}
\mu(E) e^{R_{1+\alpha}(\xi, \mu) + \beta} + \nu(E^c) e^{R_{1+\alpha}(\xi, \nu) + \beta}
\le (\mu(E) + \nu(E^c)) e^{\max \{R_{1+\alpha}(\xi, \mu), R_{1+\alpha}(\xi, \nu)\} + \beta}
\: .
\end{align*}

\end{proof}



\paragraph{Product spaces}

\begin{lemma}
  \label{lem:testing_bound_renyi_one_add_n}
  %\lean{}
  %\leanok
  \uses{def:Renyi, def:Chernoff}
  Let $\mu, \nu$ be two probability measures on $\mathcal X$, let $n \in \mathbb{N}$ and let $E$ be an event on $\mathcal X^n$. For all $\alpha > 0$,
  \begin{align*}
  \mu^{\otimes n}(E) + \nu^{\otimes n}(E^c) \ge \frac{1}{2}\exp\left( - n C_{1+\frac{\alpha}{\sqrt{n}}}(\mu, \nu) - \frac{\log 4}{\alpha}\sqrt{n}\right) \: .
  \end{align*}
\end{lemma}

\begin{proof}
\uses{lem:renyi_prod_n, lem:testing_bound_renyi_one_add}
Use Lemma~\ref{lem:renyi_prod_n} in Lemma~\ref{lem:testing_bound_renyi_one_add}. TODO: add a lemma for tensorization of the Chernoff divergence.
\end{proof}

\begin{theorem}
  \label{thm:testing_bound_chernoff}
  %\lean{}
  %\leanok
  \uses{def:Chernoff}
  Let $\mu, \nu$ be two probability measures on $\mathcal X$ and let $(E_n)_{n \in \mathbb{N}}$ be events on $\mathcal X^n$. For all $\gamma \in (0,1)$,
  \begin{align*}
  \limsup_{n \to +\infty} \frac{1}{n}\log \frac{1}{\gamma \mu^{\otimes n}(E_n) + (1 - \gamma)\nu^{\otimes n}(E_n^c)}
  \le C_1(\mu, \nu)
  \: .
  \end{align*}
\end{theorem}

\begin{proof}
\uses{cor:kl_change_measure, thm:kl_pi}
Let $\xi$ be a probability measure on $\mathcal X$ and $\beta > 0$. By Corollary~\ref{cor:kl_change_measure},
\begin{align*}
\mu^{\otimes n}(E_n) e^{\KL(\xi^{\otimes n}, \mu^{\otimes n}) + n\beta}
&\ge \xi^{\otimes n}(E_n) - \xi^{\otimes n}\left\{ \log\frac{d \xi^{\otimes n}}{d \mu^{\otimes n}} - \KL(\xi^{\otimes n}, \mu^{\otimes n}) > n\beta \right\}
\: , \\
\nu^{\otimes n}(E_n^c) e^{\KL(\xi^{\otimes n}, \nu^{\otimes n}) + n\beta}
&\ge \xi^{\otimes n}(E_n^c) - \xi^{\otimes n}\left\{ \log\frac{d \xi^{\otimes n}}{d \nu^{\otimes n}} - \KL(\xi^{\otimes n}, \nu^{\otimes n}) > n\beta \right\}
\: .
\end{align*}
We sum both inequalities with weights $\gamma$ and $1-\gamma$ respectively and use that each $\KL$ on the left is less than their max, as well as $\xi^{\otimes n}(E_n) + \xi^{\otimes n}(E_n^c) = 1$.
\begin{align*}
&e^{n\beta} (\gamma\mu^{\otimes n}(E_n) + (1-\gamma)\nu^{\otimes n}(E_n^c)) e^{\max\{\KL(\xi^{\otimes n}, \mu^{\otimes n}), \KL(\xi^{\otimes n}, \nu^{\otimes n})\}}
\\
&\ge \min\{\gamma, 1-\gamma\} - \gamma\xi^{\otimes n}\left\{ \log\frac{d \xi^{\otimes n}}{d \mu^{\otimes n}} - \KL(\xi^{\otimes n}, \mu^{\otimes n}) > n\beta \right\}
  - (1 - \gamma)\xi^{\otimes n}\left\{ \log\frac{d \xi^{\otimes n}}{d \nu^{\otimes n}} - \KL(\xi^{\otimes n}, \nu^{\otimes n}) > n\beta \right\}
\: .
\end{align*}
Let $p_{n,\mu}(\beta)$ and $p_{n, \nu}(\beta)$ be the two probabilities on the right hand side. By the law of large numbers, both tend to 0 when $n$ tends to $+\infty$.
In particular, for $n$ large enough, the right hand side is positive and we can take logarithms on both sides. We also use the tensorization of $\KL$ (Theorem~\ref{thm:kl_pi}).
\begin{align*}
& n \max\{\KL(\xi, \mu), \KL(\xi, \nu)\}
\\
&\ge \log \frac{1}{\gamma \mu^{\otimes n}(E_n) + (1 - \gamma)\nu^{\otimes n}(E_n^c)} + \log (\min\{\gamma, 1-\gamma\} - \gamma p_{n, \mu}(\beta) - (1 - \gamma) p_{n, \nu}(\beta)) - n\beta
\end{align*}
For $n \to +\infty$,
\begin{align*}
\max\{\KL(\xi, \mu), \KL(\xi, \nu)\}
\ge \limsup_{n \to + \infty}\frac{1}{n}\log \frac{1}{\gamma \mu^{\otimes n}(E_n) + (1 - \gamma)\nu^{\otimes n}(E_n^c)} - \beta
\end{align*}
Since $\beta > 0$ is arbitrary, we can take a supremum over $\beta$ on the right.
\end{proof}


\section{Sample complexity}

We study the sample complexity of binary hypothesis testing.

\begin{definition}
  \label{def:binaryPriorSampleComplexity}
  %\lean{}
  %\leanok
  \uses{def:bayesBinaryRisk, def:priorSampleComplexity, lem:bayesBinaryRisk_properties}
  The sample complexity of simple binary hypothesis testing with prior $(\pi, 1 - \pi) \in \mathcal P(\{0, 1\})$ at risk level $\delta \in \mathbb{R}_{+, \infty}$ is
  \begin{align*}
  n(\mu, \nu, \pi, \delta) = \min\{n \in \mathbb{N} \mid B_\pi(\mu^{\otimes n}, \nu^{\otimes n}) \le \delta\} \: .
  \end{align*}
  This is the sample complexity $n_\xi^P(\delta)$ of Definition~\ref{def:priorSampleComplexity} specialized to simple binary hypothesis testing.
\end{definition}


\begin{lemma}
  \label{lem:binaryPriorSampleComplexity_eq_zero}
  %\lean{}
  %\leanok
  \uses{def:binaryPriorSampleComplexity}
  For $\delta \ge \min\{\pi, 1 - \pi\}$, the sample complexity of simple binary hypothesis testing is
  $n(\mu, \nu, \pi, \delta) = 0$~.
\end{lemma}

\begin{proof}
\uses{lem:bayesBinaryRisk_le}
\end{proof}


\begin{lemma}
  \label{lem:binaryPriorSampleComplexity_le_renyi}
  %\lean{}
  %\leanok
  \uses{def:Renyi, def:binaryPriorSampleComplexity}
  The sample complexity of simple binary hypothesis testing satisfies $n(\mu, \nu, \pi, \delta) \le n_0$~, with $n_0$ the smallest natural number such that
  \begin{align*}
  \log\frac{1}{\delta}
  \le \sup_{\alpha \in (0,1)} \left( n_0 (1 - \alpha)R_\alpha(\mu, \nu) + \alpha\log\frac{1}{\pi} + (1 - \alpha)\log\frac{1}{1 - \pi} \right)
  \end{align*}
\end{lemma}

\begin{proof}%\leanok
\uses{cor:bayesBinaryRisk_le_exp_renyi, lem:renyi_prod_n}
It suffices to show that $B_\pi(\mu^{\otimes n_0}, \nu^{\otimes n_0}) \le \delta$~.
By Corollary~\ref{cor:bayesBinaryRisk_le_exp_renyi}, for all $\alpha \in (0,1)$,
\begin{align*}
B_\pi(\mu^{\otimes n_0}, \nu^{\otimes n_0})
&\le \pi^{\alpha} (1 - \pi)^{1 - \alpha}e^{- (1 - \alpha)R_\alpha(\mu^{\otimes n_0}, \nu^{\otimes n_0})}
\: .
\end{align*}
The Rényi divergence tensorizes: $R_\alpha(\mu^{\otimes n_0}, \nu^{\otimes n_0}) = n_0 R_\alpha(\mu, \nu)$~.
\begin{align*}
B_\pi(\mu^{\otimes n_0}, \nu^{\otimes n_0})
&\le \exp\left(- \left( n_0 (1 - \alpha)R_\alpha(\mu, \nu) + \alpha\log\frac{1}{\pi} + (1 - \alpha)\log\frac{1}{1 - \pi} \right) \right)
\: .
\end{align*}
By definition, $n_0$ is such that the infimum over $\alpha$ of the right hand side is less than $\delta$.
\end{proof}


\begin{lemma}
  \label{lem:binaryPriorSampleComplexity_ge_renyi}
  %\lean{}
  %\leanok
  \uses{def:binaryPriorSampleComplexity, def:Renyi}
  For $\delta \le \pi \le 1/2$, the sample complexity of simple binary hypothesis testing satisfies
  \begin{align*}
  n(\mu, \nu, \pi, \delta) \ge \frac{\log\frac{\pi}{2\delta}}{R_{1/2}(\mu, \nu)}
  \: .
  \end{align*}
\end{lemma}

\begin{proof}%\leanok
\uses{thm:log_inv_bayesBinaryRisk_le_renyi, lem:renyi_prod_n}
By Theorem~\ref{thm:log_inv_bayesBinaryRisk_le_renyi},
\begin{align*}
\log\frac{\pi}{B_\pi(\mu^{\otimes n}, \nu^{\otimes n})} \le R_{1/2}(\mu^{\otimes n}, \nu^{\otimes n}) + \log 2
\: .
\end{align*}
The Rényi divergence tensorizes (Theorem~\ref{lem:renyi_prod_n}): $R_{1/2}(\mu^{\otimes n}, \nu^{\otimes n}) = n R_{1/2}(\mu, \nu)$~.
We finally use that $B_\pi(\mu^{\otimes n}, \nu^{\otimes n}) \le \delta$ for $n = n(\mu, \nu, \pi, \delta)$~.
\end{proof}


\begin{lemma}
  \label{lem:binaryPriorSampleComplexity_ge_jensenShannon}
  %\lean{}
  %\leanok
  \uses{def:binaryPriorSampleComplexity, def:jensenShannon}
  For $\delta \le \min\{\pi, 1 - \pi\}$, the sample complexity of simple binary hypothesis testing satisfies
  \begin{align*}
  n(\mu, \nu, \pi, \delta) \ge \frac{h_2(\pi) - h_2(\delta)}{\JS_\pi(\mu, \nu)}
  \: ,
  \end{align*}
  in which $h_2: x \mapsto x\log\frac{1}{x} + (1 - x)\log\frac{1}{1 - x}$ is the binary entropy function.
\end{lemma}

\begin{proof}%\leanok
\uses{thm:sub_bayesBinaryRisk_le_jensenShannon}
We start from Theorem~\ref{thm:sub_bayesBinaryRisk_le_jensenShannon}.
\begin{align*}
\JS_\pi(\mu^{\otimes n}, \nu^{\otimes n}) \ge h_2(\pi) - h_2(B_\pi(\mu^{\otimes n}, \nu^{\otimes n})) \: .
\end{align*}
By Lemma~\ref{lem:jensenShannon_prod_n}, $\JS_\pi(\mu^{\otimes n}, \nu^{\otimes n}) \le n \JS_\pi(\mu, \nu)$. The final result is obtained by using that $B_\pi(\mu^{\otimes n}, \nu^{\otimes n}) \le \delta$ for $n = n(\mu, \nu, \pi, \delta)$.
\end{proof}
